% ===== 00_preamble.tex =====
% !TEX program = lualatex
% ============================================================
%  統計検定1級 統計応用（医薬生物学）対策 医療統計学参考書
%  単一ファイル版（LuaLaTeX + LuaTeX-ja でコンパイル）
%    lualatex 医療統計学参考書.tex   （目次・相互参照のため2〜3回）
% ============================================================
\documentclass[a4paper,11pt,openany]{ltjsbook}
\usepackage[haranoaji,deluxe]{luatexja-preset}
\usepackage[top=24mm,bottom=24mm,left=22mm,right=22mm,headsep=8mm]{geometry}

% ---------- 数式 ----------
\usepackage{lmodern}
\usepackage{amsmath,amssymb,bm,mathtools}
\allowdisplaybreaks[2]

% ---------- 表・リスト・図 ----------
\usepackage{booktabs,longtable,array,multirow,tabularx,makecell}
\usepackage{enumitem}
\setlist{itemsep=0.2\baselineskip,topsep=0.4\baselineskip}
\usepackage{graphicx}
\usepackage{xcolor}
\usepackage{tikz}
\usetikzlibrary{arrows.meta,positioning,shapes.geometric,calc,fit,backgrounds}
\usepackage{pgfplots}
\pgfplotsset{compat=1.18}
\usepgfplotslibrary{fillbetween}
\usepackage{caption}
\captionsetup{font=small,labelsep=quad}
\usepackage[most]{tcolorbox}

% ---------- 色 ----------
\definecolor{bkblue}{HTML}{1F4E79}
\definecolor{bkgreen}{HTML}{2E6B30}
\definecolor{bkred}{HTML}{A12A2A}
\definecolor{bkorange}{HTML}{B35C00}
\definecolor{bkgray}{HTML}{555555}
\definecolor{bkpurple}{HTML}{5B3A8C}

% ---------- ハイパーリンク・相互参照 ----------
\usepackage[unicode,bookmarksnumbered=true,colorlinks=true,
  linkcolor=bkblue,citecolor=bkgreen,urlcolor=bkpurple,
  pdftitle={統計検定1級 統計応用（医薬生物学）対策 医療統計学参考書},
  pdfsubject={医療統計学},pdfkeywords={統計検定1級,医薬生物学,医療統計}]{hyperref}
\usepackage{bookmark}

% ---------- 定理型環境 ----------
% cleveref は \newtheorem より前に読み込む（共有カウンタの環境も正しい名前で参照される）
\usepackage[nameinlink]{cleveref}
\newtheorem{teigi}{定義}[chapter]
\newtheorem{teiri}[teigi]{定理}
\newtheorem{meidai}[teigi]{命題}
\newtheorem{reidai}{例題}[chapter]
\tcolorboxenvironment{teigi}{enhanced,breakable,colback=bkblue!4,colframe=bkblue!70,
  boxrule=0.5pt,left=4pt,right=4pt,top=2pt,bottom=2pt,arc=1pt}
\tcolorboxenvironment{teiri}{enhanced,breakable,colback=bkgreen!4,colframe=bkgreen!70,
  boxrule=0.5pt,left=4pt,right=4pt,top=2pt,bottom=2pt,arc=1pt}
\tcolorboxenvironment{meidai}{enhanced,breakable,colback=bkgreen!4,colframe=bkgreen!70,
  boxrule=0.5pt,left=4pt,right=4pt,top=2pt,bottom=2pt,arc=1pt}
\tcolorboxenvironment{reidai}{enhanced,breakable,colback=white,colframe=bkorange!80,
  boxrule=0.5pt,left=4pt,right=4pt,top=2pt,bottom=2pt,arc=1pt,borderline west={2pt}{0pt}{bkorange}}

\crefname{teigi}{定義}{定義}
\crefname{teiri}{定理}{定理}
\crefname{meidai}{命題}{命題}
\crefname{reidai}{例題}{例題}
\crefname{table}{表}{表}
\crefname{figure}{図}{図}
\crefformat{equation}{#2式(#1)#3}
\crefrangeformat{equation}{#3式(#1)#4〜#5(#2)#6}
\crefmultiformat{equation}{#2式(#1)#3}{，#2(#1)#3}{，#2(#1)#3}{，#2(#1)#3}
\crefformat{chapter}{#2第#1章#3}
\crefformat{appendix}{#2付録#1#3}
\crefformat{section}{#2#1節#3}
\crefformat{subsection}{#2#1項#3}
\crefformat{teigi}{#2定義#1#3}
\crefformat{teiri}{#2定理#1#3}
\crefformat{meidai}{#2命題#1#3}
\crefformat{reidai}{#2例題#1#3}
\crefformat{table}{#2表#1#3}
\crefformat{figure}{#2図#1#3}
\crefname{chapter}{章}{章}
\crefname{section}{節}{節}
\crefname{subsection}{項}{項}
\crefname{appendix}{付録}{付録}
\crefname{equation}{式}{式}
\newcommand{\crefpairconjunction}{，}
\newcommand{\crefmiddleconjunction}{，}
\newcommand{\creflastconjunction}{，}
\newcommand{\crefpairgroupconjunction}{，}
\newcommand{\crefmiddlegroupconjunction}{，}
\newcommand{\creflastgroupconjunction}{，}
\crefmultiformat{chapter}{#2第#1章#3}{，#2第#1章#3}{，#2第#1章#3}{，#2第#1章#3}
\crefrangeformat{chapter}{#3第#1章#4〜#5第#2章#6}
\crefmultiformat{section}{#2#1節#3}{，#2#1節#3}{，#2#1節#3}{，#2#1節#3}
\crefmultiformat{table}{#2表#1#3}{，#2表#1#3}{，#2表#1#3}{，#2表#1#3}
\crefmultiformat{figure}{#2図#1#3}{，#2図#1#3}{，#2図#1#3}{，#2図#1#3}
\crefmultiformat{teiri}{#2定理#1#3}{，#2定理#1#3}{，#2定理#1#3}{，#2定理#1#3}
\crefmultiformat{reidai}{#2例題#1#3}{，#2例題#1#3}{，#2例題#1#3}{，#2例題#1#3}

% 定理環境の見出しの後に改行を入れない・本文は立体
\makeatletter
\def\th@plain{\normalfont}
\makeatother

% ---------- 証明環境 ----------
\newenvironment{proof}[1][証明]{\par\smallskip\noindent{\sffamily\bfseries #1}\quad\ignorespaces}{\nobreak\hfill$\square$\par\medskip}

% ---------- 解答環境 ----------
% \begin{kaitou}{reidai:label} ... \end{kaitou}
\newenvironment{kaitou}[1]{%
  \par\medskip\noindent
  \begin{tcolorbox}[enhanced,breakable,colback=bkgray!3,colframe=bkgray!60,boxrule=0.4pt,
    left=4pt,right=4pt,top=2pt,bottom=2pt,arc=1pt,
    title={\cref{#1}の解答},fonttitle=\bfseries\small,coltitle=white,colbacktitle=bkgray!80]}%
  {\end{tcolorbox}\par}

% ---------- 本書独自のボックス ----------
\newtcolorbox{goalbox}{enhanced,breakable,colback=bkblue!3,colframe=bkblue,boxrule=0.6pt,arc=1pt,
  title={この章でできるようになること},fonttitle=\bfseries,colbacktitle=bkblue,coltitle=white,
  left=4pt,right=4pt}
\newtcolorbox{exambox}{enhanced,breakable,colback=bkorange!3,colframe=bkorange,boxrule=0.6pt,arc=1pt,
  title={過去問での出題},fonttitle=\bfseries,colbacktitle=bkorange,coltitle=white,
  left=4pt,right=4pt}
\newtcolorbox{pointbox}[1][要点]{enhanced,breakable,colback=bkgreen!3,colframe=bkgreen!80,boxrule=0.5pt,arc=1pt,
  title={#1},fonttitle=\bfseries\small,colbacktitle=bkgreen!80,coltitle=white,left=4pt,right=4pt}
\newtcolorbox{warnbox}[1][よくある誤り]{enhanced,breakable,colback=bkred!3,colframe=bkred!80,boxrule=0.5pt,arc=1pt,
  title={#1},fonttitle=\bfseries\small,colbacktitle=bkred!80,coltitle=white,left=4pt,right=4pt}
\newtcolorbox{linkbox}[1][数理統計との接続]{enhanced,breakable,colback=bkpurple!3,colframe=bkpurple!70,boxrule=0.5pt,arc=1pt,
  title={#1},fonttitle=\bfseries\small,colbacktitle=bkpurple!70,coltitle=white,left=4pt,right=4pt}
% 補足・発展（過去問に直接ない内容。出典を併記する）
\newtcolorbox{supbox}[1][補足]{enhanced,breakable,colback=white,colframe=bkgray!70,boxrule=0.4pt,arc=1pt,
  borderline west={2pt}{0pt}{bkgray!70},title={#1},fonttitle=\bfseries\small,
  colbacktitle=white,coltitle=bkgray,left=4pt,right=4pt,attach title to upper={\quad}}
\newtcolorbox{summarybox}{enhanced,breakable,colback=bkblue!2,colframe=bkblue!60,boxrule=0.5pt,arc=1pt,
  title={章末まとめ},fonttitle=\bfseries,colbacktitle=bkblue!60,coltitle=white,left=4pt,right=4pt}
% 流れ図 データ構造→推定対象→手法→仮定→計算→解釈
\newcommand{\flowline}[6]{%
  \begin{center}\small
  \begin{tikzpicture}[node distance=3mm,
    box/.style={draw=bkblue,rounded corners=1pt,fill=bkblue!5,inner sep=2pt,align=center,
      text width=21mm,minimum height=11mm,font=\scriptsize}]
  \node[box](a){\textbf{データ構造}\\#1};
  \node[box,right=of a](b){\textbf{推定対象}\\#2};
  \node[box,right=of b](c){\textbf{手法}\\#3};
  \node[box,right=of c](d){\textbf{仮定}\\#4};
  \node[box,right=of d](e){\textbf{計算}\\#5};
  \node[box,right=of e](f){\textbf{解釈}\\#6};
  \foreach \x/\y in {a/b,b/c,c/d,d/e,e/f}\draw[-{Stealth[length=2mm]},bkblue](\x)--(\y);
  \end{tikzpicture}\end{center}}

% ---------- 重要度ラベル ----------
\newcommand{\tagbox}[2]{\tcbox[on line,boxsep=0.5pt,left=2pt,right=2pt,top=0.5pt,bottom=0.5pt,
  colframe=#1,colback=#1!10,boxrule=0.4pt,arc=1pt]{\scriptsize\bfseries\textcolor{#1}{#2}}}
\newcommand{\Hissu}{\tagbox{bkred}{必須}}
\newcommand{\Hinshutsu}{\tagbox{bkorange}{頻出}}
\newcommand{\Youchui}{\tagbox{bkpurple}{要注意}}
\newcommand{\Hatten}{\tagbox{bkgray}{発展}}
% 過去問参照 \kakomon{2022}{2}
\newcommand{\kakomon}[2]{\mbox{#1年 医薬 問#2}}
\newcommand{\kakomonSM}[2]{\mbox{#1年 数理 問#2}}

% ---------- 数式マクロ ----------
\newcommand{\E}{\mathrm{E}}
\newcommand{\V}{\mathrm{V}}
\newcommand{\Var}{\mathrm{Var}}
\newcommand{\Cov}{\mathrm{Cov}}
\newcommand{\Corr}{\mathrm{Corr}}
\newcommand{\Prob}{\mathrm{P}}
\newcommand{\SE}{\mathrm{SE}}
\newcommand{\SD}{\mathrm{SD}}
\newcommand{\dd}{\mathrm{d}}
\newcommand{\RR}{\mathrm{RR}}
\newcommand{\RD}{\mathrm{RD}}
\newcommand{\OR}{\mathrm{OR}}
\newcommand{\HR}{\mathrm{HR}}
\newcommand{\logit}{\operatorname{logit}}
\newcommand{\indep}{\mathop{\perp\!\!\!\perp}}
\newcommand{\iid}{\overset{\text{i.i.d.}}{\sim}}
\newcommand{\Normal}{\mathrm{N}}
\newcommand{\Bin}{\mathrm{Bin}}
\newcommand{\Po}{\mathrm{Po}}
\newcommand{\Ga}{\mathrm{Ga}}
\newcommand{\Be}{\mathrm{Be}}
\newcommand{\Ex}{\mathrm{Exp}}
\newcommand{\NB}{\mathrm{NB}}
\newcommand{\Unif}{\mathrm{U}}
\newcommand{\chisq}{\chi^2}
\newcommand{\transp}{^{\top}}
\newcommand{\ind}{\mathbb{1}}
\newcommand{\hatS}{\widehat{S}}

% ---------- その他 ----------
\setcounter{tocdepth}{1}
\setcounter{secnumdepth}{2}
\renewcommand{\arraystretch}{1.15}
\newcolumntype{L}[1]{>{\raggedright\arraybackslash}p{#1}}
\newcolumntype{C}[1]{>{\centering\arraybackslash}p{#1}}
\newcommand{\emphx}[1]{\textbf{\textcolor{bkblue}{#1}}}

% ===== 01_front.tex =====
\begin{document}

% ------------------------------------------------------------
% 表紙
% ------------------------------------------------------------
\begin{titlepage}
\begin{tikzpicture}[remember picture,overlay]
  \fill[bkblue] (current page.north west) rectangle ([yshift=-92mm]current page.north east);
  \fill[bkorange] ([yshift=-92mm]current page.north west) rectangle ([yshift=-95mm]current page.north east);
\end{tikzpicture}
\vspace*{8mm}
\begin{center}
{\color{white}\large 統計検定1級　統計応用（医薬生物学）対策\par}
\vspace{10mm}
{\color{white}\fontsize{30}{40}\selectfont\bfseries 医療統計学参考書\par}
\vspace{6mm}
{\color{white}\large データ構造から手法を選び，導出し，解釈する\par}
\end{center}
\vspace{40mm}
\begin{center}
\begin{tikzpicture}[node distance=4mm,
  b/.style={draw=bkblue,rounded corners=2pt,fill=bkblue!6,minimum width=22mm,minimum height=10mm,align=center,font=\small}]
  \node[b](a){データ構造};
  \node[b,right=of a](b){推定対象};
  \node[b,right=of b](c){手法};
  \node[b,right=of c](d){仮定};
  \node[b,right=of d](e){計算};
  \node[b,right=of e](f){解釈};
  \foreach \x/\y in {a/b,b/c,c/d,d/e,e/f}\draw[-{Stealth[length=2.5mm]},thick,bkblue](\x)--(\y);
\end{tikzpicture}
\end{center}
\vspace{30mm}
\begin{center}
\begin{tabular}{l}
過去問9回分（2016--2019，2021--2025）の分析に基づく\\
生存時間解析・分割表・診断・交絡・臨床試験デザイン・メタアナリシス・欠測・ベイズ
\end{tabular}
\end{center}
\vfill
\begin{center}
{\small 学習用電子参考書　2026年版}
\end{center}
\end{titlepage}

\frontmatter
\pagenumbering{roman}

% ------------------------------------------------------------
\chapter*{はじめに}
\addcontentsline{toc}{chapter}{はじめに}
\markboth{はじめに}{はじめに}

本書は，統計検定1級「統計応用（医薬生物学）」の受験者のための参考書である．
医療統計の入門書は多いが，1級の記述式問題は「手法の名前を知っている」だけでは解けない．
典型的な設問は次の4段で構成されている．
\begin{enumerate}[label=(\arabic*)]
\item 問題文の研究デザインとデータ構造を読み取り，
\item 推定したい量（効果指標）と，それに合う検定・推定の方法を選び，
\item 尤度・分散・漸近分布などを導出して数値を計算し，
\item 結果を医学・疫学・臨床試験の文脈で解釈する．
\end{enumerate}
本書はこの4段をそのまま章の骨格にしている．
各テーマを
\[
\boxed{\ \text{データ構造}\ \rightarrow\ \text{推定対象}\ \rightarrow\ \text{手法}\ \rightarrow\ \text{仮定}\ \rightarrow\ \text{計算}\ \rightarrow\ \text{解釈}\ }
\]
の流れで説明する．

内容は2016年から2025年までの9回分（2020年は実施なし）の医薬生物学の過去問と，
同じ年度の統計数理の過去問を分析して設計した．
設問ごとに主テーマ・副テーマ・必要能力を分類し，出題の厚い分野に紙幅を多く割いている．
ただし「過去に出ていないから不要」とは考えず，標準的な医療統計として最低限必要な内容は補った．
過去問で実際に問われた内容と，一般的に重要だが直接は問われていない内容は，
「過去問での出題」ボックスと「補足」「発展」の表示で区別している．

過去問の本文は転載せず，「\kakomon{2022}{2}」のように年と問番号で参照する．
例題はすべて，過去問と同じ能力を問うように作ったオリジナル問題である．

% ------------------------------------------------------------
\section*{本書の対象読者}
\addcontentsline{toc}{section}{本書の対象読者}
次の読者を想定している．
\begin{itemize}
\item 統計検定準1級程度の内容を習得している．
\item 『現代数理統計学の基礎』（久保川達也）第2〜7章，すなわち確率分布，多次元分布，標本分布と極限定理，推定，検定を学習済みである．
\item 最尤法，Fisher情報量，デルタ法，漸近正規性，尤度比検定，基本的な区間推定・検定を，ある程度自分で使える．
\end{itemize}
したがって「期待値とは」「最尤法とは」といった一般の数理統計は一から説明しない．
医療統計の議論で必要になった箇所だけを「数理統計との接続」として補う．
逆に，順位統計量，打ち切りのある尤度，部分尤度，多変量デルタ法，分割表固有の方法，
多重性の理論などは上記の範囲に含まれないので，本書の中で導出する．

% ------------------------------------------------------------
\section*{本書の使い方}
\addcontentsline{toc}{section}{本書の使い方}
\paragraph{章の構成}
各章は原則として次の順に並ぶ．内容に応じて一部を省略・統合している．
\begin{center}\small
\begin{tabular}{@{}ll@{}}
\toprule
1. この章でできるようになること & 8. 結果の解釈\\
2. 過去問での出題 & 9. 手法選択\\
3. 問題設定 & 10. よくある誤り\\
4. 基本概念 & 11. 数理統計との接続\\
5. 数理 & 12. オリジナル例題\\
6. 推定・検定 & 13. 解答\\
7. 仮定 & 14. 章末まとめ\\
\bottomrule
\end{tabular}
\end{center}

\paragraph{ボックスと表示}
\begin{center}\small
\begin{tabular}{@{}lL{110mm}@{}}
\toprule
表示 & 意味\\
\midrule
\Hissu & 試験で導出・計算できなければならない内容\\
\Hinshutsu & 9回の中で複数回出題された内容\\
\Youchui & 取り違えやすい概念・符号・自由度など\\
\Hatten & 試験に直接出にくいが理解を深める内容\\
\textcolor{bkorange}{\textbf{過去問での出題}} & その章のテーマが実際にどの年・問で問われたか\\
\textcolor{bkred}{\textbf{よくある誤り}} & 答案で起こりやすい誤り\\
\textcolor{bkpurple}{\textbf{数理統計との接続}} & 『現代数理統計学の基礎』や統計数理の過去問との対応\\
\textcolor{bkgray}{\textbf{補足・発展}} & 過去問・指定教材に直接ない内容（Web資料等から補った場合は出典を付記）\\
\bottomrule
\end{tabular}
\end{center}

\paragraph{記法}
本書全体で次の記法を用いる．
\begin{itemize}
\item $z_\alpha$ は標準正規分布の上側 $100\alpha\%$ 点（$\Prob(Z>z_\alpha)=\alpha$）．
      $t_\nu(\alpha)$，$\chi^2_\nu(\alpha)$ も上側点．よく使う値は $z_{0.05}=1.645$，$z_{0.025}=1.960$，$z_{0.10}=1.282$，$z_{0.20}=0.842$．
\item $\beta$ は第2種の過誤確率，検出力は $1-\beta$．
      \Youchui 『現代数理統計学の基礎』では $\beta(\theta)$ を検出力関数そのものとして用いるので逆になる．
\item 不偏分散 $s^2=\frac{1}{n-1}\sum(x_i-\bar x)^2$．
      \Youchui 同書では $S^2$ が除数 $n$，$V^2$ が除数 $n-1$ である．
\item 指数分布 $\Ex(\lambda)$ は率（ハザード）$\lambda$ による表記で，密度 $\lambda e^{-\lambda t}$，平均 $1/\lambda$．
\item ガンマ分布 $\Ga(a,b)$ は\textbf{形状 $a$・率 $b$}（密度 $\frac{b^a}{\Gamma(a)}x^{a-1}e^{-bx}$，平均 $a/b$，分散 $a/b^2$）．
      試験（\kakomon{2025}{2}，\kakomonSM{2022}{3}）と同じ流儀である．
      \Youchui 『現代数理統計学の基礎』は $\Ga(\alpha,\beta)$ を\textbf{形状・尺度}（平均 $\alpha\beta$）で定義している．
\item $\Normal(\mu,\sigma^2)$ の第2引数は分散．$\log$ は自然対数．
\item $I_n(\theta)$ は標本全体のFisher情報量，$I_1(\theta)$ は1観測あたり．
\end{itemize}

\enlargethispage{2\baselineskip}
\paragraph{数値計算について}
統計検定1級では関数電卓が使えない．本書の例題の数値計算は，四則演算・平方根と巻末付表（正規分布表・$t$ 分布表・$\chi^2$ 分布表・指数関数と対数の表）で実行できる形で示し，
付表を引いた箇所では「$t_9(0.05)=1.833$」のようにどの値を使ったかを明記する．

% ------------------------------------------------------------
\chapter*{統計検定1級・医薬生物学の出題構造}
\addcontentsline{toc}{chapter}{統計検定1級・医薬生物学の出題構造}
\markboth{出題構造}{出題構造}

\paragraph{試験の形式}
統計検定1級は「統計数理」と「統計応用」の2科目からなり，それぞれ90分の記述式試験である．
統計応用は人文科学・社会科学・理工学・医薬生物学の4分野から1つを選んで受験し，
大問5問のうち3問を選択して解答する（1問あたり約30分）\cite{web:lec,web:academaid}．
医薬生物学の大問のうち問1〜問4が分野固有の問題，問5は4分野共通の問題である
（本書で分析した9回すべてでこの構成であった）．

\paragraph{9回分の傾向}
医薬生物学の問1〜4（計36問）を主テーマで分類すると，おおよそ次のようになる
（1問が複数テーマにまたがる場合は重複して数える）．
\begin{center}\small
\begin{tabular}{@{}L{62mm}cL{80mm}@{}}
\toprule
分野 & 出題数 & 代表的な問\\
\midrule
生存時間解析（KM・ログランク・指数モデル・Cox・RMST） & 7 & \kakomon{2016}{3}，\kakomon{2019}{1}，\kakomon{2022}{1}\\
2 値データ・分割表・McNemar・MH & 7 & \kakomon{2018}{2}，\kakomon{2021}{3}，\kakomon{2022}{2}\\
診断検査・ROC & 5 & \kakomon{2018}{3}，\kakomon{2019}{3}，\kakomon{2024}{1}\\
交絡・層別・標準化・傾向スコア & 4 & \kakomon{2017}{4}，\kakomon{2019}{2}\\
中間解析・二段階・多重性・用量反応 & 5 & \kakomon{2017}{2}，\kakomon{2024}{4}，\kakomon{2025}{4}\\
縦断データ・混合モデル・ベースライン調整 & 3 & \kakomon{2022}{3}，\kakomon{2023}{4}，\kakomon{2024}{2}\\
ロジスティック回帰 & 3 & \kakomon{2018}{4}，\kakomon{2023}{1}\\
ノンパラメトリック（順位）検定 & 3 & \kakomon{2016}{1}，\kakomon{2019}{4}，\kakomon{2025}{1}\\
検出力・症例数・イベント数 & 3 & \kakomon{2018}{1}，\kakomon{2023}{2}\\
同等性・生物学的同等性 & 2 & \kakomon{2016}{4}，\kakomon{2021}{2}\\
競合リスク & 2 & \kakomon{2021}{1}，\kakomon{2025}{3}\\
メタアナリシス & 2 & \kakomon{2021}{4}，\kakomon{2023}{3}\\
ベイズ・ポアソン率 & 2 & \kakomon{2022}{4}，\kakomon{2025}{2}\\
欠測データ・多重補完 & 1 & \kakomon{2024}{3}\\
\bottomrule
\end{tabular}
\end{center}
生存時間解析は9回中7回で何らかの形で出題されており，最重要分野である．
次いで2 値データ（分割表・診断・ロジスティック）が厚い．
2021年以降は，競合リスク，メタアナリシスの変量効果，混合モデル，多重補完，多重性，ベイズ予測など，
臨床試験の方法論に寄った出題が増えている．

\enlargethispage{4\baselineskip}
\paragraph{問われる能力}
設問のほとんどは「定義を式で書く→導出する→データで計算する→解釈を述べる」の組み合わせである．
導出の多くは，尤度の最大化，分散の計算（多項分布の共分散・条件付き分散を含む），
デルタ法，正規近似から検出力式を立てる，といった数理統計の基本操作に帰着する．
統計数理の過去問は医薬固有のモデル（打ち切り，Cox，分割表）を直接は扱わないが，
指数分布・ポアソン分布・ガンマ混合・デルタ法・Fisher情報などを
医薬の問題より深く問う（詳細は各章の「数理統計との接続」）．

% ------------------------------------------------------------
\chapter*{過去問9回分の出題マップ}
\addcontentsline{toc}{chapter}{過去問9回分の出題マップ}
\markboth{出題マップ}{出題マップ}

表中の数字は医薬生物学の問番号，「5c」は4分野共通の問5で内容が関連するものを表す．
設問単位の詳細な分類（副テーマ・必要能力）は\cref{app:D}の過去問対応表にある．

\begin{center}\small
\setlength{\tabcolsep}{3.2pt}
\begin{tabular}{@{}clccccccccc@{}}
\toprule
章 & テーマ & 2016 & 2017 & 2018 & 2019 & 2021 & 2022 & 2023 & 2024 & 2025\\
\midrule
\ref{ch:01} & 問題の読み方・P 値・CI & 2,5c & & & & & 1 & & & \\
\ref{ch:02} & 研究デザイン・効果指標 & & 4,5c & & 2 & & 2 & & & \\
\ref{ch:03} & 連続・ノンパラメトリック & 1,5c & & & 4 & & & & & 1\\
\ref{ch:04} & 2 値データ・分割表 & 2 & 4 & 2 & 3 & 3 & 2 & & 1 & \\
\ref{ch:05} & 診断検査・ROC & & & 3,4 & 3 & 5c & & 1 & 1 & \\
\ref{ch:06} & ロジスティック回帰 & & & 4 & 2 & & & 1 & & \\
\ref{ch:07} & 交絡・層別・マッチング & & 4 & & 2 & 3 & 2 & & & \\
\ref{ch:08} & 縦断・ベースライン・混合 & & & & & & 3,5c & 4 & 2 & \\
\ref{ch:09} & 生存時間解析の基礎 & 3 & 1 & & 1 & 1 & 1 & & & \\
\ref{ch:10} & Cox・RMST・パラメトリック & 3 & 1 & 1 & 1 & & 1 & 2 & & \\
\ref{ch:11} & 競合リスク & & & & & 1 & & & & 3\\
\ref{ch:12} & 検出力・症例数 & & 5c & 1 & & & & 2,5c & & 1\\
\ref{ch:13} & 同等性・生物学的同等性 & 4 & & & & 2 & & & & \\
\ref{ch:14} & 中間解析・多重性・用量反応 & 2 & 2,3 & & & & & & 4 & 4\\
\ref{ch:15} & メタアナリシス & & & & & 4 & & 3 & & \\
\ref{ch:16} & 欠測データ・多重補完 & & & & 5c & & & & 3 & \\
\ref{ch:17} & ポアソン率・ベイズ & & & & & & 4 & & & 2\\
\bottomrule
\end{tabular}
\end{center}

\paragraph{読み取れること}
\begin{itemize}
\item 生存時間（第9〜11章）は2016〜2023年の毎回出題され，2025年は競合リスクの理論として出た．
\item 分割表（第4章）は形を変えて繰り返し出る．特に「対応のある2 値」（McNemar型）は2018・2019・2024年，
      「層別・MH」は2017・2021・2022年に出ている．
\item 近年は臨床試験の方法論（多重性・欠測・混合モデル・ベースライン調整）が増えている．
\end{itemize}

% ------------------------------------------------------------
\chapter*{推奨学習ロードマップ}
\addcontentsline{toc}{chapter}{推奨学習ロードマップ}
\markboth{学習ロードマップ}{学習ロードマップ}

\begin{center}
\begin{tikzpicture}[node distance=5mm and 4mm,
  st/.style={draw=bkblue,rounded corners=2pt,fill=bkblue!5,text width=62mm,align=left,font=\small,inner sep=4pt},
  lb/.style={font=\small\bfseries,text=bkorange,text width=22mm,align=right}]
\node[st](s1){第1章 読み方，第2章 デザインと効果指標，第4章 分割表，第3章 連続・順位検定};
\node[lb,left=of s1]{段階1\\基礎体力};
\node[st,below=of s1](s2){第9章 生存時間の基礎，第10章 Cox・RMST・指数モデル，第11章 競合リスク};
\node[lb,left=of s2]{段階2\\最頻出};
\node[st,below=of s2](s3){第5章 診断・ROC，第6章 ロジスティック回帰，第7章 交絡・層別・PS};
\node[lb,left=of s3]{段階3\\2 値の応用};
\node[st,below=of s3](s4){第12章 検出力・症例数，第13章 同等性，第14章 中間解析・多重性};
\node[lb,left=of s4]{段階4\\臨床試験};
\node[st,below=of s4](s5){第8章 縦断・混合，第15章 メタアナリシス，第16章 欠測，第17章 ベイズ};
\node[lb,left=of s5]{段階5\\近年の論点};
\node[st,below=of s5,fill=bkorange!8,draw=bkorange](s6){付録A 手法選択表，付録B 公式集，付録E チェックリスト，過去問演習};
\node[lb,left=of s6]{直前期};
\foreach \a/\b in {s1/s2,s2/s3,s3/s4,s4/s5,s5/s6}\draw[-{Stealth},thick,bkblue](\a)--(\b);
\end{tikzpicture}
\end{center}

\begin{itemize}
\item 段階1で，効果指標（RD・RR・OR・HR）と「独立か対応ありか」の区別を固める．以降のすべての章の前提になる．
\item 段階2は最頻出分野である．KM推定値の表の作成，ログランク検定の手計算，打ち切りのある指数分布の尤度，
      Cox部分尤度の書き下しは，時間を計って何度も練習する．
\item 段階3・4は，式の導出よりも「どの手法をなぜ使うか」「結果をどう解釈するか」の記述が得点源になる．
\item 段階5は出題回数こそ少ないが，近年の出題が多く，一度理解すれば短時間で得点できる．
\item 直前期は付録Eのチェックリストで，自力で導出できない項目を潰す．
\end{itemize}

\tableofcontents

\mainmatter
\pagenumbering{arabic}

% ===== ch01.tex =====
\chapter{医療統計の問題をどう読むか}\label{ch:01}

\begin{goalbox}
\begin{itemize}
\item 問題文から，アウトカムの型（連続・2 値・カウント・生存時間・反復測定）と，観測単位の独立性・対応を読み取れる．
\item ランダム化比較試験か観察研究かを判別し，交絡の可能性を指摘できる．
\item 推定したい量（効果指標）を式で書き，推定と検定の役割を区別できる．
\item P 値と信頼区間を正しい日本語で解釈できる（「有意差なし」と「差がない」を区別できる）．
\item 群ごとの信頼区間の重なりと2標本検定の結果の関係を導出できる．
\end{itemize}
\end{goalbox}

\begin{exambox}
\begin{itemize}
\item \kakomon{2016}{2}〔2〕〔4〕：$P=0.077$ の結果から言えること，傾向性検定で $P=0.042$ になったときの解釈．
\item \kakomon{2016}{5}（共通）：群ごとの95\%信頼区間の重なりと2標本 $t$ 検定の関係（重なりが区間幅の約29\%未満なら有意）．
\item \kakomon{2022}{1}〔5〕：ハザード比の95\%信頼区間が1を含むときの解釈．
\item \kakomon{2016}{1} と \kakomon{2019}{4}：同じ「2つの測定値の比較」でも，対応あり（治療前後）と独立（2群）で手法が変わる．
\end{itemize}
設問に「結果を解釈せよ」「この結果から言えることを述べよ」とあれば，本章の型で答える．
\end{exambox}

\section{問題設定：医療統計の設問の読み方}\label{sec:01-read}
医薬生物学の設問は，ほぼ例外なく次の情報を問題文の中に含んでいる．
答案を書き始める前に，これらを明示的に書き出す習慣をつけるとよい．
\begin{enumerate}[label=(\alph*)]
\item \textbf{研究デザイン}：ランダム化比較試験（並行群間・クロスオーバー），コホート研究，症例対照研究，マッチング，単群試験，メタアナリシスなど．
\item \textbf{観測単位とその関係}：患者，ペア，層，試験．同一患者の反復測定か，独立な2群か．
\item \textbf{アウトカムの型}：連続，2 値，カウント（人時間つき），生存時間（打ち切りつき），反復測定．
\item \textbf{比較の目的}：優越性・同等性・非劣性，関連の有無，予測，用量反応．
\item \textbf{推定対象}：平均差，リスク差，リスク比，オッズ比，ハザード比，AUC，感度など．
\item \textbf{確率モデル}：二項，多項，ポアソン，正規，指数，比例ハザードなど．問題文に明示されることが多い．
\end{enumerate}
(a)〜(c)が決まると候補となる手法はかなり絞られ，(d)〜(f)で一意に決まる．

\section{基本概念}\label{sec:01-basic}
\subsection{アウトカムの種類}\label{subsec:01-outcome}
\begin{table}[htbp]
\centering\small
\caption{アウトカムの型と代表的な要約・手法}\label{tab:01-outcome}
\begin{tabular}{@{}lL{35mm}L{35mm}L{45mm}@{}}
\toprule
型 & 例 & 要約・効果指標 & 代表的な手法（本書の章）\\
\midrule
連続 & 血圧，検査値，AUC（薬物動態） & 平均差，中央値の差，幾何平均比 & $t$ 検定，順位検定（\ref{ch:03}），ANCOVA（\ref{ch:08}），同等性（\ref{ch:13}）\\
2 値 & 有効／無効，発症の有無 & リスク差，リスク比，オッズ比 & $\chi^2$ 検定，McNemar，MH（\ref{ch:04}），ロジスティック（\ref{ch:06}）\\
カウント & 有害事象件数，死亡数 & 率（人時間あたり），率比 & ポアソンモデル（\ref{ch:17}）\\
生存時間 & 死亡・再発までの時間 & 生存関数，中央生存時間，ハザード比，RMST & KM，ログランク（\ref{ch:09}），Cox（\ref{ch:10}），競合リスク（\ref{ch:11}）\\
反復測定 & 同一患者の複数時点の値 & 時点別平均，変化量 & 混合モデル（\ref{ch:08}）\\
\bottomrule
\end{tabular}
\end{table}

生存時間は「時間という連続量」だが，観察打ち切り（censoring）があるため連続データの手法をそのまま適用できない．
これが生存時間解析という独立した分野が必要になる理由である（\cref{ch:09}）．
カウントは，観察期間（人時間）が個体や集団ごとに異なるので「件数」ではなく「率」を比較する．

\subsection{独立なデータと対応のあるデータ}\label{subsec:01-paired}
\Hissu 同じ患者から得た2つの値（治療前後，左右の眼，2種類の検査），マッチさせたペア，
クロスオーバー試験の2期の値は\emphx{対応のあるデータ}である．
対応があると2つの値は一般に正の相関をもつ．

\begin{meidai}[対応のある差の分散]\label[meidai]{prop:01-pairvar}
$(X_i,Y_i)$ $(i=1,\dots,n)$ が独立に，$\V[X_i]=\sigma_X^2$，$\V[Y_i]=\sigma_Y^2$，$\Corr(X_i,Y_i)=\rho$ に従うとする．
$D_i=Y_i-X_i$ とすると
\begin{equation}
\V[\bar D]=\frac{\sigma_X^2+\sigma_Y^2-2\rho\sigma_X\sigma_Y}{n}.\label{eq:01-pairvar}
\end{equation}
一方，大きさ $n$ の独立な2群の平均の差 $\bar Y-\bar X$ の分散は $(\sigma_X^2+\sigma_Y^2)/n$ である．
\end{meidai}
\begin{proof}
$\V[D_i]=\V[Y_i]+\V[X_i]-2\Cov(X_i,Y_i)$ と $D_i$ の独立性から直ちに従う．
\end{proof}
$\rho>0$ なら対応を利用した比較の方が分散が小さい．
逆に，対応のあるデータに独立2群の手法を使うと，分散を過大評価して検出力を失い，
独立なデータに対応のある手法は使えない（ペアが定義できない）．
2 値データでも同じで，対応があればMcNemar型（\cref{sec:04-mcnemar}），なければ2標本の比率の比較になる．

\subsection{ランダム化比較試験と観察研究}\label{subsec:01-rct}
\emphx{ランダム化}（無作為割付）は，治療の割付を患者の特性と確率的に独立にする操作である．
これにより，測定されたか否かにかかわらず，あらゆる予後因子の分布が群間で（期待値の意味で）等しくなり，
群間差を治療効果として解釈できる．
観察研究では治療の選択が患者の特性に依存するため，群間差に\emphx{交絡}が混入しうる（\cref{ch:07}）．
設問が観察研究なら「交絡の可能性」「調整の方法」に触れることが，解釈の問いでほぼ必ず求められる．

\subsection{推定したい量}\label{subsec:01-estimand}
推定したい母集団の量を先に言葉と式で定めることが，手法選択の出発点である．
たとえば2群の有効率 $\pi_1,\pi_0$ があるとき，
\[
\RD=\pi_1-\pi_0,\qquad \RR=\frac{\pi_1}{\pi_0},\qquad \OR=\frac{\pi_1/(1-\pi_1)}{\pi_0/(1-\pi_0)}
\]
はすべて「差がない」ことを同じ帰無仮説 $\pi_1=\pi_0$ で表すが，推定量・分散・解釈は異なる（\cref{ch:02}）．
臨床試験では，推定対象を母集団・治療・変数・中間事象の扱い・要約指標の5要素で規定する
「estimand」の枠組みが国際的に採用されている（\cref{sec:02-estimand}）．

\section{推定と検定}\label{sec:01-inference}
\subsection{推定と検定の役割}
\begin{itemize}
\item \textbf{点推定と区間推定}：効果の\emph{大きさ}とその不確かさを示す．臨床的判断の主役．
\item \textbf{検定}：帰無仮説（通常「効果なし」）がデータと両立するかを判定する．効果の大きさは示さない．
\end{itemize}
医学研究の報告では，効果指標の点推定値と95\%信頼区間を示し，P 値は補助的に示すのが標準である．

\subsection{P 値}\label{subsec:01-pvalue}
\begin{teigi}[P 値]\label{def:01-pvalue}
検定統計量 $T$（大きいほど帰無仮説から離れる）と観測値 $t_{\mathrm{obs}}$ に対し，
\[
p=\sup_{\theta\in\Theta_0}\Prob_\theta\bigl(T\ge t_{\mathrm{obs}}\bigr)
\]
をP 値という．「帰無仮説が正しいとしたとき，観測された結果と同じかそれ以上に極端な結果が得られる確率」である．
\end{teigi}
\Youchui P 値は「帰無仮説が正しい確率」ではなく，「効果の大きさ」でもない．
同じ効果の大きさでも標本サイズが大きければP 値は小さくなる．

\begin{pointbox}[P 値の解釈の型（有意水準5\%）]
\begin{itemize}
\item $p<0.05$：「有意水準5\%で帰無仮説は棄却され，〇〇に差があると言える．」
      さらに効果の方向と大きさ（点推定値・信頼区間）を添える．
\item $p\ge 0.05$：「有意水準5\%で帰無仮説は棄却されず，〇〇に差があるとは言えない．」
      \textbf{「差がない」「同等である」とは言わない}．検出力不足の可能性にも触れる．
\end{itemize}
\end{pointbox}

\kakomon{2016}{2} では，4群の有効率の一様性の $\chi^2$ 検定で $p=0.077$ となった．
正しい解釈は「すべての群で母有効率が等しいという帰無仮説は棄却されず，
4群間で有効率に差があるとは言えない」である．
同じデータで用量の順序を利用したCochran--Armitage傾向性検定（\cref{sec:04-trend}）は $p=0.042$ となり有意であった．
対立仮説を「用量とともに単調に変化」という方向に絞ると検出力が上がるためであり，
「どちらの検定を使うか」はデータを見る前に研究目的から決めておくべきことである．

\subsection{信頼区間}\label{subsec:01-ci}
\begin{teigi}[信頼区間]\label{def:01-ci}
データの関数 $L(\bm X),U(\bm X)$ が，すべての $\theta$ で $\Prob_\theta(L\le\theta\le U)\ge 1-\alpha$ を満たすとき，
$[L,U]$ を $\theta$ の信頼係数 $100(1-\alpha)\%$ の信頼区間という．
\end{teigi}
確率的なのは区間 $[L,U]$ の方で，$\theta$ は定数である．
「真の値が95\%の確率でこの区間に入る」は頻度論の解釈としては不正確で，
「同じ手順を繰り返すと，構成される区間の95\%が真の値を含む」が正しい．

\begin{meidai}[検定と信頼区間の双対性]\label[meidai]{prop:01-duality}
各 $\theta_0$ について有意水準 $\alpha$ の検定の受容域を $A(\theta_0)$ とすると，
$C(\bm x)=\{\theta_0:\bm x\in A(\theta_0)\}$ は信頼係数 $1-\alpha$ の信頼区間（信頼集合）である．
逆に信頼区間から検定を作ることもできる．
\end{meidai}
\begin{proof}
$\theta_0\in C(\bm X)\iff \bm X\in A(\theta_0)$ であり，$\Prob_{\theta_0}(\bm X\in A(\theta_0))\ge1-\alpha$ である．
\end{proof}
したがって，同じ統計量（例：Wald 統計量）から検定と区間を作る場合，
「$H_0:\theta=\theta_0$ の両側 $\alpha$ 検定が棄却されない」$\iff$「$100(1-\alpha)\%$ 信頼区間が $\theta_0$ を含む」．
Pearson $\chi^2$（スコア型）と Woolf の区間（Wald 型）のように作り方が異なると，境界付近で結論が一致しないことがある．
比の指標（RR，OR，HR）では $\theta_0=1$，差の指標（RD，平均差）では $\theta_0=0$ が「効果なし」である．
\Youchui 同等性の判定には90\%信頼区間を使う（\cref{ch:13}）．これは片側5\%の検定を2つ行うことに対応する．

\kakomon{2022}{1} では，Cox回帰で $\hat\beta=0.06$，$\SE=0.18$ から
$\HR=e^{0.06}\approx1.06$，95\%信頼区間 $(e^{0.06-1.96\times0.18},e^{0.06+1.96\times0.18})\approx(0.75,1.51)$ を得る．
区間が1を含むので「試験群とプラセボ群で死亡までの時間に差があるとは言えない」と解釈する．
このとき区間の幅も解釈に含めるとよい：$0.75$ から $1.51$ まで，すなわち25\%のハザード減少から51\%の増加まで，
臨床的に重要な差を否定できない程度にデータが少ない，という情報を区間は伝えている．

\subsection{片側と両側}\label{subsec:01-sided}
優越性を示す確認的試験では，片側2.5\%の検定と両側5\%の検定は，棄却限界（$z_{0.025}=1.96$）が一致する．
非劣性・同等性や群逐次デザインでは片側で考える方が自然なことが多い．
設問が片側か両側かを明示している場合はそれに従い，P 値を計算するときは
「両側P 値＝片側P 値×2」（対称な分布の場合）であることに注意する．

\section{数理：信頼区間の重なりと2標本検定}\label{sec:01-overlap}
\Hinshutsu 論文の図でよく見る「群ごとの平均と95\%信頼区間」は，群間比較の検定と同じではない．
\kakomon{2016}{5} の内容を一般化して導く．

2群の大きさがともに $n$，標本標準偏差がともに $s$ とし，$\bar x>\bar y$ とする．
各群の母平均の95\%信頼区間は $\bar x\pm c_1 s/\sqrt n$，$\bar y\pm c_1 s/\sqrt n$，$c_1=t_{n-1}(0.025)$ である．
等分散の2標本 $t$ 検定（プールした分散は $s^2$）は
\begin{equation}
\bar x-\bar y>c_2\sqrt{\tfrac{2}{n}}\,s,\qquad c_2=t_{2n-2}(0.025)\label{eq:01-ttest}
\end{equation}
のとき有意である．2つの区間の重なりの長さを $D$ とすると
\[
D=\Bigl(\bar y+c_1\frac{s}{\sqrt n}\Bigr)-\Bigl(\bar x-c_1\frac{s}{\sqrt n}\Bigr)=2c_1\frac{s}{\sqrt n}-(\bar x-\bar y).
\]
\begin{meidai}\label[meidai]{prop:01-overlap}
(1) 2つの信頼区間が重ならなければ（$D<0$），2標本 $t$ 検定は有意である．\\
(2) 重なりがあっても，
\begin{equation}
\frac{D}{2c_1 s/\sqrt n}<1-\frac{1}{\sqrt2}\cdot\frac{c_2}{c_1}\label{eq:01-overlap}
\end{equation}
すなわち重なりが1本の区間の幅に対してこの割合未満なら有意である．右辺は $n\to\infty$ で $1-1/\sqrt2\approx0.293$ に近づく．
\end{meidai}
\begin{proof}
$t$ 分布の上側点は自由度について減少するので $c_2<c_1$，よって $2c_1>\sqrt2 c_2$．
(1) $D<0$ なら $\bar x-\bar y>2c_1 s/\sqrt n>\sqrt2 c_2 s/\sqrt n$ で\cref{eq:01-ttest}が成り立つ．
(2) \cref{eq:01-ttest}は $D=2c_1s/\sqrt n-(\bar x-\bar y)<(2c_1-\sqrt2c_2)s/\sqrt n$ と同値であり，両辺を $2c_1s/\sqrt n$ で割ればよい．
\end{proof}
「区間が重なっているから有意差はない」という判断は誤りである．
差の検定の標準誤差は $\sqrt2\,s/\sqrt n$ であって，各群の標準誤差の和 $2s/\sqrt n$ ではないからである．

\section{仮定}\label{sec:01-assump}
すべての推測は仮定の上に立つ．記述式では「この検定に通常置かれる仮定を示せ」が頻出である
（\kakomon{2016}{1}，\kakomon{2019}{4}）．仮定は次の3層で整理すると漏れにくい．
\begin{enumerate}
\item \textbf{独立性}：観測単位どうしが独立（対応があるならペア単位で独立）．
\item \textbf{分布の仮定}：正規性，等分散，比例ハザード，対称性など．
\item \textbf{デザイン上の仮定}：ランダム化，無情報な打ち切り，欠測のメカニズム（MCAR/MAR），交絡の不在．
\end{enumerate}
漸近的手法ではさらに「標本サイズが十分大きい（各セルの期待度数が十分）」が加わる．

\section{結果の解釈}\label{sec:01-interp}
解釈の答案は次の3要素を含めると採点者に伝わりやすい．
\begin{enumerate}
\item \textbf{統計的結論}：有意か否か，信頼区間が帰無値を含むか．
\item \textbf{効果の方向と大きさ}：点推定値と信頼区間を医学の言葉で（例：「死亡のハザードが約30\%低い」）．
\item \textbf{限界}：観察研究なら交絡，小標本なら検出力，多重性，欠測，外的妥当性．
\end{enumerate}
\begin{center}\small
\begin{tabular}{@{}L{45mm}L{100mm}@{}}
\toprule
結果 & 解釈の例\\
\midrule
$\OR=2.0$（95\%CI 1.3--3.1） & 曝露群の発症オッズは非曝露群の約2倍で，区間は1を含まないので有意な関連がある．発症が稀ならリスク比もおよそ2と解釈できる．\\
$\HR=0.70$（95\%CI 0.55--0.89） & 試験群の死亡ハザードは対照群の0.70倍（30\%低下）．比例ハザードが成り立つなら，どの時点でもこの比が保たれる．\\
$\RD=0.05$（95\%CI $-0.02$--$0.12$） & 有効率の差の点推定は5ポイントだが区間が0を含み，差があるとは言えない．$-2$ポイントから12ポイントまでの差はデータと両立する．\\
$p=0.077$ & 5\%水準で有意ではない．「差がない」とは結論できない．\\
\bottomrule
\end{tabular}
\end{center}

\section{手法選択の基本}\label{sec:01-choice}
\begin{pointbox}[手法選択の質問リスト]
\begin{enumerate}
\item アウトカムは何型か（\cref{tab:01-outcome}）．
\item 比較する群は独立か，対応があるか（同一患者・マッチ・クロスオーバー）．
\item 群の数は2か，3以上か．群に順序（用量）があるか．
\item 調整すべき共変量・層はあるか（観察研究か，層別ランダム化か）．
\item 打ち切り・欠測・競合事象はあるか．
\item 標本サイズは大きいか（漸近近似でよいか，exactが必要か）．
\item 目的は優越性・同等性・非劣性・予測のどれか．
\end{enumerate}
\end{pointbox}
これらの答えから手法を選ぶ一覧表を\cref{app:A}に掲げた．

\section{よくある誤り}\label{sec:01-err}
\begin{warnbox}
\begin{itemize}
\item 「$p>0.05$ だから両群は同等」→ 誤り．同等性は同等性検定・信頼区間で示す（\cref{ch:13}）．
\item 「P 値が小さいほど効果が大きい」→ 誤り．P 値は効果の大きさと標本サイズの両方で決まる．
\item 群ごとの信頼区間が重なることを根拠に「有意差なし」とする → 誤り（\cref{prop:01-overlap}）．
\item 対応のあるデータに2標本 $t$ 検定を使う，またはその逆．
\item 比の指標の信頼区間を，比のスケールで $\hat\theta\pm1.96\SE$ として作る → 誤り．対数スケールで作って指数変換する．
\item データを見てから片側・両側を選ぶ，複数の検定から有意なものだけを報告する．
\end{itemize}
\end{warnbox}

\section{数理統計との接続}\label{sec:01-link}
\begin{linkbox}
\begin{itemize}
\item 本書の漸近的検定のほとんどは，Wald検定・スコア検定・尤度比検定の3つのいずれかである
      （『現代数理統計学の基礎』7.3節）．分母の分散を帰無仮説の下で推定するのがスコア型
      （例：Pearson $\chi^2$，ログランク，McNemar），推定値で評価するのがWald型（例：対数オッズ比の信頼区間）．
\item 比の指標の信頼区間は，対数変換とデルタ法（同書 定理5.20）で作る．
      多変量デルタ法は同書の範囲外なので\cref{sec:04-delta}で導出する．
\item 検定と信頼区間の双対性は，同書では明示されていないが，定義から直ちに導ける（\cref{prop:01-duality}）．
\item P 値の一様性 $p(\bm X)\sim\Unif(0,1)$（連続分布の場合）は確率積分変換による（同書 命題2.19，定理7.22）．
      多重性の議論（\cref{ch:14}）はこの性質を使う．
\end{itemize}
\end{linkbox}

\section{例題}\label{sec:01-ex}
\begin{reidai}[手法選択]\label{reidai:01-1}
次の各状況で，主解析に用いる手法を1つ挙げ，その理由を簡潔に述べよ．
\begin{enumerate}[label=(\arabic*)]
\item 高血圧患者30名に新薬を投与し，投与前と8週後の収縮期血圧を測定した．
\item 喫煙の有無で分けた2群を10年追跡し，肺がん発症までの時間を観察した．途中で転居による追跡不能がある．
\item 同一の患者150名に2種類の迅速検査を行い，それぞれ陽性・陰性を記録した．
\item 3つの病院で行った観察研究で，病院ごとに治療の有無と退院時の改善の有無を集計した．病院により患者背景が大きく異なる．
\item プラセボ，低用量，中用量，高用量の4群に無作為化し，有効・無効を判定した．用量とともに有効率が上がるかを示したい．
\end{enumerate}
\end{reidai}

\begin{reidai}[信頼区間の重なり]\label{reidai:01-2}
2群（各 $n=10$）で連続量を測定し，両群の標本標準偏差はともに $s=5.0$ であった．
$t_9(0.025)=2.262$，$t_{18}(0.025)=2.101$ を用いよ．
\begin{enumerate}[label=(\arabic*)]
\item 各群の平均の95\%信頼区間の半幅を求めよ．
\item 2標本 $t$ 検定（両側5\%）が有意になる平均の差の最小値を求めよ．
\item 平均の差が $6.0$ のとき，2つの信頼区間は重なるか．また検定は有意か．
\item 検定が有意になるための，重なりの長さの上限と，それが区間幅に占める割合を求めよ．
\end{enumerate}
\end{reidai}

\begin{reidai}[対応の利用]\label{reidai:01-3}
治療前後の値 $(X,Y)$ は
$\V[X]=\V[Y]=\sigma^2=100$，相関 $\rho=0.8$ をもつとする．
\begin{enumerate}[label=(\arabic*)]
\item 20名の前後差の平均 $\bar D$ の標準誤差を求めよ．
\item 同じ母集団から独立に20名ずつ2群を取り，平均の差をとった場合の標準誤差を求めよ．
\item (1)の標準誤差が(2)より小さくなる $\rho$ の条件を求めよ．
\end{enumerate}
\end{reidai}

\begin{reidai}[解釈]\label{reidai:01-4}
新薬とプラセボの無作為化試験で，6か月以内の再入院割合を比較した．
リスク比の点推定値は $0.80$，95\%信頼区間は $(0.62,1.03)$，両側 $p=0.08$ であった．
次の記述のうち正しいものをすべて選び，誤っているものは理由を述べよ．
\begin{enumerate}[label=(\alph*)]
\item 新薬は再入院を減らさないことが示された．
\item 有意水準5\%で，再入院割合に差があるとは言えない．
\item 新薬群の再入院リスクはプラセボ群より20\%低いと推定されたが，38\%の低下から3\%の増加までの値がデータと両立する．
\item 帰無仮説が正しい確率は8\%である．
\item 同じ効果の大きさで症例数を増やせば，有意になる可能性がある．
\end{enumerate}
\end{reidai}

\section{解答}\label{sec:01-sol}
\begin{kaitou}{reidai:01-1}
\begin{enumerate}[label=(\arabic*)]
\item \textbf{対応のある $t$ 検定}（前後差の1標本 $t$ 検定）．同一患者の前後値は対応があり，差を取ることで個人間変動を除ける．
      正規性が疑わしい，外れ値がある場合は Wilcoxon 符号付き順位検定（\cref{ch:03}）．
\item \textbf{Kaplan--Meier法とログランク検定}，調整が必要なら Cox 比例ハザードモデル（\cref{ch:09,ch:10}）．
      打ち切りがあるので発症割合の単純比較は偏る．観察研究なので交絡（年齢など）の調整を検討する．
\item \textbf{McNemar検定}（\cref{sec:04-mcnemar}）．同一患者に2検査を行った対応のある2 値データで，
      陽性率の差は不一致ペアだけから情報を得る．
\item \textbf{Mantel--Haenszel法}（層別解析，\cref{ch:07}）またはロジスティック回帰で病院を調整する．
      病院が交絡因子になりうるので，併合した表での比較は Simpson のパラドックスを起こしうる．
\item \textbf{Cochran--Armitage 傾向性検定}（\cref{sec:04-trend}）．群に順序があり，単調な用量反応を対立仮説とするので，
      順序を無視した $\chi^2$ 検定（自由度3）より検出力が高い．
\end{enumerate}
\end{kaitou}

\enlargethispage{1\baselineskip}
\begin{kaitou}{reidai:01-2}
(1) $t_9(0.025)\,s/\sqrt{10}=2.262\times5.0/3.162=3.577$．\\
(2) 差の標準誤差は $s\sqrt{2/10}=5.0\times0.4472=2.236$．有意になるのは差が $t_{18}(0.025)\times2.236=2.101\times2.236=4.698$ を超えるとき．\\
(3) 重なり $D=2\times3.577-6.0=1.154>0$ で区間は重なる．しかし $6.0>4.698$ なので検定は有意である．\\
(4) \cref{prop:01-overlap}より上限は $2\times3.577-4.698=2.456$．区間幅 $7.154$ に対する割合は
$2.456/7.154=0.343$．公式\cref{eq:01-overlap}でも $1-0.7071\times2.101/2.262=1-0.657=0.343$ となる．
自由度が小さいため漸近値 $0.293$ より大きい．
\end{kaitou}

\begin{kaitou}{reidai:01-3}
(1) $\V[D_i]=2\sigma^2(1-\rho)=2\times100\times0.2=40$ より $\SE=\sqrt{40/20}=\sqrt2=1.414$．\\
(2) $\SE=\sqrt{100/20+100/20}=\sqrt{10}=3.162$．\\
(3) (1)の分散 $2\sigma^2(1-\rho)/n$ と(2)の $2\sigma^2/n$ を比べ，$\rho>0$ のとき(1)が小さい．
ただし(1)は前後比較であり，対照群がなければ時間効果や平均への回帰（\cref{sec:08-rtm}）と治療効果を区別できない点に注意する．
\end{kaitou}

\begin{kaitou}{reidai:01-4}
正しいのは (b)，(c)，(e)．
(a) 誤り．有意でないことは効果がないことの証明ではない．区間は38\%の低下も含む．
(d) 誤り．P 値は帰無仮説の下でのデータの極端さの確率であり，帰無仮説の確率ではない．
(e) 標準誤差は症例数の平方根に反比例して小さくなるので，同じ点推定値なら区間が狭まり有意になりうる．
\end{kaitou}

\begin{summarybox}
\begin{itemize}
\item 設問を読んだら，デザイン・観測単位・アウトカムの型・目的・推定対象・確率モデルを書き出す．
\item 対応のあるデータは差（または不一致ペア）で分析する．$\V[\bar D]=(\sigma_X^2+\sigma_Y^2-2\rho\sigma_X\sigma_Y)/n$．
\item P 値は効果の大きさではない．有意でないことは「差がない」ことを意味しない．
\item 信頼区間と検定は双対：比の指標は1，差の指標は0を含むかで判定．比の指標は対数スケールで区間を作る．
\item 群ごとの95\%信頼区間が重なっていても，重なりが幅の約29\%（小標本ではもう少し大きい）未満なら2標本検定は有意．
\item 解釈の答案は「統計的結論＋効果の方向と大きさ＋限界」の3要素で書く．
\end{itemize}
\end{summarybox}

% ===== ch02.tex =====
\chapter{研究デザインと効果指標}\label{ch:02}

\begin{goalbox}
\begin{itemize}
\item ランダム化比較試験（並行群間・クロスオーバー），コホート研究，症例対照研究，マッチング研究，縦断研究，診断精度研究，メタアナリシスの特徴と，各デザインで推定できる効果指標を説明できる．
\item リスク，率，オッズを区別し，リスク差・リスク比・オッズ比・ハザード比・率比を定義できる．
\item オッズ比がリスク比を近似する条件と，その方向（OR は RR より1から遠い）を導出できる．
\item 症例対照研究で曝露オッズ比が疾病オッズ比に一致することを示せる．
\item オッズ比の非可縮性（交絡がなくても層別と併合で値が変わる）を説明できる．
\item ICH E9(R1) の estimand の5要素を挙げられる．
\end{itemize}
\end{goalbox}

\begin{exambox}
\begin{itemize}
\item \kakomon{2017}{4}：症例対照研究（後ろ向き研究）の層別データで\textbf{曝露オッズ比}を推定．
\item \kakomon{2019}{2}：コホート研究の粗リスク差と調整リスク差．
\item \kakomon{2022}{2}：マッチしたペアの死亡データで粗リスク比・オッズ比とMH推定量を比較し，「交絡の有無」と「推定されるパラメータの違い」から説明（オッズ比の非可縮性）．
\item \kakomon{2017}{5}（共通）：単純無作為化と偏りを抑える割付法の群サイズの分布．
\item HRの解釈は\kakomon{2022}{1}，RRのメタアナリシスは\kakomon{2021}{4}．
\end{itemize}
\end{exambox}

\section{問題設定}\label{sec:02-setting}
2群（曝露あり／なし，または試験治療／対照）で2 値アウトカム（発症・有効など）を比較する基本の $2\times2$ 表を
\cref{tab:02-2x2}とする．どの周辺度数が研究者によって固定されるかは研究デザインで決まり，
それによって推定可能な指標が変わる．これが本章の中心的な論点である．
\begin{table}[htbp]
\centering\small
\caption{基本の $2\times2$ 表（度数）}\label{tab:02-2x2}
\begin{tabular}{@{}lccc@{}}
\toprule
 & 発症あり $D$ & 発症なし $\bar D$ & 計\\
\midrule
曝露あり $E$ & $a$ & $b$ & $n_1=a+b$\\
曝露なし $\bar E$ & $c$ & $d$ & $n_0=c+d$\\
\midrule
計 & $m_1=a+c$ & $m_0=b+d$ & $N$\\
\bottomrule
\end{tabular}
\end{table}

\section{研究デザイン}\label{sec:02-design}
\subsection{ランダム化比較試験}\label{subsec:02-rct}
\begin{description}
\item[並行群間比較試験（parallel design）] 各被験者を1つの治療群に無作為に割り付けて並行に比較する．最も標準的．
\item[クロスオーバー試験（crossover design）] 各被験者が複数の治療を時期を変えて受ける．2剤2期（AB/BA）が基本．
      被験者内比較なので個人間変動が除かれ効率が高い．前の治療の効果が残る\emphx{持ち越し効果}（carryover）を防ぐため，
      時期の間に\emphx{休薬期間}（washout）を置く．慢性で安定した疾患や生物学的同等性試験（\cref{ch:13}）に向き，
      死亡や治癒のような不可逆なアウトカムには使えない．
\item[盲検化（blinding）] 被験者・医師・評価者が割付を知らないようにする．二重盲検では被験者と医師の両方が知らない．
      評価のバイアスとプラセボ効果の混入を防ぐ．
\item[プラセボ（placebo）] 有効成分を含まない偽薬．プラセボ対照により，自然経過や心理的効果を除いた薬効を推定する．
\end{description}

\paragraph{割付の方法}
\begin{itemize}
\item \textbf{単純無作為化}：各被験者を独立に確率 $1/2$ で割り付ける．群サイズ $X\sim\Bin(n,1/2)$，$\V[X]=n/4$ で，小標本では不均衡が起こりうる．
\item \textbf{置換ブロック法}：一定数（例：4人）ごとに各群同数になるよう割り付ける．
\item \textbf{層別無作為化}：重要な予後因子の層ごとにブロック法で割り付ける．
\item \textbf{偏りを抑える逐次割付}（最小化法，偏コイン法など）：それまでの不均衡を減らす方向の群に高い確率で割り付ける．
      \kakomon{2017}{5} の方法2は，$A$ 群 $a$ 人，$B$ 群 $b$ 人のとき次を $A$ に確率 $b/(a+b)$ で割り付ける方式で，
      群サイズの期待値は単純無作為化と同じ $n/2$ のまま分散が小さくなる．
\end{itemize}

\subsection{観察研究}\label{subsec:02-obs}
\begin{description}
\item[コホート研究] 曝露状態で分けた集団を追跡して発症を観察する（\cref{tab:02-2x2}で行和 $n_1,n_0$ を研究者が決める）．
      リスク，率，RD，RR，OR，HR をすべて推定できる．前向き（prospective）と，記録を遡る後ろ向きコホートがある．
\item[症例対照研究（case-control study）] 発症者（症例）と非発症者（対照）を集め，過去の曝露を調べる（列和 $m_1,m_0$ を研究者が決める）．
      稀な疾患に効率的．\textbf{リスクもRRも直接は推定できず，推定できるのはオッズ比}である（\cref{prop:02-orsym}）．
      「後ろ向き研究」とも呼ばれる（\kakomon{2017}{4}）．
\item[マッチング研究] 背景因子の似た症例と対照（あるいは曝露者と非曝露者）を組にする．
      ペアを層とみなす層別解析（McNemar型，条件付きオッズ比）で分析する（\cref{ch:04,ch:07}）．
\item[横断研究] ある時点で曝露と疾患を同時に調べる．有病割合（prevalence）の比較になる．
\end{description}

\subsection{その他のデザイン}
\begin{description}
\item[縦断研究（longitudinal study）] 同一個体を複数時点で測定する．個体内相関を考慮した解析が必要（\cref{ch:08}）．
\item[診断精度研究] 新しい検査と確定診断（ゴールドスタンダード）を同じ被験者に行い，感度・特異度などを推定する（\cref{ch:05}）．
      同一被験者に2検査を行えば対応のある比較になる（\kakomon{2019}{3}，\kakomon{2024}{1}）．
\item[メタアナリシス] 複数の研究の効果推定値を統合する（\cref{ch:15}）．
\end{description}

\begin{table}[htbp]
\centering\small
\caption{デザインと推定可能な指標}\label{tab:02-design}
\begin{tabular}{@{}lccccc@{}}
\toprule
デザイン & リスク & RD & RR & OR & HR・率比\\
\midrule
RCT・コホート（追跡） & ○ & ○ & ○ & ○ & ○\\
横断研究 & 有病割合 & 有病割合の差 & 有病割合の比 & ○ & ×\\
症例対照研究 & × & × & ×（稀なら OR で近似） & ○ & ×\\
\bottomrule
\end{tabular}
\end{table}

\section{効果指標}\label{sec:02-measure}
\subsection{リスク・率・オッズ}\label{subsec:02-risk}
\begin{teigi}[リスク・オッズ・率]\label{def:02-risk}
\begin{itemize}
\item \textbf{リスク}（risk，累積発生割合）：一定期間内に発症する確率 $\pi$．無次元，$0\le\pi\le1$．
\item \textbf{オッズ}（odds）：$\pi/(1-\pi)$．$0$ から $\infty$．
\item \textbf{率}（rate，発生率）：単位人時間あたりの発生数 $\lambda=(\text{発生数})/(\text{観察人時間})$．次元は（時間）$^{-1}$．
\end{itemize}
\end{teigi}
率は，人によって観察期間が異なる場合（途中参加・打ち切り）の標準的な要約である．
一定のハザード $\lambda$ をもつ指数モデルでは，期間 $t$ のリスクは $1-e^{-\lambda t}$ で，$\lambda t$ が小さいとき $\approx\lambda t$ となる．

\subsection{比較の指標}\label{subsec:02-effect}
\begin{teigi}[効果指標]\label{def:02-effect}
曝露群・非曝露群のリスクを $\pi_1,\pi_0$ とする．
\begin{align*}
\RD&=\pi_1-\pi_0\quad\text{（リスク差，risk difference）}\\
\RR&=\pi_1/\pi_0\quad\text{（リスク比，risk ratio；相対リスク relative risk）}\\
\OR&=\frac{\pi_1/(1-\pi_1)}{\pi_0/(1-\pi_0)}=\frac{\pi_1(1-\pi_0)}{\pi_0(1-\pi_1)}\quad\text{（オッズ比，odds ratio）}
\end{align*}
率 $\lambda_1,\lambda_0$ に対しては率比 $\lambda_1/\lambda_0$，率差 $\lambda_1-\lambda_0$．
ハザード関数 $h_1(t),h_0(t)$（\cref{ch:09}）に対し $h_1(t)/h_0(t)$ が $t$ によらず一定のとき，その値をハザード比 $\HR$ という．
\end{teigi}
\cref{tab:02-2x2}の標本版は $\hat\pi_1=a/n_1$，$\hat\pi_0=c/n_0$，$\widehat{\OR}=ad/(bc)$ である．

\paragraph{使い分け}
\begin{itemize}
\item \textbf{RD} は絶対的な効果で，公衆衛生・臨床的なインパクトを表す．治療が発症を減らす場合，絶対リスク減少 $\mathrm{ARR}=\pi_0-\pi_1>0$ の逆数
      $\mathrm{NNT}=1/\mathrm{ARR}$ を\emphx{NNT}（number needed to treat）といい，1人の発症を防ぐために治療が必要な人数を表す．
      治療（曝露）が有害事象を増やす場合は，絶対リスク増加 $\mathrm{ARI}=\pi_1-\pi_0>0$ の逆数 $\mathrm{NNH}=1/\mathrm{ARI}$（number needed to harm）を用い，NNT と区別する．
\item \textbf{RR} は相対的な効果で直観的に解釈しやすい．
\item \textbf{OR} は症例対照研究でも推定でき，ロジスティック回帰の係数と直結する（\cref{ch:06}）．
\item \textbf{HR} は生存時間データの標準的な指標（\cref{ch:10}）．瞬間的な発生率の比であり，一定期間のリスク比とは異なる．
\end{itemize}

\subsection{RR・OR・HRの関係}\label{subsec:02-relation}
\begin{meidai}[ORとRRの関係]\label[meidai]{prop:02-orrr}
\begin{equation}
\OR=\RR\times\frac{1-\pi_0}{1-\pi_1}.\label{eq:02-orrr}
\end{equation}
したがって (i) $\RR>1\iff\OR>1$（方向は一致），(ii) $\RR>1$ のとき $\OR>\RR$，$\RR<1$ のとき $\OR<\RR$
（OR は RR より1から遠い），(iii) $\pi_0,\pi_1$ がともに小さい（稀な疾患）とき $\OR\approx\RR$．
\end{meidai}
\begin{proof}
\cref{eq:02-orrr}は定義から直ちに従う．$\RR>1$ なら $\pi_1>\pi_0$ で $(1-\pi_0)/(1-\pi_1)>1$ となり $\OR>\RR>1$．$\RR<1$ も同様．
\end{proof}

\begin{meidai}[ORの対称性：症例対照研究でORが推定できる理由]\label[meidai]{prop:02-orsym}
曝露 $E$ と疾患 $D$ について
\[
\frac{\Prob(D\mid E)/\Prob(\bar D\mid E)}{\Prob(D\mid\bar E)/\Prob(\bar D\mid\bar E)}
=\frac{\Prob(E\mid D)/\Prob(\bar E\mid D)}{\Prob(E\mid\bar D)/\Prob(\bar E\mid\bar D)}.
\]
すなわち疾病オッズ比（コホートで定義されるOR）は曝露オッズ比（症例対照研究で推定できるOR）に等しい．
\end{meidai}
\begin{proof}
両辺とも $\dfrac{\Prob(D,E)\Prob(\bar D,\bar E)}{\Prob(\bar D,E)\Prob(D,\bar E)}$ に等しい．
症例対照研究では $\Prob(E\mid D)$ と $\Prob(E\mid\bar D)$ は推定できる（症例・対照は各疾患状態からの標本）が，
症例・対照の比率を研究者が決めるので $\Prob(D\mid E)$ は推定できない．
\end{proof}
\kakomon{2017}{4} の「曝露オッズ比」は右辺であり，標本では同じく $ad/(bc)$ で計算される．

\paragraph{HRとリスク比}
指数分布で $h_1(t)=\lambda_1$，$h_0(t)=\lambda_0$ のとき $\HR=\lambda_1/\lambda_0$ であるが，
時点 $t$ までのリスク比は
\[
\frac{1-e^{-\lambda_1t}}{1-e^{-\lambda_0t}}
\]
で，$t\to0$ のとき $\HR$ に近づき，$t\to\infty$ のとき1に近づく．
「HR $=0.7$」は「どの時点でも，まだイベントを起こしていない人の瞬間的な発生率が0.7倍」という意味であって，
「最終的にイベントを起こす人の割合が0.7倍」ではない．

\subsection{可縮性と非可縮性}\label{subsec:02-collapse}
\Youchui 層別変数 $Z$ があり，各層の効果指標が共通の値 $\theta$ をとるとする．
$Z$ が曝露と独立（例：ランダム化試験や，\kakomon{2022}{2} のように各層に曝露者・非曝露者が1人ずつのマッチング）であっても，
併合した表で計算した指標が $\theta$ に一致するとは限らない．一致する指標を\emphx{可縮}（collapsible）という．
\begin{itemize}
\item RD と RR は可縮：$Z$ と曝露が独立なら，併合リスク $\pi_x=\sum_k\Prob(Z=k)\pi_{xk}$ は層別リスクの同じ重みの平均だから，
      層共通の RD（または RR）がそのまま併合でも成り立つ．
\item OR は\textbf{非可縮}：層共通の OR があっても，併合 OR は一般に1に近い側へずれる（\cref{reidai:02-3}）．
\end{itemize}
したがって，交絡がなくても「粗OR $\ne$ 調整OR」は起こる．
\kakomon{2022}{2} では粗RR $=$ MH-RR $=2.78$ であったが，粗OR $=2.79$ に対し MH-OR $=3.98$ であった．
これは交絡バイアスではなく，推定しているパラメータ（周辺OR と層内の条件付きOR）が異なるためである．
交絡による粗・調整の差（\cref{ch:07}）と非可縮性による差を区別することが重要である．

\section{Estimand（推定対象）}\label{sec:02-estimand}
\begin{supbox}[補足（出典：ICH E9(R1) \cite{ich:e9r1}）]
ICH E9(R1)（2019年）は，臨床試験で推定すべき量（estimand）を次の5要素で規定することを求めている．
\begin{enumerate}
\item \textbf{対象集団}（population）：臨床的疑問が向けられる患者集団．
\item \textbf{治療}（treatment）：比較する治療条件．
\item \textbf{変数}（variable, endpoint）：各患者について得る評価項目．
\item \textbf{中間事象}（intercurrent events）の扱い：治療開始後に起こり，評価項目の解釈や存在に影響する事象
      （治療中止，救済治療の使用，死亡など）への対処戦略．
      治療方針戦略（treatment policy），仮想戦略（hypothetical），複合変数戦略（composite），
      治療下戦略（while on treatment），主要層戦略（principal stratum）の5つがある．
\item \textbf{集団レベルの要約}（population-level summary）：RD，RR，HR，平均差など．
\end{enumerate}
欠測データの扱い（\cref{ch:16}）や，死亡を競合事象とみなすか（\cref{ch:11}）は，この枠組みの中で estimand として先に決めるべき事柄とされる．
\end{supbox}

\section{推定と検定の概要}\label{sec:02-inference}
2つの独立な二項標本 $a\sim\Bin(n_1,\pi_1)$，$c\sim\Bin(n_0,\pi_0)$ に基づく標準的な推定を示す（導出は\cref{ch:04}）．
\begin{align}
\widehat{\V}[\widehat{\RD}]&=\frac{\hat\pi_1(1-\hat\pi_1)}{n_1}+\frac{\hat\pi_0(1-\hat\pi_0)}{n_0},\label{eq:02-varRD}\\
\widehat{\V}[\log\widehat{\RR}]&=\frac{1-\hat\pi_1}{a}+\frac{1-\hat\pi_0}{c}=\frac1a-\frac1{n_1}+\frac1c-\frac1{n_0},\label{eq:02-varRR}\\
\widehat{\V}[\log\widehat{\OR}]&=\frac1a+\frac1b+\frac1c+\frac1d.\label{eq:02-varOR}
\end{align}
比の指標は対数スケールで $\log\hat\theta\pm z_{\alpha/2}\SE$ を作り指数変換する．
帰無仮説 $\pi_1=\pi_0$ の検定は，どの指標を使っても Pearson $\chi^2$ 検定などで共通である．

\section{仮定}\label{sec:02-assump}
\begin{itemize}
\item 二項モデル：各群内で個体が独立で，発症確率が一様．クラスタ（家族・施設）内相関があれば分散は過小評価される．
\item 率の比較：発生率が観察期間中一定（ポアソン過程）であること．
\item 症例対照研究：対照が症例の発生源集団から曝露と無関係に選ばれていること（選択バイアスがないこと）．
\item 因果的解釈：交絡がない（ランダム化されている），または十分に調整されていること．
\end{itemize}

\section{結果の解釈}\label{sec:02-interp}
\begin{itemize}
\item $\RR=2$：曝露群の発症リスクは非曝露群の2倍．「2倍増える」は「100\%増加」の意味．
\item $\OR=2$：曝露群の発症オッズは2倍．リスクが小さいときのみ「リスクがおよそ2倍」と言い換えてよい．
      リスクが大きいとき（例：$\pi_0=0.4$）にORをRRのように読むと効果を過大に述べることになる．
\item $\RD=\pi_1-\pi_0=-0.05$（$\mathrm{ARR}=0.05$），$\mathrm{NNT}=20$：20人を治療すると1人の発症を防げる．逆に $\RD=+0.05$ が有害事象の差なら $\mathrm{NNH}=20$（20人に1人余分に害が生じる）．
\item $\HR=0.7$：追跡期間中どの時点でも，死亡の瞬間的な発生率が30\%低い（比例ハザードの仮定の下で）．
\end{itemize}

\section{手法選択}\label{sec:02-choice}
\begin{itemize}
\item 症例対照研究 → オッズ比（層別なら MH-OR，共変量調整ならロジスティック回帰）．
\item コホート・RCTでリスク（一定期間の割合）→ RD・RR（または OR）．
\item 観察期間が不揃い → 率（ポアソン），または生存時間解析（HR）．
\item マッチング → マッチを層とした解析（McNemar，条件付きOR，MH）．マッチを無視すると効率を失い，ORでは推定対象も変わる．
\end{itemize}

\section{よくある誤り}\label{sec:02-err}
\begin{warnbox}
\begin{itemize}
\item 症例対照研究でリスクやリスク比を計算する．
\item 発症割合が大きいのにオッズ比を「リスクが〇倍」と解釈する．
\item 「粗ORと調整ORが違う」ことを常に交絡のせいにする（非可縮性を忘れる）．
\item ハザード比を「一定期間のリスク比」と同一視する．
\item 率（人時間あたり）とリスク（割合）を混同する．
\end{itemize}
\end{warnbox}

\section{数理統計との接続}\label{sec:02-link}
\begin{linkbox}
\begin{itemize}
\item オッズの対数 $\log\{\pi/(1-\pi)\}$ は二項分布の自然母数である（『現代数理統計学の基礎』6.1.2項，例6.9）．
      ロジスティック回帰の係数がオッズ比の対数になるのはこのためである（\cref{ch:06}）．
\item $\log\hat\pi$ の漸近分散 $(1-\pi)/(n\pi)$ は同書第5章演習（2項比率の対数変換のデルタ法）の結果で，
      \cref{eq:02-varRR}はこれを2群に適用したものである．
\item 症例対照研究の尤度は曝露の条件付き分布 $\Prob(E\mid D)$ に基づくが，
      ロジスティックモデルの傾き（log OR）は前向きでも後ろ向きでも同じ値になる（Prentice--Pyke の結果；\cref{ch:06}）．
\end{itemize}
\end{linkbox}

\section{例題}\label{sec:02-ex}
\begin{reidai}[効果指標の計算]\label{reidai:02-1}
ある前向きコホート研究で，曝露群200人中40人，非曝露群400人中40人が5年以内に発症した．
\begin{enumerate}[label=(\arabic*)]
\item リスク差，リスク比，オッズ比を求めよ．
\item \cref{eq:02-orrr}が成り立つことを数値で確かめよ．
\item 曝露は有害な方向に働いている．曝露によって1人余分に発症する人数（NNH）を求めよ．
\item 対数リスク比の標準誤差と，リスク比の95\%信頼区間を求めよ．
      ただし $\log2=0.6931$，$\sqrt{0.0425}=0.2062$，$e^{0.2889}=1.335$，$e^{1.0973}=2.996$ を用いてよい．
\end{enumerate}
\end{reidai}

\begin{reidai}[症例対照研究]\label{reidai:02-2}
ある稀な疾患の症例100人と，同じ地域から選んだ対照200人について，過去の職業性曝露を調べたところ，
症例の30人，対照の20人が曝露ありであった．
\begin{enumerate}[label=(\arabic*)]
\item 曝露オッズ比を求めよ．
\item この値を「曝露者の発症リスクは非曝露者の何倍か」の推定値として用いてよい理由と，その条件を述べよ．
\item 「症例の30\%が曝露者なので，曝露者の発症リスクは30\%」という主張の誤りを指摘せよ．
\end{enumerate}
\end{reidai}

\begin{reidai}[オッズ比の非可縮性]\label{reidai:02-3}
ランダム化試験で，重症度 $Z$（軽症・重症が各50\%，治療割付とは独立）により層別すると，
改善確率は次のとおりであった：軽症層で治療 $0.8$，対照 $0.5$；重症層で治療 $0.5$，対照 $0.2$．
\begin{enumerate}[label=(\arabic*)]
\item 各層のオッズ比，リスク比，リスク差を求めよ．
\item 層を併合したときの改善確率，オッズ比，リスク差を求めよ．
\item この試験に交絡はあるか．(1)と(2)の違いを説明せよ．
\end{enumerate}
\end{reidai}

\begin{reidai}[ハザード比とリスク比]\label{reidai:02-4}
試験群・対照群の死亡までの時間がそれぞれ率 $\lambda_1=0.1$，$\lambda_0=0.2$（/年）の指数分布に従うとする．
$e^{-0.5}=0.6065$，$e^{-1}=0.3679$ を用いよ．
\begin{enumerate}[label=(\arabic*)]
\item ハザード比を求めよ．
\item 5年死亡割合の比（リスク比）を求め，ハザード比と比べよ．
\end{enumerate}
\end{reidai}

\section{解答}\label{sec:02-sol}
\begin{kaitou}{reidai:02-1}
(1) $\hat\pi_1=40/200=0.20$，$\hat\pi_0=40/400=0.10$．$\RD=0.10$，$\RR=2.0$，
$\OR=\dfrac{40\times360}{160\times40}=\dfrac{14400}{6400}=2.25$．\\
(2) $\RR\times(1-\pi_0)/(1-\pi_1)=2.0\times0.9/0.8=2.25$ で一致する．$\RR>1$ なので $\OR>\RR$．\\
(3) 絶対リスク増加 $\mathrm{ARI}=0.20-0.10=0.10$ より $\mathrm{NNH}=1/0.10=10$ 人（曝露者10人につき1人余分に発症する）．\\
(4) \cref{eq:02-varRR}より $\widehat\V=\frac1{40}-\frac1{200}+\frac1{40}-\frac1{400}=0.025-0.005+0.025-0.0025=0.0425$，$\SE=0.2062$．
$\log2=0.6931$ として $0.6931\pm1.96\times0.2062=0.6931\pm0.4042=(0.2889,\,1.0973)$．
指数変換して $(1.335,\,2.996)$，すなわちおよそ $(1.34,\,3.00)$．
\end{kaitou}

\begin{kaitou}{reidai:02-2}
(1) 症例：曝露30，非曝露70．対照：曝露20，非曝露180．$\OR=\dfrac{30\times180}{70\times20}=\dfrac{5400}{1400}=3.86$．\\
(2) \cref{prop:02-orsym}より曝露オッズ比は疾病オッズ比に等しく，\cref{prop:02-orrr}より疾患が稀（両群とも発症リスクが小さい）ならリスク比に近い．
対照が発生源集団から曝露と無関係に選ばれていることも必要である．\\
(3) 30\%は $\Prob(E\mid D)$ であり，$\Prob(D\mid E)$ ではない．症例対照研究では症例と対照の人数比を研究者が決めるので，発症リスク自体は推定できない．
\end{kaitou}

\begin{kaitou}{reidai:02-3}
(1) 軽症層：オッズは $0.8/0.2=4$，$0.5/0.5=1$ で $\OR=4$，$\RR=1.6$，$\RD=0.3$．
重症層：オッズは $0.5/0.5=1$，$0.2/0.8=0.25$ で $\OR=4$，$\RR=2.5$，$\RD=0.3$．\\
(2) 治療 $(0.8+0.5)/2=0.65$，対照 $(0.5+0.2)/2=0.35$．
$\OR=\dfrac{0.65/0.35}{0.35/0.65}=\Bigl(\dfrac{0.65}{0.35}\Bigr)^2=1.857^2=3.45$，$\RD=0.30$．\\
(3) $Z$ は割付と独立なので交絡はない．実際，層共通のRD（0.3）は併合しても0.3のまま（RDは可縮）．
一方，ORは層共通で4なのに併合では3.45となる．これはORの非可縮性による違いで，
層内の（条件付き）ORと集団全体の（周辺）ORは異なるパラメータだからである．
\end{kaitou}

\begin{kaitou}{reidai:02-4}
(1) $\HR=0.1/0.2=0.5$．\\
(2) $1-e^{-0.1\times5}=1-0.6065=0.3935$，$1-e^{-0.2\times5}=1-0.3679=0.6321$．
リスク比 $=0.3935/0.6321=0.623$．5年リスク比はハザード比0.5より1に近い．
イベントが累積すると両群のリスクはいずれも1に近づくので，長期のリスク比は1に近づく．
\end{kaitou}

\begin{summarybox}
\begin{itemize}
\item リスク（割合），オッズ，率（人時間あたり）を区別する．
\item $\OR=\RR\,(1-\pi_0)/(1-\pi_1)$：方向は一致し，OR は RR より1から遠い．稀な疾患では $\OR\approx\RR$．
\item 症例対照研究では曝露オッズ比＝疾病オッズ比のみ推定可能．RR・リスクは推定できない．
\item RD・RR は可縮，OR は非可縮．交絡がなくても粗ORと層別ORは異なりうる（\kakomon{2022}{2}）．
\item HR は瞬間的な発生率の比で，一定期間のリスク比とは異なる．
\item クロスオーバーは被験者内比較で効率的だが，持ち越し効果と休薬期間に注意．
\item estimand は対象集団・治療・変数・中間事象の扱い・要約指標の5要素で定める．
\end{itemize}
\end{summarybox}

% ===== ch03.tex =====
\chapter{連続アウトカムとノンパラメトリック法}\label{ch:03}

\begin{goalbox}
\begin{itemize}
\item 対応のある $t$ 検定と2標本 $t$ 検定を使い分け，それぞれの帰無仮説と仮定を述べられる．
\item 符号検定・Wilcoxon符号付き順位検定・Wilcoxon順位和検定（Mann--Whitney検定）の帰無仮説，仮定，検定統計量を説明できる．
\item 帰無仮説の下での順位統計量の期待値・分散を導出し，exactなP 値と正規近似のP 値を計算できる．
\item 小標本で順位検定が有意になりえない（離散性）ことを具体的に示せる．
\item 位置ずれモデルの下での順位和検定の検出力を $\Prob(X>Y)$ を用いて表せる．
\end{itemize}
\end{goalbox}

\begin{exambox}
\begin{itemize}
\item \kakomon{2016}{1}：治療前後の差（$n=7$）に対応のある $t$ 検定・符号検定・符号付き順位検定を適用．
      $\E[T\mid H_0],\V[T\mid H_0]$，符号検定のexact P 値，$\E[W^+\mid H_0],\V[W^+\mid H_0]$ の導出と正規近似（外れ値の影響は公式解答の講評で言及）．
\item \kakomon{2019}{4}：3対3の2群比較で Student $t$ 検定（仮定3つ），Wilcoxon順位和検定のexact分布（20通り），
      両側P 値，第1種の過誤確率が0になる問題点．
\item \kakomon{2025}{1}：一様分布の位置ずれモデルで $\E[R_X\mid H_0]$，$\Prob(X>Y)$，$R_X=\sum\sum I(X_i>Y_j)+n(n+1)/2$ を用いた $\E[R_X\mid H_1]$，検出力の式．
\item \kakomon{2016}{5}（共通）：2標本 $t$ 検定と信頼区間（\cref{sec:01-overlap}）．
\end{itemize}
\end{exambox}

\flowline{連続値\\対応あり／独立}{平均差・中央値\\$\Prob(X>Y)$}{$t$ 検定\\順位検定}{正規性・対称性\\独立性}{exact／正規近似}{外れ値の影響\\離散性}

\section{問題設定}\label{sec:03-setting}
\begin{itemize}
\item \textbf{1標本（対応あり）}：差 $Z_i=Y_i-X_i$ $(i=1,\dots,n)$ が独立に同一分布に従う．位置（平均 $\mu$ または中央値 $\theta$）が0かを調べる．
\item \textbf{2標本（独立）}：$X_1,\dots,X_m\iid F_X$，$Y_1,\dots,Y_n\iid F_Y$ で，両者が独立．分布の違い（特に位置のずれ）を調べる．
\end{itemize}
$t$ 検定は平均に関する正規モデルの推測であり，順位検定は分布の形をほとんど仮定しない．

\section{\texorpdfstring{$t$}{t} 検定}\label{sec:03-t}
\subsection{対応のある \texorpdfstring{$t$}{t} 検定}\label{subsec:03-pairedt}
$Z_i\iid\Normal(\mu,\sigma^2)$ とし $H_0:\mu=0$ を検定する．
\begin{equation}
t=\frac{\bar Z}{s_Z/\sqrt n}\sim t_{n-1}\quad(H_0\text{ の下})\label{eq:03-pairedt}
\end{equation}
\textbf{仮定}：(i) 差 $Z_i$ が互いに独立，(ii) 差の母集団分布が正規分布．
$n$ がある程度大きく外れ値がなければ，中心極限定理により正規性からのずれには頑健である．

\kakomon{2016}{1} では $\bar z=1.73$，$s=2.01$，$n=7$ で $t=1.73/(2.01/\sqrt7)=2.28<t_6(0.025)=2.447$ となり有意でない．
一方，同じデータの符号付き順位検定（正規近似）は $z=2.197$ で有意となった．
1人の大きな変化量（$5.6$）が $s$ を膨らませ，平均に基づく $t$ 検定がその影響を強く受けたためである．

\subsection{2標本 \texorpdfstring{$t$}{t} 検定}\label{subsec:03-twot}
$X_i\iid\Normal(\mu_1,\sigma^2)$，$Y_j\iid\Normal(\mu_2,\sigma^2)$ で $H_0:\mu_1=\mu_2$ を検定する．
\begin{equation}
t=\frac{\bar X-\bar Y}{s_p\sqrt{\frac1m+\frac1n}},\qquad s_p^2=\frac{(m-1)s_X^2+(n-1)s_Y^2}{m+n-2},\qquad t\sim t_{m+n-2}.\label{eq:03-twot}
\end{equation}
\textbf{仮定}（\kakomon{2019}{4}〔1〕）：(i) すべての測定値が互いに独立，(ii) 各群の母集団分布が正規分布，(iii) 2群の母分散が等しい．
差の $100(1-\alpha)\%$ 信頼区間は $\bar X-\bar Y\pm t_{m+n-2}(\alpha/2)\,s_p\sqrt{1/m+1/n}$．

\begin{supbox}[補足：Welch の検定]
等分散を仮定しない場合は $t_W=(\bar X-\bar Y)/\sqrt{s_X^2/m+s_Y^2/n}$ を，
自由度 $\nu=\dfrac{(s_X^2/m+s_Y^2/n)^2}{\frac{(s_X^2/m)^2}{m-1}+\frac{(s_Y^2/n)^2}{n-1}}$ の $t$ 分布で近似する（Welch--Satterthwaite）．
群の大きさが等しければ等分散からのずれの影響は小さいが，大きさも分散も異なるとStudentの検定の有意水準は大きく狂う．
\end{supbox}

\section{符号検定}\label{sec:03-sign}
\Hissu 差 $Z_i$ の分布が連続で中央値を $\theta$ とする．$H_0:\theta=0$ の下で $\Prob(Z_i>0)=1/2$ である．
\begin{teigi}[符号検定]\label{def:03-sign}
$T=\#\{i:Z_i>0\}$（差が正の個数）を検定統計量とする検定を符号検定という．$H_0$ の下で $T\sim\Bin(n,1/2)$ で
\begin{equation}
\E[T\mid H_0]=\frac n2,\qquad \V[T\mid H_0]=\frac n4.\label{eq:03-signmoments}
\end{equation}
\end{teigi}
\textbf{仮定}：観測値が互いに独立であること（と分布の連続性）だけで，分布の形（対称性・正規性）は仮定しない．
差が0の観測は除いて $n$ を数え直すのが通例である．

\paragraph{exactなP 値}
実現値 $t^*$ に対し，片側（$H_1:\theta>0$）P 値は $\sum_{i=t^*}^{n}\binom ni2^{-n}$，
\Youchui 離散分布の両側P 値には複数の定義がある．本書では $p=\min\{1,\ 2\min(P_{\mathrm{lower}},P_{\mathrm{upper}})\}$（小さい方の片側P 値の2倍，1で打ち切り）を用いる．ここで $P_{\mathrm{upper}}=\Prob(T\ge t^*\mid H_0)$，$P_{\mathrm{lower}}=\Prob(T\le t^*\mid H_0)$．このほか「観測値以下の確率をもつ値の確率の和」とする定義もあり，帰無分布が非対称なら両者は一致しない．符号検定の $\Bin(n,1/2)$ は対称なので，どちらの定義でも両側P 値は片側P 値の2倍になる．\kakomon{2016}{1} では $n=7$，$t^*=6$ で
$(\binom76+\binom77)/2^7=8/128=0.0625$，両側 $0.125$．
\paragraph{正規近似}
$z=(T-n/2)/\sqrt{n/4}$．連続修正をするなら $|T-n/2|-1/2$ を分子に用いる．

\section{Wilcoxon符号付き順位検定}\label{sec:03-signrank}
\Hissu\Hinshutsu
\begin{teigi}[符号付き順位統計量]\label{def:03-wplus}
$|Z_1|,\dots,|Z_n|$ に小さい順に順位 $1,\dots,n$ を付け，$Z_i>0$ である観測の順位の和を $W^+$ とする．
$W^-=n(n+1)/2-W^+$ は負の差の順位和である．
\end{teigi}
\textbf{帰無仮説と仮定}：$Z_i$ が互いに独立で，その分布が中央値 $\theta$ に関して\textbf{対称}な連続分布であるとし，$H_0:\theta=0$ を検定する．
（対称性を仮定しているので，平均が存在すれば $\theta$ は平均でもある．）

\begin{teiri}[$W^+$ の帰無分布の平均と分散]\label[teiri]{thm:03-wplus}
$H_0$ の下で
\begin{equation}
\E[W^+]=\frac{n(n+1)}{4},\qquad \V[W^+]=\frac{n(n+1)(2n+1)}{24}.\label{eq:03-wplus}
\end{equation}
\end{teiri}
\begin{proof}
$Z$ の分布が0に関して対称で連続なら，$\mathrm{sign}(Z)$ と $|Z|$ は独立で，$\Prob(Z>0)=1/2$ である．
したがって $|Z|$ の順位が $k$ である観測が正かどうかの指示変数 $I_k$ $(k=1,\dots,n)$ は，
順位の並びとは独立に $\mathrm{Bernoulli}(1/2)$ に従い互いに独立である．$W^+=\sum_{k=1}^n kI_k$ だから
\[
\E[W^+]=\frac12\sum_{k=1}^nk=\frac{n(n+1)}{4},\qquad
\V[W^+]=\frac14\sum_{k=1}^nk^2=\frac14\cdot\frac{n(n+1)(2n+1)}{6}.
\]
\end{proof}
\kakomon{2016}{1} のデータ（$n=7$）では，負の差は絶対値が最小の1つだけなので $W^+=28-1=27$，
$\E=14$，$\V=35$，$z=13/\sqrt{35}=2.197>1.96$ で有意となる．

\paragraph{exact分布}
$H_0$ の下で $(I_1,\dots,I_n)$ の $2^n$ 通りの値は等確率なので，
$\Prob(W^+=w)=\#\{A\subseteq\{1,\dots,n\}:\sum_{k\in A}k=w\}/2^n$ である．
分布は $n(n+1)/4$ に関して対称なので，上の定義による両側P 値は（小さい方の）片側P 値の2倍である．
小さい $n$ では，$W^-$（または $W^+$）が小さい側の部分集合を数え上げるのが速い．

\begin{supbox}[補足：同順位と0の扱い]
差が0の観測は除く．$|Z_i|$ に同順位があれば平均順位を与え，分散を
$\V[W^+]=\frac{n(n+1)(2n+1)}{24}-\frac{1}{48}\sum_g(t_g^3-t_g)$（$t_g$ は各同順位グループの大きさ）に修正する．
\end{supbox}

\section{Wilcoxon順位和検定とMann--Whitney検定}\label{sec:03-ranksum}
\subsection{定義と帰無分布}
\Hissu\Hinshutsu 2群をまとめた $N=m+n$ 個に順位 $1,\dots,N$ を付け，第1群（大きさ $m$）の順位の和を $W$ とする．
\textbf{帰無仮説}は $H_0:F_X=F_Y$（2群の分布が等しい）であり（\kakomon{2019}{4}〔2〕），
典型的な対立仮説は位置ずれ $F_X(x)=F_Y(x-\delta)$，$\delta\ne0$ である（\kakomon{2025}{1}）．
\textbf{仮定}：すべての観測が独立，分布が連続（同順位なし）．正規性や等分散は不要．

\begin{teiri}[$W$ の帰無分布]\label[teiri]{thm:03-ranksum}
$H_0$ の下で，第1群の順位の集合は $\{1,\dots,N\}$ から大きさ $m$ の部分集合を等確率 $1/\binom Nm$ で選んだものと同じ分布に従い，
\begin{equation}
\E[W]=\frac{m(N+1)}{2},\qquad \V[W]=\frac{mn(N+1)}{12}.\label{eq:03-ranksum}
\end{equation}
\end{teiri}
\begin{proof}
$H_0$ の下で $N$ 個の観測は i.i.d. の連続分布に従うので，順位ベクトルは $N!$ 通りの順列上の一様分布に従う．
よって第1群の順位 $R_1,\dots,R_m$ は $\{1,\dots,N\}$ からの非復元単純無作為抽出である．
母集団 $\{1,\dots,N\}$ の平均は $(N+1)/2$，分散は $\sigma^2=\frac1N\sum k^2-\bigl(\frac{N+1}2\bigr)^2=\frac{N^2-1}{12}$．
非復元抽出では $\V[R_i]=\sigma^2$，$\Cov(R_i,R_j)=-\sigma^2/(N-1)$ $(i\ne j)$ なので
\[
\V[W]=m\sigma^2-m(m-1)\frac{\sigma^2}{N-1}=m\sigma^2\frac{N-m}{N-1}=\frac{m n(N^2-1)}{12(N-1)}=\frac{mn(N+1)}{12}.
\]
\end{proof}
\kakomon{2025}{1} のように $m=n$ のとき $\E[W]=n(2n+1)/2$，$\V[W]=n^2(2n+1)/12\approx n^3/6$ である．

\subsection{Mann--Whitney統計量との関係}\label{subsec:03-mw}
\begin{meidai}\label[meidai]{prop:03-mw}
$U=\sum_{i=1}^m\sum_{j=1}^n I(X_i>Y_j)$ とすると
\begin{equation}
W=U+\frac{m(m+1)}{2}.\label{eq:03-wu}
\end{equation}
\end{meidai}
\begin{proof}
第1群を小さい順に $X_{(1)}<\dots<X_{(m)}$ とすると，$X_{(k)}$ の全体での順位は
「第1群で自分以下の個数 $k$」＋「$X_{(k)}$ より小さい第2群の個数 $\sum_jI(X_{(k)}>Y_j)$」である．$k$ について和をとればよい．
\end{proof}
$\theta=\Prob(X>Y)$ とおくと $\E[U]=mn\theta$ であり，$\hat\theta=U/(mn)$ は $\theta$ の不偏推定量である．
$H_0$ の下で $\theta=1/2$ なので，順位和検定は $\Prob(X>Y)=1/2$ からのずれを検出する検定である（ただし帰無分布の導出には $F_X=F_Y$ を用いる）．
$\hat\theta$ は経験ROC曲線の曲線下面積（AUC）に一致する（\cref{sec:05-auc}）．

\subsection{exact分布と離散性}\label{subsec:03-exact}
\Youchui $W$ のとる値の分布は，$\{1,\dots,N\}$ の $m$ 元部分集合の和の度数分布である．
\kakomon{2019}{4} の $m=n=3$ では $\binom63=20$ 通りで，最大値15・最小値6の確率がそれぞれ $1/20$．
両側P 値の最小値は $2/20=0.1$ なので，\textbf{有意水準5\%ではどんなデータでも棄却されず，第1種の過誤確率は0}である．
これは検出力も0であることを意味する．極端な小標本では，順位検定はそもそも有意になりえないことを設計段階で確認する必要がある．
順位和の帰無分布も $m(N+1)/2$ に関して対称なので両側P 値は片側の2倍であり，一般に最小の両側P 値は $2/\binom{N}{m}$ である．

\section{数理：位置ずれモデルでの検出力}\label{sec:03-power}
\Hinshutsu 対立仮説の下で $\E[U]=mn\theta$，$\theta=\Prob(X>Y)$ である．
分散を帰無仮説の下の値 $\V_0=mn(N+1)/12$ で近似し，$U$ を正規近似する．$Z=(U-mn/2)/\sqrt{\V_0}$ として，片側検定（$Z>z_\alpha$ で棄却）の検出力は
\begin{equation}
1-\beta\approx\Phi\!\left(\frac{mn(\theta-\frac12)}{\sqrt{mn(N+1)/12}}-z_\alpha\right).\label{eq:03-power}
\end{equation}
\paragraph{一様分布の例（\kakomon{2025}{1}）}
$Y\sim\Unif(0,1)$，$X\sim\Unif(\delta,1+\delta)$ $(0\le\delta<1)$ とする．$y\in[0,1]$ を固定すると
$\Prob(X>y)=1$（$y\le\delta$），$=1+\delta-y$（$y>\delta$）なので
\[
\theta=\int_0^1\Prob(X>y)\,\dd y=\delta+\int_\delta^1(1+\delta-y)\,\dd y=\frac{1+2\delta-\delta^2}{2}.
\]
$m=n$，$\V_0\approx n^3/6$ とすると，$n^2(\theta-\tfrac12)=n^2(\delta-\delta^2/2)$ だから検出力は
$\Phi\bigl(\sqrt{6n}(\delta-\delta^2/2)-z_\alpha\bigr)$ となる．

\paragraph{正規分布の例}
$X\sim\Normal(\mu+\delta,\sigma^2)$，$Y\sim\Normal(\mu,\sigma^2)$ なら $X-Y\sim\Normal(\delta,2\sigma^2)$ より $\theta=\Phi\bigl(\delta/(\sqrt2\sigma)\bigr)$．

\begin{supbox}[発展：漸近相対効率]
正規分布の位置ずれでは，順位和検定の $t$ 検定に対する漸近相対効率（同じ検出力を得るのに必要な標本サイズの逆比）は $3/\pi\approx0.955$ である．
すなわち正規分布でも損失は約5\%にすぎず，裾の重い分布では順位検定の方が効率がよい．
符号検定の対応のある $t$ 検定に対する効率は正規分布で $2/\pi\approx0.637$ と低い．
\end{supbox}

\section{仮定のまとめ}\label{sec:03-assump}
\begin{table}[htbp]
\centering\small
\caption{各検定の帰無仮説と仮定}\label{tab:03-assump}
\begin{tabular}{@{}lL{38mm}L{70mm}@{}}
\toprule
検定 & 帰無仮説 & 主な仮定\\
\midrule
対応のある $t$ & 差の平均 $\mu=0$ & 差が独立，差が正規分布\\
符号検定 & 差の中央値 $\theta=0$ & 差が独立（連続分布）．形の仮定なし\\
符号付き順位 & 差の中央値 $\theta=0$ & 差が独立，中央値に関して対称な連続分布\\
2標本 $t$（Student） & $\mu_1=\mu_2$ & 独立，各群正規，等分散\\
順位和（Mann--Whitney） & $F_X=F_Y$ & 独立，連続分布（位置ずれを検出したいなら形が同じ）\\
\bottomrule
\end{tabular}
\end{table}

\section{結果の解釈}\label{sec:03-interp}
\begin{itemize}
\item $t$ 検定と順位検定の結果が食い違ったら，外れ値・歪みの有無を確認する（\kakomon{2016}{1}）．
      順位検定は外れ値に頑健であり，少数の極端な値に左右されない．
\item 順位検定の有意は「分布が異なる（$X$ が確率的に大きい）」ことを示す．平均の差を示すわけではない．
      位置ずれモデルを仮定すれば中央値の差と解釈できる．
\item 小標本で順位検定が有意でないのは，離散性のため有意になりえない場合がある（\kakomon{2019}{4}）．
\end{itemize}

\section{手法選択}\label{sec:03-choice}
\begin{itemize}
\item 対応あり：差が正規 → 対応のある $t$；対称だが外れ値がある → 符号付き順位；非対称・順序尺度 → 符号検定．
\item 独立2群：正規・等分散 → Student $t$；正規・不等分散 → Welch；非正規・外れ値・順序尺度 → 順位和検定．
\item 3群以上：一元配置分散分析，順位なら Kruskal--Wallis 検定．
\item 対数正規が想定される量（薬物動態のAUCなど）→ 対数変換して $t$ 型の方法（\cref{ch:13}）．
\end{itemize}

\section{よくある誤り}\label{sec:03-err}
\begin{warnbox}
\begin{itemize}
\item 符号付き順位検定に「仮定なし」と書く → 対称性を仮定している．仮定がないのは符号検定．
\item 順位和検定の帰無仮説を「平均が等しい」と書く → 分布が等しい（$F_X=F_Y$）．
\item $\V[W]$ を独立な和として $m(N^2-1)/12$ とする → 非復元抽出の有限母集団修正 $(N-m)/(N-1)$ を忘れている．
\item 両側P 値を片側P 値のまま報告する．
\item exact分布の最小P 値を確認せずに小標本で順位検定を計画する．
\end{itemize}
\end{warnbox}

\section{数理統計との接続}\label{sec:03-link}
\begin{linkbox}
\begin{itemize}
\item $t$ 検定の根拠は，正規標本で $\bar X$ と不偏分散が独立で $(n-1)s^2/\sigma^2\sim\chi^2_{n-1}$ となること（『現代数理統計学の基礎』定理5.1）．
      同書は1標本 $t$ 検定を本文で扱っていないが，2標本 $t$ 検定（7.2.2項）と同じ構成である．
\item 順位は順序統計量（同書5.4節）とは別物である．順位検定の帰無分布は「順列の一様分布」から導かれ，
      $\V[W]$ は超幾何分布と同じ有限母集団修正を含む．\kakomonSM{2023}{5} の並べ替え検定の分散 $\frac{N-m}{N-1}\frac{\sigma_N^2}{m}$ と同型である．
\item Mann--Whitney の $U/(mn)$ は $\Prob(X>Y)$ の U 統計量であり，AUC の推定量でもある（\cref{ch:05}）．
\end{itemize}
\end{linkbox}

\section{例題}\label{sec:03-ex}
\begin{reidai}[対応のあるデータの3つの検定]\label{reidai:03-1}
8名の患者の治療後$-$治療前の変化量が
$2.1,\ -0.4,\ 3.5,\ 1.2,\ 0.8,\ 4.0,\ -1.5,\ 2.6$ であった
（$\bar z=1.5375$，$s_z=1.890$）．$t_7(0.025)=2.365$ を用いよ．
\begin{enumerate}[label=(\arabic*)]
\item 対応のある $t$ 検定（両側5\%）を行え．
\item 符号検定のexactな両側P 値を求めよ．
\item 符号付き順位統計量 $W^+$ を求め，正規近似による検定（両側5\%）を行え．
\item $W^+$ のexactな両側P 値を求めよ．
\end{enumerate}
\end{reidai}

\begin{reidai}[順位和検定のexact分布]\label{reidai:03-2}
薬剤A群3名，薬剤B群4名の測定値は A：$5.1,\ 7.3,\ 8.8$，B：$2.2,\ 3.0,\ 4.7,\ 6.1$ であった．
\begin{enumerate}[label=(\arabic*)]
\item A群の順位和 $W_A$ と Mann--Whitney 統計量 $U$，$\hat\theta=\widehat{\Prob}(X_A>X_B)$ を求めよ．
\item $H_0$ の下での $W_A$ の期待値・分散を求めよ．
\item $W_A$ のexactな片側P 値（$H_1$：Aが大きい）と両側P 値を求めよ．
\item この標本サイズで，両側5\%の順位和検定が有意になりうるかを答えよ．
\end{enumerate}
\end{reidai}

\begin{reidai}[順位和検定の検出力]\label{reidai:03-3}
2群（各 $n=50$）で，$X\sim\Normal(\mu+0.5\sigma,\sigma^2)$，$Y\sim\Normal(\mu,\sigma^2)$ とする．
$\Phi(0.354)=0.638$，$\Phi(0.42)=0.663$，$\Phi(0.54)=0.705$ を用いよ．
\begin{enumerate}[label=(\arabic*)]
\item $\theta=\Prob(X>Y)$ を求めよ．
\item \cref{eq:03-power}を両側5\%（棄却限界に $z_{0.025}=1.96$ を用い，反対側で棄却される確率は無視する近似）に適用して，順位和検定の検出力を求めよ．
\item 同じ状況の2標本 $z$ 検定（$\sigma$ 既知）の検出力と比べよ．
\end{enumerate}
\end{reidai}

\begin{reidai}[符号付き順位の分布の導出]\label{reidai:03-4}
$n=4$ のとき，$H_0$ の下での $W^+$ の確率分布を求め，その平均と分散が\cref{eq:03-wplus}に一致することを確かめよ．
\end{reidai}

\section{解答}\label{sec:03-sol}
\begin{kaitou}{reidai:03-1}
(1) $t=1.5375/(1.890/\sqrt8)=1.5375/0.6682=2.30<2.365$．有意水準5\%で有意でない．\\
(2) 正の差は6個．$\Prob(T\ge6)=(\binom86+\binom87+\binom88)/2^8=(28+8+1)/256=37/256=0.1445$，両側P 値 $0.289$．\\
(3) 絶対値の順位：$0.4(1,-)$，$0.8(2)$，$1.2(3)$，$1.5(4,-)$，$2.1(5)$，$2.6(6)$，$3.5(7)$，$4.0(8)$．
$W^+=2+3+5+6+7+8=31$（$W^-=5$）．$\E=8\cdot9/4=18$，$\V=8\cdot9\cdot17/24=51$，
$z=(31-18)/\sqrt{51}=13/7.141=1.82<1.96$ で有意でない．\\
(4) $\Prob(W^+\ge31)=\Prob(W^-\le5)$．$\{1,\dots,8\}$ の部分集合で和が5以下のものは
$\emptyset,\{1\},\{2\},\{3\},\{1,2\},\{4\},\{1,3\},\{5\},\{1,4\},\{2,3\}$ の10個．
片側 $10/256=0.039$，両側 $0.078$．
\end{kaitou}

\begin{kaitou}{reidai:03-2}
(1) 全体の順位は $2.2(1),3.0(2),4.7(3),5.1(4),6.1(5),7.3(6),8.8(7)$．$W_A=4+6+7=17$．
\cref{eq:03-wu}より $U=17-3\cdot4/2=11$，$\hat\theta=11/12=0.917$．\\
(2) $N=7$：$\E=3\times8/2=12$，$\V=3\times4\times8/12=8$．\\
(3) $\binom73=35$ 通りのうち $W_A\ge17$ は $\{5,6,7\}$（18）と $\{4,6,7\}$（17）の2通り．片側 $2/35=0.057$，両側 $0.114$．\\
(4) 最も極端な $\{5,6,7\}$ でも片側 $1/35$，両側 $2/35=0.057>0.05$ なので，両側5\%ではどのデータでも有意にならない．
\end{kaitou}

\begin{kaitou}{reidai:03-3}
(1) $X-Y\sim\Normal(0.5\sigma,2\sigma^2)$ より $\theta=\Phi(0.5/\sqrt2)=\Phi(0.354)=0.638$．\\
(2) $mn(\theta-1/2)=2500\times0.138=345$，$\sqrt{\V_0}=\sqrt{2500\times101/12}=\sqrt{21042}=145.1$．
$345/145.1=2.38$，検出力 $\approx\Phi(2.38-1.96)=\Phi(0.42)=0.663$．\\
(3) $z$ 検定では $\delta/(\sigma\sqrt{2/n})=0.5\times5=2.5$，検出力 $\Phi(2.5-1.96)=\Phi(0.54)=0.705$．
順位和検定の検出力はやや低いが近い（$\delta=0.5\sigma$ は局所対立ではないので，漸近相対効率 $3/\pi$ とは目安として比べる）．
\end{kaitou}

\begin{kaitou}{reidai:03-4}
$\{1,2,3,4\}$ の16個の部分集合の和：
0（$\emptyset$），1，2，3（$\{3\},\{1,2\}$），4（$\{4\},\{1,3\}$），5（$\{1,4\},\{2,3\}$），6（$\{1,2,3\},\{2,4\}$），
7（$\{3,4\},\{1,2,4\}$），8（$\{1,3,4\}$），9（$\{2,3,4\}$），10（全体）．
よって $\Prob(W^+=w)$ は $w=0,\dots,10$ に対し $\frac1{16}(1,1,1,2,2,2,2,2,1,1,1)$．
対称性から平均は5で $n(n+1)/4=5$ に一致．
$\E[(W^+-5)^2]=\frac1{16}\{2(25+16+9)+2\cdot2(4+1)\}=\frac{100+20}{16}=7.5$，
一方 $n(n+1)(2n+1)/24=4\cdot5\cdot9/24=7.5$ で一致する．
\end{kaitou}

\enlargethispage{3\baselineskip}
\begin{summarybox}
\begin{itemize}
\item 対応のある $t$：$t=\bar Z/(s_Z/\sqrt n)\sim t_{n-1}$．2標本 $t$：仮定は独立・正規・等分散．
\item 符号検定：$T\sim\Bin(n,1/2)$，$\E=n/2$，$\V=n/4$．仮定は独立のみ．
\item 符号付き順位：$W^+=\sum kI_k$，$\E=n(n+1)/4$，$\V=n(n+1)(2n+1)/24$．仮定は対称性．
\item 順位和：$\E[W]=m(N+1)/2$，$\V[W]=mn(N+1)/12$（非復元抽出）．$W=U+m(m+1)/2$，$\E[U]=mn\Prob(X>Y)$．
\item exact分布は部分集合の数え上げ．最小両側P 値 $2/\binom Nm$ が $\alpha$ を超えると有意になりえない．
\item 検出力 $\approx\Phi\bigl(mn(\theta-\frac12)/\sqrt{mn(N+1)/12}-z_\alpha\bigr)$．
\end{itemize}
\end{summarybox}

% ===== ch04.tex =====
\chapter{2 値データと分割表}\label{ch:04}

\begin{goalbox}
\begin{itemize}
\item 2 値データの確率モデル（二項・積二項・多項・超幾何）を研究デザインから選べる．
\item $2\times2$ 表の Pearson $\chi^2$ 検定を，比率の差の $z$ 検定・スコア検定として導出できる．
\item リスク差・リスク比・オッズ比の推定量と，デルタ法による分散・信頼区間を導出できる．
\item $K$ 群の比率の一様性の検定と Cochran--Armitage 傾向性検定を導出し，使い分けられる．
\item 対応のある2 値データで McNemar 検定・差の信頼区間・尤度比検定を導出できる．
\item 層別 $2\times2$ 表の Mantel--Haenszel 検定と共通オッズ比の推定量を書ける．
\end{itemize}
\end{goalbox}

\begin{exambox}
\begin{itemize}
\item \kakomon{2016}{2}：4群の有効率の一様性の検定統計量と漸近分布（$\chi^2_3$），最小二乗（重み $n_i$）による傾き $\hat\beta$ と $\SE[\hat\beta\mid H_0]$，Cochran--Armitage 検定．
\item \kakomon{2018}{2}：マッチドペアの多項分布．$\hat\delta=(n_{12}-n_{21})/n$ の分散，信頼区間，ラグランジュ乗数法による制約付き最尤推定，尤度比検定（自由度1）．
\item \kakomon{2019}{3}〔2〕：同一患者での2検査の感度比較（不一致ペアの二項分布 $\Bin(23,0.5)$，McNemar型）．
\item \kakomon{2021}{3}：オッズ比とWoolfの信頼区間，積二項尤度，$\hat E[n_{11}]$ と $\hat V[n_{11}-\hat E[n_{11}]]$ の導出，層別 Cochran 検定．
\item \kakomon{2024}{1}〔2〕〔3〕：$Z_i=X_i-Y_i$ の平均と分散から McNemar 統計量 $W=(n_{10}-n_{01})^2/(n_{10}+n_{01})$ を導出．
\item \kakomon{2017}{4}：層別表の超幾何分布の期待値・分散（\cref{eq:04-hyper}）と $(O-E)/V$ による1段階近似（Peto 型）の共通オッズ比．\kakomon{2022}{2}：MH 推定量（解釈は\cref{ch:07}）．
\end{itemize}
\end{exambox}

\flowline{2 値\\独立／対応／層別}{RD・RR・OR\\一様性・傾向}{$\chi^2$・MH\\McNemar・CA}{二項・多項\\超幾何}{期待度数・分散\\デルタ法}{効果の向きと大きさ}

\section{問題設定と確率モデル}\label{sec:04-model}
\cref{tab:02-2x2}の記号（$a,b,c,d$，行和 $n_1,n_0$，列和 $m_1,m_0$，総数 $N$）を用いる．
どの周辺度数が固定されているかで確率モデルが決まる．
\begin{table}[htbp]
\centering\small
\caption{分割表の確率モデル}\label{tab:04-models}
\begin{tabular}{@{}lL{42mm}L{65mm}@{}}
\toprule
モデル & 固定されるもの & 典型的なデザイン\\
\midrule
多項分布 & 総数 $N$ & 横断研究，マッチドペア（ペア数 $n$ を固定し4セルに多項）\\
積二項分布 & 行和 $n_1,n_0$ & RCT，コホート（群ごとに二項）\\
積二項分布 & 列和 $m_1,m_0$ & 症例対照研究（症例・対照ごとに曝露が二項）\\
超幾何分布 & 行和と列和 & 条件付き推測（Fisherの正確検定，MH検定）\\
\bottomrule
\end{tabular}
\end{table}
独立性（多項）と一様性（積二項）の帰無仮説は形式的に異なるが，Pearson $\chi^2$ 統計量とその漸近分布は同じになる．

\section{\texorpdfstring{$2\times2$}{2×2} 表の検定}\label{sec:04-2x2test}
\subsection{比率の差の検定と Pearson \texorpdfstring{$\chi^2$}{χ²}}\label{subsec:04-pearson}
積二項モデル $a\sim\Bin(n_1,\pi_1)$，$c\sim\Bin(n_0,\pi_0)$（独立）で $H_0:\pi_1=\pi_0\,(=\pi)$ を考える．
\paragraph{共通比率の最尤推定}
積二項尤度 $L(\pi_1,\pi_0)=\binom{n_1}{a}\pi_1^a(1-\pi_1)^{b}\binom{n_0}{c}\pi_0^c(1-\pi_0)^d$ は $H_0$ の下で
$\propto\pi^{a+c}(1-\pi)^{b+d}$ となり，$\hat\pi=m_1/N$．
\paragraph{スコア型統計量（\kakomon{2021}{3}）}
$H_0$ の下で $\E[a]=n_1\pi$ を $\hat E[a]=n_1m_1/N$ と推定すると
\[
a-\hat E[a]=\frac{Na-n_1(a+c)}{N}=\frac{n_0a-n_1c}{N}=\frac{ad-bc}{N}.
\]
$a,c$ は独立なので $\V[n_0a-n_1c]=n_0^2n_1\pi(1-\pi)+n_1^2n_0\pi(1-\pi)=n_1n_0N\pi(1-\pi)$，すなわち
\[
\V\bigl[a-\hat E[a]\bigr]=\frac{n_1n_0\pi(1-\pi)}{N},\qquad
\hat\V\bigl[a-\hat E[a]\bigr]=\frac{n_1n_0m_1m_0}{N^3}.
\]
\begin{meidai}[Pearson $\chi^2$ 統計量]\label[meidai]{prop:04-pearson}
\begin{equation}
X^2=\frac{(a-\hat E[a])^2}{\hat\V[a-\hat E[a]]}=\frac{N(ad-bc)^2}{n_1n_0m_1m_0}
=\sum_{\text{4セル}}\frac{(O-E)^2}{E}=\frac{(\hat\pi_1-\hat\pi_0)^2}{\hat\pi(1-\hat\pi)(1/n_1+1/n_0)}\label{eq:04-pearson}
\end{equation}
であり，$H_0$ の下で漸近的に $\chi^2_1$ に従う．
\end{meidai}
\begin{proof}
第1の等号の右辺は上の計算から直ちに得られる．4セルの $O-E$ はすべて $\pm(ad-bc)/N$ で，
$\sum 1/E=\frac{N}{n_1m_1}+\frac{N}{n_1m_0}+\frac{N}{n_0m_1}+\frac{N}{n_0m_0}=\frac{N^3}{n_1n_0m_1m_0}$ より第2の等号．
$\hat\pi_1-\hat\pi_0=(ad-bc)/(n_1n_0)$ を代入すれば最後の等号を得る．漸近分布は中心極限定理とSlutskyの定理による．
\end{proof}
最後の形は比率の差をプールした分散で標準化した $z$ 統計量の2乗である．

\begin{supbox}[補足：Yates補正と Fisher の正確検定]
度数が小さいときは連続修正 $X^2_Y=N(|ad-bc|-N/2)^2/(n_1n_0m_1m_0)$ を用いることがある（保守的になる）．
周辺度数をすべて固定すると，$H_0$ の下で $a$ は超幾何分布
$\Prob(a)=\binom{n_1}{a}\binom{n_0}{m_1-a}\big/\binom{N}{m_1}$ に従う．これに基づく検定が Fisher の正確検定である．
例：治療群5人中4人，対照群5人中0人が有効なら，$m_1=4$ で $a\ge4$ となる確率は
$\binom54\binom50/\binom{10}{4}=5/210=0.024$（片側）．
\end{supbox}

\subsection{尤度比検定}\label{subsec:04-lrt}
尤度比統計量 $G^2=2\sum O\log(O/E)$ も $H_0$ の下で $\chi^2_1$ に従い，$X^2$ と漸近同等である
（パラメータは積二項で2個，$H_0$ で1個．差が自由度）．

\section{効果指標の推定とデルタ法}\label{sec:04-delta}
\subsection{デルタ法}
\begin{teiri}[多変量デルタ法]\label[teiri]{thm:04-delta}
$\sqrt n(\hat{\bm\theta}-\bm\theta)\xrightarrow{d}\Normal_k(\bm0,\bm\Sigma)$ で $g:\mathbb R^k\to\mathbb R$ が $\bm\theta$ で微分可能，
$\nabla g(\bm\theta)\ne\bm0$ なら
\begin{equation}
\sqrt n\bigl(g(\hat{\bm\theta})-g(\bm\theta)\bigr)\xrightarrow{d}\Normal\bigl(0,\ \nabla g(\bm\theta)\transp\bm\Sigma\,\nabla g(\bm\theta)\bigr).\label{eq:04-delta}
\end{equation}
\end{teiri}
\begin{proof}[証明の概略]
$g(\hat{\bm\theta})=g(\bm\theta)+\nabla g(\bm\theta)\transp(\hat{\bm\theta}-\bm\theta)+o_p(n^{-1/2})$ と1次のテイラー展開し，
線形変換された正規ベクトルの分布と Slutsky の定理を用いる．
\end{proof}
実用上は「$\V[g(\hat{\bm\theta})]\approx\nabla g\transp\,\V[\hat{\bm\theta}]\,\nabla g$」と覚え，
独立な推定量なら各成分の寄与の和 $\sum_j(\partial g/\partial\theta_j)^2\V[\hat\theta_j]$ になる．

\subsection{各指標の分散}\label{subsec:04-var}
\begin{meidai}\label[meidai]{prop:04-var}
独立な $a\sim\Bin(n_1,\pi_1)$，$c\sim\Bin(n_0,\pi_0)$ について
\begin{align}
\V[\hat\pi_1-\hat\pi_0]&=\frac{\pi_1(1-\pi_1)}{n_1}+\frac{\pi_0(1-\pi_0)}{n_0},\\
\V[\log\hat\pi_1]&\approx\frac{1-\pi_1}{n_1\pi_1},\qquad
\V[\log\widehat{\RR}]\approx\frac{1-\pi_1}{n_1\pi_1}+\frac{1-\pi_0}{n_0\pi_0},\label{eq:04-varlogRR}\\
\V\Bigl[\log\frac{\hat\pi_1}{1-\hat\pi_1}\Bigr]&\approx\frac{1}{n_1\pi_1(1-\pi_1)},\qquad
\V[\log\widehat{\OR}]\approx\frac{1}{n_1\pi_1(1-\pi_1)}+\frac{1}{n_0\pi_0(1-\pi_0)}.\label{eq:04-varlogOR}
\end{align}
推定量は $\pi$ を $\hat\pi$ に置き換えて，$\widehat\V[\log\widehat{\RR}]=\frac1a-\frac1{n_1}+\frac1c-\frac1{n_0}$，
$\widehat\V[\log\widehat{\OR}]=\frac1a+\frac1b+\frac1c+\frac1d$（Woolf）．
\end{meidai}
\begin{proof}
$g(\pi)=\log\pi$ なら $g'=1/\pi$ で $\V\approx\pi^{-2}\pi(1-\pi)/n_1$．
$g(\pi)=\log\{\pi/(1-\pi)\}$ なら $g'=1/\{\pi(1-\pi)\}$ で $\V\approx\{\pi(1-\pi)\}^{-2}\pi(1-\pi)/n_1$．
2群は独立なので分散は和．$n_1\hat\pi_1(1-\hat\pi_1)=ab/n_1$ より $1/\{n_1\hat\pi_1(1-\hat\pi_1)\}=1/a+1/b$．
\end{proof}
信頼区間は $\log\hat\theta\pm z_{\alpha/2}\widehat{\SE}$ を指数変換する．
\kakomon{2021}{3} では $\widehat{\OR}=\frac{59\times65}{71\times60}=0.90$，$\widehat\V=\frac1{59}+\frac1{71}+\frac1{60}+\frac1{65}=0.063$ で，
95\%信頼区間は $(0.55,1.47)$ となる．

\section{\texorpdfstring{$K$}{K} 群の比率の一様性の検定}\label{sec:04-homog}
$K$ 群で $y_i\sim\Bin(n_i,\pi_i)$（独立），$p_i=y_i/n_i$，$p=\sum y_i/N$ とする．$H_0:\pi_1=\dots=\pi_K$ に対し
\begin{equation}
X^2=\sum_{i=1}^K\Bigl(\frac{p_i-p}{\sqrt{p(1-p)/n_i}}\Bigr)^2=\frac{1}{p(1-p)}\sum_{i=1}^Kn_i(p_i-p)^2\ \xrightarrow{d}\ \chi^2_{K-1}.\label{eq:04-homog}
\end{equation}
これは $K\times2$ 表の Pearson $\chi^2$ 統計量と一致する（\kakomon{2016}{2}〔1〕）．
自由度は，$K$ 個の比率から共通の $p$ を推定して1つ失うので $K-1$．
対立仮説は「少なくとも1つの等号が成り立たない」で，方向を持たない．

\section{Cochran--Armitage 傾向性検定}\label{sec:04-trend}
\Hinshutsu 群に用量 $x_1<\dots<x_K$ の順序があるとき，単調な用量反応を検出したい．
\kakomon{2016}{2}〔3〕の定式化に従う．

\paragraph{重み付き最小二乗}
$\pi_i=\alpha+\beta x_i$ を，$(x_i,p_i)$ に重み $n_i$ で当てはめる：
$Q=\sum n_i(p_i-\alpha-\beta x_i)^2$ を最小化すると，$\bar x=\sum n_ix_i/N$，$S_{xx}=\sum n_i(x_i-\bar x)^2$ として
\begin{equation}
\hat\beta=\frac{\sum_in_i(p_i-p)(x_i-\bar x)}{S_{xx}}=\frac{\sum_iy_i(x_i-\bar x)}{S_{xx}}.\label{eq:04-cabeta}
\end{equation}
（$\partial Q/\partial\alpha=0$ から $\hat\alpha=p-\hat\beta\bar x$．これを $\partial Q/\partial\beta=0$ に代入し，
$\sum n_i(x_i-\bar x)=0$ を用いると上式になる．）

\paragraph{帰無仮説の下での分散}
$H_0$ の下で $p_i$ は独立に $\V[p_i]=\pi(1-\pi)/n_i$ なので，
\[
\V[\hat\beta\mid H_0]=\frac{\sum_in_i^2(x_i-\bar x)^2\,\pi(1-\pi)/n_i}{S_{xx}^2}=\frac{\pi(1-\pi)}{S_{xx}},\qquad
\SE[\hat\beta\mid H_0]=\sqrt{\frac{p(1-p)}{S_{xx}}}.
\]
\begin{teiri}[Cochran--Armitage 傾向性検定]\label[teiri]{thm:04-ca}
\begin{equation}
Y=\frac{\hat\beta^2}{p(1-p)/S_{xx}}=\frac{\bigl\{\sum_iy_i(x_i-\bar x)\bigr\}^2}{p(1-p)\sum_in_i(x_i-\bar x)^2}\ \xrightarrow{d}\ \chi^2_1\quad(H_0\text{ の下}).\label{eq:04-ca}
\end{equation}
さらに $X^2=Y+(X^2-Y)$ と分解でき，$X^2-Y$ は直線からのずれの検定統計量で，$H_0$ の下で漸近的に $\chi^2_{K-2}$ に従う．
\end{teiri}
一様性の検定は $K-1$ 自由度すべての方向の差を検出しようとするのに対し，傾向性検定は用量反応の方向に1自由度を集中させる．
真の関係がスコアに対して直線に近ければ $Y$ は $X^2$ のほぼすべてを占め，自由度が小さい分だけ有意になりやすい．
これが \kakomon{2016}{2} で $X^2$ は $p=0.077$，$Y$ は $p=0.042$ となった理由である．
ただし単調でない（例：中用量で最大）関係には傾向性検定は検出力が低い．

\section{対応のある2 値データと McNemar 検定}\label{sec:04-mcnemar}
\Hissu\Hinshutsu
\subsection{設定}
$n$ 組のペア（または $n$ 人の被験者に2つの処置・検査）について，\cref{tab:04-paired}のように集計する．
セル度数 $(n_{11},n_{12},n_{21},n_{22})$ は多項分布 $\mathrm{Mult}(n;\pi_{11},\pi_{12},\pi_{21},\pi_{22})$ に従う．
\begin{table}[htbp]
\centering\small
\caption{対応のある2 値データ（行：処置1，列：処置2）}\label{tab:04-paired}
\begin{tabular}{@{}lccc@{}}
\toprule
 & 処置2 有効 & 処置2 無効 & 計\\
\midrule
処置1 有効 & $n_{11}$ & $n_{12}$ & $n_{1\cdot}$\\
処置1 無効 & $n_{21}$ & $n_{22}$ & $n_{2\cdot}$\\
\midrule
計 & $n_{\cdot1}$ & $n_{\cdot2}$ & $n$\\
\bottomrule
\end{tabular}
\end{table}
有効率の差は $\pi_{1\cdot}-\pi_{\cdot1}=(\pi_{11}+\pi_{12})-(\pi_{11}+\pi_{21})=\pi_{12}-\pi_{21}=:\delta$ で，
一致したペア（$n_{11},n_{22}$）は差に寄与しない．

\subsection{差の分散と信頼区間}
\begin{meidai}[\kakomon{2018}{2}〔1〕〔2〕]\label[meidai]{prop:04-mcvar}
$\hat\delta=(n_{12}-n_{21})/n$ について
\begin{equation}
\V[\hat\delta]=\frac{1}{n}\bigl\{\pi_{12}+\pi_{21}-(\pi_{12}-\pi_{21})^2\bigr\}.\label{eq:04-mcvar}
\end{equation}
正規近似による $100(1-\alpha)\%$ 信頼区間は
$\hat\delta\pm z_{\alpha/2}\sqrt{\dfrac{n_{12}+n_{21}}{n^2}-\dfrac{(n_{12}-n_{21})^2}{n^3}}$．
\end{meidai}
\begin{proof}
多項分布で $\V[n_{ij}]=n\pi_{ij}(1-\pi_{ij})$，$\Cov(n_{12},n_{21})=-n\pi_{12}\pi_{21}$ なので
\[
\V[\hat\delta]=\frac{1}{n^2}\{n\pi_{12}(1-\pi_{12})+n\pi_{21}(1-\pi_{21})+2n\pi_{12}\pi_{21}\}
=\frac1n\{\pi_{12}+\pi_{21}-(\pi_{12}-\pi_{21})^2\}.
\]
\end{proof}
別の見方（\kakomon{2024}{1}）：被験者ごとに処置1・2の結果を $X_i,Y_i\in\{0,1\}$（1＝有効）とし，$\pi_{jk}=\Prob(X_i=j,Y_i=k)$ と書くと $\pi_{10}=\pi_{12}$，$\pi_{01}=\pi_{21}$ に対応する．$Z_i=X_i-Y_i\in\{1,0,-1\}$ とすると
$\E[Z_i]=\pi_{10}-\pi_{01}$，$\E[Z_i^2]=\pi_{10}+\pi_{01}$ より $\V[Z_i]=(\pi_{10}+\pi_{01})-(\pi_{10}-\pi_{01})^2$ で，
$\hat\pi_X-\hat\pi_Y=\bar Z$ の分散は $\V[Z_i]/n$ となり同じ結果を得る．

\subsection{McNemar 検定}
\begin{teiri}[McNemar 検定]\label[teiri]{thm:04-mcnemar}
$H_0:\pi_{12}=\pi_{21}$（周辺確率が等しい）の下で $\V[Z_i]=\pi_{12}+\pi_{21}$ であり，これを $(n_{12}+n_{21})/n$ で推定すると
\begin{equation}
W=\frac{n\,\bar Z^2}{\hat\V_0[Z_i]}=\frac{(n_{12}-n_{21})^2}{n_{12}+n_{21}}\ \xrightarrow{d}\ \chi^2_1.\label{eq:04-mcnemar}
\end{equation}
\end{teiri}
\paragraph{条件付き二項検定}
不一致ペア数 $s=n_{12}+n_{21}$ を与えると，$n_{12}\mid s\sim\Bin\bigl(s,\ \pi_{12}/(\pi_{12}+\pi_{21})\bigr)$ であり，
$H_0$ の下で $\Bin(s,1/2)$．正規近似 $z=(n_{12}-s/2)/\sqrt{s/4}$ の2乗が $W$ である．
\kakomon{2019}{3}〔2〕の感度の比較は，罹患者における2検査の不一致ペア $s=23$（$n_{+-}=7$，$n_{-+}=16$）で
$z=(7-11.5)/\sqrt{5.75}=-1.88$，有意でない．小標本では $\Bin(s,1/2)$ で exact に検定する．

\subsection{尤度比検定（\kakomon{2018}{2}〔3〕〔4〕）}\label{subsec:04-mclrt}
対数尤度 $\ell=\sum n_{ij}\log\pi_{ij}$ を，制約 $\sum\pi_{ij}=1$ の下で最大化すると $\hat\pi_{ij}=n_{ij}/n$．
$H_0:\pi_{12}=\pi_{21}$ の制約も加え，ラグランジュ関数
$\ell-\lambda(\sum\pi_{ij}-1)-\phi(\pi_{12}-\pi_{21})$ を微分して0とおくと
$n_{11}=\lambda\pi_{11}$，$n_{12}=(\lambda+\phi)\pi_{12}$，$n_{21}=(\lambda-\phi)\pi_{21}$，$n_{22}=\lambda\pi_{22}$．
辺々加えて $\lambda=n$，さらに $\pi_{12}=\pi_{21}$ から
\[
\hat\pi_{11}^{(0)}=\frac{n_{11}}{n},\quad \hat\pi_{12}^{(0)}=\hat\pi_{21}^{(0)}=\frac{n_{12}+n_{21}}{2n},\quad \hat\pi_{22}^{(0)}=\frac{n_{22}}{n}.
\]
一致セルの推定値は両モデルで同じなので打ち消し合い，
\begin{equation}
G^2=2\log\Lambda=2\Bigl[n_{12}\log\frac{2n_{12}}{n_{12}+n_{21}}+n_{21}\log\frac{2n_{21}}{n_{12}+n_{21}}\Bigr]\ \xrightarrow{d}\ \chi^2_1.\label{eq:04-mclrt}
\end{equation}
自由度は対立仮説の自由パラメータ数3と帰無仮説の2の差で1である．

\subsection{条件付きオッズ比}
ペアごとに異なるベースラインを許すロジスティックモデル $\logit\Prob(\text{有効})=\alpha_k+\beta\cdot(\text{処置})$ の下で，
$\alpha_k$ を消去した条件付き尤度から $e^{\hat\beta}=n_{12}/n_{21}$ となる．
これはペアを層とみなした MH オッズ比（\cref{sec:04-mh}）とも一致する（\kakomon{2022}{2}：同乗者の運転者に対する $\widehat{\OR}_{\mathrm{MH}}=2632/662=3.98$，すなわち不一致ペアの比）．

\section{層別 \texorpdfstring{$2\times2$}{2×2} 表と Mantel--Haenszel 法}\label{sec:04-mh}
\Hinshutsu 層 $k=1,\dots,K$ ごとの表（$a_k,b_k,c_k,d_k$，行和 $n_{1k},n_{0k}$，列和 $m_{1k},m_{0k}$，総数 $N_k$）を考える．
交絡の調整としての意味は\cref{ch:07}で詳しく扱い，ここでは統計量をまとめる．
\paragraph{MH 検定}
各層で周辺度数を固定すると，$H_0$（全層で OR $=1$）の下で $a_k$ は超幾何分布に従い
\begin{equation}
\E[a_k]=\frac{n_{1k}m_{1k}}{N_k},\qquad \V[a_k]=\frac{n_{1k}n_{0k}m_{1k}m_{0k}}{N_k^2(N_k-1)}.\label{eq:04-hyper}
\end{equation}
層は独立なので
\begin{equation}
\chi^2_{\mathrm{MH}}=\frac{\bigl\{\sum_k(a_k-\E[a_k])\bigr\}^2}{\sum_k\V[a_k]}\ \xrightarrow{d}\ \chi^2_1.\label{eq:04-mhtest}
\end{equation}
\kakomon{2021}{3} の Cochran 検定は，分散に条件なし（積二項）版 $n_{1k}n_{0k}m_{1k}m_{0k}/N_k^3$ を用いたもので，
$N_k$ が大きければ両者はほぼ等しい．MH 版は各層が小さい（マッチドペア $N_k=2$ など）場合にも使える．
\paragraph{MH 推定量}
\begin{equation}
\widehat{\OR}_{\mathrm{MH}}=\frac{\sum_ka_kd_k/N_k}{\sum_kb_kc_k/N_k},\qquad
\widehat{\RR}_{\mathrm{MH}}=\frac{\sum_ka_kn_{0k}/N_k}{\sum_kc_kn_{1k}/N_k}.\label{eq:04-mhest}
\end{equation}
層内で比が共通なら一致推定量であり，層が多数で各層が小さくても一致性を保つ．
このほか，$\widehat{\log\psi}=\sum_k(a_k-\E[a_k])/\sum_k\V[a_k]$（Peto 型の1段階推定；\cref{subsec:07-peto}）も用いられる．

\section{仮定}\label{sec:04-assump}
\begin{itemize}
\item 独立な2群の比較：個体間の独立性，各群内で発症確率が一定．
\item $\chi^2$ 近似：期待度数が十分大きい（目安として全セル5以上）．満たさなければ Fisher の正確検定・exact二項検定．
\item McNemar：ペア間の独立性．ペア内の相関は許容される（むしろそれを利用する）．
\item 傾向性検定：用量のスコア $x_i$ の選び方が結果を左右する（等間隔スコアが一般的）．
\item MH：層内でORが共通（均一性）であることを推定の前提とする．検定自体は均一でなくても有効だが，効果が逆向きの層があると検出力を失う．
\end{itemize}

\section{結果の解釈}\label{sec:04-interp}
\begin{itemize}
\item 一様性の検定が非有意でも傾向性検定が有意なら，「用量とともに有効率が増加する傾向がある」と述べられる．
      ただし，どちらを主解析とするかは事前に決めておくべきである（\kakomon{2016}{2}〔4〕）．
\item McNemar 検定で有意なら，2処置（2検査）の周辺陽性率が異なる．一致ペアは情報を持たないことを答案で示すとよい．
\item OR の信頼区間が1を含まなければ関連あり．RR が推定できるデザインでは RR も併記すると解釈しやすい．
\end{itemize}

\section{手法選択}\label{sec:04-choice}
\begin{table}[htbp]
\centering\small
\caption{2 値データの手法選択}\label{tab:04-choice}
\begin{tabular}{@{}L{55mm}L{85mm}@{}}
\toprule
状況 & 手法\\
\midrule
独立2群，大標本 & Pearson $\chi^2$（＝比率の差の $z$ 検定），RD・RR・OR の区間推定\\
独立2群，小標本 & Fisher の正確検定\\
独立 $K$ 群，順序なし & $K\times2$ 表の $\chi^2$（自由度 $K-1$）\\
独立 $K$ 群，用量の順序あり & Cochran--Armitage 傾向性検定（自由度1）\\
対応あり（前後，マッチドペア，同一患者に2検査） & McNemar 検定，差の区間推定，条件付きOR $n_{12}/n_{21}$\\
層別（交絡の調整） & MH 検定，MH 推定量（\cref{ch:07}）\\
連続共変量を含めて調整 & ロジスティック回帰（\cref{ch:06}）\\
\bottomrule
\end{tabular}
\end{table}

\section{よくある誤り}\label{sec:04-err}
\begin{warnbox}
\begin{itemize}
\item 対応のある2 値データを独立2群の $2\times2$ 表として $\chi^2$ 検定する．
\item McNemar 統計量の分母に一致ペアを含める．
\item Woolf の分散 $1/a+1/b+1/c+1/d$ を RR に使う（RR は $1/a-1/n_1+1/c-1/n_0$）．
\item 一様性の検定の自由度を $K$ とする（正しくは $K-1$）．傾向性検定の自由度を $K-1$ とする（正しくは1）．
\item 傾向性検定の $\V[\hat\beta]$ を対立仮説の下の分散で計算する（検定では $H_0$ の下の共通 $p$ を用いる）．
\end{itemize}
\end{warnbox}

\section{数理統計との接続}\label{sec:04-link}
\begin{linkbox}
\begin{itemize}
\item 分割表の独立性の Pearson $\chi^2$ 検定と自由度 $(r-1)(c-1)$ の導出は『現代数理統計学の基礎』7.4節（例7.12は薬と疾患の $2\times2$ 表）．
      Yates補正・Fisherの正確検定・McNemar・MH・傾向性検定は同書の範囲外であり，本章で導出した．
\item 多項分布の共分散 $\Cov(n_{ij},n_{kl})=-n\pi_{ij}\pi_{kl}$（同書 定義4.23，第4章演習）は McNemar の分散の導出の要である．
\item 制約付き最尤推定のラグランジュ乗数法は\kakomonSM{2025}{1}（ヒストグラム型多項モデル）でも問われた．
\item 超幾何分布の平均・分散（同書 式(3.9)(3.10)）が MH 検定の分散の源である．\kakomonSM{2018}{2} では有限母集団修正と超幾何分布の分散が導出された．
\item $X^2$ はスコア検定，$G^2$ は尤度比検定，Woolf の区間は Wald 型である（同書7.3節）．
\end{itemize}
\end{linkbox}

\section{例題}\label{sec:04-ex}
\begin{reidai}[$2\times2$ 表の総合問題]\label{reidai:04-1}
無作為化試験で，治療群50人中30人，対照群50人中18人が有効であった．
\begin{enumerate}[label=(\arabic*)]
\item Pearson $\chi^2$ 統計量を2通り（\cref{eq:04-pearson}の第2式と最後の式）で計算し，有意水準5\%で検定せよ（$\chi^2_1(0.05)=3.841$）．
\item リスク差の95\%信頼区間を求めよ．
\item オッズ比とその95\%信頼区間を求めよ（$\log(8/3)=0.9808$，$e^{0.1724}=1.188$，$e^{1.7893}=5.985$）．
\item リスク比とその95\%信頼区間を求めよ（$\log(5/3)=0.5108$，$e^{0.0775}=1.081$，$e^{0.9442}=2.571$）．
\end{enumerate}
\end{reidai}

\begin{reidai}[一様性と傾向性]\label{reidai:04-2}
プラセボと3用量（スコア $x=0,1,2,3$）の各40人で，有効数は $8,12,15,19$ であった．
\begin{enumerate}[label=(\arabic*)]
\item 一様性の $\chi^2$ 統計量 $X^2$ を求め，$\chi^2_3(0.05)=7.815$ と比較せよ．
\item $\hat\beta$，$\SE[\hat\beta\mid H_0]$，傾向性検定の統計量 $Y$ を求め，$\chi^2_1(0.05)=3.841$ と比較せよ．
\item 直線からのずれの統計量 $X^2-Y$ を求め，結果を解釈せよ．
\end{enumerate}
\end{reidai}

\begin{reidai}[McNemar]\label{reidai:04-3}
患者120人に2種類の治療法を左右の患部に1つずつ（左右は無作為）行い，有効・無効を判定した．
両方有効60，治療1のみ有効20，治療2のみ有効8，両方無効32であった．
\begin{enumerate}[label=(\arabic*)]
\item 有効率の差 $\hat\delta$ と，その95\%信頼区間を求めよ．
\item McNemar 検定（有意水準5\%）を行え．
\item 尤度比統計量 $G^2$ を求めよ（$\log(10/7)=0.3567$，$\log(4/7)=-0.5596$）．
\item この120組を独立な2群とみなして $\chi^2$ 検定を行うことの問題点を述べよ．
\end{enumerate}
\end{reidai}

\begin{reidai}[超幾何分布の分散の導出]\label{reidai:04-4}
周辺度数 $n_1,n_0,m_1,m_0$ を固定した $2\times2$ 表で，$H_0$ の下で $a$ は，$N$ 人から $m_1$ 人（発症者）を非復元で選ぶとき
その中に含まれる曝露者（全体で $n_1$ 人）の数とみなせる．
$\E[a]=n_1m_1/N$，$\V[a]=n_1n_0m_1m_0/\{N^2(N-1)\}$ を導け．
\end{reidai}

\section{解答}\label{sec:04-sol}
\begin{kaitou}{reidai:04-1}
$a=30,b=20,c=18,d=32$，$n_1=n_0=50$，$m_1=48$，$m_0=52$，$N=100$．\\
(1) $ad-bc=960-360=600$．$X^2=\frac{100\times600^2}{50\cdot50\cdot48\cdot52}=\frac{36{,}000{,}000}{6{,}240{,}000}=5.77$．
また $\hat\pi_1=0.60$，$\hat\pi_0=0.36$，$\hat\pi=0.48$ で
$X^2=\frac{0.24^2}{0.48\times0.52\times(2/50)}=\frac{0.0576}{0.009984}=5.77$．$5.77>3.841$ で有意．\\
(2) $\SE=\sqrt{0.6\times0.4/50+0.36\times0.64/50}=\sqrt{0.0048+0.004608}=0.0970$．
$0.24\pm1.96\times0.0970=(0.050,\,0.430)$．\\
(3) $\widehat{\OR}=960/360=2.667$．$\SE=\sqrt{1/30+1/20+1/18+1/32}=\sqrt{0.1701}=0.4125$．
$0.9808\pm0.8085=(0.1724,\,1.7893)$ → $(1.19,\,5.99)$．\\
(4) $\widehat{\RR}=0.6/0.36=1.667$．$\SE=\sqrt{1/30-1/50+1/18-1/50}=\sqrt{0.04889}=0.2211$．
$0.5108\pm0.4334=(0.0775,\,0.9442)$ → $(1.08,\,2.57)$．
OR の区間は RR の区間より1から離れた位置にあり，$\RR=1.67$ に対し $\OR=2.67$ と大きい（有効率が高いので OR は RR を近似しない）．
\end{kaitou}

\begin{kaitou}{reidai:04-2}
$N=160$，$y=54$，$p=54/160=0.3375$，$p(1-p)=0.22359$．$\bar x=1.5$，$S_{xx}=40\{(1.5)^2+(0.5)^2+(0.5)^2+(1.5)^2\}=200$．\\
(1) $p_i=0.200,0.300,0.375,0.475$．$\sum n_i(p_i-p)^2=40\{0.018906+0.001406+0.001406+0.018906\}=1.6250$．
$X^2=1.6250/0.22359=7.27<7.815$ で有意でない．\\
(2) $\sum y_i(x_i-\bar x)=8(-1.5)+12(-0.5)+15(0.5)+19(1.5)=-12-6+7.5+28.5=18$．
$\hat\beta=18/200=0.090$．$\SE[\hat\beta\mid H_0]=\sqrt{0.22359/200}=0.0334$．
$Y=0.090^2/(0.22359/200)=0.0081/0.0011180=7.25>3.841$ で有意．\\
(3) $X^2-Y=0.02$（自由度2）で，直線からのずれはほとんどない．
一様性の検定は自由度3のすべての方向の差を検出しようとするため有意にならなかったが，
差のほぼすべてが直線的な用量反応で説明され，傾向性検定では有意となる．
\end{kaitou}

\begin{kaitou}{reidai:04-3}
(1) $\hat\delta=(20-8)/120=0.100$．
$\widehat\V=\frac{28}{120^2}-\frac{12^2}{120^3}=0.0019444-0.0000833=0.0018611$，$\SE=0.0431$．
$0.100\pm1.96\times0.0431=(0.015,\,0.185)$．\\
(2) $W=(20-8)^2/(20+8)=144/28=5.14>3.841$ で有意．治療1の有効率が高い．\\
(3) $s=28$．$G^2=2\{20\log(40/28)+8\log(16/28)\}=2\{20\times0.3567+8\times(-0.5596)\}=2(7.134-4.477)=5.31$．\\
(4) 同一患者の左右は相関するので独立性の仮定が破れる．独立2群として扱うと，
$\hat\pi_1-\hat\pi_2$ の分散を $\pi_1(1-\pi_1)/n+\pi_2(1-\pi_2)/n$ と見積もり，正の相関による分散の減少を無視する（通常は保守的になり検出力を失う）．
\end{kaitou}

\begin{kaitou}{reidai:04-4}
発症者として選ばれる $m_1$ 人に番号を付け，$j$ 番目が曝露者なら1となる指示変数を $I_j$ とする．$a=\sum_{j=1}^{m_1}I_j$．
$\Prob(I_j=1)=n_1/N$ なので $\E[a]=m_1n_1/N$．
$\V[I_j]=\frac{n_1}{N}\frac{n_0}{N}$，$\Cov(I_j,I_l)=\frac{n_1(n_1-1)}{N(N-1)}-\frac{n_1^2}{N^2}=-\frac{n_1n_0}{N^2(N-1)}$ $(j\ne l)$．
\[
\V[a]=m_1\frac{n_1n_0}{N^2}-m_1(m_1-1)\frac{n_1n_0}{N^2(N-1)}=\frac{n_1n_0m_1}{N^2}\cdot\frac{N-m_1}{N-1}=\frac{n_1n_0m_1m_0}{N^2(N-1)}.
\]
\end{kaitou}

\begin{summarybox}
\begin{itemize}
\item $X^2=N(ad-bc)^2/(n_1n_0m_1m_0)$ ＝ プール分散による比率の差の $z^2$ ＝ スコア検定．
\item $\V[\log\widehat{\RR}]\approx\frac1a-\frac1{n_1}+\frac1c-\frac1{n_0}$，$\V[\log\widehat{\OR}]\approx\frac1a+\frac1b+\frac1c+\frac1d$（デルタ法）．
\item $K$ 群一様性：$X^2=\sum n_i(p_i-p)^2/\{p(1-p)\}\sim\chi^2_{K-1}$．
      傾向性：$\hat\beta=\sum y_i(x_i-\bar x)/S_{xx}$，$\V_0=p(1-p)/S_{xx}$，$Y\sim\chi^2_1$．
\item 対応あり：$\V[\hat\delta]=\{\pi_{12}+\pi_{21}-(\pi_{12}-\pi_{21})^2\}/n$，McNemar $W=(n_{12}-n_{21})^2/(n_{12}+n_{21})$，
      LRT $2[n_{12}\log\frac{2n_{12}}{s}+n_{21}\log\frac{2n_{21}}{s}]$，条件付き OR $=n_{12}/n_{21}$．
\item MH：超幾何の $\E[a_k]=n_{1k}m_{1k}/N_k$，$\V[a_k]=n_{1k}n_{0k}m_{1k}m_{0k}/\{N_k^2(N_k-1)\}$，$\widehat{\OR}_{\mathrm{MH}}=\sum a_kd_k/N_k\big/\sum b_kc_k/N_k$．
\end{itemize}
\end{summarybox}

% ===== ch05.tex =====
\chapter{診断検査とROC解析}\label{ch:05}

\begin{goalbox}
\begin{itemize}
\item 感度・特異度・陽性的中率（PPV）・陰性的中率（NPV）・有病率・尤度比を条件付き確率として定義し，Bayes の定理で相互に変換できる．
\item PPV・NPV が有病率に依存すること，罹患者・非罹患者の人数を固定した研究では標本の PPV がそのまま使えないことを説明できる．
\item カットオフ値を変えたときの感度と特異度のトレードオフを，正規モデルで計算できる．
\item 経験ROC曲線を描き，台形公式で AUC を計算し，Youden 指数を最大にするカットオフを選べる．
\item AUC $=\Prob(X_1>X_0)$ を示し，経験AUCが Mann--Whitney 統計量に一致することを導出できる．
\item 同一被験者に行った2検査の感度（McNemar型）と PPV（多項分布＋デルタ法）を比較できる．
\end{itemize}
\end{goalbox}

\begin{exambox}
\begin{itemize}
\item \kakomon{2018}{3}：ロジスティック回帰の予測値による判定の真陽性率・偽陽性率，経験ROC曲線の描画，「真陽性率$-$偽陽性率」最大のカットオフ，
      ROC上の長方形の面積と AUC が $\frac{1}{8\times8}\sum\sum I(x_{1i}>x_{0j})$ で表されることの証明．
\item \kakomon{2018}{4}〔5〕：リスクスコアの AUC $0.70$（95\%CI 0.67--0.72）と医師の予測の ROC の比較の解釈．
\item \kakomon{2019}{3}：同一被験者200名での2検査の感度・PPV，感度の比較（不一致ペアの二項検定），PPV の比較（8セル多項分布の $\bm D-\bm\pi\bm\pi\transp$ とデルタ法）．
\item \kakomon{2021}{5}（共通）：正規モデルの感度・特異度，有病率1\%での PPV・NPV，感度を固定したカットオフ，「感度1で陰性的中率1」の検査が有効とは限らない理由．
\item \kakomon{2023}{1}〔5〕：ロジスティックモデルの予測確率 $\ge0.6$ による判定の PPV・NPV．
\item \kakomon{2024}{1}：4つのカットオフでの感度・特異度・PPV・NPV の表の完成，折れ線 ROC と台形公式による AUC，対応のある2検査の McNemar 検定．
\end{itemize}
\end{exambox}

\flowline{検査値と\\確定診断}{感度・特異度\\PPV・AUC}{2×2表\\ROC・McNemar}{参照標準が正確\\独立な被験者}{Bayes・台形\\デルタ法}{有病率への依存\\カットオフの選択}

\section{問題設定}\label{sec:05-setting}
各被験者について，確定診断（参照標準，ゴールドスタンダード）による疾患の有無 $D\in\{1,0\}$ と，
検査値 $X$（連続）またはその判定結果（陽性 $+$／陰性 $-$）が得られる．
$X\ge c$ を陽性とする判定を考え，罹患者・非罹患者の $X$ の分布関数を $F_1,F_0$ とする．
\begin{table}[htbp]
\centering\small
\caption{診断検査の $2\times2$ 表}\label{tab:05-2x2}
\begin{tabular}{@{}lccc@{}}
\toprule
 & 罹患 $D=1$ & 非罹患 $D=0$ & 計\\
\midrule
陽性 $+$ & TP（真陽性） & FP（偽陽性） & TP$+$FP\\
陰性 $-$ & FN（偽陰性） & TN（真陰性） & FN$+$TN\\
\midrule
計 & TP$+$FN & FP$+$TN & $N$\\
\bottomrule
\end{tabular}
\end{table}

\section{基本概念}\label{sec:05-basic}
\begin{teigi}[診断性能の指標]\label{def:05-measures}
\begin{center}\small
\renewcommand{\arraystretch}{1.9}
\begin{tabular}{@{}ll@{\quad}ll@{}}
感度（TPF） & $\Prob(+\mid D=1)\approx\dfrac{\mathrm{TP}}{\mathrm{TP}+\mathrm{FN}}$ &
特異度 & $\Prob(-\mid D=0)\approx\dfrac{\mathrm{TN}}{\mathrm{FP}+\mathrm{TN}}$\\
偽陽性率（FPF） & $1-\text{特異度}=\Prob(+\mid D=0)$ &
有病率 & $p=\Prob(D=1)$\\
陽性的中率（PPV） & $\Prob(D=1\mid +)\approx\dfrac{\mathrm{TP}}{\mathrm{TP}+\mathrm{FP}}$ &
陰性的中率（NPV） & $\Prob(D=0\mid -)\approx\dfrac{\mathrm{TN}}{\mathrm{FN}+\mathrm{TN}}$\\
陽性尤度比 $\mathrm{LR}^+$ & $\dfrac{\text{感度}}{1-\text{特異度}}$ &
陰性尤度比 $\mathrm{LR}^-$ & $\dfrac{1-\text{感度}}{\text{特異度}}$
\end{tabular}
\end{center}
\end{teigi}
\Youchui 感度・特異度は「疾患の状態を条件とした検査結果の確率」，PPV・NPV は「検査結果を条件とした疾患の確率」である．
臨床で患者が知りたいのは後者だが，後者は有病率に依存する．

\begin{meidai}[Bayes の定理による的中率]\label[meidai]{prop:05-bayes}
感度 $Se$，特異度 $Sp$，有病率 $p$ とすると
\begin{equation}
\mathrm{PPV}=\frac{Se\,p}{Se\,p+(1-Sp)(1-p)},\qquad
\mathrm{NPV}=\frac{Sp\,(1-p)}{Sp\,(1-p)+(1-Se)\,p}.\label{eq:05-ppv}
\end{equation}
オッズで書くと $\dfrac{\mathrm{PPV}}{1-\mathrm{PPV}}=\mathrm{LR}^+\times\dfrac{p}{1-p}$（検査後オッズ＝尤度比×検査前オッズ）．
\end{meidai}
\begin{proof}
$\Prob(D=1\mid+)=\Prob(+\mid D=1)\Prob(D=1)/\{\Prob(+\mid D=1)\Prob(D=1)+\Prob(+\mid D=0)\Prob(D=0)\}$．NPV も同様．
\end{proof}
\paragraph{有病率への依存}
$p\to0$ なら $\mathrm{PPV}\to0$，$\mathrm{NPV}\to1$ である．
\kakomon{2021}{5} では $Se=0.841$，$Sp=0.933$ でも $p=0.01$ だと $\mathrm{PPV}\approx0.113$ にすぎない．
「感度がほぼ1で，陰性者に感染者がいない（NPV $\approx1$）」という主張は，
閾値を下げて何でも陽性にすれば達成でき，偽陽性が増えて PPV が下がるので，それだけで有効な検査とは言えない．

\paragraph{研究デザインと的中率}
\Youchui 罹患者20人・非罹患者20人のように\textbf{疾患の状態ごとに人数を固定}した研究（\kakomon{2024}{1}）では，
感度・特異度は推定できるが，標本の PPV・NPV は標本内の有病率（ここでは0.5）を反映した値にすぎない．
実際の対象集団の有病率 $p$ での的中率は\cref{eq:05-ppv}で計算し直す必要がある．
\kakomon{2019}{3} のように被験者全員に検査と確定診断を行い，被験者が対象集団を代表する標本であれば，標本の PPV は母集団の PPV の推定量になる．

\section{数理：カットオフと ROC 曲線}\label{sec:05-roc}
\subsection{正規モデルとカットオフ}\label{subsec:05-binormal}
$X\mid D=1\sim\Normal(\mu_1,\sigma_1^2)$，$X\mid D=0\sim\Normal(\mu_0,\sigma_0^2)$，$\mu_1>\mu_0$ とする（2正規モデル）．$X\ge c$ を陽性とすると
\[
Se(c)=1-\Phi\Bigl(\frac{c-\mu_1}{\sigma_1}\Bigr),\qquad Sp(c)=\Phi\Bigl(\frac{c-\mu_0}{\sigma_0}\Bigr).
\]
$c$ を上げると特異度は上がり感度は下がる．
感度を $s$ に固定するなら $c=\mu_1-z_{1-s}\sigma_1$（例：$s=0.95$ なら $c=\mu_1-1.645\sigma_1$）．
両方を同時に $0.95$ にするには $(\mu_1-c)/\sigma_1=(c-\mu_0)/\sigma_0=1.645$ が必要で，
等分散なら $c=(\mu_0+\mu_1)/2$，$\sigma=(\mu_1-\mu_0)/(2\times1.645)$（\kakomon{2021}{5}〔4〕）．

\subsection{ROC 曲線}\label{subsec:05-roccurve}
\begin{teigi}[ROC 曲線]\label{def:05-roc}
カットオフ $c$ を $+\infty$ から $-\infty$ まで動かしたときの点 $(\mathrm{FPF}(c),\mathrm{TPF}(c))=(1-F_0(c),\,1-F_1(c))$ の軌跡を
ROC（receiver operating characteristic）曲線という．$(0,0)$ から $(1,1)$ に至る非減少曲線である．
\end{teigi}
対角線 $\mathrm{TPF}=\mathrm{FPF}$ は判別力のない検査（コイン投げ）に対応し，左上 $(0,1)$ に近いほど良い．

\paragraph{経験 ROC 曲線の描き方（\kakomon{2018}{3}〔2〕）}
罹患者 $n_1$ 人，非罹患者 $n_0$ 人の検査値（または予測確率）を大きい順に並べ，
原点から出発して，罹患者なら上へ $1/n_1$，非罹患者なら右へ $1/n_0$ 進む（同じ値の罹患者と非罹患者があれば斜めに進む）．
最終的に $(1,1)$ に到達する階段状の曲線が経験ROC曲線である．
カットオフが有限個しか与えられていない場合（\kakomon{2024}{1}）は，与えられた点を $(0,0)$，$(1,1)$ と折れ線で結ぶ．

\subsection{Youden 指数}\label{subsec:05-youden}
\begin{teigi}[Youden 指数]\label{def:05-youden}
$J(c)=Se(c)+Sp(c)-1=\mathrm{TPF}(c)-\mathrm{FPF}(c)$．$J$ を最大にするカットオフを最適カットオフの1つの基準とする．
\end{teigi}
ROC 上の点 $(u,v)$ から対角線への距離は $(v-u)/\sqrt2$ なので，$J$ の最大化は ROC 曲線上で対角線から最も離れた点を選ぶことである
（\kakomon{2018}{3}〔3〕）．2正規・等分散モデルでは $J$ を最大にする $c$ は $(\mu_0+\mu_1)/2$ である
（$\frac{\dd}{\dd c}J=-f_1(c)+f_0(c)=0$（$f_1,f_0$ は罹患者・非罹患者の密度）より2つの密度が交わる点）．

\section{AUC とその確率的解釈}\label{sec:05-auc}
\Hissu\Hinshutsu
\begin{teiri}[AUC $=\Prob(X_1>X_0)$]\label[teiri]{thm:05-auc}
$X_1\sim F_1$，$X_0\sim F_0$ が独立な連続確率変数なら，ROC 曲線下面積は
\begin{equation}
\mathrm{AUC}=\Prob(X_1>X_0).\label{eq:05-auc}
\end{equation}
\end{teiri}
\begin{proof}
$u=\mathrm{FPF}(c)=1-F_0(c)$ とおくと $\dd u=-f_0(c)\,\dd c$ で，$c:+\infty\to-\infty$ が $u:0\to1$ に対応する．
\[
\mathrm{AUC}=\int_0^1\mathrm{TPF}\,\dd u=\int_{-\infty}^{\infty}\{1-F_1(c)\}f_0(c)\,\dd c=\int\Prob(X_1>c)f_0(c)\,\dd c=\Prob(X_1>X_0).
\]
\end{proof}
2正規モデルでは $X_1-X_0\sim\Normal(\mu_1-\mu_0,\sigma_1^2+\sigma_0^2)$ より
$\mathrm{AUC}=\Phi\bigl((\mu_1-\mu_0)/\sqrt{\sigma_1^2+\sigma_0^2}\bigr)$．

\begin{meidai}[経験 AUC と Mann--Whitney 統計量（\kakomon{2018}{3}〔4〕〔5〕）]\label[meidai]{prop:05-empauc}
罹患者の値 $x_{11},\dots,x_{1n_1}$，非罹患者の値 $x_{01},\dots,x_{0n_0}$ に同順位がなければ，経験ROC曲線の下の面積は
\begin{equation}
\widehat{\mathrm{AUC}}=\frac{1}{n_1n_0}\sum_{i=1}^{n_1}\sum_{j=1}^{n_0}I(x_{1i}>x_{0j})=\frac{U}{n_1n_0}.\label{eq:05-empauc}
\end{equation}
同順位がある場合は $I(x_{1i}>x_{0j})+\frac12I(x_{1i}=x_{0j})$ を用いる（斜めの線分の下の台形が $1/2$ を与える）．
\end{meidai}
\begin{proof}
経験ROC曲線の右方向の各ステップは1人の非罹患者 $j$ に対応し，幅 $1/n_0$，そのときの高さは
「$x_{0j}$ より大きい値をもつ罹患者の割合」$\frac{1}{n_1}\sum_iI(x_{1i}>x_{0j})$ である．
曲線下面積はこれらの長方形の和であり，\cref{eq:05-empauc}を得る．
（曲線の上側を罹患者ごとの長方形に分けると $\frac{1}{n_1n_0}\sum\sum I(x_{1i}<x_{0j})$ となり，$1$ からこれを引いても同じ結果になる．）
\end{proof}
したがって経験AUCは Mann--Whitney 統計量 $U$ の定数倍（Wilcoxon 順位和 $W$ の1次関数；\cref{sec:03-ranksum}）であり，
「AUC $=0.5$」の検定は順位和検定そのものである．

\paragraph{台形公式}
与えられた点 $(u_0,v_0)=(0,0),(u_1,v_1),\dots,(u_K,v_K)=(1,1)$ を折れ線で結ぶと
$\mathrm{AUC}=\sum_{k=1}^K\frac{(v_{k-1}+v_k)(u_k-u_{k-1})}{2}$．
\kakomon{2024}{1} では点 $(0,0),(0.15,0.60),(0.60,0.90),(1,1)$ から $\mathrm{AUC}=0.045+0.3375+0.38=0.7625$．

\begin{figure}[htbp]
\centering
\begin{tikzpicture}
\begin{axis}[width=62mm,height=62mm,xmin=0,xmax=1,ymin=0,ymax=1,
  xlabel={偽陽性率 FPF},ylabel={真陽性率 TPF},xtick={0,0.2,0.4,0.6,0.8,1},ytick={0,0.25,0.5,0.75,1},
  grid=major,grid style={gray!20},tick label style={font=\scriptsize},label style={font=\small}]
\addplot[dashed,gray] coordinates {(0,0)(1,1)};
\addplot[very thick,bkblue] coordinates {(0,0)(0,0.25)(0.2,0.25)(0.2,0.75)(0.4,0.75)(0.4,1)(1,1)};
\addplot[only marks,mark=*,bkorange,mark size=2pt] coordinates {(0.4,1)};
\node[font=\scriptsize,anchor=west] at (axis cs:0.42,0.93) {$J$ 最大（$X\ge4$）};
\end{axis}
\end{tikzpicture}
\caption{\cref{reidai:05-2}の経験ROC曲線（AUC $=0.80$）}\label{fig:05-roc}
\end{figure}

\section{推定・検定：2検査の比較}\label{sec:05-compare}
\subsection{対応のある感度の比較}\label{subsec:05-sens}
同一の罹患者に検査A・Bを行った場合，感度の差は罹患者での対応のある2 値データの周辺比率の差である．
不一致ペア（A陽性・B陰性の $n_{+-}$，A陰性・B陽性の $n_{-+}$）だけが情報を持ち，McNemar 検定（\cref{sec:04-mcnemar}）
\[
z=\frac{n_{+-}-s/2}{\sqrt{s/4}},\quad s=n_{+-}+n_{-+}
\]
で検定する（\kakomon{2019}{3}〔2〕，\kakomon{2024}{1}〔3〕）．特異度は非罹患者で同様に比較する．

\subsection{対応のある PPV の比較}\label{subsec:05-ppv}
\Youchui PPV は「検査陽性者」を分母とするので，罹患・非罹患の両方の表にまたがる．
\kakomon{2019}{3} では，罹患者・非罹患者それぞれの (A,B) の $2\times2$ 表の計8セルを，総数 $n$ を固定した多項分布とみなす．
セル確率ベクトル $\bm\pi$，標本比率 $\bm p=\bm x/n$ に対し，多変量中心極限定理から
\[
\sqrt n(\bm p-\bm\pi)\xrightarrow{d}\Normal_8\bigl(\bm0,\ \bm D_{\bm\pi}-\bm\pi\bm\pi\transp\bigr),\qquad \bm D_{\bm\pi}=\mathrm{diag}(\bm\pi).
\]
$\mathrm{PPV}_A=\dfrac{\pi_{+\bullet D}}{\pi_{+\bullet D}+\pi_{+\bullet\bar D}}$，$\mathrm{PPV}_B=\dfrac{\pi_{\bullet+D}}{\pi_{\bullet+D}+\pi_{\bullet+\bar D}}$，
$f(\bm\pi)=\mathrm{PPV}_A-\mathrm{PPV}_B$ とすると，デルタ法（\cref{thm:04-delta}）より
\[
\V\bigl[\sqrt n f(\bm p)\bigr]\approx\sigma^2=\nabla f\transp(\bm D_{\bm\pi}-\bm\pi\bm\pi\transp)\nabla f.
\]
成分まで計算すると，$q_A=\pi_{+\bullet D}+\pi_{+\bullet\bar D}$（Aが陽性の確率），$q_B=\pi_{\bullet+D}+\pi_{\bullet+\bar D}$ として
\begin{equation}
\begin{split}
\sigma^2={}&\frac{\mathrm{PPV}_A(1-\mathrm{PPV}_A)}{q_A}+\frac{\mathrm{PPV}_B(1-\mathrm{PPV}_B)}{q_B}\\
&-2\,\frac{\pi_{++D}(1-\mathrm{PPV}_A)(1-\mathrm{PPV}_B)+\pi_{++\bar D}\,\mathrm{PPV}_A\mathrm{PPV}_B}{q_Aq_B}.
\end{split}\label{eq:05-ppvvar}
\end{equation}
第3項は両検査が陽性の被験者が両方の PPV の分母に入ることによる共分散である．
検定統計量は $z=\sqrt n f(\bm p)/\hat\sigma$ で，\kakomon{2019}{3} の数値では $f(\bm p)=-0.1$，$\hat\sigma^2=0.42$，$n=200$ から $z=-2.18$（有意）となる．

\begin{supbox}[発展：AUC の比較]
2つのマーカーの AUC の比較には，Mann--Whitney 型推定量の分散・共分散をU統計量の理論で推定する方法（DeLong ら \cite{delong1988}）が標準的である．
同一被験者で測定した2マーカーの AUC は正の相関をもつので，独立とみなすと検定は保守的になる．
\end{supbox}

\section{仮定}\label{sec:05-assump}
\begin{itemize}
\item 参照標準が正確であること（誤分類があると感度・特異度は偏る）．
\item 検査を受けた全員に参照標準が実施されていること（一部だけだと検証バイアス）．
\item 被験者が独立で，対象集団を代表していること．PPV・NPV を標本から直接推定するには，有病率も代表的である必要がある．
\item 予測モデルの性能は，モデル作成に使ったデータで評価すると楽観的に偏る．独立な検証コホートで評価する（\kakomon{2018}{4}）．
\end{itemize}

\section{結果の解釈}\label{sec:05-interp}
\begin{itemize}
\item 感度が高い検査は，陰性なら疾患を否定しやすい（除外診断，SnNout）．特異度が高い検査は，陽性なら疾患を確定しやすい（確定診断，SpPin）．
      ただしこれらは経験的な覚え方（heuristic mnemonic）にすぎない．検査結果による判断の変化の大きさは尤度比 $\mathrm{LR}^+=Se/(1-Sp)$，$\mathrm{LR}^-=(1-Se)/Sp$ で決まり，      どちらも感度と特異度の両方に依存する（感度が高くても特異度が低ければ $\mathrm{LR}^-$ は1から大きく離れない）．さらに検査前確率（有病率）にも依存する．
\item PPV は有病率が低い集団（検診）では小さくなる．同じ検査でも専門外来（有病率が高い）では PPV が高い．
\item AUC は「無作為に選んだ罹患者の検査値が，無作為に選んだ非罹患者の値より大きい確率」．
      $0.5$ は判別力なし，$1$ は完全判別．$0.70$ 程度は中等度の判別能であり（\kakomon{2018}{4}），信頼区間が0.5を含まないことと臨床的に十分であることは別である．
\item ROC 曲線が交差する2検査では，AUC だけでなく，臨床的に重要な FPF 範囲での感度を比べる．
\end{itemize}

\section{手法選択}\label{sec:05-choice}
\begin{itemize}
\item 単一検査・単一カットオフ → $2\times2$ 表から感度・特異度（二項の区間推定），有病率がわかれば\cref{eq:05-ppv}．
\item 連続マーカー → ROC 曲線と AUC．カットオフは Youden 指数や臨床的要件（感度 $\ge0.95$ など）で決める．
\item 同一被験者で2検査 → 感度・特異度は McNemar，PPV・NPV は多項分布＋デルタ法，AUC は DeLong 法．
\item 独立な被験者群で2検査 → 2群の比率の比較（\cref{ch:04}）．
\end{itemize}

\section{よくある誤り}\label{sec:05-err}
\begin{warnbox}
\begin{itemize}
\item 感度と PPV を取り違える（条件付けの向きの誤り）．
\item 罹患者・非罹患者の人数を固定した研究で，標本の PPV を実際の集団の PPV として報告する．
\item 経験ROC曲線を描くとき，罹患者で右，非罹患者で上に進む（逆）．
\item AUC の Mann--Whitney 表現で不等号の向きを逆にする（ROC の上側の面積になる）．
\item 同一被験者の2検査の感度を，独立2群の $\chi^2$ 検定で比べる．
\end{itemize}
\end{warnbox}

\section{数理統計との接続}\label{sec:05-link}
\begin{linkbox}
\begin{itemize}
\item \cref{thm:05-auc}は確率積分変換と同じ置換積分である．AUC の推定量は U 統計量で，順位和検定の統計量と一致する（\cref{sec:03-ranksum}）．
\item 多項分布の共分散行列 $\bm D_{\bm\pi}-\bm\pi\bm\pi\transp$（『現代数理統計学の基礎』4.4.2項）と多変量デルタ法（同書は1次元のみ；\cref{thm:04-delta}）の組み合わせは，
      比率の関数の分散を求める万能の方法である．
\item \textbf{等分散の}2正規モデルでは，尤度比 $f_1(x)/f_0(x)$ が $x$ の単調増加関数であることから，Neyman--Pearson の補題（同書 定理7.16）により
      「$X\ge c$ で陽性」がある特異度の下で感度を最大にする判定になる（$\sigma_1\ne\sigma_0$ では尤度比が単調でなく，この判定は最適とは限らない）．ROC 曲線は検定の「サイズ（FPF）に対する検出力（TPF）」の曲線にほかならない．
\end{itemize}
\end{linkbox}

\section{例題}\label{sec:05-ex}
\begin{reidai}[正規モデルと有病率]\label{reidai:05-1}
罹患者の検査値は $\Normal(12,2^2)$，非罹患者は $\Normal(8,2^2)$ に従い，$X\ge11$ を陽性とする．
$\Phi(0.5)=0.6915$，$\Phi(1)=0.8413$，$\Phi(1.5)=0.9332$，$\Phi(1.414)=0.9214$ を用いよ．
\begin{enumerate}[label=(\arabic*)]
\item 感度・特異度・陽性尤度比を求めよ．
\item 有病率5\%の集団での PPV・NPV を求めよ．
\item Youden 指数を最大にするカットオフとそのときの $J$ を求めよ．
\item AUC を求めよ．
\end{enumerate}
\end{reidai}

\begin{reidai}[経験 ROC と AUC]\label{reidai:05-2}
罹患者4人の検査値が $9,7,6,4$，非罹患者5人の値が $8,5,3,2,1$ である．値 $x$ 以上を陽性とする．
\begin{enumerate}[label=(\arabic*)]
\item 経験ROC曲線の頂点の座標を列挙せよ．
\item AUC を面積として計算し，$\frac{1}{n_1n_0}\sum\sum I(x_{1i}>x_{0j})$ と一致することを確かめよ．
\item Youden 指数を最大にするカットオフを求めよ．
\end{enumerate}
\end{reidai}

\begin{reidai}[研究デザインと PPV]\label{reidai:05-3}
罹患者100人と非罹患者100人を集めた評価研究で，新検査の陽性は罹患者90人，非罹患者20人であった．
\begin{enumerate}[label=(\arabic*)]
\item 感度，特異度，この標本での陽性的中率を求めよ．
\item この検査を有病率2\%の検診集団に用いるときの PPV を求めよ．
\item (1)と(2)の PPV が大きく異なる理由を述べよ．
\end{enumerate}
\end{reidai}

\begin{reidai}[対応のある感度の比較]\label{reidai:05-4}
確定診断で罹患と判定された80人に検査A・Bを行った．両方陽性50，Aのみ陽性14，Bのみ陽性4，両方陰性12であった．
\begin{enumerate}[label=(\arabic*)]
\item 検査A・Bの感度を求めよ．
\item 感度が等しいという帰無仮説を有意水準5\%で検定せよ．
\item 感度の差の95\%信頼区間を求めよ．
\end{enumerate}
\end{reidai}

\section{解答}\label{sec:05-sol}
\begin{kaitou}{reidai:05-1}
(1) $Se=\Prob(Z>(11-12)/2)=\Phi(0.5)=0.6915$，$Sp=\Phi((11-8)/2)=\Phi(1.5)=0.9332$．
$\mathrm{LR}^+=0.6915/0.0668=10.35$．\\
(2) $\mathrm{PPV}=\dfrac{0.6915\times0.05}{0.6915\times0.05+0.0668\times0.95}=\dfrac{0.03458}{0.03458+0.06346}=0.353$．
$\mathrm{NPV}=\dfrac{0.9332\times0.95}{0.9332\times0.95+0.3085\times0.05}=\dfrac{0.8865}{0.8865+0.0154}=0.983$．\\
(3) 等分散なので $c=(8+12)/2=10$．$Se=Sp=\Phi(1)=0.8413$，$J=0.683$．\\
(4) $\mathrm{AUC}=\Phi(4/\sqrt{8})=\Phi(1.414)=0.921$．
\end{kaitou}

\begin{kaitou}{reidai:05-2}
(1) 値の大きい順に $9(D),8(\bar D),7(D),6(D),5(\bar D),4(D),3,2,1(\bar D)$．
頂点は $(0,0)\to(0,0.25)\to(0.2,0.25)\to(0.2,0.75)\to(0.4,0.75)\to(0.4,1)\to(1,1)$（\cref{fig:05-roc}）．\\
(2) 面積：$0.2\times0.25+0.2\times0.75+0.6\times1=0.05+0.15+0.60=0.80$．
対の数え上げ：$9$ は5人全員より大，$7$ と $6$ はそれぞれ4人より大，$4$ は3人より大で，計 $16$．$16/20=0.80$ で一致．\\
(3) 各頂点の $J=\mathrm{TPF}-\mathrm{FPF}$ は $0.25,\ 0.05,\ 0.55,\ 0.35,\ 0.60$．最大は $(0.4,1)$ で，カットオフ「$4$ 以上を陽性」．
\end{kaitou}

\begin{kaitou}{reidai:05-3}
(1) $Se=0.90$，$Sp=80/100=0.80$，標本の PPV $=90/110=0.818$．\\
(2) $\mathrm{PPV}=\dfrac{0.9\times0.02}{0.9\times0.02+0.2\times0.98}=\dfrac{0.018}{0.214}=0.084$．\\
(3) 標本では有病率が $0.5$ に固定されている．PPV は有病率に強く依存し，有病率2\%では陽性者の大部分が偽陽性になる．
尤度比 $\mathrm{LR}^+=4.5$ を使えば，検査後オッズ $=4.5\times(0.02/0.98)=0.0918$，PPV $=0.0918/1.0918=0.084$ とも計算できる．
\end{kaitou}

\begin{kaitou}{reidai:05-4}
(1) $Se_A=(50+14)/80=0.800$，$Se_B=(50+4)/80=0.675$．\\
(2) 不一致ペア $s=18$．$W=(14-4)^2/18=5.56>3.841$ で有意．検査Aの感度が高い．\\
(3) $\hat\delta=10/80=0.125$．$\widehat\V=\frac{18}{80^2}-\frac{10^2}{80^3}=0.0028125-0.0001953=0.0026172$，$\SE=0.0512$．
$0.125\pm1.96\times0.0512=(0.025,\,0.225)$．
\end{kaitou}

\begin{summarybox}
\begin{itemize}
\item 感度・特異度は $D$ を条件，PPV・NPV は検査結果を条件とする確率．PPV $=Se\,p/\{Se\,p+(1-Sp)(1-p)\}$．
\item 検査後オッズ＝尤度比×検査前オッズ．PPV は有病率が低いと小さい．
\item 疾患状態ごとに人数を固定した研究では標本の PPV・NPV は使えない（有病率で計算し直す）．
\item ROC：$(\mathrm{FPF},\mathrm{TPF})$ の軌跡．経験ROCは罹患者で上，非罹患者で右．台形公式で AUC．
\item AUC $=\Prob(X_1>X_0)$，経験 AUC $=\frac{1}{n_1n_0}\sum\sum I(x_{1i}>x_{0j})$（Mann--Whitney）．2正規なら $\Phi(\Delta\mu/\sqrt{\sigma_1^2+\sigma_0^2})$．
\item Youden $J=Se+Sp-1$ は対角線から最も離れた点．
\item 同一被験者の2検査：感度は McNemar，PPV は8セル多項分布＋デルタ法．
\end{itemize}
\end{summarybox}

% ===== ch06.tex =====
\chapter{ロジスティック回帰}\label{ch:06}

\begin{goalbox}
\begin{itemize}
\item ロジスティック回帰モデルを書き，回帰係数を（調整）対数オッズ比として解釈できる．
\item 個別データ・集計データの対数尤度，スコア，Fisher情報行列を導出できる．
\item 2群比較のロジスティック回帰の最尤推定値と標準誤差が $2\times2$ 表の対数オッズ比と Woolf の分散に一致することを示せる．
\item Wald 検定・尤度比検定・スコア検定を使い分け，尤度比検定の自由度を正しく数えられる．
\item 交互作用項のあるモデルでオッズ比を計算できる．予測確率から判定表・PPV・NPV・ROC を作れる．
\item AIC の定義と Kullback--Leibler 情報量との関係を説明できる．
\end{itemize}
\end{goalbox}

\begin{exambox}
\begin{itemize}
\item \kakomon{2018}{4}：ロジスティックモデルの対数尤度，二値共変量の係数が調整オッズ比の対数であることの証明，AIC の式，Bernoulli分布間の KL 情報量，リスクスコアの解釈（150字）．
\item \kakomon{2023}{1}：治療群の対数オッズとオッズ比の表現，集計データの対数尤度 $l=c_1\alpha+c_2\beta+c_3\log(1+e^\alpha)+c_4\log(1+e^{\alpha+\beta})$ の係数，
      2群モデルの最尤推定値 $\hat\alpha=\log(4/6)$，$\hat\beta=\log(7/2)$，年齢を加えたモデルとの尤度比検定，予測確率 $\ge0.6$ による判定の PPV・NPV．
\item \kakomon{2019}{2}〔3〕：傾向スコアを飽和ロジスティックモデル（交互作用項つき）で推定．
\item \kakomon{2018}{3}：ロジスティックモデルの予測値による ROC 解析（\cref{ch:05}）．
\end{itemize}
\end{exambox}

\flowline{2 値アウトカム\\＋共変量}{調整OR\\予測確率}{ロジスティック\\回帰}{logit 線形\\独立}{最尤法\\Wald・LRT}{調整ORの解釈\\判別・較正}

\section{問題設定}\label{sec:06-setting}
被験者 $i=1,\dots,n$ について，2 値アウトカム $y_i\in\{0,1\}$ と共変量ベクトル $\bm x_i=(1,x_{i1},\dots,x_{ip})\transp$ が得られる．
目的は (a) 他の共変量を調整した曝露・治療の効果（調整オッズ比）の推定，(b) 発症確率の予測，のいずれかである．
2群比較や層別表では足りない「複数の共変量（連続量を含む）の同時調整」を行うのがロジスティック回帰である．

\section{モデルと係数の解釈}\label{sec:06-model}
\begin{teigi}[ロジスティック回帰モデル]\label{def:06-logit}
$\pi_i=\Prob(Y_i=1\mid\bm x_i)$ について
\begin{equation}
\logit\pi_i=\log\frac{\pi_i}{1-\pi_i}=\bm x_i\transp\bm\beta=\beta_0+\beta_1x_{i1}+\dots+\beta_px_{ip},\qquad
\pi_i=\frac{e^{\bm x_i\transp\bm\beta}}{1+e^{\bm x_i\transp\bm\beta}}.\label{eq:06-logit}
\end{equation}
\end{teigi}
\begin{meidai}[係数は調整対数オッズ比（\kakomon{2018}{4}〔2〕）]\label[meidai]{prop:06-coef}
$x_1$ が2 値（0/1）のとき，他の共変量 $x_2,\dots,x_p$ を固定すると
\[
\log\left\{\frac{\Prob(Y=1\mid x_1=1,\cdot)/\Prob(Y=0\mid x_1=1,\cdot)}{\Prob(Y=1\mid x_1=0,\cdot)/\Prob(Y=0\mid x_1=0,\cdot)}\right\}=\beta_1.
\]
$x_1$ が連続なら，$x_1$ が1単位増えるごとのオッズ比が $e^{\beta_1}$，$k$ 単位なら $e^{k\beta_1}$．
\end{meidai}
\begin{proof}
各条件での対数オッズは $\bm x\transp\bm\beta$ なので，その差は $(\beta_0+\beta_1+\sum_{j\ge2}\beta_jx_j)-(\beta_0+\sum_{j\ge2}\beta_jx_j)=\beta_1$．
\end{proof}
「調整」とは「他の共変量が同じ値の人どうしで比べた」という意味であり，
このモデルは他の共変量の値によらずオッズ比が一定（交互作用なし）であることを仮定している．

\paragraph{カテゴリ変数}
$L$ 水準のカテゴリ変数は参照水準を除く $L-1$ 個のダミー変数で表す．各係数は参照水準に対する対数オッズ比である
（\kakomon{2018}{4} の年齢区分 [8--12]/[5--7]，[13--18]/[5--7]）．
このときカテゴリ変数全体の効果の検定は自由度 $L-1$ になる．

\paragraph{交互作用}
$\logit\pi=\beta_0+\beta_1x+\beta_2z+\beta_3xz$ では，$z$ を固定したときの $x$ のオッズ比は $e^{\beta_1+\beta_3z}$ で，$z$ によって変わる．
$\beta_3\ne0$ は「$z$ による効果修飾」（\cref{sec:07-modification}）をオッズ比の尺度で表したものである．

\section{数理：尤度・スコア・情報行列}\label{sec:06-likelihood}
\subsection{対数尤度}
\begin{meidai}[対数尤度（\kakomon{2018}{4}〔1〕）]\label[meidai]{prop:06-loglik}
$Y_i$ が独立に $\mathrm{Bernoulli}(\pi_i)$ に従うとき，
\begin{equation}
\ell(\bm\beta)=\sum_{i=1}^n\{y_i\log\pi_i+(1-y_i)\log(1-\pi_i)\}=\sum_{i=1}^n\Bigl\{y_i\bm x_i\transp\bm\beta-\log\bigl(1+e^{\bm x_i\transp\bm\beta}\bigr)\Bigr\}.\label{eq:06-loglik}
\end{equation}
\end{meidai}
\begin{proof}
$\log\pi_i=\bm x_i\transp\bm\beta-\log(1+e^{\bm x_i\transp\bm\beta})$，$\log(1-\pi_i)=-\log(1+e^{\bm x_i\transp\bm\beta})$ を代入する．
\end{proof}
\paragraph{集計データ}
共変量パターン $g$ ごとに $n_g$ 人中 $y_g$ 人がイベントを起こしたとき（二項），定数を除いて
$\ell=\sum_g\{y_g\,\bm x_g\transp\bm\beta-n_g\log(1+e^{\bm x_g\transp\bm\beta})\}$．
\kakomon{2023}{1} の2群（$T$：10人中7人，$C$：10人中4人，$\logit\pi_C=\alpha$，$\logit\pi_T=\alpha+\beta$）では
\[
\ell(\alpha,\beta)=7(\alpha+\beta)+4\alpha-10\log(1+e^{\alpha+\beta})-10\log(1+e^{\alpha})=11\alpha+7\beta-10\log(1+e^\alpha)-10\log(1+e^{\alpha+\beta}).
\]

\subsection{スコアと情報行列}
\begin{meidai}\label[meidai]{prop:06-score}
\begin{align}
\bm U(\bm\beta)&=\frac{\partial\ell}{\partial\bm\beta}=\sum_i\bm x_i(y_i-\pi_i)=\bm X\transp(\bm y-\bm\pi),\label{eq:06-score}\\
\bm I(\bm\beta)&=-\frac{\partial^2\ell}{\partial\bm\beta\partial\bm\beta\transp}=\sum_i\pi_i(1-\pi_i)\bm x_i\bm x_i\transp=\bm X\transp\bm W\bm X,\label{eq:06-info}
\end{align}
$\bm W=\mathrm{diag}\{\pi_i(1-\pi_i)\}$．$\bm I$ は $\bm y$ に依存しないので観測情報＝期待情報である．
\end{meidai}
\begin{proof}
$\frac{\partial}{\partial\bm\beta}\log(1+e^{\bm x\transp\bm\beta})=\bm x\,\pi$，$\frac{\partial\pi}{\partial\bm\beta\transp}=\pi(1-\pi)\bm x\transp$ による．
\end{proof}
尤度方程式 $\bm X\transp(\bm y-\hat{\bm\pi})=\bm0$ は一般に陽に解けず，Newton--Raphson 法
$\bm\beta^{(t+1)}=\bm\beta^{(t)}+\bm I(\bm\beta^{(t)})^{-1}\bm U(\bm\beta^{(t)})$（反復重み付き最小二乗法；IRLS）で解く．
切片を含むモデルでは第1成分から $\sum_iy_i=\sum_i\hat\pi_i$（観測イベント数＝予測イベント数）が成り立つ．
MLE は漸近的に $\hat{\bm\beta}\approx\Normal(\bm\beta,\bm I(\bm\beta)^{-1})$ に従う．

\subsection{2群の場合：\texorpdfstring{$2\times2$}{2×2} 表との一致}\label{subsec:06-2group}
\begin{meidai}\label[meidai]{prop:06-2group}
$\logit\pi=\alpha+\beta x$（$x\in\{0,1\}$）を $2\times2$ 表（\cref{tab:02-2x2}）に当てはめると
\[
\hat\alpha=\log\frac cd,\qquad \hat\beta=\log\frac{ad}{bc}=\log\widehat{\OR},\qquad
\widehat\V[\hat\beta]=\frac1a+\frac1b+\frac1c+\frac1d.
\]
\end{meidai}
\begin{proof}
パラメータ数（2）と群の数（2）が等しい飽和モデルなので，MLE は各群の観測割合を再現する：$\hat\pi_1=a/n_1$，$\hat\pi_0=c/n_0$
（スコア方程式 $\sum y=\sum\hat\pi$，$\sum_{x=1}y=\sum_{x=1}\hat\pi$ から）．
よって $\hat\alpha=\logit(c/n_0)=\log(c/d)$，$\hat\alpha+\hat\beta=\log(a/b)$．
情報行列は $w_1=n_1\pi_1(1-\pi_1)$，$w_0=n_0\pi_0(1-\pi_0)$ として
$\bm I=\begin{pmatrix}w_0+w_1&w_1\\ w_1&w_1\end{pmatrix}$，$\det\bm I=w_0w_1$，
$(\bm I^{-1})_{22}=(w_0+w_1)/(w_0w_1)=1/w_0+1/w_1$．
$\hat w_1=ab/n_1$ より $1/\hat w_1=1/a+1/b$，同様に $1/\hat w_0=1/c+1/d$．
\end{proof}
\kakomon{2023}{1} では $\hat\alpha=\log(4/6)=-0.405$，$\hat\beta=\log\{(7/3)/(4/6)\}=\log(7/2)=1.253$．
年齢を加えたモデルで治療の係数が $1.1$ に変わったのは，ランダム化にもかかわらず偶然に年齢分布が群間で異なった（$C$ 群に70歳が多い）ことによる調整の効果である（OR の非可縮性の影響も含む）．

\section{推定と検定}\label{sec:06-test}
\subsection{3つの検定}
$H_0:\beta_j=0$ に対して
\begin{itemize}
\item \textbf{Wald 検定}：$z=\hat\beta_j/\SE(\hat\beta_j)$，$z^2\sim\chi^2_1$．信頼区間 $\exp\{\hat\beta_j\pm z_{\alpha/2}\SE\}$ と双対．
\item \textbf{尤度比検定}：入れ子のモデルの最大対数尤度 $\ell_1\ge\ell_0$ に対し $G^2=2(\ell_1-\ell_0)\sim\chi^2_{d}$，$d$ は自由パラメータ数の差．
\item \textbf{スコア検定}：制約モデルの MLE で評価したスコアを使う．2群の $\beta=0$ のスコア検定は Pearson $\chi^2$ である．
\end{itemize}
\kakomon{2023}{1} では $G^2=2\{-11.79-(-12.84)\}=2.10<\chi^2_1(0.05)=3.84$ で，年齢の効果は有意でない．
3つの検定は漸近的に同等だが，小標本や係数が大きい場合は値が異なり，尤度比検定が最も信頼できることが多い．
\Youchui Wald 検定は $|\hat\beta|$ が大きいと標準誤差も大きくなり検出力が落ちることがある（Hauck--Donner 効果）．

\subsection{完全分離}
\begin{supbox}[補足：完全分離]
ある共変量の閾値でアウトカムが完全に分かれる（例：治療群全員が有効）と，尤度は $|\beta|\to\infty$ で単調に増加し MLE が存在しない．
$2\times2$ 表にゼロのセルがあると $\log\widehat{\OR}$ が定義できないのと同じ現象で，exact 法や罰則付き尤度（Firth 法）を用いる．
\end{supbox}

\section{共変量調整と交絡}\label{sec:06-adjust}
\begin{itemize}
\item 観察研究では，交絡因子を共変量に加えることで調整オッズ比を推定する（\cref{ch:07}）．
\item 層別変数を共変量に含めたロジスティック回帰は，層共通ORの MH 推定と同じ対象を推定する．
\item RCT でも予後因子で調整すると OR の値は変わる（\cref{subsec:02-collapse}の非可縮性）．
      調整 OR は「同じ予後の患者どうしの条件付き効果」で，粗 OR は「集団全体の周辺効果」である．
      ランダム化が成功していれば両者の違いは交絡ではない．
\end{itemize}

\section{予測と判別}\label{sec:06-predict}
予測確率 $\hat\pi(\bm x)=\mathrm{expit}(\bm x\transp\hat{\bm\beta})$，$\mathrm{expit}(u)=e^u/(1+e^u)$．
カットオフ $c$ で「$\hat\pi\ge c$ なら陽性」と判定すれば，感度・特異度・PPV・NPV（\cref{ch:05}）が計算できる．
\kakomon{2023}{1}〔5〕では，予測確率 $\ge0.6$ の群（9人）のうち実際に効果ありが6人で $\mathrm{PPV}=2/3$，
$<0.6$ の群（11人）のうち効果なしが6人で $\mathrm{NPV}=6/11$．
\begin{itemize}
\item \textbf{判別能}（discrimination）：AUC（c統計量）．
\item \textbf{較正}（calibration）：予測確率と観測割合の一致．予測値を10分位に分けて比べる Hosmer--Lemeshow 検定など．
\end{itemize}
予測モデルは作成データで評価すると過大評価になるので，独立な検証コホートで評価する（\kakomon{2018}{4}）．

\section{モデル選択：AIC と KL 情報量}\label{sec:06-aic}
\begin{teigi}[AIC]\label{def:06-aic}
自由パラメータ数 $k$ のモデルの最大対数尤度を $\ell(\hat{\bm\beta})$ として $\mathrm{AIC}=-2\ell(\hat{\bm\beta})+2k$．
切片と $p$ 個の説明変数のロジスティックモデルでは $k=p+1$（\kakomon{2018}{4}〔3〕）．小さいほど良い．
\end{teigi}
\begin{teigi}[Kullback--Leibler 情報量]\label{def:06-kl}
真の分布 $g$ に対するモデル $f$ の KL 情報量は $I(g;f)=\E_g[\log\{g(X)/f(X)\}]\ge0$．
$g=\mathrm{Bern}(\theta)$，$f=\mathrm{Bern}(\pi)$ なら
\begin{equation}
I(g;f)=\theta\log\frac\theta\pi+(1-\theta)\log\frac{1-\theta}{1-\pi}.\label{eq:06-kl}
\end{equation}
\end{teigi}
$I(g;f)=\E_g[\log g]-\E_g[\log f]$ の第1項はモデルによらないので，KL の最小化は期待平均対数尤度 $\E_g[\log f]$ の最大化と同値である．
最大対数尤度はこの量を楽観的に（同じデータで推定と評価をするため）過大推定し，その偏りが漸近的に $k$ であることから
$-2\times(\ell-k)$ が AIC となる（\kakomonSM{2023}{4} で線形回帰の場合の罰則 $2(p+1)$ が導出された）．

\section{症例対照研究とロジスティック回帰}\label{sec:06-cc}
\Hatten 症例対照研究で，症例は確率 $s_1$，対照は確率 $s_0$ で抽出される（$\bm x$ によらない）とする．
抽出された ($S=1$) 被験者での発症確率は Bayes の定理より
\begin{gather*}
\Prob(D=1\mid\bm x,S=1)=\frac{s_1\pi(\bm x)}{s_1\pi(\bm x)+s_0\{1-\pi(\bm x)\}},\\
\logit\Prob(D=1\mid\bm x,S=1)=\log\frac{s_1}{s_0}+\beta_0+\beta_1x_1+\dots.
\end{gather*}
したがって症例対照データにロジスティック回帰を当てはめると，\textbf{傾き（対数OR）のパラメータは抽出によって変わらず（保存され），切片だけが $\log(s_1/s_0)$ ずれる}．
したがって通常の条件（抽出確率が $\bm x$ によらない，十分な標本サイズ）の下で，症例対照データの最尤推定量は対数 OR を一致推定する（小標本での不偏性までは保証されない）．
症例対照研究でも調整ORが推定できるのはこのためで，一方で発症確率（リスク）そのものは推定できない．

\section{仮定}\label{sec:06-assump}
\begin{itemize}
\item 観測の独立性（クラスタや反復測定があれば GEE・混合モデル）．
\item 連続共変量について logit が線形（必要ならカテゴリ化・スプライン）．
\item 交互作用を入れていなければ，オッズ比が他の共変量の値によらず一定．
\item 十分なイベント数．「変数（推定するパラメータ）1つあたりイベント10以上」（EPV $\ge10$）は古典的な粗い目安であり，絶対的な条件ではない．イベントが少ないと推定が不安定で過学習する．
\end{itemize}

\section{結果の解釈}\label{sec:06-interp}
\begin{itemize}
\item 「年齢・性別を調整した喫煙の OR は 2.1（95\%CI 1.4--3.2）」：同じ年齢・性別の人どうしで比べて，喫煙者の発症オッズは約2倍．
\item 連続変数の OR は単位に依存する（年齢1歳あたり1.04なら10歳あたり $1.04^{10}\approx1.48$）．
\item リスクスコアの評価（\kakomon{2018}{4}〔5〕）：どの因子が独立した予測因子か（調整ORと信頼区間），判別能（AUC とその区間），
      既存の方法（医師の判断）との比較，検証コホートでの評価であることを述べる．
\end{itemize}

\section{手法選択}\label{sec:06-choice}
\begin{itemize}
\item 調整したい変数が少数のカテゴリ変数 → MH 法でも可．連続変数や多数の変数 → ロジスティック回帰．
\item リスク比を直接推定したい（イベントが稀でない） → log-二項モデルや修正ポアソン回帰．
\item 時間の情報（打ち切り）がある → Cox 回帰（\cref{ch:10}）．
\item マッチした症例対照研究 → 条件付きロジスティック回帰（ペアの切片を条件付けで消去）．
\end{itemize}

\section{よくある誤り}\label{sec:06-err}
\begin{warnbox}
\begin{itemize}
\item 係数 $\beta$ をそのままオッズ比と書く（オッズ比は $e^\beta$）．
\item イベントが稀でないのに調整 OR をリスク比として解釈する．
\item 尤度比検定の自由度をモデル全体のパラメータ数とする（差が自由度）．
\item AIC のパラメータ数に切片を数え忘れる．
\item 交互作用項があるモデルで $e^{\beta_1}$ を「$x$ のオッズ比」と書く（$z=0$ のときのオッズ比にすぎない）．
\end{itemize}
\end{warnbox}

\section{数理統計との接続}\label{sec:06-link}
\begin{linkbox}
\begin{itemize}
\item ロジットは二項分布（指数型分布族）の自然母数である（『現代数理統計学の基礎』6.1.2項，例6.9）．
      自然母数に線形モデルを置くのが正準リンクで，このとき $\bm X\transp\bm y$ が十分統計量，観測情報と期待情報が一致する．
\item MLE の漸近正規性と情報行列の逆行列（同書 6.4.3項）がそのまま Wald 型の標準誤差になる．
\item 尤度比検定の $\chi^2_{k-r}$ 近似は同書 命題7.8．AIC の罰則の導出は\kakomonSM{2023}{4}，多項モデルの AIC は\kakomonSM{2025}{1}．
\end{itemize}
\end{linkbox}

\section{例題}\label{sec:06-ex}
\begin{reidai}[2群のロジスティック回帰]\label{reidai:06-1}
治療群30人中18人，対照群30人中10人が改善した．$\logit\pi=\alpha+\beta x$（治療 $x=1$）を当てはめる．
\begin{enumerate}[label=(\arabic*)]
\item 対数尤度を $\alpha,\beta$ の関数として書け．
\item $\hat\alpha,\hat\beta$ とオッズ比の推定値を求めよ．
\item $\SE(\hat\beta)$ を求め，Wald 検定を行え．オッズ比の95\%信頼区間も求めよ（$e^{0.0451}=1.046$，$e^{2.1521}=8.60$）．
\item 最大対数尤度は，モデル全体で $-39.29$，$\beta=0$ のモデルで $-41.46$ である．尤度比検定を行え．
\item Pearson $\chi^2$ 統計量を求め，(3)(4)と比べよ．
\end{enumerate}
\end{reidai}

\begin{reidai}[係数の解釈と予測]\label{reidai:06-2}
ある疾患の発症について $\logit\pi=-5.0+0.04\,(\text{年齢})+0.9\,(\text{喫煙})+0.5\,(\text{男性})$ が得られた．
$e^{0.9}=2.460$，$e^{0.4}=1.492$，$e^{1.2}=3.320$ を用いよ．
\begin{enumerate}[label=(\arabic*)]
\item 喫煙の調整オッズ比と，年齢10歳あたりのオッズ比を求め，解釈せよ．
\item 60歳の喫煙男性の発症確率の予測値を求めよ．
\item 喫煙と性別の交互作用項 $0.3\,(\text{喫煙}\times\text{男性})$ を加えたモデルでは，男性での喫煙のオッズ比は何を用いて表されるか．
\end{enumerate}
\end{reidai}

\begin{reidai}[尤度比検定と AIC]\label{reidai:06-3}
同じデータに次の入れ子のモデルを当てはめた．
M0：切片と治療（パラメータ2個），最大対数尤度 $-120.4$．
M1：M0 に年齢（3個），$-117.9$．
M2：M1 に3水準の重症度（ダミー2個を追加，計5個），$-116.8$．
\begin{enumerate}[label=(\arabic*)]
\item M0 対 M1，M1 対 M2 の尤度比検定を行え（$\chi^2_1(0.05)=3.841$，$\chi^2_2(0.05)=5.991$）．
\item 各モデルの AIC を求め，選ばれるモデルを答えよ．
\end{enumerate}
\end{reidai}

\section{解答}\label{sec:06-sol}
\begin{kaitou}{reidai:06-1}
(1) $\ell=18(\alpha+\beta)+10\alpha-30\log(1+e^{\alpha+\beta})-30\log(1+e^\alpha)=28\alpha+18\beta-30\log(1+e^\alpha)-30\log(1+e^{\alpha+\beta})$．\\
(2) 飽和モデルなので $\hat\pi_1=0.6$，$\hat\pi_0=1/3$．$\hat\alpha=\log(10/20)=-0.693$，$\hat\alpha+\hat\beta=\log(18/12)=0.405$，
$\hat\beta=\log3=1.099$，$\widehat{\OR}=3$．\\
(3) $\SE=\sqrt{1/18+1/12+1/10+1/20}=\sqrt{0.2889}=0.5375$．$z=1.099/0.5375=2.04>1.96$ で有意．
区間 $1.099\pm1.96\times0.5375=(0.045,\,2.152)$ → OR の区間 $(1.05,\,8.60)$．\\
(4) $G^2=2\{-39.29-(-41.46)\}=4.34>3.841$ で有意．\\
(5) $X^2=\frac{60(18\cdot20-12\cdot10)^2}{30\cdot30\cdot28\cdot32}=\frac{60\times57600}{806400}=4.29$．
Wald $z^2=4.18$，尤度比 $4.34$，スコア（Pearson）$4.29$ と，3つの検定は近いが一致はしない．いずれも5\%で有意．
\end{kaitou}

\begin{kaitou}{reidai:06-2}
(1) 喫煙の調整OR $=e^{0.9}=2.46$：同じ年齢・性別の人どうしで比べて，喫煙者の発症オッズは非喫煙者の約2.5倍．
年齢10歳あたり $e^{0.4}=1.49$：同じ喫煙状況・性別なら，10歳年上だとオッズが約1.5倍．\\
(2) 線形予測子 $-5.0+2.4+0.9+0.5=-1.2$．$\hat\pi=1/(1+e^{1.2})=1/4.320=0.231$．\\
(3) $e^{0.9+0.3}=e^{1.2}=3.32$（女性では $e^{0.9}=2.46$）．喫煙の効果が性別で異なる（効果修飾）ことを表す．
\end{kaitou}

\begin{kaitou}{reidai:06-3}
(1) M0 対 M1：$G^2=2(120.4-117.9)=5.0>3.841$（自由度1）で年齢は有意．
M1 対 M2：$G^2=2(117.9-116.8)=2.2<5.991$（自由度2）で重症度の追加は有意でない．\\
(2) $\mathrm{AIC}_0=240.8+4=244.8$，$\mathrm{AIC}_1=235.8+6=241.8$，$\mathrm{AIC}_2=233.6+10=243.6$．M1 が選ばれる．
\end{kaitou}

\begin{summarybox}
\begin{itemize}
\item $\logit\pi=\bm x\transp\bm\beta$．$e^{\beta_j}$ は他の共変量を固定した調整オッズ比．交互作用があれば $e^{\beta_1+\beta_3z}$．
\item $\ell=\sum\{y_i\bm x_i\transp\bm\beta-\log(1+e^{\bm x_i\transp\bm\beta})\}$，$\bm U=\bm X\transp(\bm y-\bm\pi)$，$\bm I=\bm X\transp\bm W\bm X$．
\item 2群では $\hat\beta=\log(ad/bc)$，$\V=1/a+1/b+1/c+1/d$（$2\times2$ 表と一致）．
\item 検定：Wald $\hat\beta/\SE$，尤度比 $2(\ell_1-\ell_0)\sim\chi^2_{\text{パラメータ数の差}}$，スコア（2群なら Pearson）．
\item $\mathrm{AIC}=-2\ell+2(p+1)$．KL：$\theta\log(\theta/\pi)+(1-\theta)\log\{(1-\theta)/(1-\pi)\}$．
\item 症例対照研究では傾き（対数OR）は抽出で保存され一致推定できる．切片は $\log(s_1/s_0)$ だけずれる．
\end{itemize}
\end{summarybox}

% ===== ch07.tex =====
\chapter{疫学研究・交絡・層別・マッチング}\label{ch:07}

\begin{goalbox}
\begin{itemize}
\item 交絡の定義と条件を述べ，粗（crude）効果と調整（adjusted）効果の違いを説明できる．
\item 交絡と効果修飾を区別し，層間の均一性を検定できる．
\item 層別 $2\times2$ 表に Mantel--Haenszel 法・逆分散重み法・Peto の1段階推定を適用し，共通オッズ比とその信頼区間を計算できる．
\item 直接標準化による調整リスクを計算できる．
\item マッチングの目的と，マッチしたデータの正しい解析法を説明できる．
\item 傾向スコアの定義とバランス特性を証明し，層別・重み付けによる調整を実行できる．
\item Simpson のパラドックスを例示し，因果解釈上の注意を述べられる．
\end{itemize}
\end{goalbox}

\begin{exambox}
\begin{itemize}
\item \kakomon{2017}{4}：年齢3層の症例対照データ．層別の曝露OR，超幾何分布の $\E[y_{11k}],\V[y_{11k}]$，
      $\widehat{\log\psi_k}=(y_{11k}-\E)/\V$ による推定，逆分散重みの共通 OR（$\hat\psi=6.22$）と分散 $1/\sum V_k$，95\%CI，均一性の確認方法（Breslow--Day）．
\item \kakomon{2019}{2}：乳がん化学療法のコホート．粗リスク差の CI，4層の直接標準化による調整リスク差，
      飽和ロジスティックモデルによる傾向スコアとバランス特性の確認，傾向スコア層別による調整リスク差，バランス特性の証明（繰り返し期待値）．
\item \kakomon{2021}{3}〔3〕〔4〕：層別 Cochran 検定の導出と，性別で層別すると結論が変わる（Simpson のパラドックス）．
\item \kakomon{2022}{2}：マッチしたペアでの粗 RR・OR と MH 推定量の比較，「交絡の有無」と「推定されるパラメータの違い」，間接比較によるシートベルトの効果．
\end{itemize}
\end{exambox}

\flowline{層別2×2表\\共変量}{調整RR・OR\\標準化リスク}{MH・標準化\\PS・回帰}{均一性\\未測定交絡なし}{重み付き和\\超幾何分散}{粗と調整の差\\の理由}

\section{問題設定}\label{sec:07-setting}
観察研究では治療・曝露 $X$ の選択が患者の特性 $Z$ に依存する．$Z$ が結果 $Y$ にも影響するなら，
$X$ と $Y$ の粗い関連には $Z$ を介した関連が混入する．
本章では，層（$Z$ の値）ごとの $2\times2$ 表 $k=1,\dots,K$（$a_k,b_k,c_k,d_k$，行和 $n_{1k},n_{0k}$，列和 $m_{1k},m_{0k}$，総数 $N_k$）を
主なデータ形式とする．

\section{基本概念}\label{sec:07-basic}
\subsection{交絡}\label{subsec:07-confound}
\begin{teigi}[交絡因子]\label{def:07-confounder}
次の3条件を満たす変数 $Z$ を交絡因子という．
(i) 曝露 $X$ と関連する．(ii) 曝露がない集団でも結果 $Y$ の危険因子である．(iii) 曝露から結果への経路上の中間変数ではない．
\end{teigi}
\Youchui 中間変数（治療→血圧低下→脳卒中予防の「血圧」など）で調整すると，治療効果の一部を消してしまう．
治療後に測定された変数で調整してはいけない．

\begin{supbox}[補足：潜在結果による定式化]
各個体に，治療を受けた場合の結果 $Y(1)$ と受けなかった場合の結果 $Y(0)$ を考える．
平均因果効果は $\E[Y(1)]-\E[Y(0)]$ である．観測できるのは $Y=XY(1)+(1-X)Y(0)$ だけなので，
$\E[Y\mid X=1]-\E[Y\mid X=0]$ が因果効果に等しいためには $\{Y(1),Y(0)\}\indep X$（交換可能性）が必要で，ランダム化はこれを保証する．
観察研究では $Z$ を条件として $\{Y(1),Y(0)\}\indep X\mid Z$（未測定の交絡がない）と $0<\Prob(X=1\mid Z)<1$（正値性）を仮定し，
$\E[Y(x)]=\E_Z\bigl[\E[Y\mid X=x,Z]\bigr]$ によって因果効果を識別する（標準化の公式）．
\end{supbox}

\subsection{粗効果と調整効果・Simpson のパラドックス}\label{subsec:07-simpson}
\kakomon{2021}{3} では，併合表の OR は $0.90$（有意差なし）だったが，性別で層別すると男性 $4.09$，女性 $2.06$ で，Cochran 検定は $X_C=8.94$ と有意になった．
新規治療群に女性（有効率が低い層）が多く，既存治療群に男性（有効率が高い層）が多かったためである．
このように，併合した関連の向きや大きさが層別の関連と食い違う現象を Simpson のパラドックスという．

\subsection{効果修飾}\label{sec:07-modification}
\begin{teigi}[効果修飾]\label{def:07-modification}
曝露の効果の大きさが $Z$ の値によって異なることを，$Z$ による効果修飾（effect modification；交互作用）という．
\end{teigi}
交絡は推定の\textbf{偏り}であり調整によって除くべきもの，効果修飾は\textbf{自然の性質}であり報告すべきものである．
効果修飾があれば共通 OR を1つの数値で要約することは適切でなく，層ごとの効果を示す．
また効果修飾の有無は尺度に依存する：RD が層間で一定なら，ベースラインが異なる層では RR や OR は一定でない．

\section{数理：層別解析}\label{sec:07-strat}
\subsection{Mantel--Haenszel 法}\label{subsec:07-mh}
\Hissu\Hinshutsu MH 検定統計量と MH 推定量は\cref{sec:04-mh}で定義した．ここでは性質をまとめる．
\begin{itemize}
\item $\chi^2_{\mathrm{MH}}=\{\sum(a_k-\E[a_k])\}^2/\sum\V[a_k]$（超幾何分散）は，全層で $\OR=1$ の帰無仮説の検定．
      各層の $O-E$ を足してから2乗するので，層間で効果の向きが揃っているときに検出力が高い．
\item \kakomon{2021}{3} の Cochran 検定は分散に $n_{1k}n_{0k}m_{1k}m_{0k}/N_k^3$（積二項の条件なし分散）を用い，
      各 $Y_k$ が漸近的に正規で層間独立なので $Y=\sum Y_k$ が漸近的に $\Normal(0,\sum\V[Y_k])$ に従うことから導かれる．
\item $\widehat{\OR}_{\mathrm{MH}}=\sum(a_kd_k/N_k)/\sum(b_kc_k/N_k)$ は，層内の推定値 $a_kd_k/(b_kc_k)$ を重み $b_kc_k/N_k$ で平均したものとみなせる．
      ゼロのセルがあっても計算でき，層が多数で各層が小さいときも一致推定量である．
\end{itemize}
\begin{supbox}[補足：MH オッズ比の分散（Robins--Breslow--Greenland \cite{rbg1986}）]
$P_k=(a_k+d_k)/N_k$，$Q_k=(b_k+c_k)/N_k$，$R_k=a_kd_k/N_k$，$S_k=b_kc_k/N_k$，$R=\sum R_k$，$S=\sum S_k$ として
\[
\widehat\V[\log\widehat{\OR}_{\mathrm{MH}}]=\frac{\sum P_kR_k}{2R^2}+\frac{\sum(P_kS_k+Q_kR_k)}{2RS}+\frac{\sum Q_kS_k}{2S^2}.
\]
\end{supbox}

\subsection{逆分散重み法と Peto の1段階推定}\label{subsec:07-peto}
各層の対数 OR の推定値 $\hat\theta_k$ とその分散 $v_k$ があれば，共通の $\theta$ の推定量として
\begin{equation}
\hat\theta_w=\frac{\sum_kw_k\hat\theta_k}{\sum_kw_k},\qquad w_k=\frac1{v_k},\qquad \V[\hat\theta_w]=\frac{1}{\sum_kw_k}\label{eq:07-iv}
\end{equation}
が得られる（固定効果メタアナリシスと同じ；\cref{ch:15}）．$\hat\theta_k=\log(a_kd_k/b_kc_k)$，$v_k=1/a_k+1/b_k+1/c_k+1/d_k$ とするのが Woolf 法である．

\paragraph{Peto の1段階推定（\kakomon{2017}{4}）}
各層で周辺度数を固定した条件付き尤度（非心超幾何分布）の，$\theta=\log\OR$ に関するスコアと情報量は，$\theta=0$ で
$U_k(0)=a_k-\E[a_k]$，$I_k(0)=\V[a_k]$（\cref{eq:04-hyper}）となる．$\theta=0$ から Newton 法を1回だけ行うと
\begin{equation}
\widehat{\log\psi_k}=\frac{a_k-\E[a_k]}{\V[a_k]},\qquad
\widehat{\log\psi}=\frac{\sum_kV_k\cdot\frac{a_k-E_k}{V_k}}{\sum_kV_k}=\frac{\sum_k(a_k-E_k)}{\sum_kV_k},\qquad
\V[\widehat{\log\psi}]=\frac{1}{\sum_kV_k}.\label{eq:07-peto}
\end{equation}
重みは $1/\V[\widehat{\log\psi_k}]=V_k$ である．分子は MH 検定の分子と同じで，$\widehat{\log\psi}^2\sum V_k=\chi^2_{\mathrm{MH}}$ となる．
真の OR が1に近く，群の大きさが釣り合っているときに良い近似である．OR が1から遠い場合や群の大きさが大きく異なる場合には偏りが大きくなり，その向き（1に近づくか遠ざかるか）は一定しない（\kakomon{2017}{4} の第1層では1段階推定値 $31.9$ が層の OR $7.71$ を大きく上回る）．
\kakomon{2017}{4} では $\sum(O-E)=96-48.18=47.82$，$\sum V=26.17$ から $\widehat{\log\psi}=1.827$，$\hat\psi=6.22$，
95\%CI は $\exp(1.827\pm1.96/\sqrt{26.17})=(4.24,\,9.12)$．

\subsection{均一性の検定}\label{subsec:07-homog}
共通 OR の仮定（\kakomon{2017}{4}〔5〕）は，次のような方法で確認する．
\begin{itemize}
\item \textbf{Woolf の均一性検定}：$Q=\sum_kw_k(\hat\theta_k-\hat\theta_w)^2\approx\chi^2_{K-1}$（\cref{sec:15-q}の Cochran の $Q$ と同じ）．
\item \textbf{Breslow--Day 検定}：共通 $\widehat{\OR}_{\mathrm{MH}}$ の下での各層の期待度数と観測度数を比べる $\chi^2_{K-1}$ 検定．
\item ロジスティック回帰で「曝露×層」の交互作用項を加えた尤度比検定（自由度 $K-1$）．
\end{itemize}
層数が少なく各層が小さいとこれらの検定は検出力が低いので，有意でないことは均一性の証明にならない．

\section{標準化}\label{sec:07-standard}
\subsection{直接標準化}
\begin{teigi}[直接標準化]\label{def:07-direct}
層 $k$ の曝露群 $x$ のリスクを $p_{xk}$，標準集団の層の割合を $w_k$（$\sum w_k=1$）として，
$R_{xw}=\sum_kw_kp_{xk}$ を標準化リスクという．$R_{1w}-R_{0w}$，$R_{1w}/R_{0w}$ が標準化（調整）RD・RR．
\end{teigi}
両群に同じ重みを用いるので，$Z$ の分布の違いによる交絡が除かれる．
標準集団に研究全体（両群を合わせた集団）を用いると，推定対象は（未測定の交絡がなければ）集団全体での平均因果効果 $\E[Y(1)]-\E[Y(0)]$ になる（前掲の補足の標準化公式）．
\kakomon{2019}{2} では $w_k=(N_{0k}+N_{1k})/(N_0+N_1)$ として，粗 RD $0.27$ に対し標準化 RD は $0.026$ にまで小さくなった．
化学療法は閉経後・進行がん（再発リスクが高い層）に多く行われており，粗 RD の大部分が交絡によるものだったことを示す．

\subsection{間接標準化}
\begin{supbox}[補足：標準化死亡比（SMR）]
対象集団の層別人数 $n_k$ と参照集団の層別率 $\lambda_k^{\mathrm{ref}}$ から期待死亡数 $E=\sum n_k\lambda_k^{\mathrm{ref}}$ を求め，
$\mathrm{SMR}=O/E$（観測死亡数/期待死亡数）とする．対象集団の層別率が不安定（小集団）なときに用いる．
$O$ をポアソン分布とみなして検定・区間推定する（\cref{ch:17}）．
\end{supbox}

\section{マッチング}\label{sec:07-match}
\Hinshutsu マッチングは，交絡因子の分布を群間でそろえるために，背景の似た個体を組にするデザインである．
\begin{itemize}
\item \textbf{マッチした症例対照研究}：症例と対照がペア．マッチングで対照のマッチ因子の分布が症例に近づき，マッチ因子が曝露と関連していれば対照の曝露分布も症例に近づくので，マッチを無視した粗 OR は通常1の方向に偏る．
      ペアを層とする MH-OR（＝不一致ペアの比 $n_{10}/n_{01}$），McNemar 検定，条件付きロジスティック回帰で解析する．
\item \textbf{マッチしたコホート}（\kakomon{2022}{2}）：同じ事故の運転者（着用）と同乗者（非着用）がペア．
      事故状況 $Z$ の分布は両群で完全に等しく，$Z$ と曝露は独立なので，$Z$ による交絡はない．
      それでも粗 OR $2.79$ と MH-OR $3.98$ が異なるのは OR の非可縮性（\cref{subsec:02-collapse}）による．
      RR は可縮なので粗 RR $=$ MH-RR $=3074/1104=2.78$．
\end{itemize}
\paragraph{間接比較（\kakomon{2022}{2}〔5〕）}
「着用の運転者 vs 非着用の同乗者」の RR（2.78）には座席の違いと着用の違いが混ざっている．
両者とも非着用の事故データから座席の効果（同乗者/運転者の RR $=1.01$）を推定し，
$2.78/1.01\approx2.75$（丸めずに計算すると $2.77$）を着用の効果（非着用/着用）の推定値とする．これは座席の効果と着用の効果が乗法的に独立（交互作用なし）という仮定に基づく．

\section{傾向スコア}\label{sec:07-ps}
\Hinshutsu
\begin{teigi}[傾向スコア]\label{def:07-ps}
共変量 $\bm Z$ をもつ個体が治療を受ける確率 $e(\bm z)=\Prob(X=1\mid\bm Z=\bm z)=\E[X\mid\bm Z=\bm z]$ を傾向スコア（propensity score）という．
\end{teigi}
\begin{teiri}[バランス特性（\kakomon{2019}{2}〔5〕）]\label[teiri]{thm:07-balance}
\[
\Prob\bigl(X=1\mid\bm Z,e(\bm Z)\bigr)=\Prob\bigl(X=1\mid e(\bm Z)\bigr)=e(\bm Z).
\]
すなわち $X\indep\bm Z\mid e(\bm Z)$：傾向スコアが同じ個体の集団では，$\bm Z$ の分布が治療群と対照群で等しい．
\end{teiri}
\begin{proof}
$e(\bm Z)$ は $\bm Z$ の関数なので $\Prob(X=1\mid\bm Z,e(\bm Z))=\Prob(X=1\mid\bm Z)=e(\bm Z)$．
一方，繰り返し期待値の法則より
\[
\Prob(X=1\mid e(\bm Z))=\E[X\mid e(\bm Z)]=\E\bigl[\E[X\mid\bm Z,e(\bm Z)]\bigm| e(\bm Z)\bigr]=\E\bigl[e(\bm Z)\mid e(\bm Z)\bigr]=e(\bm Z).
\]
両者が等しいので，$e(\bm Z)$ を与えたとき $X$ の条件付き分布は $\bm Z$ に依存しない．
\end{proof}
さらに，$\bm Z$ を与えて未測定の交絡がない（$\{Y(1),Y(0)\}\indep X\mid\bm Z$）なら，$e(\bm Z)$ を与えても同じことが成り立つ．
多次元の $\bm Z$ の調整を1次元のスコアの調整に縮約できるのが傾向スコアの利点である．

\paragraph{推定と利用}
$e(\bm z)$ は通常ロジスティック回帰で推定する．\kakomon{2019}{2} では2 値の $\bm Z=(\text{閉経},\text{進行度})$ と交互作用を含む飽和モデルなので，
推定値は各層の治療割合 $0.143,0.550,0.550,0.853$ に一致した．層2と層3で PS が等しく，この2層を合わせた集団では $\bm Z$ の分布が両群で一致する（$0.513:0.487$）．
利用法には次がある．
\begin{itemize}
\item \textbf{層別}：PS の値（または5分位など）で層別し，標準化する．\kakomon{2019}{2}〔4〕では調整 RD $=0.026$ で，$\bm Z$ による直接標準化と一致した．
\item \textbf{マッチング}：PS の近い治療者と非治療者を組にする．
\item \textbf{逆確率重み付け（IPW）}：$\widehat{\E}[Y(1)]=\frac1n\sum\frac{X_iY_i}{\hat e(\bm Z_i)}$，$\widehat{\E}[Y(0)]=\frac1n\sum\frac{(1-X_i)Y_i}{1-\hat e(\bm Z_i)}$．
      $XY=XY(1)$ と $Y(1)\indep X\mid\bm Z$ より $\E\bigl[XY/e(\bm Z)\bigr]=\E\bigl[\E[X\mid\bm Z]\E[Y(1)\mid\bm Z]/e(\bm Z)\bigr]=\E[Y(1)]$ による．
\item \textbf{共変量としての調整}：回帰モデルに $\hat e(\bm Z)$ を共変量として加える．
\end{itemize}

\section{仮定}\label{sec:07-assump}
\begin{itemize}
\item \textbf{未測定の交絡がない}：調整した変数以外に交絡因子がない．データからは検証できない最大の仮定．
\item \textbf{正値性}：どの層にも治療者と非治療者がいる（$0<e(\bm z)<1$）．PS が0や1に近い個体があると IPW の重みが不安定になる．
\item \textbf{モデルの正しさ}：PS モデル・結果モデルの特定．
\item MH 推定：層間で効果が均一．
\end{itemize}

\section{結果の解釈}\label{sec:07-interp}
\begin{itemize}
\item 粗効果と調整効果が異なる理由として，(a) 交絡（\kakomon{2019}{2}，\kakomon{2021}{3}），(b) OR の非可縮性（\kakomon{2022}{2}），(c) 効果修飾による要約の違い，を区別して述べる．
\item 調整しても，因果効果と言えるのは未測定の交絡がない場合に限られる．観察研究の結論には「測定されていない交絡因子の影響は否定できない」を添える．
\item 層別の結果が逆向きなら共通の要約はせず，層別に報告する．
\end{itemize}

\section{手法選択}\label{sec:07-choice}
\begin{table}[htbp]
\centering\small
\caption{交絡の調整法}\label{tab:07-choice}
\begin{tabular}{@{}L{40mm}L{95mm}@{}}
\toprule
状況 & 方法\\
\midrule
デザイン段階 & ランダム化（最善），層別ランダム化，マッチング，限定（対象を1層に絞る）\\
少数のカテゴリ交絡因子 & MH 検定・MH 推定量，直接標準化\\
層が多く各層が小さい（ペアなど） & MH 推定量，条件付きロジスティック回帰（Woolf 法は不適）\\
連続・多数の交絡因子 & ロジスティック回帰，傾向スコア（層別・マッチング・IPW）\\
稀な疾患の症例対照研究 & マッチング＋条件付き解析，ロジスティック回帰による調整OR\\
\bottomrule
\end{tabular}
\end{table}

\section{よくある誤り}\label{sec:07-err}
\begin{warnbox}
\begin{itemize}
\item 粗 OR と調整 OR の差をすべて交絡と解釈する（非可縮性を無視）．
\item 効果修飾がある（層で OR が逆向き）のに共通 OR で要約する．
\item 中間変数・治療後の変数で調整する．
\item 傾向スコアで調整すれば未測定の交絡も除けると考える．
\item バランス特性の確認で，周辺分布だけを比べる（本来は同時分布．\kakomon{2019}{2} では両者が同値になる特殊な場合であった）．
\item 標準化の重みを群ごとに別々に取る（同じ重みを両群に使わないと調整にならない）．
\end{itemize}
\end{warnbox}

\section{数理統計との接続}\label{sec:07-link}
\begin{linkbox}
\begin{itemize}
\item バランス特性の証明は繰り返し期待値の法則（『現代数理統計学の基礎』式(4.4)）そのものである．
\item Peto の1段階推定は，条件付き尤度のスコアと情報量による1回の Newton 法であり，MLE の1段階近似（one-step estimator）の典型例である．
\item 逆分散重み付き平均が線形不偏推定量の中で分散最小であることは同書第6章演習（問9）．
\item 超幾何分布の分散（同書 式(3.10)）が $\V[a_k]$ の源であり，条件付けによって局外母数（各層のベースラインリスク）が消えることが MH 法の要点である．
\item \kakomonSM{2022}{2} では「条件付け（選択）によって見かけの相関が生じる」ことが問われた．これは選択バイアスや合流点での条件付けの数理的な原型である．
\end{itemize}
\end{linkbox}

\section{例題}\label{sec:07-ex}
\begin{reidai}[Simpson のパラドックスと MH 法]\label{reidai:07-1}
2病院で治療の有無と改善の有無を集計した．
病院A（軽症中心）：治療100人中90人改善，無治療500人中400人改善．
病院B（重症中心）：治療500人中300人改善，無治療100人中50人改善．
\begin{enumerate}[label=(\arabic*)]
\item 各病院の OR と，併合した表の OR を求めよ．
\item MH 共通 OR を求めよ．
\item MH 検定統計量を求め，有意水準5\%で検定せよ．各層の $\E[a_k]$，$\V[a_k]$ も示せ．
\item Peto の1段階推定による共通 OR を求め，(2)と比べよ（$e^{0.508}=1.66$）．
\item 併合した OR が層別の OR と逆向きになった理由を説明せよ．
\end{enumerate}
\end{reidai}

\begin{reidai}[直接標準化と IPW]\label{reidai:07-2}
2 値の交絡因子 $Z$ で層別したコホートで，
$Z=0$：治療20人中4人発症，無治療80人中8人発症；$Z=1$：治療60人中30人発症，無治療40人中16人発症．
\begin{enumerate}[label=(\arabic*)]
\item 粗リスク差を求めよ．
\item 研究集団全体を標準集団とする直接標準化リスク差を求めよ．
\item 各層の傾向スコアを求め，IPW 推定量によるリスク差を求めよ．(2)と一致することを確かめよ．
\end{enumerate}
\end{reidai}

\begin{reidai}[バランス特性の数値確認]\label{reidai:07-3}
2 値の共変量 $Z_1,Z_2$ による4層で，治療者数/非治療者数が $(Z_1,Z_2)=(0,0)$：$20/80$，$(0,1)$：$30/20$，$(1,0)$：$60/40$，$(1,1)$：$45/5$ であった．
飽和モデルによる傾向スコアを求め，傾向スコアが等しい層を併合した集団で，$(Z_1,Z_2)$ の同時分布が治療群と非治療群で等しいことを示せ．
\end{reidai}

\section{解答}\label{sec:07-sol}
\begin{kaitou}{reidai:07-1}
(1) A：$\frac{90\times100}{10\times400}=2.25$．B：$\frac{300\times50}{200\times50}=1.50$．
併合：治療600人中390人，無治療600人中450人で $\frac{390\times150}{210\times450}=0.62$．\\
(2) A：$a_kd_k/N_k=9000/600=15$，$b_kc_k/N_k=4000/600=6.667$．B：$15000/600=25$，$10000/600=16.667$．
$\widehat{\OR}_{\mathrm{MH}}=40/23.33=1.71$．\\
(3) A：$E=100\times490/600=81.67$，$V=\frac{100\cdot500\cdot490\cdot110}{600^2\cdot599}=12.50$．
B：$E=500\times350/600=291.67$，$V=\frac{500\cdot100\cdot350\cdot250}{600^2\cdot599}=20.29$．
$\sum(O-E)=(90-81.67)+(300-291.67)=16.67$，$\sum V=32.79$，$\chi^2_{\mathrm{MH}}=16.67^2/32.79=8.47>3.841$ で有意．\\
(4) $\widehat{\log\psi}=16.67/32.79=0.508$，$\hat\psi=1.66$．MH の1.71（条件付き最尤推定値は約1.69）より1にやや近い．
Peto 法は OR が1に近く群の大きさが釣り合うときに良い近似で，本例は層内の群の大きさが100対500，500対100と大きく不均衡なため，他の推定量とずれが生じた（ずれの向きは一般には決まらない）．\\
(5) 治療は重症（改善率の低い）病院Bに多く（500人），無治療は軽症の病院Aに多い．病院（重症度）が治療と改善の両方に関連する交絡因子で，
併合すると「治療群は重症が多いから改善率が低い」という効果が治療の効果に混入し，向きが逆転した．
\end{kaitou}

\begin{kaitou}{reidai:07-2}
(1) 治療 $34/80=0.425$，無治療 $24/120=0.200$，粗 RD $=0.225$．\\
(2) 層別リスク：$Z=0$ で $0.2$ と $0.1$，$Z=1$ で $0.5$ と $0.4$．標準集団の重みは各層 $100/200=0.5$．
$R_{1w}=0.5\times0.2+0.5\times0.5=0.35$，$R_{0w}=0.5\times0.1+0.5\times0.4=0.25$，標準化 RD $=0.10$．\\
(3) $e(0)=20/100=0.2$，$e(1)=60/100=0.6$．
$\widehat\E[Y(1)]=\frac{1}{200}\bigl(\frac{4}{0.2}+\frac{30}{0.6}\bigr)=\frac{20+50}{200}=0.35$，
$\widehat\E[Y(0)]=\frac{1}{200}\bigl(\frac{8}{0.8}+\frac{16}{0.4}\bigr)=\frac{10+40}{200}=0.25$．RD $=0.10$ で(2)と一致する．
（飽和 PS モデルの IPW は，研究集団を標準とする直接標準化と一致する．）
\end{kaitou}

\begin{kaitou}{reidai:07-3}
PS は $(0,0)$：$0.2$，$(0,1)$：$30/50=0.6$，$(1,0)$：$60/100=0.6$，$(1,1)$：$0.9$．
PS $=0.6$ の集団（層 $(0,1)$ と $(1,0)$）では，治療群 $30+60=90$ 人の内訳が $(0,1):(1,0)=30:60=1/3:2/3$，
非治療群 $20+40=60$ 人の内訳も $20:40=1/3:2/3$ で一致し，$(0,0)$，$(1,1)$ はどちらの群でも0である．
PS $=0.2$，$0.9$ の集団はそれぞれ1層のみなので，同時分布は自明に一致する．
\end{kaitou}

\begin{summarybox}
\begin{itemize}
\item 交絡因子：曝露と関連，結果の危険因子，中間変数でない．交絡は偏り，効果修飾は自然の性質．
\item MH 検定 $\{\sum(a_k-E_k)\}^2/\sum V_k$，MH-OR $\sum a_kd_k/N_k\big/\sum b_kc_k/N_k$．
\item Peto：$\widehat{\log\psi}=\sum(O-E)/\sum V$，$\V=1/\sum V$（OR が1に近いとき良い近似）．
\item 均一性：Woolf の $Q$，Breslow--Day（自由度 $K-1$）．有意でなくても均一の証明ではない．
\item 直接標準化 $R_{xw}=\sum w_kp_{xk}$（同じ重みを両群に）．
\item 傾向スコア $e(\bm z)=\Prob(X=1\mid\bm z)$．バランス特性 $X\indep\bm Z\mid e(\bm Z)$ は繰り返し期待値で証明．
\item 粗と調整の差の原因：交絡，非可縮性，効果修飾．調整しても未測定交絡は残る．
\end{itemize}
\end{summarybox}

% ===== ch08.tex =====
\chapter[縦断データ・ベースライン調整・混合モデル]{縦断データ・ベースライン調整・\\混合モデル}\label{ch:08}

\begin{goalbox}
\begin{itemize}
\item 最終値・変化量・共分散分析（ANCOVA）による治療効果の推定量の分散を導出し，相関 $\rho$ による優劣を説明できる．
\item 平均への回帰を2変量正規分布と真値モデルで説明し，前後比較の解釈の注意を述べられる．
\item 反復測定データの平均ベクトルの推定量の分布，変化量の標準誤差を計算できる．
\item 複合対称・AR(1)・Toeplitz・無構造の共分散構造のパラメータ数を数え，AIC で選択できる．
\item 線形混合効果モデル $\bm Y=\bm X\bm\beta+\bm Z\bm\gamma+\bm\varepsilon$ の周辺分散 $\bm V=\bm Z\bm G\bm Z\transp+\bm R$ と一般化最小二乗推定量を導出できる．
\item 固定効果と変量効果の違いを説明し，変量切片が複合対称構造を生むことを示せる．
\end{itemize}
\end{goalbox}

\begin{exambox}
\begin{itemize}
\item \kakomon{2024}{2}：2群のベースライン $X$ と1年後 $Y$（2変量正規，$\sigma^2,\rho$ 既知）．
      最終値の差 $\hat\delta_d$，変化量の差 $\hat\delta_c$，回帰モデル $Y=\alpha+\delta_az+\rho X+\varepsilon$ の最小二乗推定 $\hat\delta_a$ の係数と分散，
      3推定量の期待値が一致する条件，相対分散のグラフ．
\item \kakomon{2023}{4}：3時点の反復測定（3変量正規）．$\hat{\bm\mu}$ の分布，Toeplitz と無構造の AIC，相関行列，8週の変化量の平均と標準誤差，$Z$ 検定．
\item \kakomon{2022}{3}：血中鉛濃度の線形混合モデル．$\bm V=\bm Z\bm G\bm Z\transp+\bm R$，変量効果なしの場合の $\hat{\bm\beta}=(\bm X\transp\bm X)^{-1}\bm X\transp\bm Y$ と不偏性，数値計算，
      ブロック対角の複合対称構造を生む $\bm Z,\bm G$，被験者内相関を考慮する意義．
\item \kakomon{2022}{5}（共通）：血圧の平均への回帰（2変量正規の条件付き分布と真値モデル）．
\end{itemize}
\end{exambox}

\flowline{同一個体の\\反復測定}{平均の差\\時点別・変化}{ANCOVA\\混合モデル}{多変量正規\\共分散構造}{GLS・AIC\\分散比較}{個体内相関\\平均への回帰}

\section{問題設定}\label{sec:08-setting}
同一の患者から複数時点で測定した値は，患者固有の水準（個体差）を共有するため正の相関をもつ．
この相関を (a) ベースライン値で調整して治療効果の精度を上げる，(b) 反復測定全体の共分散としてモデル化する，の2つの方向で扱う．
相関を無視すると，推定量の標準誤差を誤る（多くの場合過小評価する）．

\section{ベースライン調整}\label{sec:08-baseline}
\Hissu\Hinshutsu
\subsection{3つの推定量}\label{subsec:08-three}
\kakomon{2024}{2} の設定：各群 $n$ 人，被験者 $i$ のベースライン $X_i$ と最終値 $Y_i$ が
$\V[X_i]=\V[Y_i]=\sigma^2$，$\Corr(X_i,Y_i)=\rho$ の2変量正規分布に従う．群 $g\in\{T,C\}$ の平均を $\mu_g^X,\mu_g^Y$ とする．
\begin{align*}
\text{最終値：}&\ \hat\delta_d=\bar Y_T-\bar Y_C,\\
\text{変化量：}&\ \hat\delta_c=(\bar Y_T-\bar X_T)-(\bar Y_C-\bar X_C),\\
\text{ANCOVA：}&\ \hat\delta_a=(\bar Y_T-\bar Y_C)-\rho(\bar X_T-\bar X_C).
\end{align*}
$\hat\delta_a$ は $Y_i=\alpha+\delta_az_i+\rho X_i+\varepsilon_i$（傾き $\rho$ は既知）の最小二乗推定量である．
実際，正規方程式 $\sum(Y_i-\alpha-\delta_az_i-\rho X_i)=0$，$\sum z_i(Y_i-\alpha-\delta_az_i-\rho X_i)=0$ から
$\hat\alpha=\bar Y_C-\rho\bar X_C$，$\hat\delta_a=\bar Y_T-\bar Y_C-\rho(\bar X_T-\bar X_C)$ となり，係数は $(c_1,c_2,c_3,c_4)=(1,-1,-\rho,\rho)$．

\begin{meidai}[3推定量の分散]\label[meidai]{prop:08-var}
\begin{equation}
\V[\hat\delta_d]=\frac{2\sigma^2}{n},\qquad
\V[\hat\delta_c]=\frac{4\sigma^2(1-\rho)}{n},\qquad
\V[\hat\delta_a]=\frac{2\sigma^2(1-\rho^2)}{n}.\label{eq:08-var}
\end{equation}
\end{meidai}
\begin{proof}
群は独立．$\V[\bar Y_g-\bar X_g]=\frac{\sigma^2+\sigma^2-2\rho\sigma^2}{n}=\frac{2\sigma^2(1-\rho)}{n}$ を2群分足して $\V[\hat\delta_c]$．
$\V[\bar Y_g-\rho\bar X_g]=\frac{\sigma^2+\rho^2\sigma^2-2\rho\cdot\rho\sigma^2}{n}=\frac{\sigma^2(1-\rho^2)}{n}$ から $\V[\hat\delta_a]$．
\end{proof}
\begin{meidai}[期待値と比較]\label[meidai]{prop:08-compare}
(1) $\E[\hat\delta_d]=\mu_T^Y-\mu_C^Y$，$\E[\hat\delta_c]=(\mu_T^Y-\mu_C^Y)-(\mu_T^X-\mu_C^X)$，$\E[\hat\delta_a]=(\mu_T^Y-\mu_C^Y)-\rho(\mu_T^X-\mu_C^X)$ であり，
3つが一致するのはベースラインの母平均が等しい $\mu_T^X=\mu_C^X$ ときである（ランダム化試験では成り立つ）．\\
(2) 相対分散は $\V[\hat\delta_c]/\V[\hat\delta_d]=2(1-\rho)$，$\V[\hat\delta_a]/\V[\hat\delta_d]=1-\rho^2$．
変化量が最終値より有利なのは $\rho>1/2$ のときに限る．ANCOVA は常に両者以下である．
\end{meidai}
\begin{proof}
(2) $1-\rho^2\le1$ は明らか．$2(1-\rho)-(1-\rho^2)=(1-\rho)^2\ge0$ より $1-\rho^2\le2(1-\rho)$．
\end{proof}

\begin{figure}[htbp]
\centering
\begin{tikzpicture}
\begin{axis}[width=80mm,height=58mm,xmin=0,xmax=1,ymin=0,ymax=2.05,
  xlabel={相関係数 $\rho$},ylabel={相対分散（最終値$=1$）},grid=major,grid style={gray!20},
  legend style={font=\scriptsize,at={(0.98,0.98)},anchor=north east},tick label style={font=\scriptsize}]
\addplot[thick,dashdotted,bkgray,domain=0:1] {1};
\addplot[thick,dashed,bkred,domain=0:1] {2*(1-x)};
\addplot[very thick,bkblue,domain=0:1,samples=60] {1-x^2};
\legend{最終値 $\hat\delta_d$,変化量 $\hat\delta_c$：$2(1-\rho)$,ANCOVA $\hat\delta_a$：$1-\rho^2$}
\end{axis}
\end{tikzpicture}
\caption{ベースライン調整の3推定量の相対分散（\kakomon{2024}{2}〔5〕）}\label{fig:08-relvar}
\end{figure}

\paragraph{実際の ANCOVA}
実際には傾きは未知で，群内のプールした回帰係数 $\hat b=\sum_g\sum_i(X_i-\bar X_g)(Y_i-\bar Y_g)/\sum_g\sum_i(X_i-\bar X_g)^2$ を用いる．
$\hat b$ の推定による分散の増加は $O(1/n^2)$ で，大標本では\cref{eq:08-var}の $\V[\hat\delta_a]$ が成り立つ．
ランダム化試験ではベースラインの群間差は偶然によるものだけなので，ANCOVA は偏りなく精度を上げる．
観察研究でベースラインが群間で系統的に異なる場合は，変化量と ANCOVA で推定対象が異なり結論が食い違うことがある（Lord のパラドックス）．

\subsection{平均への回帰}\label{sec:08-rtm}
\Youchui $(X,Y)$ が平均 $\mu$，分散 $\sigma^2$，相関 $\rho$ の2変量正規分布に従うとき，
\[
\E[Y\mid X=x]=\mu+\rho(x-\mu),\qquad \V[Y\mid X=x]=\sigma^2(1-\rho^2).
\]
$\rho<1$ なら，1回目に高かった人の2回目の期待値は平均の方へ戻る．これを平均への回帰（regression to the mean）という．
\paragraph{真値モデル}
$X=\theta+\varepsilon_1$，$Y=\theta+\varepsilon_2$，$\theta\sim\Normal(\mu,\tau^2)$，$\varepsilon_j\sim\Normal(0,\psi^2)$（すべて独立）とすると
$\V[X]=\tau^2+\psi^2$，$\Cov(X,Y)=\tau^2$，$\rho=\tau^2/(\tau^2+\psi^2)$（信頼性係数）であり，
\[
\E[\theta\mid X=x]=\mu+\rho(x-\mu),\qquad \V[\theta\mid X=x]=\tau^2(1-\rho).
\]
1回目の値が高い人は「真値が高い」ことに加えて「たまたま正の測定誤差が乗った」人が多いので，真値の条件付き平均は $x$ より平均寄りになる．
\kakomon{2022}{5} では $\mu=120$，$\sigma=12$，$\rho=0.75$ で，$X=132$ の人の2回目の期待値は $120+0.75\times12=129$．
4 mmHg の低下のうち3 mmHg は平均への回帰で説明される．
\paragraph{臨床試験への含意}
「ベースラインが基準値以上の患者を登録する」試験では，無治療でも改善して見える．
単群の前後比較ではこれと治療効果を区別できないので，並行対照群を置き，ANCOVA で調整する．

\section{反復測定データ}\label{sec:08-repeated}
\subsection{平均ベクトルの推定}\label{subsec:08-mean}
$N$ 人の $T$ 時点の測定値 $\bm Y_i=(Y_{i1},\dots,Y_{iT})\transp\iid\Normal_T(\bm\mu,\bm\Sigma)$ とする（欠測なし）．
$\hat{\bm\mu}=\frac1N\sum\bm Y_i\sim\Normal_T(\bm\mu,\bm\Sigma/N)$ であり（\kakomon{2023}{4}〔1〕），
時点 $s$ から $t$ への変化量の推定量 $\hat\mu_t-\hat\mu_s$ の分散は
\begin{equation}
\V[\hat\mu_t-\hat\mu_s]=\frac{\sigma_{tt}+\sigma_{ss}-2\sigma_{st}}{N}.\label{eq:08-change}
\end{equation}
\kakomon{2023}{4} では Toeplitz 構造の推定値 $\hat\sigma_0=2.56$，$\hat\sigma_2=1.48$，$N=40$ から
$\widehat\V=(2.56+2.56-2\times1.48)/40=0.054$，$\SE=0.23$，変化量 $0.54$ に対し $Z=0.54/0.23\approx2.35$（丸めずに計算すると $2.32$）で有意となる．

\subsection{共分散構造}\label{subsec:08-cov}
\begin{table}[htbp]
\centering\small
\caption{主な共分散構造（$T$ 時点）}\label{tab:08-cov}
\begin{tabular}{@{}lL{62mm}c@{}}
\toprule
構造 & $(s,t)$ 成分 & パラメータ数\\
\midrule
独立 & $\sigma^2\ind(s=t)$ & 1\\
複合対称（CS，交換可能） & $\sigma_b^2+\sigma^2\ind(s=t)$ & 2\\
AR(1) & $\sigma^2\phi^{|s-t|}$ & 2\\
Toeplitz & $\sigma_{|s-t|}$（時点差ごとに共通） & $T$\\
無構造（UN） & $\sigma_{st}$（すべて自由） & $T(T+1)/2$\\
\bottomrule
\end{tabular}
\end{table}
構造の選択には AIC（$=-2\ell+2\times$（平均と分散のパラメータ数））を使う．
\kakomon{2023}{4}（$T=3$）では Toeplitz が平均3＋分散3で6個，無構造が3＋6で9個であり，
$\mathrm{AIC}=412.84+12=424.84$ 対 $409.74+18=427.74$ で Toeplitz が選ばれる．
ただし差は2.9と小さく，決定的とは言えない．
\begin{supbox}[補足：REML]
分散成分の最尤推定量は平均の推定による自由度の損失を考慮しないため下方に偏る（$\hat\sigma^2=\frac1n\sum(x_i-\bar x)^2$ と同じ現象）．
混合モデルの実務では制限付き最尤法（REML）が標準である．固定効果が異なるモデルの比較を REML 尤度の AIC で行ってはいけない．
\end{supbox}

\section{線形混合効果モデル}\label{sec:08-lmm}
\subsection{モデル}\label{subsec:08-lmmdef}
\begin{teigi}[線形混合効果モデル]\label{def:08-lmm}
\begin{equation}
\bm Y=\bm X\bm\beta+\bm Z\bm\gamma+\bm\varepsilon,\qquad \bm\gamma\sim\Normal(\bm0,\bm G),\quad \bm\varepsilon\sim\Normal(\bm0,\bm R),\quad \bm\gamma\indep\bm\varepsilon.\label{eq:08-lmm}
\end{equation}
$\bm\beta$ は\textbf{固定効果}（集団平均の構造），$\bm\gamma$ は\textbf{変量効果}（個体ごとの確率的なずれ）．
\end{teigi}
\begin{meidai}[周辺分布（\kakomon{2022}{3}〔1〕）]\label[meidai]{prop:08-marginal}
$\E[\bm Y]=\bm X\bm\beta$，$\V[\bm Y]=\bm V=\bm Z\bm G\bm Z\transp+\bm R$．
\end{meidai}
\begin{proof}
$\bm X\bm\beta$ は定数，$\bm\gamma$ と $\bm\varepsilon$ は独立なので
$\V[\bm Y]=\V[\bm Z\bm\gamma]+\V[\bm\varepsilon]=\bm Z\E[\bm\gamma\bm\gamma\transp]\bm Z\transp+\bm R=\bm Z\bm G\bm Z\transp+\bm R$．
\end{proof}

\subsection{一般化最小二乗推定}\label{subsec:08-gls}
\begin{teiri}[GLS 推定量]\label[teiri]{thm:08-gls}
$\bm V$ が既知なら $\bm\beta$ の最尤推定量（＝一般化最小二乗推定量）は
\begin{equation}
\hat{\bm\beta}=(\bm X\transp\bm V^{-1}\bm X)^{-1}\bm X\transp\bm V^{-1}\bm Y,\qquad
\E[\hat{\bm\beta}]=\bm\beta,\qquad \V[\hat{\bm\beta}]=(\bm X\transp\bm V^{-1}\bm X)^{-1}.\label{eq:08-gls}
\end{equation}
\end{teiri}
\begin{proof}
対数尤度の $\bm\beta$ を含む部分は $-\frac12(\bm Y-\bm X\bm\beta)\transp\bm V^{-1}(\bm Y-\bm X\bm\beta)$ で，これを最大化すると正規方程式
$\bm X\transp\bm V^{-1}\bm X\bm\beta=\bm X\transp\bm V^{-1}\bm Y$．$\E[\bm Y]=\bm X\bm\beta$ から不偏性，
$\V[\hat{\bm\beta}]=\bm A\bm V\bm A\transp$，$\bm A=(\bm X\transp\bm V^{-1}\bm X)^{-1}\bm X\transp\bm V^{-1}$ を計算すればよい．
\end{proof}
実際は $\bm V$ を推定値 $\hat{\bm V}$ で置き換える．
変量効果がなく $\bm R=\sigma^2\bm I$ なら $\hat{\bm\beta}=(\bm X\transp\bm X)^{-1}\bm X\transp\bm Y$（通常の最小二乗；\kakomon{2022}{3}〔2〕）．
\kakomon{2022}{3} の2群（各2人×3時点）では，切片はプラセボ群の6測定値の平均 $121.0/6=20.17$，
治療効果はサクシマー群の平均 $70.3/6=11.72$ との差 $-8.45$ になる．

\subsection{変量切片と複合対称}\label{subsec:08-ri}
被験者 $i$ の $j$ 回目の測定を $Y_{ij}=\bm x_{ij}\transp\bm\beta+b_i+\varepsilon_{ij}$，$b_i\sim\Normal(0,\sigma_b^2)$，$\varepsilon_{ij}\sim\Normal(0,\sigma^2)$ とする（変量切片モデル）．
被験者ごとのブロックで $\bm Z_i=\bm 1_{m}$，$G=\sigma_b^2$ だから
\[
\bm V_i=\sigma_b^2\bm 1\bm 1\transp+\sigma^2\bm I=
\begin{pmatrix}\sigma^2+\sigma_b^2&\sigma_b^2&\cdots\\ \sigma_b^2&\sigma^2+\sigma_b^2&\cdots\\ \vdots&&\ddots\end{pmatrix},
\qquad \Corr(Y_{ij},Y_{ik})=\frac{\sigma_b^2}{\sigma_b^2+\sigma^2}\ (j\ne k).
\]
これが複合対称構造であり，相関（級内相関係数，ICC）は時点差によらない．
\kakomon{2022}{3}〔3〕のように被験者ごとに異なる変量切片の分散 $\sigma_i$（過去問の記法で，ここでの $\sigma_i$ は分散）を許すなら，$\bm Z$ を被験者の指示行列（$12\times4$），$\bm G=\mathrm{diag}(\sigma_1,\dots,\sigma_4)$ とすればよい．
変量傾きモデル $Y_{ij}=\bm x_{ij}\transp\bm\beta+b_{0i}+b_{1i}t_j+\varepsilon_{ij}$ では $\bm Z_i=(\bm1,\bm t)$，$\bm G$ は $2\times2$ で，分散が時間とともに変化する構造になる．

\paragraph{相関を無視したときの誤り}
各被験者 $m$ 回測定の変量切片モデルで被験者レベルの共変量（治療群）の効果を推定するとき，
各群 $n$ 人の群平均の分散は $\frac{\sigma_b^2+\sigma^2/m}{n}=\frac{(\sigma_b^2+\sigma^2)\{1+(m-1)\mathrm{ICC}\}}{nm}$ である．
独立を仮定した OLS の分散は真の分散の $1/\{1+(m-1)\mathrm{ICC}\}$ となり（$1+(m-1)\mathrm{ICC}$ がデザイン効果），過小評価されて有意になりやすくなる．
変量効果を入れたモデル（\kakomon{2022}{3}〔3〕）は，この被験者内相関を正しく考慮した推測を可能にする．

\subsection{固定効果と変量効果}\label{subsec:08-fixrand}
\begin{itemize}
\item \textbf{固定効果}：関心のある特定の水準の効果（治療，時点，性別）．水準の効果そのものを推定・比較する．
\item \textbf{変量効果}：母集団から抽出されたとみなす水準（被験者，施設，試験）の効果．個々の値ではなく分散 $\sigma_b^2$ に関心がある．
\end{itemize}
メタアナリシスの変量効果モデル（\cref{ch:15}）も，試験を変量効果とする混合モデルの一種である．

\section{仮定}\label{sec:08-assump}
\begin{itemize}
\item 個体間の独立性（個体内の相関は共分散構造でモデル化）．
\item 多変量正規性，平均構造と共分散構造の正しい特定．
\item ANCOVA：ベースラインと最終値の関係が線形で，傾きが群間で共通（交互作用なし）．
\item 欠測がある場合，尤度に基づく混合モデルは欠測が MAR なら妥当（\cref{ch:16}）．完全ケース解析は MCAR を要する．
\end{itemize}

\section{結果の解釈}\label{sec:08-interp}
\begin{itemize}
\item ランダム化試験では，最終値・変化量・ANCOVA はいずれも同じ治療効果を不偏に推定するが，精度が異なる．$\rho$ が大きいほど ANCOVA の利得が大きい．
\item 前後比較の改善は，治療効果・自然経過・平均への回帰・プラセボ効果の合計である．
\item 共分散構造の選択は推定値そのものより標準誤差に影響する．
\end{itemize}

\section{手法選択}\label{sec:08-choice}
\begin{itemize}
\item 前後2時点の RCT → ANCOVA（ベースラインを共変量）．
\item 多時点・欠測あり → 混合モデル（MMRM：時点×群を固定効果，無構造共分散）．
\item 個体差の構造そのものに関心 → 変量切片・変量傾きモデル．
\item 2 値・カウントの反復測定 → GEE または一般化線形混合モデル．
\end{itemize}

\section{よくある誤り}\label{sec:08-err}
\begin{warnbox}
\begin{itemize}
\item 変化量を使えば常に精度が上がると考える（$\rho<1/2$ ではかえって悪化）．
\item $\V[\hat\delta_c]$ で群内の共分散の項を2群分数え忘れる．
\item 反復測定の全観測を独立とみなして $t$ 検定・OLS を行う．
\item 無構造のパラメータ数を $T^2$ とする（正しくは $T(T+1)/2$）．AIC で平均のパラメータを数え忘れる．
\item 前後比較の改善をそのまま治療効果とする（平均への回帰）．
\end{itemize}
\end{warnbox}

\section{数理統計との接続}\label{sec:08-link}
\begin{linkbox}
\begin{itemize}
\item 2変量正規の条件付き分布（『現代数理統計学の基礎』第4章演習）が ANCOVA と平均への回帰の基礎である．
      真値モデルの $\E[\theta\mid X]$ は正規--正規モデルのベイズ事後平均（同書 例6.13）と同じ形の縮小である．
\item 多変量正規の線形変換（同書 命題4.25）から $\hat{\bm\mu}$ と GLS 推定量の分布が得られる．
\item 条件付き分散・共分散の公式（同書 命題4.11，第4章演習）で $\V[\bm Y]=\E[\V[\bm Y\mid\bm\gamma]]+\V[\E[\bm Y\mid\bm\gamma]]=\bm R+\bm Z\bm G\bm Z\transp$ とも導ける．
\item 分散の MLE の除数 $n$ による偏り（同書 式(5.5)，\kakomonSM{2017}{1}）が REML の動機である．
\item \kakomonSM{2022}{5} は対応のある比較を二元配置分散分析として扱い，被験者を固定効果とみるか変量効果とみるかの考え方に通じる．
\end{itemize}
\end{linkbox}

\section{例題}\label{sec:08-ex}
\begin{reidai}[ベースライン調整の効率]\label{reidai:08-1}
各群50人のランダム化試験で，ベースラインと最終値の標準偏差がともに10，相関が $\rho=0.6$ とする．
\begin{enumerate}[label=(\arabic*)]
\item $\hat\delta_d,\hat\delta_c,\hat\delta_a$ の標準誤差を求めよ．
\item ANCOVA で最終値の比較と同じ精度を得るには，各群何人でよいか．
\item $\rho=0.4$ なら変化量と最終値のどちらが精度が高いか．
\end{enumerate}
\end{reidai}

\begin{reidai}[平均への回帰]\label{reidai:08-2}
ある検査値は集団で平均130，標準偏差15，2回の測定の相関0.7の2変量正規分布に従う．
1回目が160以上の人だけを対象に無治療で2回目を測定する．
\begin{enumerate}[label=(\arabic*)]
\item 1回目が160だった人の2回目の条件付き期待値と条件付き標準偏差を求めよ（$\sqrt{0.51}=0.714$）．
\item 真値モデルで考えたとき，測定誤差の分散 $\psi^2$ と真値の分散 $\tau^2$ を求めよ．
\item この対象者で2回目の平均が下がったことを「自然に改善した」と解釈してよいか．
\end{enumerate}
\end{reidai}

\begin{reidai}[反復測定]\label{reidai:08-3}
30人の3時点（0，4，8週）の測定値の平均ベクトルの推定値が $(10.2,9.1,8.3)$，
共分散行列の推定値が $\begin{pmatrix}4.0&2.4&2.0\\2.4&4.4&2.6\\2.0&2.6&4.8\end{pmatrix}$ である．
\begin{enumerate}[label=(\arabic*)]
\item 0週から8週の変化量の推定値と標準誤差を求め，$H_0:\mu_8=\mu_0$ を有意水準5\%で検定せよ．
\item 0週と8週の相関係数を求めよ（$\sqrt{19.2}=4.382$）．
\item 4時点の研究で，CS，AR(1)，Toeplitz，無構造の共分散パラメータ数を答えよ．
\item 3時点で，無構造（最大対数尤度 $-150.2$）と AR(1)（$-153.9$）の AIC を比較せよ．平均は各時点で自由とする．
\end{enumerate}
\end{reidai}

\begin{reidai}[変量切片モデル]\label{reidai:08-4}
$Y_{ij}=\mu+\delta z_i+b_i+\varepsilon_{ij}$（$j=1,\dots,m$，$b_i\sim\Normal(0,\sigma_b^2)$，$\varepsilon_{ij}\sim\Normal(0,\sigma^2)$，独立）で，各群 $n$ 人とする．
\begin{enumerate}[label=(\arabic*)]
\item 被験者 $i$ の $m$ 個の測定値の共分散行列を書け．
\item 群平均の差 $\hat\delta$ の分散を求め，全観測を独立とみなした場合の分散との比（デザイン効果）を求めよ．
\item $\sigma_b^2=3$，$\sigma^2=1$，$m=4$ のときデザイン効果はいくらか．
\end{enumerate}
\end{reidai}

\section{解答}\label{sec:08-sol}
\begin{kaitou}{reidai:08-1}
(1) $\V[\hat\delta_d]=2\times100/50=4$（SE 2.00），$\V[\hat\delta_c]=4\times100\times0.4/50=3.2$（SE 1.79），
$\V[\hat\delta_a]=2\times100\times0.64/50=2.56$（SE 1.60）．\\
(2) 分散は $n$ に反比例するので $50\times(1-\rho^2)=50\times0.64=32$ 人．\\
(3) $2(1-0.4)=1.2>1$ なので最終値の方が精度が高い．
\end{kaitou}

\begin{kaitou}{reidai:08-2}
(1) $\E=130+0.7\times30=151$，$\SD=15\sqrt{1-0.49}=15\times0.714=10.7$．\\
(2) $\rho=\tau^2/(\tau^2+\psi^2)=0.7$，$\tau^2+\psi^2=225$ より $\tau^2=157.5$，$\psi^2=67.5$．\\
(3) 解釈できない．1回目が高い人には正の測定誤差が乗った人が多く，2回目は平均の方へ戻る（1回目160の人でも期待値151）．
対照群なしでは平均への回帰と自然経過を区別できない．
\end{kaitou}

\begin{kaitou}{reidai:08-3}
(1) $8.3-10.2=-1.9$．$\V=(4.0+4.8-2\times2.0)/30=4.8/30=0.16$，$\SE=0.40$．$Z=-1.9/0.40=-4.75$，$|Z|>1.96$ で有意．\\
(2) $2.0/\sqrt{4.0\times4.8}=2.0/4.382=0.456$．\\
(3) CS 2，AR(1) 2，Toeplitz 4，無構造 $4\cdot5/2=10$．\\
(4) 無構造：$300.4+2\times(3+6)=318.4$．AR(1)：$307.8+2\times(3+2)=317.8$．AR(1) の AIC が小さい．
\end{kaitou}

\begin{kaitou}{reidai:08-4}
(1) 対角成分 $\sigma_b^2+\sigma^2$，非対角成分 $\sigma_b^2$ の $m\times m$ 行列（複合対称）．\\
(2) 被験者平均 $\bar Y_{i\cdot}$ の分散は $\sigma_b^2+\sigma^2/m$ で，群平均はその $n$ 人の平均なので
$\V[\hat\delta]=2(\sigma_b^2+\sigma^2/m)/n$．独立とみなすと $2(\sigma_b^2+\sigma^2)/(nm)$．
比は $\dfrac{m\sigma_b^2+\sigma^2}{\sigma_b^2+\sigma^2}=1+(m-1)\mathrm{ICC}$．\\
(3) $\mathrm{ICC}=3/4$，$1+3\times0.75=3.25$．独立を仮定すると分散を約3分の1に過小評価する．
\end{kaitou}

\begin{summarybox}
\begin{itemize}
\item $\V[\hat\delta_d]=2\sigma^2/n$，$\V[\hat\delta_c]=4\sigma^2(1-\rho)/n$，$\V[\hat\delta_a]=2\sigma^2(1-\rho^2)/n$．変化量は $\rho>1/2$ のときのみ有利，ANCOVA は常に最良．
\item 3推定量の期待値の一致はベースラインの母平均が等しいとき（ランダム化）．
\item 平均への回帰：$\E[Y\mid X=x]=\mu+\rho(x-\mu)$．単群前後比較の改善には回帰分が含まれる．
\item $\hat{\bm\mu}\sim\Normal(\bm\mu,\bm\Sigma/N)$，変化量の分散 $(\sigma_{tt}+\sigma_{ss}-2\sigma_{st})/N$．
\item 共分散構造のパラメータ数：CS 2，AR(1) 2，Toeplitz $T$，UN $T(T+1)/2$．AIC で選択．
\item LMM：$\bm V=\bm Z\bm G\bm Z\transp+\bm R$，$\hat{\bm\beta}=(\bm X\transp\bm V^{-1}\bm X)^{-1}\bm X\transp\bm V^{-1}\bm Y$．変量切片 → 複合対称，ICC $=\sigma_b^2/(\sigma_b^2+\sigma^2)$．
\end{itemize}
\end{summarybox}

% ===== ch09.tex =====
\chapter{生存時間解析の基礎}\label{ch:09}

\begin{goalbox}
\begin{itemize}
\item 生存関数・密度・ハザード・累積ハザードの関係 $h=f/S$，$H=-\log S$，$S=e^{-H}$，$f=hS$ を導出できる．
\item $H(T)\sim\Ex(1)$ を示し，比例ハザードモデルからの乱数生成に使える．
\item 打ち切りの種類と「独立な打ち切り」の仮定を説明できる．
\item リスク集合を数え，Kaplan--Meier 推定値の表を作成し，生存曲線を描き，中央生存時間を読み取れる．
\item Greenwood の分散公式をデルタ法で導出できる．
\item Nelson--Aalen 推定量を計算し，$\tilde S(t)\ge\hat S(t)$ を示せる．
\item ログランク検定の $O-E$ と分散を時点ごとの $2\times2$ 表から導出し，手計算できる．
\end{itemize}
\end{goalbox}

\begin{exambox}
\begin{itemize}
\item \kakomon{2016}{3}：2群各5例の KM 曲線，ログランク検定のスコア $u=\sum(d_t-n_{1t}/n_t)$ と分散 $v=\sum\frac{n_{1t}}{n_t}\frac{n_t-n_{1t}}{n_t}$ から $\chi^2=5.62$，$P=0.018$．
\item \kakomon{2017}{1}〔1〕〜〔3〕：$\lambda=f/S$，$\Lambda=-\log S$，$\Prob(\Lambda(T)>t)=e^{-t}$ の証明．
\item \kakomon{2019}{1}〔1〕〔2〕：5例の KM 推定値の表と曲線，Nelson--Aalen 推定値が KM 以上であることの証明．
\item \kakomon{2021}{1}〔1〕〔2〕：$f=\lambda S$，8例の KM から $1-\hat S(t)$ のグラフ（競合リスクとの比較は\cref{ch:11}）．
\item \kakomon{2022}{1}〔2〕：$S(t;x)=S_0(t)^{\exp(\beta x)}$ と $S_0(t)=\exp(-\int_0^t\lambda_0)$．
\end{itemize}
Cox モデル・指数モデル・RMST は\cref{ch:10}，競合リスクは\cref{ch:11}で扱う．
\end{exambox}

\flowline{時間＋打ち切り\\$(Y_i,\delta_i)$}{$S(t)$・中央値\\群間差}{KM・NA\\ログランク}{独立打ち切り\\比例ハザード}{リスク集合\\$O-E$・分散}{生存曲線\\の比較}

\section{問題設定}\label{sec:09-setting}
被験者 $i$ のイベント（死亡・再発など）までの時間を $T_i$，打ち切り時間を $C_i$ とすると，観測されるのは
\[
Y_i=\min(T_i,C_i),\qquad \delta_i=\ind(T_i\le C_i)\quad(1\text{：イベント},\ 0\text{：打ち切り})
\]
である．打ち切りのある観測 $(Y_i,\delta_i=0)$ は「$T_i>Y_i$」という情報しか与えない．
打ち切りを除外したり，打ち切り時点をイベント時点とみなしたりすると偏った推定になる．

\begin{warnbox}[打ち切り指示変数の向きに注意]
過去問では年度によって指示変数の向きが逆である．
\kakomon{2016}{3}，\kakomon{2019}{1} は「1：打ち切り，0：イベント」，\kakomon{2018}{1}，\kakomon{2021}{1} は「0：打ち切り，1：イベント」．
表を読む前に必ず確認すること．
\end{warnbox}

\section{生存時間の分布を表す関数}\label{sec:09-functions}
\Hissu
\begin{teigi}[生存関数・ハザード関数]\label{def:09-hazard}
$T\ge0$ を連続な確率変数，分布関数 $F$，密度 $f$ とする．
\begin{align*}
S(t)&=\Prob(T>t)=1-F(t)\quad\text{（生存関数）},\\
h(t)&=\lim_{\Delta\to0}\frac{\Prob(t\le T<t+\Delta\mid T\ge t)}{\Delta}\quad\text{（ハザード関数）},\\
H(t)&=\int_0^th(u)\,\dd u\quad\text{（累積ハザード関数）}.
\end{align*}
\end{teigi}
ハザードは「時点 $t$ まで生存した人が，次の瞬間にイベントを起こす率」であり，確率ではない（1を超えうる．単位は時間$^{-1}$）．

\begin{teiri}[基本関係式]\label[teiri]{thm:09-relations}
\begin{equation}
h(t)=\frac{f(t)}{S(t)}=-\frac{\dd}{\dd t}\log S(t),\qquad H(t)=-\log S(t),\qquad S(t)=e^{-H(t)},\qquad f(t)=h(t)S(t).\label{eq:09-relations}
\end{equation}
\end{teiri}
\begin{proof}
条件付き確率の定義と $\{t\le T<t+\Delta\}\subset\{T\ge t\}$ より
\[
h(t)=\lim_{\Delta\to0}\frac{1}{\Delta}\frac{\Prob(t\le T<t+\Delta)}{\Prob(T\ge t)}=\frac{1}{S(t)}\lim_{\Delta\to0}\frac{F(t+\Delta)-F(t)}{\Delta}=\frac{f(t)}{S(t)}.
\]
$f=-S'$ なので $h=-S'/S=-(\log S)'$．$S(0)=1$ より $H(t)=\int_0^t-(\log S)'\,\dd u=-\log S(t)+\log S(0)=-\log S(t)$．
残りの2式はこれを書き直したもの．
\end{proof}

\begin{meidai}[$H(T)$ は標準指数分布（\kakomon{2017}{1}〔3〕）]\label[meidai]{prop:09-HT}
$T$ が連続なら $\Prob(H(T)>t)=e^{-t}$，すなわち $H(T)\sim\Ex(1)$．同値に $S(T)\sim\Unif(0,1)$．
\end{meidai}
\begin{proof}
確率積分変換より $F(T)\sim\Unif(0,1)$．
$\Prob(H(T)>t)=\Prob(-\log S(T)>t)=\Prob(S(T)<e^{-t})=\Prob(F(T)>1-e^{-t})=e^{-t}$．
\end{proof}
\paragraph{乱数生成への応用}
$V\sim\Unif(0,1)$ なら $T=H^{-1}(-\log V)$ は累積ハザード $H$ をもつ．
たとえば $h(t)=\frac32\sqrt t\,e^{z/2}$（\kakomon{2017}{1}〔6〕）なら $H(t)=t^{3/2}e^{z/2}$ より $T=\{-e^{-z/2}\log V\}^{2/3}$．

\subsection{代表的なパラメトリック分布}\label{subsec:09-param}
\begin{table}[htbp]
\centering\small
\caption{代表的な生存時間分布}\label{tab:09-dist}
\begin{tabular}{@{}lllll@{}}
\toprule
分布 & $h(t)$ & $H(t)$ & $S(t)$ & 特徴\\
\midrule
指数 $\Ex(\lambda)$ & $\lambda$ & $\lambda t$ & $e^{-\lambda t}$ & ハザード一定，無記憶性，平均 $1/\lambda$\\
Weibull & $\alpha\gamma(\gamma t)^{\alpha-1}$ & $(\gamma t)^{\alpha}$ & $e^{-(\gamma t)^\alpha}$ & $\alpha>1$ 増加，$\alpha<1$ 減少\\
Gompertz & $ae^{bt}$ & $\frac ab(e^{bt}-1)$ & $e^{-H(t)}$ & 成人の死亡率\\
\bottomrule
\end{tabular}
\end{table}
Weibull 型の原因別ハザードは\kakomon{2025}{3}で用いられた．
中央生存時間 $t_{0.5}$ は $S(t_{0.5})=0.5$，指数分布なら $t_{0.5}=\log2/\lambda$．
平均生存時間は $\E[T]=\int_0^\infty S(t)\,\dd t$（生存曲線下の面積；\cref{sec:10-rmst}）．

\section{打ち切り}\label{sec:09-censor}
\begin{itemize}
\item \textbf{右側打ち切り}：イベントが観察終了時点より後に起こることだけがわかる．最も一般的．
      試験終了による\textbf{管理的打ち切り}，転居などによる追跡不能がある．
\item \textbf{左側打ち切り}：観察開始時にすでにイベントが起きていた．\textbf{区間打ち切り}：2回の受診の間に起きたことだけがわかる．
\item \textbf{独立（無情報）な打ち切り}：打ち切られた人のその後のハザードが，同じ時点でリスク集合に残っている人と同じであること．
      「状態が悪化したから来院しなくなった」といった打ち切りは，この仮定を破る（情報のある打ち切り）．
\end{itemize}
本章の方法（KM，NA，ログランク，Cox）はすべて独立な打ち切りを仮定する．
競合する別のイベント（他因死）を打ち切りとみなすと，この仮定の解釈が変わる（\cref{ch:11}）．

\section{Kaplan--Meier 推定量}\label{sec:09-km}
\Hissu\Hinshutsu
\subsection{定義と導出}
相異なるイベント時点を $t_{(1)}<\dots<t_{(r)}$ とし，
$n_j$ を時点 $t_{(j)}$ の直前にイベントも打ち切りも起きていない人数（\emphx{リスク集合}の大きさ），$d_j$ を $t_{(j)}$ のイベント数とする．
\begin{teigi}[Kaplan--Meier 推定量]\label{def:09-km}
\begin{equation}
\hat S(t)=\prod_{j:\,t_{(j)}\le t}\Bigl(1-\frac{d_j}{n_j}\Bigr).\label{eq:09-km}
\end{equation}
\end{teigi}
\paragraph{積極限としての導出}
$t_{(j)}\le t<t_{(j+1)}$ なら，条件付き確率の連鎖で
\[
S(t)=\Prob(T>t_{(1)})\,\Prob(T>t_{(2)}\mid T>t_{(1)})\cdots\Prob(T>t_{(j)}\mid T>t_{(j-1)})
\]
（イベントのない区間では生存確率は減らない）．各因子 $\Prob(T>t_{(k)}\mid T\ge t_{(k)})$ を，
時点 $t_{(k)}$ にリスク集合にいた $n_k$ 人のうち生き残った割合 $1-d_k/n_k$ で推定すると\cref{eq:09-km}を得る．
打ち切られた人は打ち切り時点まではリスク集合に含まれ，その後の分母から外れる．こうして打ち切りの情報を無駄なく使う．
\begin{supbox}[補足：ノンパラメトリック最尤推定量]
各時点のハザード $q_j$（離散）を未知母数とすると尤度は $\prod_jq_j^{d_j}(1-q_j)^{n_j-d_j}$ に比例し，
MLE $\hat q_j=d_j/n_j$ から $\hat S=\prod(1-\hat q_j)$ となる．KM はノンパラメトリック最尤推定量である．
\end{supbox}

\subsection{表の作り方（手順）}
\begin{enumerate}
\item 観察時間を昇順に並べる．同じ時点にイベントと打ち切りがあるときは，\textbf{イベントが先}（打ち切られた人はその時点のリスク集合に含める）．
\item 各イベント時点 $t_{(j)}$ で $n_j$（その時点以上の観察時間をもつ人数）と $d_j$ を数える．
\item $\hat S$ を前の値に $(1-d_j/n_j)$ を掛けて更新する．打ち切り時点では $\hat S$ は変わらず，次の $n$ が1減る．
\end{enumerate}
\kakomon{2019}{1} の5例（10，13，$18^+$，19，$23^+$）では，
$\hat S(10)=4/5=0.8$，$\hat S(13)=0.8\times3/4=0.6$，$\hat S(19)=0.6\times1/2=0.3$．
最長の観察が打ち切り（$23^+$）なので曲線は0に到達しない．

\subsection{曲線と中央生存時間}
KM 曲線は右連続な階段関数で，イベント時点でのみ下がる．打ち切り時点は曲線上に目印（$+$）を付けることが多い．
中央生存時間の推定値は $\hat S(t)\le0.5$ となる最小の $t$ である．曲線が0.5まで下がらなければ中央値は推定できない．

\begin{figure}[htbp]
\centering
\begin{tikzpicture}
\begin{axis}[width=90mm,height=55mm,xmin=0,xmax=16,ymin=0,ymax=1.05,
  xlabel={時間（週）},ylabel={$\hat S(t)$},xtick={0,3,6,10,15},ytick={0,0.25,0.5,0.75,1},
  grid=major,grid style={gray!20},tick label style={font=\scriptsize}]
\addplot[very thick,bkblue,const plot] coordinates {(0,1)(3,0.875)(6,0.5833)(10,0.3889)(15,0)(16,0)};
\addplot[only marks,mark=+,mark size=3pt,thick,bkorange] coordinates {(5,0.875)(8,0.5833)(12,0.3889)};
\addplot[dashed,gray] coordinates {(0,0.5)(16,0.5)};
\end{axis}
\end{tikzpicture}
\caption{\cref{reidai:09-2}の KM 曲線（$+$ は打ち切り）．中央生存時間は10週．}\label{fig:09-km}
\end{figure}

\subsection{Greenwood の分散公式}\label{subsec:09-greenwood}
\begin{teiri}[Greenwood の公式]\label[teiri]{thm:09-greenwood}
\begin{equation}
\widehat\V[\hat S(t)]=\hat S(t)^2\sum_{j:\,t_{(j)}\le t}\frac{d_j}{n_j(n_j-d_j)}.\label{eq:09-greenwood}
\end{equation}
\end{teiri}
\begin{proof}[導出（デルタ法）]
$\log\hat S(t)=\sum_j\log\hat p_j$，$\hat p_j=1-d_j/n_j$ とする．$n_j$ を与えると $n_j-d_j\sim\Bin(n_j,p_j)$ で，
異なる時点の $\hat p_j$ は（条件付きで）無相関とみなせる．デルタ法より
\[
\V[\log\hat p_j]\approx\frac{\V[\hat p_j]}{p_j^2}=\frac{p_j(1-p_j)}{n_jp_j^2}=\frac{1-p_j}{n_jp_j}\approx\frac{d_j}{n_j(n_j-d_j)}.
\]
よって $\V[\log\hat S]\approx\sum_jd_j/\{n_j(n_j-d_j)\}$．さらに $\V[\hat S]\approx\hat S^2\V[\log\hat S]$（$g(x)=e^x$ のデルタ法）．
\end{proof}
\kakomon{2016}{3} の解答の「標準偏差」欄はこの公式による．たとえば群2の時点1（$n=5,d=1$）で
$\sqrt{0.8^2\times\frac{1}{5\cdot4}}=0.179$，時点6で $\sqrt{0.3^2(\frac1{20}+\frac1{12}+\frac12)}=0.239$．
最後のイベントでリスク集合が尽きる（$n_j=d_j$）と公式は定義できない．
\begin{supbox}[補足：信頼区間]
$\hat S\pm1.96\SE$ は $[0,1]$ をはみ出すことがある．$\log(-\log\hat S)$ のスケールで区間を作り逆変換する方法がよく用いられる：
$\V[\log(-\log\hat S)]\approx\frac{\sum d_j/\{n_j(n_j-d_j)\}}{(\log\hat S)^2}$．
\end{supbox}

\section{Nelson--Aalen 推定量}\label{sec:09-na}
\begin{teigi}[Nelson--Aalen 推定量]\label{def:09-na}
\begin{equation}
\hat H(t)=\sum_{j:\,t_{(j)}\le t}\frac{d_j}{n_j},\qquad \tilde S(t)=\exp\{-\hat H(t)\}=\prod_{j:\,t_{(j)}\le t}\exp\Bigl(-\frac{d_j}{n_j}\Bigr).\label{eq:09-na}
\end{equation}
分散の推定量は $\widehat\V[\hat H(t)]=\sum d_j/n_j^2$．
\end{teigi}
各時点のハザードの増分 $d_j/n_j$ を足し上げたものが累積ハザードの推定量で，$S=e^{-H}$ の関係から生存関数の推定量が得られる．
\begin{meidai}[$\tilde S\ge\hat S$（\kakomon{2019}{1}〔2〕）]\label[meidai]{prop:09-nage}
任意の $t$ で $\tilde S(t)\ge\hat S(t)$．
\end{meidai}
\begin{proof}
$g(x)=e^{-x}-(1-x)$ は $g'(x)=1-e^{-x}$ より $x=0$ で最小値0をとるので $e^{-x}\ge1-x$．
$x=d_j/n_j\in[0,1]$ に適用して各因子で $\exp(-d_j/n_j)\ge1-d_j/n_j\ge0$．非負の数の積なので $\tilde S\ge\hat S$．
\end{proof}
$d_j/n_j$ が小さい（リスク集合が大きい）とき両者はほぼ一致する．
Nelson--Aalen 推定量の対数を $\log t$ に対してプロットすると，Weibull 性や比例ハザード性の確認に使える（\cref{sec:10-phcheck}）．

\section{ログランク検定}\label{sec:09-logrank}
\Hissu\Hinshutsu
\subsection{導出}
2群の生存関数の同一性 $H_0:S_1(t)=S_0(t)$（すべての $t$）を検定する．
各イベント時点 $t_{(j)}$ で，リスク集合を群×（イベントあり／なし）の $2\times2$ 表にまとめる．
\begin{center}\small
\begin{tabular}{@{}lccc@{}}
\toprule
時点 $t_{(j)}$ & イベント & イベントなし & リスク集合\\
\midrule
群1 & $d_{1j}$ & $n_{1j}-d_{1j}$ & $n_{1j}$\\
群0 & $d_{0j}$ & $n_{0j}-d_{0j}$ & $n_{0j}$\\
\midrule
計 & $d_j$ & $n_j-d_j$ & $n_j$\\
\bottomrule
\end{tabular}
\end{center}
$H_0$ の下で周辺度数を固定すると，$d_{1j}$ は超幾何分布に従い
\begin{equation}
E_{1j}=\E[d_{1j}]=d_j\frac{n_{1j}}{n_j},\qquad V_j=\V[d_{1j}]=\frac{d_j\,n_{1j}\,n_{0j}\,(n_j-d_j)}{n_j^2(n_j-1)}.\label{eq:09-lrmoments}
\end{equation}
\begin{teiri}[ログランク検定]\label[teiri]{thm:09-logrank}
\begin{equation}
U=\sum_j(d_{1j}-E_{1j})=O_1-E_1,\qquad V=\sum_jV_j,\qquad \chi^2_{\mathrm{LR}}=\frac{U^2}{V}\ \xrightarrow{d}\ \chi^2_1.\label{eq:09-logrank}
\end{equation}
\end{teiri}
時点ごとの表を層とみなした Mantel--Haenszel 検定（\cref{sec:04-mh}）そのものである．
異なる時点の $d_{1j}-E_{1j}$ は独立ではないが，過去を条件とした条件付き平均が0の和（マルチンゲール差）になり，分散は足し合わせられる．
\paragraph{各時点のイベント数が1の場合}
$d_j=1$ なら $V_j=\frac{n_{1j}n_{0j}}{n_j^2}=\frac{n_{1j}}{n_j}\Bigl(1-\frac{n_{1j}}{n_j}\Bigr)$ で，
「確率 $n_{1j}/n_j$ の Bernoulli 分布の分散」である（\kakomon{2016}{3}〔2〕の説明のとおり）．
\kakomon{2016}{3} では群1について $U=-2.3222$，$V=0.9591$，$\chi^2=5.62$，$\sqrt{5.62}=2.37$ から両側 $P=0.018$ となる．

\subsection{性質}
\begin{itemize}
\item ログランク検定は Cox 比例ハザードモデルの $H_0:\beta=0$ のスコア検定と一致する（同順位がなければ厳密に；\cref{sec:10-cox}）．
      したがって\textbf{比例ハザード型の局所対立仮説（HR が1に近い）の下で効率的}（局所的に最も検出力が高い）である．
\item 生存曲線が交差する（ハザード比が途中で逆転する）と，前半と後半の $O-E$ が打ち消し合い検出力を失う．
\item 両群の $O-E$ の和は0（$\sum_j(d_j-d_j)=0$）なので，どちらの群で計算しても $\chi^2$ は同じ．
\end{itemize}
\begin{supbox}[補足：重み付きログランク検定と層別]
時点ごとに重み $w_j$ を付けた $U_w=\sum w_j(d_{1j}-E_{1j})$，$V_w=\sum w_j^2V_j$ も検定になる．
$w_j=n_j$（一般化 Wilcoxon，Gehan）や $w_j=\hat S(t_{(j)}-)$，それに近い Peto--Prentice 重み $\prod_{i\le j}\{1-d_i/(n_i+1)\}$ は早期の差を重視する．
層別ログランク検定は層ごとの $U_s,V_s$ を足して $(\sum U_s)^2/\sum V_s$ とする．
\end{supbox}

\section{仮定}\label{sec:09-assump}
\begin{itemize}
\item 独立（無情報）な打ち切り．打ち切りが予後と関連すると KM は偏る．
\item 被験者の独立性．
\item 登録時期によって生存の分布が変わらないこと（登録が長期に及ぶ場合）．
\item ログランク検定の最適性は比例ハザードを前提とする（検定の妥当性自体は $H_0$ の下で保たれる）．
\end{itemize}

\section{結果の解釈}\label{sec:09-interp}
\begin{itemize}
\item KM 曲線は右端ほどリスク集合が小さく不安定である．リスク集合の人数を併記するのが望ましい．
\item 「5年生存率 $\hat S(5)$」と中央生存時間は KM から直接読める要約であり，比較にはログランク検定や HR を用いる．
\item ログランク検定が有意でも差の大きさはわからない．HR（\cref{ch:10}）や RMST の差（\cref{sec:10-rmst}）で大きさを示す．
\end{itemize}

\section{手法選択}\label{sec:09-choice}
\begin{itemize}
\item 生存曲線の推定 → KM（累積ハザードなら NA）．
\item 2群以上の比較 → ログランク検定（比例ハザードを想定）．早期の差を重視 → 一般化 Wilcoxon．
\item 共変量の調整・効果の大きさ → Cox 回帰．分布の仮定が妥当なら指数・Weibull モデル（\cref{ch:10}）．
\item 他因死などの競合イベント → 累積発生関数（\cref{ch:11}）．
\end{itemize}

\section{よくある誤り}\label{sec:09-err}
\begin{warnbox}
\begin{itemize}
\item 打ち切り例を分母から最初から除く，または打ち切り時点で死亡とみなす．
\item 同じ時点の打ち切りをイベントより先に処理してリスク集合を1減らす．
\item 打ち切り時点で KM 曲線を下げる．
\item ログランクの分散を $\sum E_{1j}$（ポアソン近似）で代用し，出題の指示（二項分散・超幾何分散）と異なる計算をする．
\item ハザードを「確率」と説明する．
\end{itemize}
\end{warnbox}

\section{数理統計との接続}\label{sec:09-link}
\begin{linkbox}
\begin{itemize}
\item ハザード関数と $F(x)=1-\exp\{-\int_0^x\lambda\}$ は『現代数理統計学の基礎』式(3.24)(3.25)，Weibull 分布は式(3.26)
      （同書は $f=abx^{b-1}e^{-ax^b}$ と書き，$a$ は累積ハザードの係数）．打ち切り・KM・ログランクは同書の範囲外である．
\item $E[X]=\int_0^\infty\{1-F(x)\}\dd x$（同書第2章演習，\kakomonSM{2024}{4}）は平均生存時間＝生存曲線下面積の根拠．
\item 確率積分変換（同書 命題2.19）が $H(T)\sim\Ex(1)$ の証明の要．\kakomonSM{2021}{1} は $h(x)=1-e^{-x}$ による $\Ex(1)$ と $\Unif(0,1)$ の対応を扱った．
\item 最小値のハザードは各ハザードの和（同書第5章演習）：競合リスク（\cref{ch:11}）で $\lambda=\sum\lambda_j$ となる理由．
\item Greenwood の公式は，二項比率の対数のデルタ法（同書第5章演習）を時点ごとに足し合わせたものである．
\end{itemize}
\end{linkbox}

\section{例題}\label{sec:09-ex}
\begin{reidai}[関数の関係]\label{reidai:09-1}
ハザード関数が $h(t)=2t/\theta^2$ $(t\ge0,\ \theta>0)$ の生存時間 $T$ について，
$H(t)$，$S(t)$，$f(t)$，中央生存時間を求めよ．また $\Unif(0,1)$ 乱数 $V$ から $T$ を生成する式を示せ．
\end{reidai}

\begin{reidai}[KM・Greenwood・NA]\label{reidai:09-2}
8例の観察時間（週，$+$ は打ち切り）が $3,\ 5^+,\ 6,\ 6,\ 8^+,\ 10,\ 12^+,\ 15$ であった．
\begin{enumerate}[label=(\arabic*)]
\item KM 推定値の表（$t_{(j)},n_j,d_j,\hat S$）を作れ．
\item 中央生存時間を求めよ．
\item $t=10$ における Greenwood の標準誤差を求めよ（$\sqrt{0.04051}=0.2013$）．
\item $t=10$ における Nelson--Aalen 推定値 $\hat H$ と $\tilde S=e^{-\hat H}$ を求め，$\hat S(10)$ と比べよ（$e^{-0.7917}=0.453$）．
\end{enumerate}
\end{reidai}

\begin{reidai}[ログランク検定の手計算]\label{reidai:09-3}
群A：$2,\ 4,\ 5^+,\ 7,\ 9^+$，群B：$6,\ 8,\ 10^+,\ 12,\ 14^+$（月）．
\begin{enumerate}[label=(\arabic*)]
\item 各イベント時点のリスク集合と群Aの期待イベント数，分散（\cref{eq:09-lrmoments}）を表にせよ．
\item ログランク統計量を求め，有意水準5\%で検定せよ．
\end{enumerate}
\end{reidai}

\begin{reidai}[概念]\label{reidai:09-4}
次の各記述の正誤を答え，誤りなら理由を述べよ．
\begin{enumerate}[label=(\alph*)]
\item ハザード関数の値は0から1の間にある．
\item 追跡不能の患者は，その時点で死亡したとみなすと保守的な推定になるので望ましい．
\item KM 曲線の最後の観察が打ち切りなら，曲線は0にならない．
\item 生存曲線が途中で交差する場合でも，ログランク検定は最も検出力が高い．
\end{enumerate}
\end{reidai}

\section{解答}\label{sec:09-sol}
\begin{kaitou}{reidai:09-1}
$H(t)=\int_0^t2u/\theta^2\,\dd u=t^2/\theta^2$，$S(t)=e^{-t^2/\theta^2}$，$f(t)=hS=\frac{2t}{\theta^2}e^{-t^2/\theta^2}$（Weibull，形状2）．
中央値：$t^2/\theta^2=\log2$ より $t_{0.5}=\theta\sqrt{\log2}$．
$H(T)\sim\Ex(1)$ と $-\log V\sim\Ex(1)$ より $T^2/\theta^2=-\log V$，$T=\theta\sqrt{-\log V}$．
\end{kaitou}

\begin{kaitou}{reidai:09-2}
(1)
\begin{center}\small
\begin{tabular}{@{}ccccc@{}}
\toprule
$t_{(j)}$ & $n_j$ & $d_j$ & $1-d_j/n_j$ & $\hat S$\\
\midrule
3 & 8 & 1 & 7/8 & 0.875\\
6 & 6 & 2 & 4/6 & 0.583\\
10 & 3 & 1 & 2/3 & 0.389\\
15 & 1 & 1 & 0 & 0\\
\bottomrule
\end{tabular}
\end{center}
（5週と8週の打ち切りでリスク集合が1ずつ減る．12週の打ち切りの後，15週のリスク集合は1人．）\\
(2) $\hat S(t)\le0.5$ となる最小の $t$ は10週（\cref{fig:09-km}）．\\
(3) $\sum d_j/\{n_j(n_j-d_j)\}=\frac{1}{8\cdot7}+\frac{2}{6\cdot4}+\frac{1}{3\cdot2}=0.0179+0.0833+0.1667=0.2679$．
$\widehat\V=0.389^2\times0.2679=0.0405$，$\SE=0.201$．\\
(4) $\hat H(10)=\frac18+\frac26+\frac13=0.792$，$\tilde S=e^{-0.792}=0.453>\hat S(10)=0.389$（\cref{prop:09-nage}のとおり）．
\end{kaitou}

\begin{kaitou}{reidai:09-3}
(1)
\begin{center}\small
\begin{tabular}{@{}cccccccc@{}}
\toprule
時点 & $n_{Aj}$ & $n_{Bj}$ & $n_j$ & $d_{Aj}$ & $d_j$ & $E_{Aj}=n_{Aj}/n_j$ & $V_j=n_{Aj}n_{Bj}/n_j^2$\\
\midrule
2 & 5 & 5 & 10 & 1 & 1 & 0.500 & 0.250\\
4 & 4 & 5 & 9 & 1 & 1 & 0.444 & 0.247\\
6 & 2 & 5 & 7 & 0 & 1 & 0.286 & 0.204\\
7 & 2 & 4 & 6 & 1 & 1 & 0.333 & 0.222\\
8 & 1 & 4 & 5 & 0 & 1 & 0.200 & 0.160\\
12 & 0 & 2 & 2 & 0 & 1 & 0 & 0\\
\midrule
計 & & & & 3 & & 1.764 & 1.083\\
\bottomrule
\end{tabular}
\end{center}
（群Aは5月の打ち切り後，6月には2人．）\\
(2) $U=3-1.764=1.236$，$\chi^2=1.236^2/1.083=1.41<3.841$ で有意でない．
群Aは期待より多くのイベントを起こしているが，標本が小さく差があるとは言えない．
\end{kaitou}

\begin{kaitou}{reidai:09-4}
(a) 誤り．ハザードは率で，単位時間あたりの値なので1を超えうる（時間の単位を変えると値が変わる）．\\
(b) 誤り．打ち切りを死亡とみなすと生存を過小評価する偏った推定になる．独立な打ち切りの下ではKMで正しく扱える．\\
(c) 正しい．最後の観察が打ち切りなら，その人はすべてのイベント時点のリスク集合に含まれるので $n_j>d_j$ となり，各因子 $1-d_j/n_j$ が正で積が0にならないため．\\
(d) 誤り．ログランク検定が効率的なのは比例ハザード型の（局所）対立仮説のとき．交差すると前後の差が打ち消し合う．
\end{kaitou}

\begin{summarybox}
\begin{itemize}
\item $h=f/S=-(\log S)'$，$H=-\log S$，$S=e^{-H}$，$f=hS$．$H(T)\sim\Ex(1)$ で乱数生成 $T=H^{-1}(-\log V)$．
\item 打ち切り指示変数の向きを確認．独立な打ち切りが全手法の前提．
\item KM：$\hat S(t)=\prod_{t_{(j)}\le t}(1-d_j/n_j)$．同時点ではイベントを先に．中央値は $\hat S\le0.5$ の最小時点．
\item Greenwood：$\hat S^2\sum d_j/\{n_j(n_j-d_j)\}$（デルタ法）．
\item NA：$\hat H=\sum d_j/n_j$，$e^{-\hat H}\ge\hat S$（$e^{-x}\ge1-x$）．
\item ログランク：$E_{1j}=d_jn_{1j}/n_j$，$V_j=d_jn_{1j}n_{0j}(n_j-d_j)/\{n_j^2(n_j-1)\}$（$d_j=1$ なら $n_{1j}n_{0j}/n_j^2$），$\chi^2=(O-E)^2/V$．比例ハザード型の局所対立仮説で効率的．
\end{itemize}
\end{summarybox}

% ===== ch10.tex =====
\chapter[Cox比例ハザード・RMST・パラメトリック生存解析]{Cox比例ハザード・RMST・\\パラメトリック生存解析}\label{ch:10}

\begin{goalbox}
\begin{itemize}
\item Cox 比例ハザードモデルを書き，回帰係数を対数ハザード比として解釈し，比例ハザード性の意味を説明できる．
\item $S(t\mid x)=S_0(t)^{\exp(\beta x)}$ と $\log\{-\log S(t\mid x)\}=\log H_0(t)+\beta x$ を導き，2重対数プロットで比例ハザード性を確認できる．
\item 部分尤度を小さなデータで書き下し，スコア検定がログランク検定に一致することを示せる．
\item 打ち切りのある指数分布の尤度・最尤推定量・Fisher情報量を導出し，デルタ法でハザード比の検定・区間を作れる．
\item RMST（境界内平均生存時間）が生存曲線下の面積であることを示し，指数モデルと KM で計算できる．
\end{itemize}
\end{goalbox}

\begin{exambox}
\begin{itemize}
\item \kakomon{2016}{3}〔3〕：2群のデータの Cox 部分尤度 $L(\beta)$ の書き下し．
\item \kakomon{2017}{1}〔4〕〜〔6〕：Cox モデルで $\log(-\log S(t\mid\bm Z))=\log\Lambda_0(t)+\bm\beta\transp\bm Z$，
      $\epsilon(T)$ が極値分布 $1-\exp(-e^x)$ に従うこと，年齢が一様分布の共変量のときの乱数生成．
\item \kakomon{2018}{1}：打ち切りのある指数分布の尤度，$\hat\lambda=\sum\delta_i/\sum t_i$，$I(\lambda)=\delta/\lambda^2$，
      デルタ法による $\V[\log\hat\lambda_1-\log\hat\lambda_2]=1/d_1+1/d_2$，$Z$ 検定，必要イベント数．
\item \kakomon{2019}{1}〔3〕〜〔5〕：RMST $=\int_0^\tau S$ の証明，指数分布での $\V[\min(T,\tau)]$，$\hat\lambda=3/83$ と $\tau=20$ の RMST・分散．
\item \kakomon{2022}{1}：比例ハザード性，$S_0(t)$ の表現，2重対数プロットが平行になる理由，$\hat\beta=0.06$（SE 0.18）の HR と CI の解釈．
\item \kakomon{2023}{2}：打ち切りなしの指数分布の尤度・Fisher情報量と症例数設計（\cref{ch:12}）．
\end{itemize}
\end{exambox}

\flowline{時間＋打ち切り\\＋共変量}{ハザード比\\RMST}{Cox・指数\\モデル}{比例ハザード\\独立打ち切り}{部分尤度\\デルタ法}{HR・RMST差\\の臨床的意味}

\section{Cox 比例ハザードモデル}\label{sec:10-cox}
\Hissu\Hinshutsu
\subsection{モデル}
\begin{teigi}[Cox 比例ハザードモデル]\label{def:10-cox}
共変量 $\bm x$ をもつ個体のハザードを
\begin{equation}
h(t\mid\bm x)=h_0(t)\exp(\bm\beta\transp\bm x)\label{eq:10-cox}
\end{equation}
とする．$h_0(t)$ はベースラインハザード（$\bm x=\bm0$ のハザード）で，形を特定しない（セミパラメトリックモデル）．
\end{teigi}
2つの個体のハザード比は $h(t\mid\bm x_1)/h(t\mid\bm x_2)=\exp\{\bm\beta\transp(\bm x_1-\bm x_2)\}$ で $t$ によらない．
これが\emphx{比例ハザード性}である（\kakomon{2022}{1}〔1〕）．
2 値の共変量なら $e^{\beta}$ が HR，連続なら1単位増加あたりの HR である．

\begin{meidai}[生存関数と2重対数変換]\label[meidai]{prop:10-loglog}
$H_0(t)=\int_0^th_0$，$S_0(t)=e^{-H_0(t)}$ とすると
\begin{equation}
S(t\mid\bm x)=S_0(t)^{\exp(\bm\beta\transp\bm x)},\qquad
\log\{-\log S(t\mid\bm x)\}=\log H_0(t)+\bm\beta\transp\bm x.\label{eq:10-loglog}
\end{equation}
\end{meidai}
\begin{proof}
$H(t\mid\bm x)=\int_0^th_0(u)e^{\bm\beta\transp\bm x}\dd u=e^{\bm\beta\transp\bm x}H_0(t)$．
$S=e^{-H}$ より $S(t\mid\bm x)=\{e^{-H_0(t)}\}^{\exp(\bm\beta\transp\bm x)}$．両辺の $-\log$ をとってから $\log$ をとれば第2式．
\end{proof}
\paragraph{2重対数プロット}\label{sec:10-phcheck}
群ごとの KM 推定値 $\hat S_g(t)$ について $\log\{-\log\hat S_g(t)\}$ を $\log t$（または $t$）に対してプロットすると，
比例ハザードの下では群間の差が $\beta$ で一定，すなわち曲線が\textbf{平行}になる（\kakomon{2022}{1}〔3〕）．
交差したり間隔が大きく変わったりすれば比例ハザード性が疑われる．
さらに直線になれば Weibull 分布（$\log H=\alpha\log\gamma+\alpha\log t$），傾き1の直線なら指数分布である．
\begin{supbox}[補足：その他の確認法]
Schoenfeld 残差を時間に対してプロットする（傾向があれば非比例），共変量と時間の関数の交互作用項 $\beta_2\,x\log t$ を加えて検定する，などがある．
非比例のときは，層別 Cox モデル（非比例の変数で層別し $h_{0s}(t)$ を層ごとに変える），時間依存係数，RMST による要約を検討する．
\end{supbox}

\paragraph{線形変換モデル（\kakomon{2017}{1}〔5〕）}
\cref{eq:10-loglog}を $t=T$ で評価すると $\log H_0(T)=-\bm\beta\transp\bm x+\epsilon$，$\epsilon=\log H(T\mid\bm x)$．
$H(T\mid\bm x)\sim\Ex(1)$（\cref{prop:09-HT}）なので
\[
\Prob(\epsilon\le u)=\Prob\bigl(H(T\mid\bm x)\le e^u\bigr)=1-\exp(-e^u),
\]
すなわち $\epsilon$ は（最小値の）極値分布に従う．
$H_0(t)=t^{3/2}$，$\beta=1/2$ なら $\frac32\log T=-\frac Z2+\epsilon$，$\epsilon=\log(-\log V)$ より
$T=\exp\{-Z/3+\frac23\log(-\log V)\}$ で乱数を作れる（$Z=40+30U$）．

\subsection{部分尤度}\label{subsec:10-partial}
\Hissu 同順位のないイベント時点 $t_{(1)}<\dots<t_{(r)}$，時点 $t_{(j)}$ でイベントを起こした個体を $(j)$，リスク集合を $R_j$ とする．
\begin{teiri}[Cox の部分尤度]\label[teiri]{thm:10-partial}
\begin{equation}
L(\bm\beta)=\prod_{j=1}^r\frac{\exp(\bm\beta\transp\bm x_{(j)})}{\sum_{l\in R_j}\exp(\bm\beta\transp\bm x_l)}.\label{eq:10-partial}
\end{equation}
\end{teiri}
\begin{proof}[導出]
時点 $t_{(j)}$ で $R_j$ の中からちょうど1個体がイベントを起こしたという条件の下で，それが個体 $(j)$ である条件付き確率は
\[
\frac{h(t_{(j)}\mid\bm x_{(j)})}{\sum_{l\in R_j}h(t_{(j)}\mid\bm x_l)}=\frac{h_0(t_{(j)})e^{\bm\beta\transp\bm x_{(j)}}}{\sum_{l\in R_j}h_0(t_{(j)})e^{\bm\beta\transp\bm x_l}}
=\frac{e^{\bm\beta\transp\bm x_{(j)}}}{\sum_{l\in R_j}e^{\bm\beta\transp\bm x_l}}.
\]
$h_0$ が約分されて消える．これをすべてのイベント時点について掛け合わせたものが部分尤度である．
打ち切り例は分子には現れず，打ち切り時点まで分母のリスク集合に含まれる．
\end{proof}
\kakomon{2016}{3} では群2（$x=1$）の4つのイベント時点で分子 $e^\beta$，分母は（群1の人数）$+$（群2の人数）$e^\beta$ で，
\[
L(\beta)=\frac{e^\beta}{5+5e^\beta}\cdot\frac{e^\beta}{5+4e^\beta}\cdot\frac{e^\beta}{3+2e^\beta}\cdot\frac{e^\beta}{2+e^\beta}\cdot\frac12\cdot\frac11.
\]

\begin{meidai}[スコアと情報量]\label[meidai]{prop:10-score}
1次元の $\beta$ について，$\bar x_j(\beta)=\sum_{l\in R_j}x_le^{\beta x_l}/\sum_{l\in R_j}e^{\beta x_l}$（リスク集合での重み付き平均）とすると
\[
U(\beta)=\frac{\partial\log L}{\partial\beta}=\sum_j\{x_{(j)}-\bar x_j(\beta)\},\qquad
I(\beta)=-\frac{\partial^2\log L}{\partial\beta^2}=\sum_j\Bigl\{\frac{\sum_{R_j}x_l^2e^{\beta x_l}}{\sum_{R_j}e^{\beta x_l}}-\bar x_j(\beta)^2\Bigr\}.
\]
2 値の $x$ で $\beta=0$ とすると $U(0)=\sum_j(d_{1j}-n_{1j}/n_j)$，$I(0)=\sum_jn_{1j}n_{0j}/n_j^2$ で，
スコア検定 $U(0)^2/I(0)$ はログランク検定（\cref{thm:09-logrank}，同順位なし）に一致する．
\end{meidai}
部分尤度は通常の尤度と同様に扱える：$\hat\beta$ は漸近正規で $\SE(\hat\beta)=I(\hat\beta)^{-1/2}$，
Wald 検定 $\hat\beta/\SE$，部分尤度比検定 $2\{\log L(\hat\beta)-\log L(0)\}$ はいずれも $\chi^2_1$ に近似できる．
\begin{supbox}[補足：同順位とベースラインハザード]
同時点に複数のイベントがあるときは，Breslow 近似 $\prod_j\exp(\bm\beta\transp\bm s_j)/\{\sum_{R_j}e^{\bm\beta\transp\bm x_l}\}^{d_j}$（$\bm s_j$ はイベント個体の共変量の和）や Efron 近似を用いる．
ベースラインの累積ハザードは Breslow 推定量 $\hat H_0(t)=\sum_{t_{(j)}\le t}d_j/\sum_{l\in R_j}e^{\hat{\bm\beta}\transp\bm x_l}$ で推定する．
\end{supbox}

\subsection{推定結果の解釈}\label{subsec:10-interp}
\kakomon{2022}{1} の数値では，$\hat\beta=0.06$，$\SE=0.18$ から
\[
\widehat{\HR}=e^{0.06}=1.06,\qquad 95\%\text{CI}=(e^{0.06-0.353},\,e^{0.06+0.353})=(0.75,\,1.51).
\]
HR $<1$ なら試験群のハザードが低い（有効），$>1$ なら高い．点推定値はプラセボ群がわずかに良い方向だが，
区間が1を含むので，死亡までの時間に差があるとは言えない．

\section{パラメトリックモデル：指数分布}\label{sec:10-exp}
\Hissu\Hinshutsu
\subsection{打ち切りのある尤度}
\begin{teiri}[打ち切りのある尤度（\kakomon{2018}{1}〔2〕）]\label[teiri]{thm:10-censlik}
独立な右側打ち切りの下で，観測 $(t_i,\delta_i)$ の尤度は
\begin{equation}
L=\prod_{i=1}^nf(t_i)^{\delta_i}S(t_i)^{1-\delta_i}=\prod_{i=1}^nh(t_i)^{\delta_i}S(t_i).\label{eq:10-censlik}
\end{equation}
\end{teiri}
イベント例は「その時点で起きた」（密度），打ち切り例は「その時点より後に起きる」（生存関数）を寄与する．
打ち切りの分布はパラメータを含まないとして尤度から除いている（独立で無情報な打ち切り）．

\begin{meidai}[指数モデルの MLE と情報量]\label[meidai]{prop:10-expmle}
$h(t)=\lambda$ のとき，$d=\sum\delta_i$（イベント数），$W=\sum t_i$（総観察時間，人時間）として
\begin{equation}
\log L=d\log\lambda-\lambda W,\qquad \hat\lambda=\frac dW,\qquad I(\lambda)=-\frac{\partial^2\log L}{\partial\lambda^2}=\frac{d}{\lambda^2}.\label{eq:10-expmle}
\end{equation}
\end{meidai}
\begin{proof}
$f=\lambda e^{-\lambda t}$，$S=e^{-\lambda t}$ を\cref{eq:10-censlik}に代入すると $L=\prod\lambda^{\delta_i}e^{-\lambda t_i}=\lambda^de^{-\lambda W}$．
$\partial\log L/\partial\lambda=d/\lambda-W=0$ から $\hat\lambda$．2階微分は $-d/\lambda^2$．
\end{proof}
$\hat\lambda$ は「イベント数/人時間」の発生率そのものである．打ち切りがなければ $d=n$，$\hat\lambda=1/\bar T$，$I_n(\lambda)=n/\lambda^2$（\kakomon{2023}{2}）．
\Youchui 打ち切りがあるとき，期待情報量は $\E[d]/\lambda^2$ で，$\E[d]$ は打ち切り分布に依存する．実用上は観測情報 $d/\hat\lambda^2$ を用いる．
\kakomon{2019}{1} のデータ（イベント3，総観察時間83週）では $\hat\lambda=3/83=0.036$/週．

\subsection{デルタ法とハザード比}\label{subsec:10-expdelta}
$\V[\hat\lambda]\approx\lambda^2/d$ なので，デルタ法で $\V[\log\hat\lambda]\approx\lambda^{-2}\cdot\lambda^2/d=1/d$．
独立な2群では
\begin{equation}
\V[\log\hat\lambda_1-\log\hat\lambda_2]\approx\frac{1}{d_1}+\frac{1}{d_2},\qquad
Z=\frac{\log\hat\lambda_1-\log\hat\lambda_2}{\sqrt{1/d_1+1/d_2}}.\label{eq:10-expz}
\end{equation}
$\log$ 変換後の分散が\textbf{イベント数だけ}で決まることが重要である．生存時間試験の精度は症例数ではなくイベント数で決まる．
指数モデルの HR $=\lambda_1/\lambda_2$ の95\%CI は $\exp\{\log(\hat\lambda_1/\hat\lambda_2)\pm1.96\sqrt{1/d_1+1/d_2}\}$．

\paragraph{必要イベント数（\kakomon{2018}{1}〔5〕）}
各群のイベント数を $r$ とすると，対立仮説の下で $Z$ はおよそ $\Normal\bigl(\frac{\log\lambda_1-\log\lambda_2}{\sqrt{2/r}},1\bigr)$ で，
両側 $\alpha$・検出力 $1-\beta$ の条件 $\frac{|\log\lambda_1-\log\lambda_2|}{\sqrt{2/r}}=z_{\alpha/2}+z_\beta$ から
\begin{equation}
r=\frac{2(z_{\alpha/2}+z_\beta)^2}{(\log\lambda_1-\log\lambda_2)^2}\quad\text{（各群）}.\label{eq:10-events}
\end{equation}
導出の一般論は\cref{ch:12}．

\begin{supbox}[補足：Weibull モデルと加速モデル]
Weibull モデル $h(t)=\alpha\gamma(\gamma t)^{\alpha-1}$ の打ち切りつき対数尤度は
$\sum_i\{\delta_i\log(\alpha\gamma^\alpha t_i^{\alpha-1})-(\gamma t_i)^\alpha\}$ で，数値的に最大化する．
Weibull（指数を含む）は比例ハザードモデルであると同時に，$\log T=\mu+\bm\eta\transp\bm x+\sigma\epsilon$ の形の加速故障時間（AFT）モデルでもある唯一の分布族である．
\end{supbox}

\section{RMST（境界内平均生存時間）}\label{sec:10-rmst}
\Hinshutsu
\subsection{定義と性質}
KM 曲線が0に達しない（最長観察が打ち切り）と，平均生存時間 $\int_0^\infty S$ は推定できない．そこで境界時間 $\tau$ を定めて
\begin{teigi}[RMST]\label{def:10-rmst}
$X(\tau)=\min(T,\tau)$ として $\mu(\tau)=\E[X(\tau)]$ を境界内平均生存時間（restricted mean survival time）という．
\end{teigi}
\begin{teiri}[\kakomon{2019}{1}〔3〕〔4〕]\label[teiri]{thm:10-rmst}
\begin{equation}
\E[X(\tau)]=\int_0^\tau S(t)\,\dd t,\qquad \E[X(\tau)^2]=2\int_0^\tau tS(t)\,\dd t.\label{eq:10-rmst}
\end{equation}
\end{teiri}
\begin{proof}
$\E[X(\tau)]=\int_0^\tau tf(t)\dd t+\tau S(\tau)$．部分積分 $\int_0^\tau tf\,\dd t=[-tS(t)]_0^\tau+\int_0^\tau S\,\dd t=-\tau S(\tau)+\int_0^\tau S$ を代入すると第1式．
同様に $\int_0^\tau t^2f\,\dd t=-\tau^2S(\tau)+2\int_0^\tau tS$，$\E[X^2]=\int_0^\tau t^2f+\tau^2S(\tau)$ から第2式．
（別証：$X(\tau)=\int_0^\tau\ind(T>t)\dd t$ の期待値をとり，積分の順序を交換する．）
\end{proof}
\paragraph{指数分布の場合}
$S(t)=e^{-\lambda t}$ なら
\begin{equation}
\mu(\tau)=\frac{1-e^{-\lambda\tau}}{\lambda},\qquad
\V[X(\tau)]=\frac{1-2\lambda\tau e^{-\lambda\tau}-e^{-2\lambda\tau}}{\lambda^2}.\label{eq:10-rmstexp}
\end{equation}
（$2\int_0^\tau te^{-\lambda t}\dd t=\frac{2}{\lambda^2}\{1-(1+\lambda\tau)e^{-\lambda\tau}\}$ から $\mu(\tau)^2$ を引く．）
\kakomon{2019}{1} では $\hat\lambda=3/83$，$\tau=20$ で $\hat\mu(20)\approx14.2$週，$\widehat\V\approx48$．

\paragraph{KM による推定}
$\hat\mu(\tau)=\int_0^\tau\hat S(t)\dd t$ は階段関数の面積で，$0=t_0<t_1<\dots<t_D\le\tau$ を $\tau$ 以前のイベント時点，$t_{D+1}=\tau$ として
$\hat\mu(\tau)=\sum_{k=0}^{D}(t_{k+1}-t_k)\hat S(t_k)$．\kakomon{2019}{1} では $10\times1+3\times0.8+6\times0.6+1\times0.3=16.3$ 週．
指数モデルによる値（14.2）との違いは，このデータで指数分布の仮定が合っていない可能性を示唆する（ただし5例と小標本なので断定はできない）．

\subsection{効果指標としての RMST}
2群の RMST の差 $\mu_1(\tau)-\mu_0(\tau)$ は「$\tau$ までの平均的なイベントのない時間の延長」で，
比例ハザード性を仮定せず，時間の単位で臨床的に解釈しやすい．
HR が非比例ハザードで解釈しにくいときの代替指標として推奨されている．$\tau$ は十分な人数がリスク集合に残る範囲で事前に決める．
\begin{supbox}[補足：KM による RMST の分散]
$A_j=\int_{t_{(j)}}^\tau\hat S(t)\dd t$ として $\widehat\V[\hat\mu(\tau)]=\sum_{t_{(j)}\le\tau}A_j^2\frac{d_j}{n_j(n_j-d_j)}$（Greenwood 型）．
\end{supbox}

\section{仮定}\label{sec:10-assump}
\begin{itemize}
\item Cox：比例ハザード（HR が時間によらず一定），共変量の効果が $\exp(\bm\beta\transp\bm x)$ の形（対数線形），独立な打ち切り．
\item 指数モデル：ハザード一定．Weibull：$\log H$ が $\log t$ の1次式．
\item RMST：独立な打ち切り（比例ハザードは不要）．$\tau$ まで十分な追跡．
\end{itemize}

\section{結果の解釈}\label{sec:10-interp}
\begin{itemize}
\item 「HR $=0.70$（95\%CI 0.55--0.89）」：追跡期間中どの時点でも，生存している人の死亡の瞬間的な発生率が試験群で30\%低い．
      「死亡割合が30\%低い」ではない（\cref{subsec:02-relation}）．
\item 共変量で調整した HR は「他の共変量が同じ人どうしの比較」．Cox の HR も OR と同様に非可縮であり，RCT でも調整で値が変わりうる．
\item RMST の差「$\tau=5$年で0.4年」：5年間で平均して約5か月長くイベントのない期間を過ごせる．
\end{itemize}

\section{手法選択}\label{sec:10-choice}
\begin{itemize}
\item 2群比較の検定 → ログランク．効果の大きさ・共変量調整 → Cox 回帰．
\item ハザード一定がもっともらしく小標本・外挿が必要 → 指数（または Weibull）モデル．
\item 比例ハザードが疑わしい → 層別 Cox，RMST の差，時点ごとの生存率の差．
\item 競合リスクがある → 原因別ハザードの Cox か Fine--Gray モデル（\cref{ch:11}）．
\end{itemize}

\section{よくある誤り}\label{sec:10-err}
\begin{warnbox}
\begin{itemize}
\item 部分尤度の分母を，その時点のリスク集合（両群を含め，直前までイベントも打ち切りも起きていない全員）ではなく同じ群の人だけとする．打ち切り例を分母から早く外しすぎる．
\item 打ち切り例の尤度の寄与を $f(t_i)$ とする（正しくは $S(t_i)$）．
\item 指数モデルの分散 $\V[\log\hat\lambda]$ を $1/n$ とする（正しくは $1/d$，イベント数）．
\item 2重対数プロットで「直線であること」を比例ハザードの条件とする（条件は群間で平行であること）．
\item RMST の $\tau$ をデータを見てから選ぶ．
\end{itemize}
\end{warnbox}

\section{数理統計との接続}\label{sec:10-link}
\begin{linkbox}
\begin{itemize}
\item 指数分布の MLE・Fisher情報（打ち切りなし）は『現代数理統計学の基礎』第6・7章演習（指数分布の片側検定と検出力，2群の指数分布の比較）と同じ計算で，
      打ち切りがあると $n$ が $d$ に，$\sum T_i$ が人時間 $W$ に置き換わる．
      \kakomonSM{2016}{2} は指数分布の MLE・不変性・$\sum X_i\sim\Ga(n,\lambda)$（形状・率）・生存関数の UMVU を問うた．
\item 指数モデルの尤度 $\lambda^de^{-\lambda W}$ は，平均 $\lambda W$ のポアソン分布で $d$ を観測した尤度と $\lambda$ について同じ形である（\cref{ch:17}）．
\item デルタ法（同書 定理5.20）で $\V[\log\hat\lambda]=1/d$．$\log$ 変換は分散を $\lambda$ に依存しなくする分散安定化変換でもある．
\item 部分尤度は条件付き確率の積であり，局外母数 $h_0$ を消去する点で MH 法の条件付けと同じ発想である．
\item $E[\min(T,\tau)]=\int_0^\tau S$ は\kakomonSM{2024}{4}（$E[X]=\int(1-F)$）の打ち切り版である．
\end{itemize}
\end{linkbox}

\section{例題}\label{sec:10-ex}
\begin{reidai}[部分尤度とスコア検定]\label{reidai:10-1}
5例のデータ（時間，イベント有無，群 $x$）：$(2,\text{イベント},1)$，$(3,\text{イベント},0)$，$(4,\text{打ち切り},1)$，$(5,\text{イベント},0)$，$(6,\text{イベント},1)$．
\begin{enumerate}[label=(\arabic*)]
\item Cox モデル $h(t\mid x)=h_0(t)e^{\beta x}$ の部分尤度 $L(\beta)$ を書け．
\item $\beta=0$ でのスコア $U(0)$ と情報量 $I(0)$ を求め，スコア検定統計量を計算せよ．
\item (2)がログランク検定統計量に一致することを，時点ごとの $O-E$ と分散から確かめよ．
\end{enumerate}
\end{reidai}

\begin{reidai}[指数モデルによる2群比較]\label{reidai:10-2}
群Aは総観察時間240人月でイベント12件，群Bは200人月で20件であった．指数分布を仮定する．
\begin{enumerate}[label=(\arabic*)]
\item 各群のハザードの MLE と中央生存時間を求めよ（$\log2=0.6931$）．
\item ハザード比 $\lambda_A/\lambda_B$ の推定値と，\cref{eq:10-expz}の $Z$ 検定（両側5\%）を行え．
\item HR の95\%信頼区間を求めよ（$e^{-1.4088}=0.244$，$e^{0.0225}=1.023$）．
\end{enumerate}
\end{reidai}

\begin{reidai}[RMST]\label{reidai:10-3}
\cref{reidai:09-2}のデータ（$3,5^+,6,6,8^+,10,12^+,15$）について，$\tau=12$ とする．
\begin{enumerate}[label=(\arabic*)]
\item KM に基づく RMST を求めよ．
\item 指数モデルの MLE $\hat\lambda$ と，それに基づく RMST を求めよ（$e^{-0.9231}=0.3973$）．
\item 2つの値の違いについて考察せよ．
\end{enumerate}
\end{reidai}

\begin{reidai}[2重対数プロット]\label{reidai:10-4}
対照群の生存関数が $S_0(t)=\exp\{-(t/10)^{1.5}\}$ で，試験群のハザード比が $0.6$ の比例ハザードモデルに従うとする．
\begin{enumerate}[label=(\arabic*)]
\item 試験群の生存関数を書け．
\item 両群の $\log\{-\log S(t)\}$ を $\log t$ の関数として書き，プロットが平行な直線になることを示せ．
\item $t=10$ での両群の生存確率を求めよ（$e^{-1}=0.368$，$e^{-0.6}=0.549$）．
\end{enumerate}
\end{reidai}

\section{解答}\label{sec:10-sol}
\begin{kaitou}{reidai:10-1}
(1) リスク集合：$t=2$：全5例（$x=1,0,1,0,1$），$t=3$：4例（$0,1,0,1$），$t=5$：2例（$0,1$），$t=6$：1例（$1$）．
\[
L(\beta)=\frac{e^\beta}{3e^\beta+2}\cdot\frac{1}{2e^\beta+2}\cdot\frac{1}{e^\beta+1}\cdot\frac{e^\beta}{e^\beta}
=\frac{e^\beta}{(3e^\beta+2)(2e^\beta+2)(e^\beta+1)}.
\]
(2) $\beta=0$ で $\bar x_j=3/5,\ 2/4,\ 1/2,\ 1$．$U(0)=(1-0.6)+(0-0.5)+(0-0.5)+(1-1)=-0.6$．
$I(0)=\sum\bar x_j(1-\bar x_j)=0.24+0.25+0.25+0=0.74$．スコア統計量 $0.36/0.74=0.49$．\\
(3) 群1（$x=1$）の観測イベント2，期待 $0.6+0.5+0.5+1=2.6$，$O-E=-0.6$．分散 $\sum n_{1j}n_{0j}/n_j^2=0.74$．一致する．
\end{kaitou}

\begin{kaitou}{reidai:10-2}
(1) $\hat\lambda_A=12/240=0.05$/月，$\hat\lambda_B=20/200=0.10$/月．中央値 $\log2/\hat\lambda$：A 13.9か月，B 6.9か月．\\
(2) $\widehat{\HR}=0.5$，$\log0.5=-0.693$，$\SE=\sqrt{1/12+1/20}=0.365$，$Z=-1.90$，$|Z|<1.96$ で有意でない．\\
(3) $-0.693\pm1.96\times0.365=(-1.409,\,0.023)$ → $(0.24,\,1.02)$．区間が1をわずかに含む．
\end{kaitou}

\begin{kaitou}{reidai:10-3}
(1) $\hat S$ は $[0,3)$ で1，$[3,6)$ で0.875，$[6,10)$ で0.583，$[10,12]$ で0.389．
$\hat\mu(12)=3\times1+3\times0.875+4\times0.583+2\times0.389=3+2.625+2.333+0.778=8.74$ 週．\\
(2) $d=5$，$W=65$ で $\hat\lambda=0.0769$．$\hat\mu(12)=(1-e^{-0.9231})/0.0769=0.6027/0.0769=7.84$ 週．\\
(3) KM による値の方が大きい．KM では最初の3週にイベントがなく，その後もゆるやかに減るのに対し，指数モデルは初期からハザード一定を仮定するため，
早期の生存を過小評価している可能性がある．どちらが妥当かはハザードが一定かどうか（2重対数プロットなど）で判断する．
\end{kaitou}

\begin{kaitou}{reidai:10-4}
(1) $S_1(t)=S_0(t)^{0.6}=\exp\{-0.6(t/10)^{1.5}\}$．\\
(2) $\log\{-\log S_0\}=1.5\log t-1.5\log10$，$\log\{-\log S_1\}=\log0.6+1.5\log t-1.5\log10$．
どちらも $\log t$ の傾き1.5の直線で，縦方向の差は $\log0.6=\beta$ で一定（平行）．\\
(3) $S_0(10)=e^{-1}=0.368$，$S_1(10)=e^{-0.6}=0.549$．
\end{kaitou}

\begin{summarybox}
\begin{itemize}
\item Cox：$h(t\mid\bm x)=h_0(t)e^{\bm\beta\transp\bm x}$，HR $=e^\beta$ は時間によらず一定．$S=S_0^{\exp(\bm\beta\transp\bm x)}$，$\log(-\log S)=\log H_0+\bm\beta\transp\bm x$（平行性で確認）．
\item 部分尤度 $\prod_j e^{\bm\beta\transp\bm x_{(j)}}/\sum_{R_j}e^{\bm\beta\transp\bm x_l}$．$\beta=0$ のスコア検定＝ログランク．
\item 打ち切り尤度 $\prod f^{\delta}S^{1-\delta}$．指数：$\hat\lambda=d/W$，$I=d/\lambda^2$，$\V[\log\hat\lambda]\approx1/d$．
\item 2群：$\V[\log\hat\lambda_1-\log\hat\lambda_2]\approx1/d_1+1/d_2$，必要イベント数 $r=2(z_{\alpha/2}+z_\beta)^2/(\log\HR)^2$（各群）．
\item RMST $=\int_0^\tau S$，$\E[X(\tau)^2]=2\int_0^\tau tS$．指数：$(1-e^{-\lambda\tau})/\lambda$．KM：階段の面積．
\end{itemize}
\end{summarybox}

% ===== ch11.tex =====
\chapter{競合リスク}\label{ch:11}

\begin{goalbox}
\begin{itemize}
\item 競合リスク（competing risks）の状況を識別し，競合イベントと打ち切りの違いを説明できる．
\item 原因別ハザード $\lambda_j(t)$，部分密度 $f_j(t)$，累積発生関数（CIF）$F_j(t)$ を定義し，$\lambda=\sum_j\lambda_j$，$f_j=\lambda_jS$，$F_j(t)=\int_0^t\lambda_jS$ を導出できる．
\item $F_j(\infty)=\Prob(J=j)<1$ であること，$1-F_j(t)>S_j(t)$ を示せる．
\item Aalen--Johansen 型の推定量で CIF を手計算できる．
\item 競合イベントを打ち切りとみなした $1-\mathrm{KM}$ が CIF を過大推定することを示し，理由を説明できる．
\item 原因別ハザードの Cox モデルと Fine--Gray モデル（部分分布ハザード）の解釈の違いを説明できる．
\end{itemize}
\end{goalbox}

\begin{exambox}
\begin{itemize}
\item \kakomon{2021}{1}〔3〕〜〔5〕：心血管死（イベント1）と他因死（イベント2）．$F_j(t)=\int_0^t\lambda_j(u)S_2(u)\dd u$ の導出（試験の $S_2$ は本章の全生存関数 $S$），
      $\hat F_j(t)=\sum_{t_{ji}\le t}\frac{d_{ji}}{n_{ji}}\hat S_2(t_{ji}^-)$ による推定表（$\hat F_1(18)=0.75$，$\hat F_2=0.25$），
      $1-\hat S_1(t)$（他因死を打ち切り扱い）との比較と過大推定の理由．
\item \kakomon{2025}{3}：$m$ 種類の競合する死因．$\lambda_j=f_j/S$，$1-F_j(t)>S_j(t)$，$\lambda=\sum\lambda_j$ の証明，
      Weibull 型原因別ハザードの下での $S(t)=\exp\{-(\gamma_1t)^{\alpha_1}-(\gamma_2t)^{\alpha_2}\}$，
      原因別ハザードの比例ハザードモデルの $\beta$ で $F_1(t)$ への効果を評価することの適否．
\end{itemize}
\end{exambox}

\flowline{時間＋\\イベントの種類}{CIF $F_j(t)$\\原因別ハザード}{Aalen--Johansen\\Fine--Gray}{独立打ち切り\\（競合は打ち切りでない）}{$\sum\frac{d_j}{n}\hat S(t^-)$}{$1-$KM の過大推定}

\section{問題設定}\label{sec:11-setting}
がん死を主要評価項目とする研究で，他因死が起こると以後がん死は観察できない．
このように，あるイベントの発生によって別のイベントが起こりえなくなる（あるいは意味が変わる）とき，それらを\emphx{競合リスク}という．
観測されるのは最初のイベントまでの時間 $T$ と，その種類 $J\in\{1,\dots,m\}$，および打ち切り時間 $C$ である．
\[
Y=\min(T,C),\qquad J^*=\begin{cases}J&(T\le C)\\0&(T>C)\ \text{（打ち切り）}\end{cases}
\]
\Youchui 打ち切りは「イベントはいずれ起こるが観察できなかった」，競合イベントは「そのイベントは（少なくとも最初のイベントとしては）二度と起こらない」である．
両者を同じに扱うことが本章で扱う誤りの根源である．

\section{基本概念}\label{sec:11-basic}
\Hissu
\begin{teigi}[競合リスクの関数]\label{def:11-functions}
\begin{align*}
&\text{全生存関数} & S(t)&=\Prob(T>t),\\
&\text{原因別ハザード} & \lambda_j(t)&=\lim_{\Delta\to0}\frac{\Prob(t\le T<t+\Delta,\ J=j\mid T\ge t)}{\Delta},\\
&\text{部分密度} & f_j(t)&=\lim_{\Delta\to0}\frac{\Prob(t\le T<t+\Delta,\ J=j)}{\Delta},\\
&\text{累積発生関数（CIF）} & F_j(t)&=\Prob(T\le t,\ J=j)=\int_0^tf_j(u)\,\dd u.
\end{align*}
\end{teigi}
原因別ハザードは「時点 $t$ までにどのイベントも起こしていない人が，次の瞬間に原因 $j$ のイベントを起こす率」である．

\begin{teiri}[基本関係式]\label[teiri]{thm:11-relations}
\begin{align}
\lambda_j(t)&=\frac{f_j(t)}{S(t)},\qquad f_j(t)=\lambda_j(t)S(t),\label{eq:11-fj}\\
\lambda(t)&=\sum_{j=1}^m\lambda_j(t),\qquad S(t)=\exp\Bigl\{-\sum_{j=1}^m\Lambda_j(t)\Bigr\},\ \ \Lambda_j(t)=\int_0^t\lambda_j,\label{eq:11-sum}\\
F_j(t)&=\int_0^t\lambda_j(u)S(u)\,\dd u,\qquad \sum_{j=1}^mF_j(t)=1-S(t).\label{eq:11-cif}
\end{align}
\end{teiri}
\begin{proof}
\cref{eq:11-fj}：$\{t\le T<t+\Delta,J=j\}\subset\{T\ge t\}$ なので条件付き確率の定義から
$\lambda_j(t)=\lim\frac{\Prob(t\le T<t+\Delta,J=j)}{\Delta\,\Prob(T\ge t)}=f_j(t)/S(t)$（$T$ は連続で $\Prob(T\ge t)=S(t)$）．\\
\cref{eq:11-sum}：イベントは同時に起こらないので，$\{t\le T<t+\Delta\}$ は互いに排反な $\{t\le T<t+\Delta,J=j\}$ $(j=1,\dots,m)$ の和．
両辺を $\Delta\,\Prob(T\ge t)$ で割って極限をとれば $\lambda=\sum\lambda_j$．$S=\exp(-\int_0^t\lambda)$（\cref{thm:09-relations}）に代入する．\\
\cref{eq:11-cif}：$F_j=\int f_j$ に \cref{eq:11-fj} を代入．$\sum_jF_j(t)=\Prob(T\le t)=1-S(t)$．
\end{proof}
CIF は原因 $j$ のハザードだけでなく，$S(u)$ を通じて\textbf{他のすべての原因のハザードに依存する}．
他の原因のハザードが大きいと，原因 $j$ のイベントを起こす前に他のイベントが起きやすいので $F_j$ は小さくなる．

\begin{meidai}[\kakomon{2025}{3}〔2〕]\label[meidai]{prop:11-Fj}
$m\ge2$ で各原因の確率が正なら，$F_j(\infty)=\Prob(J=j)<1$ であり，$S_j(t)=\Prob(T>t,J=j)$ について $1-F_j(t)>S_j(t)$．
\end{meidai}
\begin{proof}
$\{J=j\}=\{T\le t,J=j\}\cup\{T>t,J=j\}$（排反）より $\Prob(J=j)=F_j(t)+S_j(t)$．
$\Prob(J=j)=1-\sum_{k\ne j}\Prob(J=k)<1$ なので $S_j(t)=\Prob(J=j)-F_j(t)<1-F_j(t)$．
\end{proof}
したがって CIF は「部分分布関数」であり，$t\to\infty$ で1ではなく $\Prob(J=j)$ に収束する．
通常の分布関数のように $1-F_j$ を「生存関数」と解釈することはできない．

\paragraph{例：原因別ハザードが一定の場合}
$\lambda_j(t)=\lambda_j$ なら $S(t)=e^{-(\lambda_1+\lambda_2)t}$，
\[
F_1(t)=\int_0^t\lambda_1e^{-(\lambda_1+\lambda_2)u}\dd u=\frac{\lambda_1}{\lambda_1+\lambda_2}\bigl\{1-e^{-(\lambda_1+\lambda_2)t}\bigr\},\qquad F_1(\infty)=\frac{\lambda_1}{\lambda_1+\lambda_2}.
\]
（独立な $\Ex(\lambda_1)$，$\Ex(\lambda_2)$ の最小値とどちらが先かの同時分布と同じ．）
\paragraph{例：Weibull 型（\kakomon{2025}{3}〔4〕）}
$\lambda_k(t)=\alpha_k\gamma_k(\gamma_kt)^{\alpha_k-1}$ なら $\Lambda_k(t)=(\gamma_kt)^{\alpha_k}$ で $S(t)=\exp\{-(\gamma_1t)^{\alpha_1}-(\gamma_2t)^{\alpha_2}\}$．

\section{推定}\label{sec:11-est}
\subsection{Aalen--Johansen 型推定量}\label{subsec:11-aj}
\Hissu\cref{eq:11-cif}の $\lambda_j(u)\dd u$ を時点ごとの原因別ハザードの推定値 $d_{ji}/n_i$ に，$S(u)$ をその直前の全イベントの KM 推定値に置き換えると
\begin{equation}
\hat F_j(t)=\sum_{i:\,t_i\le t}\frac{d_{ji}}{n_i}\hat S(t_i^-),\label{eq:11-aj}
\end{equation}
ここで $t_i$ は（どれかの原因の）イベント時点，$n_i$ は $t_i$ の直前にどのイベントも起こさず観察中の人数，$d_{ji}$ は $t_i$ での原因 $j$ のイベント数，
$\hat S$ は\textbf{すべての原因をイベントとした} KM 推定量である（\kakomon{2021}{1}〔4〕）．
各時点で「そこまで無事だった確率」×「そこで原因 $j$ が起こる条件付き確率」を足し上げている．
$\sum_j\hat F_j(t)=1-\hat S(t)$ が成り立つ．

\subsection{\texorpdfstring{$1-\mathrm{KM}$}{1−KM} による過大推定}\label{subsec:11-bias}
\Hinshutsu 競合イベントを打ち切りとみなして原因 $j$ だけの KM $\hat S_j^{\mathrm{KM}}$ を計算すると，
各因子は $1-d_{ji}/n_i$ で，$1-\hat S_j^{\mathrm{KM}}(t)$ は $1-\exp\{-\Lambda_j(t)\}$ を推定する．
\begin{meidai}\label[meidai]{prop:11-overestimate}
$F_j(t)\le1-e^{-\Lambda_j(t)}$．等号は $[0,t]$ 上で（$\lambda_j>0$ となる範囲で）他の原因の累積ハザードが0のときに限る．
\end{meidai}
\begin{proof}
$S(u)=e^{-\Lambda_j(u)}e^{-\sum_{k\ne j}\Lambda_k(u)}\le e^{-\Lambda_j(u)}$ より
$F_j(t)=\int_0^t\lambda_j(u)S(u)\dd u\le\int_0^t\lambda_j(u)e^{-\Lambda_j(u)}\dd u=1-e^{-\Lambda_j(t)}$．
\end{proof}
$1-\mathrm{KM}$ は「競合イベントが起こらない仮想的な世界での原因 $j$ の累積リスク」に相当し，
しかもそれは潜在時間の独立性という検証不能な仮定の下でしか意味を持たない．
競合イベントを起こした人は二度と原因 $j$ のイベントを起こさないのに，打ち切り扱いは
「その人も，リスク集合に残った人と同じ率で将来イベントを起こす」とみなすので，原因 $j$ のイベントを起こさない確率 $1-F_j$ を過小に，累積発生を過大に推定する（\kakomon{2021}{1}〔5〕）．

\begin{figure}[htbp]
\centering
\begin{tikzpicture}
\begin{axis}[width=90mm,height=58mm,xmin=0,xmax=13,ymin=0,ymax=0.65,
  xlabel={時間},ylabel={累積発生},grid=major,grid style={gray!20},tick label style={font=\scriptsize},
  legend style={font=\scriptsize,at={(0.02,0.98)},anchor=north west}]
\addplot[very thick,bkblue,const plot] coordinates {(0,0)(2,0.1)(5,0.2143)(7,0.3286)(9,0.4810)(13,0.4810)};
\addplot[thick,dashed,bkred,const plot] coordinates {(0,0)(2,0.1)(5,0.2286)(7,0.3829)(9,0.5886)(13,0.5886)};
\legend{CIF $\hat F_1(t)$（Aalen--Johansen）,$1-\mathrm{KM}$（原因2を打ち切り扱い）}
\end{axis}
\end{tikzpicture}
\caption{\cref{reidai:11-2}のデータでの累積発生の推定．$1-\mathrm{KM}$ は過大になる．}\label{fig:11-cif}
\end{figure}

\section{回帰モデル}\label{sec:11-reg}
\subsection{原因別ハザードの Cox モデル}
原因 $j$ の原因別ハザードに比例ハザードを仮定する：$\lambda_j(t\mid\bm x)=\lambda_{j0}(t)\exp(\bm\beta_j\transp\bm x)$．
推定は他の原因のイベントを打ち切りとみなした通常の Cox 回帰で行える（原因別ハザードのリスク集合は「どのイベントも起きていない人」だから）．
\textbf{ただし $\bm\beta_j$ は原因別ハザードへの効果であり，CIF への効果ではない}（\kakomon{2025}{3}〔5〕）．
\cref{eq:11-cif}により $F_1(t\mid\bm x)=\int_0^t\lambda_1(u\mid\bm x)\exp\{-\Lambda_1(u\mid\bm x)-\Lambda_2(u\mid\bm x)\}\dd u$ は $\lambda_2$ にも依存し，
$\bm x$ が $\lambda_2$ を大きくするなら，$\bm\beta_1$ の値にかかわらず $F_1$ は小さくなりうる．
原因別ハザードモデルは病因（etiology）の研究，すなわち「まだ生きている人でその原因の発生率が変わるか」に向く．

\subsection{Fine--Gray モデル（部分分布ハザード）}
\Hatten
\begin{teigi}[部分分布ハザード]\label{def:11-subdist}
$\tilde\lambda_j(t)=\lim_{\Delta\to0}\frac{1}{\Delta}\Prob\bigl(t\le T<t+\Delta,J=j\bigm|T\ge t\ \text{または}\ (T<t,\ J\ne j)\bigr)
=\frac{f_j(t)}{1-F_j(t)}=-\frac{\dd}{\dd t}\log\{1-F_j(t)\}$．
\end{teigi}
リスク集合に「すでに他の原因のイベントを起こした人」も含めるのが特徴で，$1-F_j(t)=\exp\{-\int_0^t\tilde\lambda_j\}$ となる．
Fine--Gray モデル \cite{finegray1999} は $\tilde\lambda_1(t\mid\bm x)=\tilde\lambda_{10}(t)\exp(\bm\gamma\transp\bm x)$ とし，
\[
F_1(t\mid\bm x)=1-\{1-F_{10}(t)\}^{\exp(\bm\gamma\transp\bm x)}
\]
から，$\bm\gamma>0$ は「CIF を大きくする」と直接解釈できる．予後予測（ある期間内に原因 $j$ のイベントを経験する確率）に向く．
\begin{supbox}[補足：CIF の群間比較]
2群の CIF の比較には Gray の検定（部分分布ハザードに基づく重み付きログランク型の検定）が用いられる．
原因別ハザードのログランク検定とは帰無仮説が異なる．
\end{supbox}

\section{仮定}\label{sec:11-assump}
\begin{itemize}
\item 打ち切り（追跡不能・試験終了）は独立であること．\textbf{競合イベントは打ち切りではない}．
\item CIF・原因別ハザードは観測データから識別できる．一方，「競合イベントがない世界での原因 $j$ の分布」（潜在時間の周辺分布）は，潜在時間の独立性を仮定しない限り識別できない．
\item 回帰モデル：原因別ハザード（または部分分布ハザード）の比例性．
\end{itemize}

\section{結果の解釈}\label{sec:11-interp}
\begin{itemize}
\item 「5年累積がん死亡率 $\hat F_1(5)=0.20$」：5年以内にがんで亡くなる人の割合の推定値（他因死が起こる現実の世界で）．
\item 高齢者で他因死が多いと，がん死の原因別ハザードは同じでも CIF は低くなる．
\item 治療が原因1の原因別ハザードを下げても，生存が延びた結果として原因2のイベントが増えることがある．両方の原因について結果を示すのが望ましい．
\item estimand の観点では，競合イベントを複合エンドポイント（全死亡）に含める，while alive 戦略で CIF を推定する，などを事前に決める（\cref{sec:02-estimand}）．
\end{itemize}

\section{手法選択}\label{sec:11-choice}
\begin{table}[htbp]
\centering\small
\caption{競合リスクの手法選択}\label{tab:11-choice}
\begin{tabular}{@{}L{55mm}L{80mm}@{}}
\toprule
目的 & 手法\\
\midrule
原因 $j$ の累積発生の記述 & Aalen--Johansen 推定量（$1-\mathrm{KM}$ は不可）\\
CIF の群間比較 & Gray 検定\\
共変量が原因 $j$ の発生率に与える影響（病因） & 原因別ハザードの Cox（競合を打ち切り扱い）\\
共変量が原因 $j$ の累積発生に与える影響（予測） & Fine--Gray モデル\\
全イベントの記述 & 全イベントを1つとした KM・Cox\\
\bottomrule
\end{tabular}
\end{table}

\section{よくある誤り}\label{sec:11-err}
\begin{warnbox}
\begin{itemize}
\item 競合イベントを打ち切りとして $1-\mathrm{KM}$ で累積発生を報告する．
\item Aalen--Johansen 推定量の $\hat S(t^-)$ に，原因 $j$ だけの KM を用いる（全原因の KM を用いる）．
\item 原因別ハザードの HR を「累積発生の比」と解釈する．
\item $F_j(t)\to1$ と考える（極限は $\Prob(J=j)$）．
\end{itemize}
\end{warnbox}

\section{数理統計との接続}\label{sec:11-link}
\begin{linkbox}
\begin{itemize}
\item 独立な $\Ex(\lambda)$ と $\Ex(\mu)$ の最小値 $Z$ と，どちらが最小かの指示変数 $W$ の同時分布（『現代数理統計学の基礎』第4章演習）は，
      一定の原因別ハザードをもつ競合リスクそのものである：$Z\sim\Ex(\lambda+\mu)$，$\Prob(W=1)=\lambda/(\lambda+\mu)$，$Z\indep W$．
\item 最小値のハザードは各ハザードの和（同書第5章演習）が $\lambda=\sum\lambda_j$ に対応する．
\item \kakomonSM{2025}{2} は最小値の分布 $1-\{S(t)\}^n$ を扱った（直列系・最初のイベントの時間）．
\end{itemize}
\end{linkbox}

\section{例題}\label{sec:11-ex}
\begin{reidai}[一定の原因別ハザード]\label{reidai:11-1}
原因1・原因2の原因別ハザードがそれぞれ $0.02$，$0.03$（/年）で一定とする．$e^{-0.5}=0.6065$，$e^{-0.2}=0.8187$ を用いよ．
\begin{enumerate}[label=(\arabic*)]
\item 10年間いずれのイベントも起こさない確率を求めよ．
\item $F_1(10)$，$F_2(10)$，$F_1(\infty)$ を求めよ．
\item 原因2を打ち切りとみなした場合に $1-\mathrm{KM}$ が推定する量 $1-e^{-\Lambda_1(10)}$ を求め，$F_1(10)$ と比べよ．
\end{enumerate}
\end{reidai}

\begin{reidai}[Aalen--Johansen 推定]\label{reidai:11-2}
10例の観察時間とイベントの種類（1：原因1，2：原因2，0：打ち切り）が
$(2,1),(3,2),(4,0),(5,1),(6,2),(7,1),(8,0),(9,1),(10,0),(12,2)$ であった．
\begin{enumerate}[label=(\arabic*)]
\item 各イベント時点の $n_i$，$\hat S(t_i^-)$，$\hat F_1(t)$，$\hat F_2(t)$ を表にせよ．
\item 原因2を打ち切りとみなした $1-\mathrm{KM}$ を $t=9$ で求め，$\hat F_1(9)$ と比べよ．
\item $\hat F_1(12)+\hat F_2(12)+\hat S(12)=1$ を確かめよ．
\end{enumerate}
\end{reidai}

\begin{reidai}[モデルの解釈]\label{reidai:11-3}
高齢がん患者で，治療Aは対照に比べ，がん死の原因別ハザード比が $0.8$，他因死の原因別ハザード比が $1.5$ であった（ともに比例ハザード，原因別ハザードは一定とする）．
対照群のがん死・他因死の原因別ハザードをそれぞれ $0.10$，$0.10$ とする．
\begin{enumerate}[label=(\arabic*)]
\item 各群の $F_{\text{がん}}(\infty)$ を求めよ．
\item 「治療Aはがん死を減らす」という主張を，原因別ハザードと CIF の観点から評価せよ．
\end{enumerate}
\end{reidai}

\section{解答}\label{sec:11-sol}
\begin{kaitou}{reidai:11-1}
(1) $S(10)=e^{-0.05\times10}=e^{-0.5}=0.607$．\\
(2) $F_1(10)=\frac{0.02}{0.05}(1-0.607)=0.4\times0.393=0.157$，$F_2(10)=0.6\times0.393=0.236$，$F_1(\infty)=0.4$．\\
(3) $1-e^{-0.2}=0.181>0.157$．$1-\mathrm{KM}$ は約15\%過大になる．
\end{kaitou}

\begin{kaitou}{reidai:11-2}
(1)
\begin{center}\small
\begin{tabular}{@{}cccccccc@{}}
\toprule
$t_i$ & 種類 & $n_i$ & $\hat S(t_i^-)$ & $\hat S(t_i)$ & $\hat F_1(t_i)$ & $\hat F_2(t_i)$ & $1-\mathrm{KM}_1$\\
\midrule
2 & 1 & 10 & 1.000 & 0.900 & 0.100 & 0.000 & 0.100\\
3 & 2 & 9 & 0.900 & 0.800 & 0.100 & 0.100 & 0.100\\
5 & 1 & 7 & 0.800 & 0.686 & 0.214 & 0.100 & 0.229\\
6 & 2 & 6 & 0.686 & 0.571 & 0.214 & 0.214 & 0.229\\
7 & 1 & 5 & 0.571 & 0.457 & 0.329 & 0.214 & 0.383\\
9 & 1 & 3 & 0.457 & 0.305 & 0.481 & 0.214 & 0.589\\
12 & 2 & 1 & 0.305 & 0.000 & 0.481 & 0.519 & 0.589\\
\bottomrule
\end{tabular}
\end{center}
例：$t=5$ では $\hat F_1=0.1+\frac17\times0.8=0.214$．\\
(2) 原因1だけの KM は $(1-\frac1{10})(1-\frac17)(1-\frac15)(1-\frac13)=0.9\times0.857\times0.8\times0.667=0.411$，
$1-\mathrm{KM}=0.589>\hat F_1(9)=0.481$（\cref{fig:11-cif}）．\\
(3) $0.481+0.519+0=1$．
\end{kaitou}

\begin{kaitou}{reidai:11-3}
(1) 対照：$\frac{0.10}{0.20}=0.50$．治療A：がん $0.08$，他因 $0.15$ で $\frac{0.08}{0.23}=0.348$．\\
(2) がん死の原因別ハザードは20\%低下しており，「生存している人でのがん死の発生率を下げる」という意味では正しい．
CIF も低下しているが，その一部は他因死のハザードが上がって「がん死する前に他因で亡くなる人」が増えた結果である．
他因死の CIF は $0.15/0.23=0.652$（対照0.50）に増え，全死亡のハザードも $0.23>0.20$ と悪化している．
原因別ハザード・CIF・全死亡をあわせて評価すべきで，がん死の CIF の低下だけを根拠に治療を有益とは言えない．
\end{kaitou}

\enlargethispage{3\baselineskip}
\begin{summarybox}
\begin{itemize}
\item 競合イベントは打ち切りではない．観測は $(T,J)$．
\item $\lambda_j=f_j/S$，$\lambda=\sum\lambda_j$，$S=\exp(-\sum\Lambda_j)$，$F_j(t)=\int_0^t\lambda_jS$，$\sum F_j=1-S$．
\item $F_j(\infty)=\Prob(J=j)<1$，$1-F_j(t)>S_j(t)$．
\item 推定：$\hat F_j(t)=\sum_{t_i\le t}(d_{ji}/n_i)\hat S(t_i^-)$（$\hat S$ は全原因の KM）．
\item $1-\mathrm{KM}$（競合を打ち切り扱い）は $1-e^{-\Lambda_j}\ge F_j$ を推定し過大．
\item 原因別 Cox の $\beta$ はハザードへの効果（病因），Fine--Gray の $\gamma$ は CIF への効果（予測）．
\end{itemize}
\end{summarybox}

% ===== ch12.tex =====
\chapter{検出力・症例数・イベント数設計}\label{ch:12}

\begin{goalbox}
\begin{itemize}
\item 第1種の過誤・第2種の過誤・検出力・効果量の関係を説明できる．
\item 「帰無仮説で $\Normal(0,1)$，対立仮説で $\Normal(\mu,1)$」の枠組みから $\mu=z_\alpha+z_\beta$ を導出し，任意の検定の症例数の式を作れる．
\item 平均の差，比率の差，単群の比率，生存時間（イベント数）の症例数を計算できる．
\item 生存時間試験では必要な\textbf{イベント数}が精度を決めることを説明し，症例数に換算できる．
\item 分散を帰無仮説の下で評価するか対立仮説の下で評価するかで結果が変わることを理解する．
\end{itemize}
\end{goalbox}

\begin{exambox}
\begin{itemize}
\item \kakomon{2018}{1}〔5〕：指数モデルの2群比較で，検出力 $1-\beta$ を満たす各群のイベント数 $r=2(z_{\alpha/2}+z_\beta)^2/(\log\lambda_1-\log\lambda_2)^2$ の導出．
\item \kakomon{2023}{2}：単群試験（指数分布，打ち切りなし）．$H_0:Z\sim\Normal(0,1)$，$H_1:Z\sim\Normal(\mu,1)$ から $\mu=z_\alpha+z_\beta$，
      Wald 型統計量 $Z=(\hat\lambda-\lambda_0)/\sqrt{I_n(\hat\lambda)^{-1}}$ による $n=\lambda_1^2(z_\alpha+z_\beta)^2/(\lambda_1-\lambda_0)^2$，回復率から $\lambda$ を定めて $n=32$．
\item \kakomon{2025}{1}〔4〕：順位和検定の検出力 $\Phi(\sqrt{6n}(\delta-\delta^2/2)-z_\alpha)$（\cref{sec:03-power}）．
\item \kakomon{2017}{5}〔5〕（共通）：信頼区間幅を一定以下にする確率が0.8以上となる $n$（二項分布の正規近似）．
\item \kakomon{2017}{3}：二段階デザインの $\alpha,\beta$ を二項分布で正確に計算（\cref{ch:14}）．
\end{itemize}
\end{exambox}

\flowline{デザイン\\主要評価項目}{効果量 $\Delta$\\分散}{$Z$ 検定の\\正規近似}{$\alpha$・$1-\beta$\\分散の見積もり}{$n=(z+z)^2v/\Delta^2$}{脱落・多重性\\の上乗せ}

\section{問題設定と基本概念}\label{sec:12-basic}
\begin{center}\small
\begin{tabular}{@{}lcc@{}}
\toprule
 & $H_0$ が真 & $H_1$ が真\\
\midrule
$H_0$ を棄却 & 第1種の過誤（確率 $\alpha$） & 正しい（検出力 $1-\beta$）\\
$H_0$ を採択 & 正しい（$1-\alpha$） & 第2種の過誤（確率 $\beta$）\\
\bottomrule
\end{tabular}
\end{center}
症例数設計では，(i) 有意水準 $\alpha$（片側・両側），(ii) 検出したい効果量 $\Delta$（臨床的に意味のある最小の差，または期待される差），
(iii) 分散などの局外母数の見積もり，(iv) 目標検出力 $1-\beta$（通常0.8または0.9）を定め，必要な $n$ を求める．
\Youchui 『現代数理統計学の基礎』では $\beta(\theta)$ を検出力関数そのものとして用いる．本書（と医療統計の慣習）では $\beta$ は第2種の過誤確率である．

\section{数理：\texorpdfstring{$\mu=z_\alpha+z_\beta$}{μ = zα + zβ} の導出}\label{sec:12-core}
\Hissu
\begin{teiri}[基本公式]\label[teiri]{thm:12-core}
検定統計量 $Z$ が $H_0$ の下で $\Normal(0,1)$，$H_1$ の下で $\Normal(\mu,1)$（$\mu>0$）に従い，$Z>z_\alpha$ で棄却するとき，
検出力は $\Phi(\mu-z_\alpha)$ であり，検出力 $1-\beta$ を達成する条件は
\begin{equation}
\mu=z_\alpha+z_\beta.\label{eq:12-core}
\end{equation}
\end{teiri}
\begin{proof}
$1-\beta=\Prob(Z>z_\alpha\mid H_1)=\Prob(Z-\mu>z_\alpha-\mu)=1-\Phi(z_\alpha-\mu)=\Phi(\mu-z_\alpha)$．
一方 $\Phi(z_\beta)=1-\beta$ なので，$\Phi$ の単調性から $\mu-z_\alpha=z_\beta$．
\end{proof}
\begin{figure}[htbp]
\centering
\begin{tikzpicture}
\begin{axis}[width=110mm,height=45mm,xmin=-3.5,xmax=6.5,ymin=0,ymax=0.45,axis lines=middle,
  xtick={0,1.96,2.8},xticklabels={$0$,$z_\alpha$,$\mu$},ytick=\empty,tick label style={font=\scriptsize},
  samples=120,domain=-3.5:6.5]
\addplot[name path=h0,thick,bkblue] {exp(-x^2/2)/sqrt(2*pi)};
\addplot[name path=h1,thick,bkred] {exp(-(x-2.8)^2/2)/sqrt(2*pi)};
\addplot[name path=axis,draw=none] {0};
\addplot[fill=bkblue!30,draw=none] fill between[of=h0 and axis,soft clip={domain=1.96:6.5}];
\addplot[fill=bkred!20,draw=none] fill between[of=h1 and axis,soft clip={domain=-3.5:1.96}];
\draw[dashed] (axis cs:1.96,0)--(axis cs:1.96,0.42);
\node[font=\scriptsize,bkblue] at (axis cs:-1.4,0.36) {$H_0:\Normal(0,1)$};
\node[font=\scriptsize,bkred] at (axis cs:4.6,0.36) {$H_1:\Normal(\mu,1)$};
\node[font=\scriptsize] at (axis cs:1.35,0.05) {$\beta$};
\node[font=\scriptsize] at (axis cs:2.45,0.03) {$\alpha$};
\draw[<->] (axis cs:0,0.41)--(axis cs:1.96,0.41) node[midway,above,font=\scriptsize]{$z_\alpha$};
\draw[<->] (axis cs:1.96,0.41)--(axis cs:2.8,0.41) node[midway,above,font=\scriptsize]{$z_\beta$};
\end{axis}
\end{tikzpicture}
\caption{$\mu=z_\alpha+z_\beta$ の意味：棄却限界 $z_\alpha$ が，$H_1$ の分布の下側 $100\beta\%$ 点 $\mu-z_\beta$ に一致する．}\label{fig:12-core}
\end{figure}
両側検定（$|Z|>z_{\alpha/2}$）では，反対側の棄却確率を無視して $\mu=z_{\alpha/2}+z_\beta$ とする．

\paragraph{症例数の一般公式}
推定量 $\hat\theta$ が $\hat\theta\approx\Normal(\theta,v/n)$（$v$ は1単位あたりの分散）で，$H_0:\theta=\theta_0$，対立仮説の値 $\theta_1$ に対して
$Z=(\hat\theta-\theta_0)/\sqrt{v/n}$ を用いるなら $\mu=(\theta_1-\theta_0)\sqrt{n/v}$ だから
\begin{equation}
n=\frac{(z_\alpha+z_\beta)^2\,v}{(\theta_1-\theta_0)^2}.\label{eq:12-general}
\end{equation}
分散が $H_0$ と $H_1$ で異なり（$v_0,v_1$），棄却限界を $H_0$ の分散で決めるなら
\begin{equation}
n=\frac{\bigl(z_\alpha\sqrt{v_0}+z_\beta\sqrt{v_1}\bigr)^2}{(\theta_1-\theta_0)^2}.\label{eq:12-general2}
\end{equation}
（$\Prob(\hat\theta>\theta_0+z_\alpha\sqrt{v_0/n}\mid\theta_1)=1-\beta$ を解く．）

\section{主な症例数の式}\label{sec:12-formulas}
\subsection{連続アウトカム}
\begin{itemize}
\item 2群の平均の差（各群 $n$，共通 SD $\sigma$，差 $\Delta$）：$v=2\sigma^2$ より
\begin{equation}
n=\frac{2\sigma^2(z_{\alpha/2}+z_\beta)^2}{\Delta^2}\quad\text{（各群）}.\label{eq:12-means}
\end{equation}
\item 対応のある比較（差の SD $\sigma_D$）：$n=\sigma_D^2(z_{\alpha/2}+z_\beta)^2/\Delta^2$．
\item ベースラインで調整する ANCOVA：\cref{eq:12-means}に $1-\rho^2$ を掛ける（\cref{ch:08}）．
\item 小標本では $z$ を $t$ に置き換えて反復計算する（数人増える程度）．
\end{itemize}
\subsection{2 値アウトカム}
比率 $\pi_1,\pi_0$，$\bar\pi=(\pi_1+\pi_0)/2$，各群 $n$ 人で，プール分散による検定（Pearson $\chi^2$）を用いるなら $v_0=2\bar\pi(1-\bar\pi)$，$v_1=\pi_1(1-\pi_1)+\pi_0(1-\pi_0)$ で
\begin{equation}
n=\frac{\bigl\{z_{\alpha/2}\sqrt{2\bar\pi(1-\bar\pi)}+z_\beta\sqrt{\pi_1(1-\pi_1)+\pi_0(1-\pi_0)}\bigr\}^2}{(\pi_1-\pi_0)^2}.\label{eq:12-prop}
\end{equation}
簡便には $n=(z_{\alpha/2}+z_\beta)^2\{\pi_1(1-\pi_1)+\pi_0(1-\pi_0)\}/(\pi_1-\pi_0)^2$．

\subsection{生存時間：イベント数}\label{subsec:12-events}
\Hinshutsu
\paragraph{指数モデル（\kakomon{2018}{1}）}
$\V[\log\hat\lambda_g]\approx1/d_g$ より，各群 $r$ イベントなら $\log\widehat{\HR}$ の分散は $2/r$．\cref{eq:12-general}で $v/n\to2/r$ とおけば
\begin{equation}
r=\frac{2(z_{\alpha/2}+z_\beta)^2}{(\log\HR)^2}\quad\text{（各群）}.\label{eq:12-exp}
\end{equation}
\paragraph{ログランク検定（Schoenfeld の式）}
割付比 $p:(1-p)$，総イベント数 $D$ のとき，$H_0$ 近傍で $\V[\log\widehat{\HR}]\approx1/\{Dp(1-p)\}$ なので
\begin{equation}
D=\frac{(z_{\alpha/2}+z_\beta)^2}{p(1-p)(\log\HR)^2},\qquad p=\tfrac12\ \text{なら}\ D=\frac{4(z_{\alpha/2}+z_\beta)^2}{(\log\HR)^2}.\label{eq:12-schoenfeld}
\end{equation}
1:1 割付なら $2r=D$ で\cref{eq:12-exp}と一致する．
\paragraph{イベント数から症例数へ}
必要症例数は $N=D/\Prob(\text{イベント})$ で，イベント確率は登録期間・追跡期間・ハザードから求める．
\begin{supbox}[補足：一様登録の下でのイベント確率]
登録期間 $a$ に一様に登録し，最終登録後さらに $f$ だけ追跡する．ハザード $\lambda$ の指数分布なら，1人がイベントを観察される確率は
$\Prob=1-\dfrac{e^{-\lambda f}-e^{-\lambda(a+f)}}{\lambda a}$（$S(t)$ を追跡期間 $[f,a+f]$ で平均して1から引く）．
\end{supbox}
生存時間試験で「症例数を増やす」だけでなく「追跡を延ばす」ことで検出力が上がるのは，精度を決めるのがイベント数だからである．

\subsection{単群試験（\kakomon{2023}{2}）}\label{subsec:12-single}
打ち切りなしの指数分布で $H_0:\lambda=\lambda_0$ 対 $H_1:\lambda>\lambda_0$．
$\hat\lambda=1/\bar T$，$I_n(\lambda)=n/\lambda^2$ で，$Z=(\hat\lambda-\lambda_0)/\sqrt{\lambda_1^2/n}$（情報量を $\lambda_1$ で評価）とすると
$\mu=\sqrt n(\lambda_1-\lambda_0)/\lambda_1$，\cref{eq:12-core}から
\begin{equation}
n=\frac{\lambda_1^2(z_\alpha+z_\beta)^2}{(\lambda_1-\lambda_0)^2}.\label{eq:12-single}
\end{equation}
回復率 $F(5)$ が $H_0$ で0.5，$H_1$ で0.75なら $\lambda_0=\log2/5$，$\lambda_1=\log4/5=2\lambda_0$，
$z_{0.025}=1.96$，$z_{0.2}=0.84$ から $n=4\times2.8^2=31.4$，すなわち32人．
\begin{supbox}[発展：尺度の選び方で症例数が変わる]
同じ問題でも，$\V[\log\hat\lambda]=1/n$ を使った $Z=\sqrt n(\log\hat\lambda-\log\lambda_0)$ なら $n=(z_\alpha+z_\beta)^2/(\log\lambda_1/\lambda_0)^2=7.84/0.480=16.3$，すなわち17人となる．
$2\lambda\sum T_i\sim\chi^2_{2n}$ を用いた正確な検定では19人で検出力0.82となり，
\cref{eq:12-single}の32人は保守的（検出力0.97），対数尺度の17人はやや楽観的（検出力0.77）である．
正規近似の精度は統計量の尺度に依存するので，試験では「どの検定を用いるか」を先に決め，その検定に合わせた式を使う．
\end{supbox}

\section{その他の考慮}\label{sec:12-other}
\begin{itemize}
\item \textbf{脱落}：脱落率 $w$ を見込むなら $n/(1-w)$．
\item \textbf{不均等割付}：$k:1$ 割付では，総数は均等割付の $(1+k)^2/(4k)$ 倍必要．
\item \textbf{多重性}：複数の主要比較があれば $\alpha$ を分割して $z_{\alpha/(2m)}$ を用いる（\cref{ch:14}）．
\item \textbf{中間解析}：群逐次デザインでは最大症例数がわずかに増える（\cref{ch:14}）．
\item \textbf{同等性・非劣性}：対立仮説の側が変わるので式が異なる（\cref{ch:13}）．
\end{itemize}

\section{仮定}\label{sec:12-assump}
\begin{itemize}
\item 検定統計量の正規近似が成り立つ程度の標本サイズ．
\item 分散・対照群のイベント率・相関などの見積もりが正しいこと（外れると検出力が大きく変わる）．
\item 生存時間：比例ハザード，登録・追跡・脱落の見込み．
\end{itemize}

\section{結果の解釈}\label{sec:12-interp}
\begin{itemize}
\item 「$n=122$ で検出力90\%」は，「真の差が $\Delta$ なら90\%の確率で有意になる」の意味．真の差が小さければ検出力は下がる．
\item 試験後に観測された効果量で「事後の検出力」を計算しても，P 値と1対1に対応するだけで新しい情報は与えない．非有意の解釈は信頼区間で行う．
\end{itemize}

\section{手法選択}\label{sec:12-choice}
主解析で用いる検定と同じ統計量・同じ尺度で症例数を計算する．
平均の比較なら\cref{eq:12-means}，比率なら\cref{eq:12-prop}，生存時間ならイベント数\cref{eq:12-schoenfeld}から症例数へ，
順位検定なら $\Prob(X>Y)$ を用いた\cref{eq:03-power}，二段階デザインなら二項分布による正確計算．

\section{よくある誤り}\label{sec:12-err}
\begin{warnbox}
\begin{itemize}
\item 両側検定なのに $z_\alpha$（片側）を使う，またはその逆．$z_\beta$ と $z_{\beta/2}$ を混同する．
\item 平均の差の式で係数2（2群分の分散）を落とす．結果を「総数」と「各群」で取り違える．
\item 生存時間で症例数とイベント数を混同する．
\item 端数を切り捨てる（必ず切り上げる）．
\end{itemize}
\end{warnbox}

\section{数理統計との接続}\label{sec:12-link}
\begin{linkbox}
\begin{itemize}
\item 検出力関数（『現代数理統計学の基礎』定義7.13）と Neyman--Pearson の補題（定理7.16）．同書の $\beta(\theta)$ は検出力である．
\item \kakomonSM{2024}{1} では片側 $z$ 検定の検出力 $1-\Phi(z_\alpha-\beta_1/\SE)$ を2つの推定量で比べ，分散の小さい（有効な）推定量ほど検出力が高いことが問われた．
      症例数設計は「分散を小さくする」ことと「$n$ を増やす」ことが同じ効果をもつことの定量化である．
\item Wald 型の検定では情報量 $I_n(\theta)$ をどの $\theta$ で評価するかが近似の質を左右する（\cref{subsec:12-single}の発展）．
\end{itemize}
\end{linkbox}

\section{例題}\label{sec:12-ex}
\begin{reidai}[平均の差]\label{reidai:12-1}
主要評価項目の SD を12，臨床的に意味のある差を5とし，両側 $\alpha=0.05$，検出力90\%（$z_{0.10}=1.282$）で各群の症例数を求めよ．
ベースラインとの相関0.5を利用して ANCOVA を行う場合はどうか．
\end{reidai}

\begin{reidai}[比率の差]\label{reidai:12-2}
対照群の有効率0.45，試験群0.60を想定し，両側 $\alpha=0.05$，検出力80\%（$z_{0.20}=0.842$）とする．
\cref{eq:12-prop}と簡便式で各群の症例数を求めよ（$\sqrt{0.49875}=0.7062$，$\sqrt{0.4875}=0.6982$）．
\end{reidai}

\begin{reidai}[イベント数]\label{reidai:12-3}
HR $=0.7$ を両側 $\alpha=0.05$，検出力90\%で検出したい（1:1割付，$\log0.7=-0.3567$）．
\begin{enumerate}[label=(\arabic*)]
\item 必要な総イベント数を求めよ．
\item 追跡期間中にイベントを起こす確率が両群平均で0.6と見込まれるとき，必要な総症例数を求めよ．
\item 脱落を10\%見込むとどうなるか．
\end{enumerate}
\end{reidai}

\begin{reidai}[検出力の計算]\label{reidai:12-4}
各群50人，効果量 $\Delta/\sigma=0.5$ の2群比較を片側 $\alpha=0.025$ で行うときの検出力を求めよ（$\Phi(0.54)=0.705$）．
検出力を90\%にするには各群何人必要か．
\end{reidai}

\section{解答}\label{sec:12-sol}
\begin{kaitou}{reidai:12-1}
$n=2\times144\times(1.960+1.282)^2/25=288\times10.51/25=121.1$，各群122人．
ANCOVA では $1-\rho^2=0.75$ 倍で $90.8$，各群91人．
\end{kaitou}

\begin{kaitou}{reidai:12-2}
$\bar\pi=0.525$．$1.96\times0.7062=1.384$，$0.842\times0.6982=0.588$，和 $1.972$，2乗 $3.889$．$n=3.889/0.0225=172.8$ → 各群173人．
簡便式：$(2.802)^2\times0.4875/0.0225=7.851\times21.67=170.1$ → 171人．
\end{kaitou}

\begin{kaitou}{reidai:12-3}
(1) $D=4\times(1.960+1.282)^2/0.3567^2=4\times10.51/0.1272=330.5$ → 331イベント．\\
(2) $331/0.6=551.7$ → 552人．\\
(3) $331/(0.6\times0.9)=613.0$ → 613人（途中で切り上げた $552/0.9$ を使えば614人．脱落が打ち切りとしてイベントを減らすことを見込む）．
\end{kaitou}

\begin{kaitou}{reidai:12-4}
$\mu=0.5\sqrt{50/2}=2.5$，検出力 $\Phi(2.5-1.96)=\Phi(0.54)=0.705$．
90\%には $n=2(1.960+1.282)^2/0.25=84.1$ → 各群85人．
\end{kaitou}

\begin{summarybox}
\begin{itemize}
\item $H_0:\Normal(0,1)$，$H_1:\Normal(\mu,1)$ なら検出力 $\Phi(\mu-z_\alpha)$，条件は $\mu=z_\alpha+z_\beta$（両側は $z_{\alpha/2}$）．
\item 一般：$n=(z_\alpha+z_\beta)^2v/\Delta^2$，分散が異なれば $(z_\alpha\sqrt{v_0}+z_\beta\sqrt{v_1})^2/\Delta^2$．
\item 平均：各群 $2\sigma^2(z_{\alpha/2}+z_\beta)^2/\Delta^2$．比率：\cref{eq:12-prop}．
\item 生存：総イベント数 $4(z_{\alpha/2}+z_\beta)^2/(\log\HR)^2$（指数なら各群 $2(\cdot)^2/(\log\HR)^2$）．症例数 $=$ イベント数/イベント確率．
\item 単群指数：$n=\lambda_1^2(z_\alpha+z_\beta)^2/(\lambda_1-\lambda_0)^2$（\kakomon{2023}{2}）．尺度により結果が変わる．
\end{itemize}
\end{summarybox}

% ===== ch13.tex =====
\chapter{同等性・非劣性・生物学的同等性}\label{ch:13}

\begin{goalbox}
\begin{itemize}
\item 優越性・同等性・非劣性の帰無仮説と対立仮説を書き分けられる．
\item 「有意差なし」と「同等」が異なることを，信頼区間を用いて説明できる．
\item 2つの片側検定（TOST）の第1種の過誤確率が $\alpha$ 以下であることを示し，$100(1-2\alpha)\%$ 信頼区間による判定と同値であることを説明できる．
\item 対数正規分布の中央値・密度・平均・分散・変動係数を導出できる．
\item 2剤2期クロスオーバー試験のデータから幾何平均比の90\%信頼区間を計算し，生物学的同等性を判定できる．
\end{itemize}
\end{goalbox}

\begin{exambox}
\begin{itemize}
\item \kakomon{2016}{4}：AUC の対数正規分布．中央値 $e^\mu$，密度の導出，MGF による平均・分散，$\mathrm{CV}=1$ から $\sigma^2=\log2$，
      10名のクロスオーバー試験（$\bar Y=0.1$，$s=0.5$）で中央値の比の90\%CI $(0.83,1.48)$ → 同等域 $(0.8,1.25)$ に含まれず生物学的同等でない．
\item \kakomon{2021}{2}：$\mu$ の信頼区間，片側仮説 $H_{01}:\mu\le-\Delta$，$H_{02}:\mu\ge\Delta$ の検定統計量と棄却域，
      $H_0=H_{01}\cup H_{02}$ の検定の第1種の過誤確率が $\alpha$ 以下であること，20名で平均 $-0.1$，SD $0.26$，$\Delta=\log1.25$ の判定（90\%CI $[-0.201,0.001]$ → 同等）．
\end{itemize}
\end{exambox}

\flowline{2群または\\クロスオーバー}{平均差・\\幾何平均比}{TOST\\$(1-2\alpha)$CI}{対数正規\\持ち越しなし}{$t$ 区間\\指数変換}{同等域に\\含まれるか}

\section{問題設定：3種類の比較}\label{sec:13-setting}
試験治療 $T$ と対照 $C$ の差（または対数比）を $\theta=\mu_T-\mu_C$（大きいほど $T$ が良い）とし，臨床的に許容できる差（マージン）を $\Delta>0$ とする．
\begin{table}[htbp]
\centering\small
\caption{3種類の比較の仮説}\label{tab:13-hyp}
\begin{tabular}{@{}lll@{}}
\toprule
目的 & 帰無仮説 $H_0$ & 対立仮説 $H_1$\\
\midrule
優越性（superiority） & $\theta\le0$ & $\theta>0$\\
非劣性（non-inferiority） & $\theta\le-\Delta$ & $\theta>-\Delta$\\
同等性（equivalence） & $\theta\le-\Delta$ または $\theta\ge\Delta$ & $-\Delta<\theta<\Delta$\\
\bottomrule
\end{tabular}
\end{table}
同等性・非劣性では「示したいこと（同等・劣らない）」を対立仮説に置く．
\begin{warnbox}[「有意差なし」と「同等」]
優越性の検定で $H_0:\theta=0$ が棄却されないことは，差がないことの証拠ではない．症例数が少なければ大きな差があっても有意にならない．
同等性は，信頼区間全体が $(-\Delta,\Delta)$ に含まれること（差が大きい可能性をデータが否定すること）で示す．
\cref{reidai:13-2}は，$t$ 検定で有意差がないのに同等性も示せない例である．
\end{warnbox}

\section{数理：2つの片側検定（TOST）}\label{sec:13-tost}
\Hissu
$Y_1,\dots,Y_n\iid\Normal(\theta,\sigma^2)$（差の標本，$\sigma^2$ 未知），$\bar Y$，不偏分散 $s^2$ とする．
\begin{meidai}[片側検定（\kakomon{2021}{2}〔1-2〕）]\label[meidai]{prop:13-onesided}
$H_{01}:\theta\le-\Delta$ 対 $H_{A1}:\theta>-\Delta$ は $T_1=\sqrt n(\bar Y+\Delta)/s>t_{n-1}(\alpha)$ で，
$H_{02}:\theta\ge\Delta$ 対 $H_{A2}:\theta<\Delta$ は $T_2=\sqrt n(\bar Y-\Delta)/s<-t_{n-1}(\alpha)$ で棄却する．いずれも有意水準 $\alpha$ の検定である．
\end{meidai}
\begin{proof}
$\theta=-\Delta$ で $T_1\sim t_{n-1}$ なので棄却確率は $\alpha$．$\theta<-\Delta$ なら $T_1$ は分布が左にずれ，棄却確率は $\alpha$ より小さい．$T_2$ も同様．
\end{proof}
\begin{teiri}[交差和検定（IUT）の有意水準（\kakomon{2021}{2}〔1-3〕）]\label[teiri]{thm:13-iut}
$H_0=H_{01}\cup H_{02}$ を，「$H_{01}$ と $H_{02}$ がともに有意水準 $\alpha$ で棄却されたときに棄却する」検定の第1種の過誤確率は $\alpha$ 以下である．
\end{teiri}
\begin{proof}
$H_0$ が真なら $H_{01}$ と $H_{02}$ の少なくとも一方が真である．$H_{01}$ が真なら
$\Prob(\text{両方棄却})\le\Prob(T_1\in R_1)\le\alpha$．$H_{02}$ が真の場合も同様．
\end{proof}
\Youchui 各検定を $\alpha/2$ に調整する必要はない．$H_0$ の下で「両方棄却」が起こるには，真である方の帰無仮説を誤って棄却しなければならないからである．

\begin{meidai}[信頼区間による判定]\label[meidai]{prop:13-ci}
TOST（各片側 $\alpha$）で $H_0$ を棄却することは，$100(1-2\alpha)\%$ 信頼区間
$\bar Y\pm t_{n-1}(\alpha)s/\sqrt n$ が $(-\Delta,\Delta)$ に含まれることと同値である．
\end{meidai}
\begin{proof}
$T_1>t_{n-1}(\alpha)\iff\bar Y-t_{n-1}(\alpha)s/\sqrt n>-\Delta$（下限 $>-\Delta$），
$T_2<-t_{n-1}(\alpha)\iff\bar Y+t_{n-1}(\alpha)s/\sqrt n<\Delta$（上限 $<\Delta$）．
\end{proof}
片側5\%の TOST は\textbf{90\%信頼区間}で判定する．
非劣性（片側 $\alpha=2.5\%$）は95\%信頼区間の下限が $-\Delta$ を上回るかで判定する．

\section{非劣性試験}\label{sec:13-ni}
\begin{itemize}
\item 標準治療より副作用が少ない・投与が簡便などの利点がある新治療で，「有効性が許容範囲を超えて劣らない」ことを示す．
\item 判定：$\theta$ の95\%CI の下限 $>-\Delta$．下限が0も上回れば，続けて優越性も主張できる（閉手順なので多重性の調整は不要）．
\item マージン $\Delta$ はデータを見る前に，臨床的判断と過去のプラセボ対照試験の効果の大きさから決める．
\end{itemize}
\begin{supbox}[補足（出典：ICH E10 \cite{ich:e10}，FDA 非劣性ガイダンス \cite{fda:ni}）]
試験薬の対照薬に対する劣り幅が $M_1$（対照薬のプラセボに対する効果を過去の試験から保守的に見積もった値）を超えないこと（＝試験薬がプラセボより有効であること）を最低条件とし，さらに臨床的判断で $M_1$ の一部（例えば50\%）を保持するよう $M_2\,(\le M_1)$ を非劣性マージンとして設定するのが一般的である．
非劣性試験では，(i) 分析感度（assay sensitivity：試験が差を検出できる能力），(ii) 恒常性（過去の試験と現在の試験で対照薬の効果が同じ）が前提となる．
質の低い試験（脱落が多い，遵守が悪い）は両群の差を小さくし，非劣性を示しやすくしてしまう．このため ITT 解析と治験実施計画書に適合した対象集団の解析の両方で一貫した結果が求められる．
\end{supbox}

\section{対数正規分布と生物学的同等性}\label{sec:13-be}
\subsection{対数正規分布}\label{subsec:13-lognormal}
\Hissu 薬物動態パラメータ（AUC，$C_{\max}$）は正の値をとり右に歪むので，$X\sim LN(\mu,\sigma^2)$，すなわち $Y=\log X\sim\Normal(\mu,\sigma^2)$ とモデル化する．
\begin{meidai}[\kakomon{2016}{4}〔1〕〜〔4〕]\label[meidai]{prop:13-ln}
\begin{enumerate}[label=(\arabic*)]
\item 中央値 $e^\mu$（$\log$ は単調なので $\Prob(X\le e^\mu)=\Prob(Y\le\mu)=1/2$）．
\item 密度 $f(x)=\dfrac{1}{x\sqrt{2\pi}\sigma}\exp\Bigl\{-\dfrac{(\log x-\mu)^2}{2\sigma^2}\Bigr\}$（$x>0$）．
\item $\E[X]=e^{\mu+\sigma^2/2}$，$\V[X]=e^{2\mu+\sigma^2}(e^{\sigma^2}-1)$．
\item 変動係数 $\mathrm{CV}=\sqrt{\V[X]}/\E[X]=\sqrt{e^{\sigma^2}-1}$，逆に $\sigma^2=\log(1+\mathrm{CV}^2)$．
\end{enumerate}
\end{meidai}
\begin{proof}
(2) $\Prob(X\le x)=\Prob(Y\le\log x)=\Phi\bigl((\log x-\mu)/\sigma\bigr)$ を $x$ で微分し，$(\log x)'=1/x$．
(3) $M_Y(t)=\E[e^{tY}]=e^{\mu t+\sigma^2t^2/2}$ で，$\E[X]=M_Y(1)$，$\E[X^2]=M_Y(2)=e^{2\mu+2\sigma^2}$．
(4) (3)から直ちに従う．
\end{proof}
\kakomon{2016}{4} では $\mathrm{CV}=1$ から $\sigma^2=\log2$．CV が小さいとき $\sigma\approx\mathrm{CV}$ である．

\subsection{幾何平均比}
対数スケールでの平均の差 $\mu_T-\mu_R$ を逆変換した $e^{\mu_T-\mu_R}$ は中央値の比（幾何平均の比，GMR）である（\kakomon{2021}{2}〔2-1〕）．
\Youchui 算術平均の比 $e^{\mu_T+\sigma_T^2/2}/e^{\mu_R+\sigma_R^2/2}$ とは，分散が等しいときに限り一致する．

\subsection{生物学的同等性の判定}
\begin{supbox}[補足（出典：後発医薬品の生物学的同等性試験ガイドライン \cite{jp:be}）]
日本のガイドラインでは，AUC と $C_{\max}$ の対数値の平均の差の90\%信頼区間が $\log(0.80)$ から $\log(1.25)$ の範囲にあるとき生物学的に同等と判定する．
これは片側5\%の2つの片側検定に相当し，同等でない製剤が合格する確率（消費者危険）を5\%以下に抑える．
$\log1.25=0.2231=-\log0.80$ なので，同等域は対数スケールで $0$ に関して対称である．
\end{supbox}
\paragraph{判定の手順}
\begin{enumerate}
\item AUC を対数変換し，被験者ごとの差 $d_i=\log X_i^T-\log X_i^R$（またはクロスオーバーのモデルに基づく推定値）を求める．
\item $\bar d\pm t_\nu(0.05)\,\SE$ を計算する．
\item 指数変換して GMR の90\%CI を得て，$(0.80,1.25)$ に含まれるかを判定する．
\end{enumerate}
\kakomon{2016}{4}〔5〕：$\bar Y=0.1$，$s=0.5$，$n=10$，$t_9(0.05)=1.833$ から $0.1\pm0.29=(-0.19,0.39)$，
GMR（既存薬/後発品）の区間 $(0.83,1.48)$ は上限が1.25を超え，同等とは言えない．\\
\kakomon{2021}{2}〔2-2〕：$\bar d=-0.1$，$s=0.26$，$n=20$，$t_{19}(0.05)=1.729$ から $[-0.201,0.001]\subset(-0.223,0.223)$ で同等．

\subsection{2剤2期クロスオーバー試験}\label{subsec:13-crossover}
\Hatten 系列 TR（第1期T，第2期R）に $n_1$ 人，RT に $n_2$ 人を割り付け，被験者 $i$ の第 $p$ 期の対数値を
\[
Y_{ip}=\mu+\xi_i+\pi_p+\tau_{d(i,p)}+\varepsilon_{ip}
\]
とする（$\xi_i$：被験者効果，$\pi_p$：時期効果，$\tau_d$：製剤効果，$\varepsilon$：被験者内誤差 $\Normal(0,\sigma_w^2)$，持ち越し効果なし）．
被験者ごとの期間差の半分 $c_i=(Y_{i2}-Y_{i1})/2$ は $\xi_i$ を消去し，
系列 RT では $\E[c_i]=\frac{\pi_2-\pi_1}{2}+\frac{\tau_T-\tau_R}{2}$，系列 TR では $\frac{\pi_2-\pi_1}{2}-\frac{\tau_T-\tau_R}{2}$．よって
\[
\hat\tau=\bar c_{RT}-\bar c_{TR},\qquad \V[\hat\tau]=\frac{\sigma_w^2}{2}\Bigl(\frac{1}{n_1}+\frac{1}{n_2}\Bigr),
\]
自由度 $n_1+n_2-2$ の $t$ 分布で区間を作る．時期効果が相殺される点が，単純な対応のある差 $\log X^T-\log X^R$ の $t$ 区間（試験問題で用いられた簡略化）との違いである．
持ち越し効果の検定は被験者間比較（被験者ごとの和の系列間比較）になり検出力が低いので，十分な休薬期間を設けてデザインで防ぐ．

\section{症例数の考え方}\label{sec:13-n}
\begin{supbox}[補足：同等性試験の症例数]
真の差が0のとき，TOST の検出力はおよそ $2\Phi(\Delta/\SE-z_\alpha)-1$ である．これを $1-\beta$ とおくと $\Delta/\SE=z_\alpha+z_{\beta/2}$．
差の標本の SD を $\sigma$ とすれば $n=(z_\alpha+z_{\beta/2})^2\sigma^2/\Delta^2$．真の差が0でないと，必要数は大きく増える．
\end{supbox}

\section{仮定}\label{sec:13-assump}
\begin{itemize}
\item 対数変換後の正規性（対数正規性），被験者の独立性．
\item クロスオーバー：持ち越し効果がない（十分な休薬期間），時期と製剤の交互作用がない．
\item 非劣性：分析感度・恒常性，マージンの事前設定．
\end{itemize}

\section{結果の解釈}\label{sec:13-interp}
\begin{itemize}
\item 「GMR の90\%CI $(0.98,1.13)$ が $(0.80,1.25)$ に含まれるので生物学的に同等」．
\item 「90\%CI $(0.86,1.29)$ は上限が1.25を超え，同等性は示されない」．この区間は1を含むので両側5\%の優越性検定も有意ではないが，それは同等の根拠にならない．
\item 非劣性：「95\%CI の下限 $-7.9\%$ がマージン $-10\%$ を上回るので非劣性が示された．上限が0を含むので優越性は示されない．」
\end{itemize}

\section{手法選択}\label{sec:13-choice}
\begin{itemize}
\item 後発品の生物学的同等性 → 対数変換，クロスオーバー，90\%CI を $(0.80,1.25)$ と比較．
\item 新治療の有効性が標準治療に劣らないこと → 非劣性（95\%CI の下限とマージン）．
\item 差が両方向に小さいこと → 同等性（TOST，$(1-2\alpha)$ CI）．
\end{itemize}

\section{よくある誤り}\label{sec:13-err}
\begin{warnbox}
\begin{itemize}
\item 優越性検定の非有意を同等とみなす．
\item 同等性の判定に95\%CI を使う（片側5\%の TOST に対応するのは90\%CI）．あるいは TOST の各検定を $\alpha/2$ にする．
\item AUC を対数変換せずに比の区間を作る，または対数スケールの区間を指数変換し忘れる．
\item 対数正規分布の中央値を $e^{\mu+\sigma^2/2}$ とする（それは平均）．
\end{itemize}
\end{warnbox}

\section{数理統計との接続}\label{sec:13-link}
\begin{linkbox}
\begin{itemize}
\item 対数正規分布は『現代数理統計学の基礎』第3章演習．$\E[e^{\bar X}]$ の偏りと不偏化 $\exp(\bar X-\frac1{2n})$ は\kakomonSM{2016}{1}，
      平均と分散の MLE $\exp(\hat\mu+\hat\sigma^2/2)$ は\kakomonSM{2025}{5}（不変性）．
\item $t$ 区間の根拠は正規標本の $\bar X$ と不偏分散の独立性（同書 定理5.1）．
\item 片側検定の有意水準が境界 $\theta=-\Delta$ で最大になることは，単調尤度比をもつ分布族の性質（同書 定理7.18）．
\item \kakomonSM{2022}{5} の「対応のある2標本＝二元配置」はクロスオーバーの被験者効果の消去と同じ考え方である．
\end{itemize}
\end{linkbox}

\section{例題}\label{sec:13-ex}
\begin{reidai}[生物学的同等性]\label{reidai:13-1}
24名のクロスオーバー試験で，被験者ごとの $\log\mathrm{AUC}_T-\log\mathrm{AUC}_R$ の平均が $0.05$，標準偏差が $0.20$ であった．
$t_{23}(0.05)=1.714$，$e^{-0.020}=0.980$，$e^{0.120}=1.127$ を用いよ．
\begin{enumerate}[label=(\arabic*)]
\item GMR の点推定値と90\%信頼区間を求め，生物学的同等性を判定せよ．
\item 同じ判定を TOST の2つの $t$ 統計量で示せ（$\log1.25=0.2231$）．
\end{enumerate}
\end{reidai}

\begin{reidai}[有意差なしと同等]\label{reidai:13-2}
8名の同様の試験で，平均 $0.05$，標準偏差 $0.30$ であった（$t_7(0.05)=1.895$，$t_7(0.025)=2.365$，$e^{-0.151}=0.860$，$e^{0.251}=1.285$）．
\begin{enumerate}[label=(\arabic*)]
\item $H_0:\mu_T=\mu_R$ の両側5\%の $t$ 検定を行え．
\item 生物学的同等性を判定せよ．
\item (1)(2)の結果から何が言えるか．
\end{enumerate}
\end{reidai}

\begin{reidai}[対数正規分布]\label{reidai:13-3}
AUC の変動係数が30\%と見込まれる．$\log(1.09)=0.0862$，$e^{0.0431}=1.044$ を用いよ．
\begin{enumerate}[label=(\arabic*)]
\item 対数スケールの分散 $\sigma^2$ を求めよ．
\item 中央値に対する平均の比を求めよ．
\end{enumerate}
\end{reidai}

\begin{reidai}[非劣性]\label{reidai:13-4}
有効率の差（新治療$-$標準治療）の推定値が $-0.02$，標準誤差 $0.03$，非劣性マージンが $0.10$ である．非劣性と優越性を判定せよ．
\end{reidai}

\section{解答}\label{sec:13-sol}
\begin{kaitou}{reidai:13-1}
(1) $\SE=0.20/\sqrt{24}=0.0408$．$0.05\pm1.714\times0.0408=0.05\pm0.070=(-0.020,0.120)$．
GMR $=e^{0.05}=1.05$，90\%CI $(0.980,1.127)\subset(0.80,1.25)$ で生物学的に同等．\\
(2) $T_1=(0.05+0.2231)/0.0408=6.69>1.714$，$T_2=(0.05-0.2231)/0.0408=-4.24<-1.714$．両方の片側帰無仮説が棄却され同等．
\end{kaitou}

\begin{kaitou}{reidai:13-2}
(1) $\SE=0.30/\sqrt8=0.106$，$t=0.05/0.106=0.47$，$|t|<2.365$ で有意差なし．\\
(2) 90\%CI $0.05\pm1.895\times0.106=(-0.151,0.251)$ → $(0.860,1.285)$．上限が1.25を超え同等とは言えない．\\
(3) 有意差がないのは症例数が少なく推定が不正確だからであり，同等であることは示されない．「差があるとも同等とも言えない」が正しい結論．
\end{kaitou}

\begin{kaitou}{reidai:13-3}
(1) $\sigma^2=\log(1+0.3^2)=\log1.09=0.0862$．\\
(2) $\E[X]/e^\mu=e^{\sigma^2/2}=e^{0.0431}=1.044$．平均は中央値より約4\%大きい．
\end{kaitou}

\begin{kaitou}{reidai:13-4}
95\%CI $-0.02\pm1.96\times0.03=(-0.079,\,0.039)$．下限 $-0.079>-0.10$ で非劣性が示される．区間が0を含むので優越性は示されない．
\end{kaitou}

\begin{summarybox}
\begin{itemize}
\item 優越性 $H_0:\theta\le0$，非劣性 $H_0:\theta\le-\Delta$，同等性 $H_0:|\theta|\ge\Delta$．示したいことを対立仮説に．
\item TOST：各片側 $\alpha$ で両方棄却 → 全体の第1種の過誤 $\le\alpha$（IUT）．$\iff$ $(1-2\alpha)$CI $\subset(-\Delta,\Delta)$．
\item 非劣性：95\%CI 下限 $>-\Delta$．
\item 対数正規：中央値 $e^\mu$，平均 $e^{\mu+\sigma^2/2}$，$\V=e^{2\mu+\sigma^2}(e^{\sigma^2}-1)$，$\mathrm{CV}=\sqrt{e^{\sigma^2}-1}$．
\item BE：$\log$ 差の90\%CI を指数変換して $(0.80,1.25)$ に含まれれば同等．
\item 「有意差なし」≠「同等」．
\end{itemize}
\end{summarybox}

% ===== ch14.tex =====
\chapter{中間解析・二段階試験・多重性・用量反応}\label{ch:14}

\begin{goalbox}
\begin{itemize}
\item 繰り返し検定による第1種の過誤の増大を説明し，情報時間と Brown 運動の構造から中間解析の検定統計量の相関を求められる．
\item $\alpha$ 消費関数から各解析で消費する $\alpha$ と棄却限界値を求め，2変量正規分布を用いて最終解析の限界値を定める方法を説明できる．
\item 第II相の二段階デザインの第1種・第2種の過誤確率，早期中止確率，期待症例数を二項分布で計算できる．
\item FWER を定義し，Bonferroni・Šidák・Holm・Simes 法を比較できる．相関が FWER に与える影響を示せる．
\item Dunnett 型の多対1比較の相関，閉検定手順・ゲートキーピングの考え方を説明できる．
\item 用量反応の対比（contrast）を設計し，対比検定どうしの相関を計算できる．
\end{itemize}
\end{goalbox}

\begin{exambox}
\begin{itemize}
\item \kakomon{2017}{2}：ログランク検定の群逐次デザイン．$\alpha(t)=0.05t^2$，$t=0.4$ で $\alpha=0.008$，$z_1=2.41$，
      $\Cov[Z(0.4),Z(1)]=\sqrt{0.4}$，$z_2$ を $\int_{-\infty}^{z_1}\int_{z_2}^\infty\phi_2(x_1,x_2;\sqrt{0.4})=0.042$ で定める方法，解析の追加（$t=0.7$）．
\item \kakomon{2017}{3}：Simon の二段階デザイン（$\pi_0=0.1$，$\pi_1=0.5$，$n_1=5,r_1=0,n_2=7,r=3$）．$\alpha\approx0.024$，$\beta\approx0.089$，PET，$\E[N\mid\pi_0]=7.87$．
\item \kakomon{2024}{4}：プラセボと3用量．独立3検定の FWER，Bonferroni，分散分析をゲートとした $z$ 検定で FWER が制御される条件，
      $Z_j$ と $Z_k$ の相関 $\sqrt{n_jn_k/\{(n_j+n_0)(n_k+n_0)\}}$，直線的用量反応の対比 $(-3,-1,1,3)/\sqrt{20}$，対比検定の相関 $3/\sqrt{15}$．
\item \kakomon{2025}{4}：2つの主要評価項目．独立な $p$ 値の FWER $0.0975$，Bonferroni の FWER $0.0494$，Šidák 型の制約，Simes と Bonferroni の棄却域の差，
      $\mathrm{FWER}(\rho)=2\alpha-\phi(\rho)\le\mathrm{FWER}(0)$．
\item \kakomon{2016}{2}：Cochran--Armitage 傾向性検定（\cref{sec:04-trend}）．
\end{itemize}
\end{exambox}

\section{問題設定}\label{sec:14-setting}
1つの試験の中で複数回・複数の仮説を検定すると，「どれか1つでも誤って有意になる」確率が名目の $\alpha$ を超える．
臨床試験では，(a) 時間的な繰り返し（中間解析），(b) 複数の群・用量，(c) 複数の評価項目，(d) 部分集団，で多重性が生じる．
本章では，これらを制御する標準的な方法を導く．

\section{群逐次デザインと \texorpdfstring{$\alpha$}{α} 消費関数}\label{sec:14-gsd}
\Hinshutsu
\subsection{情報時間と検定統計量の相関}
最終解析の情報量（ログランクならイベント数 $D$）に対する時点の情報量の割合を\emphx{情報時間} $t\in(0,1]$ とする（$t=L/D$）．
情報時間 $t$ の標準化統計量を $Z(t)$，$B(t)=\sqrt t\,Z(t)$ とすると，$H_0$ の下で $B(t)$ は近似的に標準 Brown 運動になり，
$\E[B(t)]=0$，$\Cov[B(t_1),B(t_2)]=\min(t_1,t_2)$（\kakomon{2017}{2}）．
\begin{meidai}\label[meidai]{prop:14-corr}
$t_1<t_2$ のとき $\Corr[Z(t_1),Z(t_2)]=\sqrt{t_1/t_2}$．
\end{meidai}
\begin{proof}
$\Cov[Z(t_1),Z(t_2)]=\Cov[B(t_1),B(t_2)]/\sqrt{t_1t_2}=t_1/\sqrt{t_1t_2}=\sqrt{t_1/t_2}$，分散はともに1．
\end{proof}
直観的には，$Z(t_2)$ は $Z(t_1)$ の情報に新しい情報を足したもの（独立増分）で，情報の重なりの割合が相関になる．

\paragraph{繰り返し検定の問題}
$t=0.5$ と $t=1$ の2回，それぞれ片側 $z_{0.025}=1.96$ で検定すると，全体の第1種の過誤確率は
$1-\Prob(Z(0.5)\le1.96,Z(1)\le1.96)\approx0.042$（相関 $\sqrt{0.5}$ の2変量正規）となり，名目の0.025を大きく超える．

\subsection{\texorpdfstring{$\alpha$}{α} 消費関数}\label{subsec:14-spending}
\begin{teigi}[$\alpha$ 消費関数]\label{def:14-spending}
$\alpha(0)=0$，$\alpha(1)=\alpha$ の非減少関数 $\alpha(t)$ を定め，情報時間 $t_k$ までに累積で $\alpha(t_k)$ の第1種の過誤確率を「消費」するように棄却限界値 $c_k$ を決める
（Lan--DeMets の方法 \cite{landemets1983}）．
\end{teigi}
\begin{enumerate}
\item 第1回：$\Prob(Z(t_1)>c_1)=\alpha(t_1)$ から $c_1=z_{\alpha(t_1)}$．
\item 第 $k$ 回：$\Prob(Z(t_1)\le c_1,\dots,Z(t_{k-1})\le c_{k-1},Z(t_k)>c_k)=\alpha(t_k)-\alpha(t_{k-1})$ を満たす $c_k$ を，多変量正規分布の数値積分で求める．
\end{enumerate}
\kakomon{2017}{2} では $\alpha(t)=0.05t^2$ で $t_1=0.4$ なら $\alpha(0.4)=0.008$，$c_1=z_{0.008}=2.41$，最終解析で消費するのは $0.042$．
最終の限界値 $c_2$ は
\[
\int_{-\infty}^{c_1}\int_{c_2}^{\infty}\phi_2\bigl(x_1,x_2;\sqrt{0.4}\bigr)\,\dd x_2\,\dd x_1=0.042
\]
を満たす値で，数値計算すると $c_2\approx1.68$（名目の $z_{0.05}=1.645$ よりわずかに大きい）．
$\alpha$ 消費関数法の利点は，解析の回数や時期をあらかじめ固定しなくてよいことで，
途中で $t=0.7$ の解析を追加しても，消費する $\alpha$ は $\alpha(0.7)-\alpha(0.4)=0.0165$，最終は $0.05-\alpha(0.7)=0.0255$ と自動的に決まる．

\begin{table}[htbp]
\centering\small
\caption{代表的な $\alpha$ 消費関数（$\alpha$ は全体の片側有意水準）}\label{tab:14-spending}
\begin{tabular}{@{}llL{62mm}@{}}
\toprule
名称 & $\alpha(t)$ & 特徴\\
\midrule
O'Brien--Fleming 型 & $2-2\Phi\bigl(z_{\alpha/2}/\sqrt t\bigr)$ & 初期はほとんど消費しない．最終の限界値は名目値に近い\\
Pocock 型 & $\alpha\log\{1+(e-1)t\}$ & 各回ほぼ同じ限界値．早期中止しやすいが最終が厳しい\\
べき型 & $\alpha t^\rho$ & $\rho=1$ で一様，$\rho$ が大きいほど初期の消費が少ない（2017年は $\rho=2$）\\
\bottomrule
\end{tabular}
\end{table}
\Youchui O'Brien--Fleming 型の式に $z_{\alpha/2}$ が現れるのは両側検定だからではない．O'Brien--Fleming 型の境界は $Z(t)>c/\sqrt t$，すなわち $B(t)>c$ であり，Brown 運動の鏡像原理から，時点 $t$ までに片側の水準 $c$ を越える確率は$\Prob(\max_{s\le t}B(s)\ge c)=2\{1-\Phi(c/\sqrt t)\}$ である．$t=1$ でこれを片側水準 $\alpha$ とおくと $c=z_{\alpha/2}$ となり，$\alpha(t)=2-2\Phi(z_{\alpha/2}/\sqrt t)$ を得る（$\alpha(1)=\alpha$）．たとえば片側 $\alpha=0.025$ なら $z_{\alpha/2}=z_{0.0125}=2.241$ を用いる．両側水準 $\alpha$ で対称な境界を設計する場合は，片側 $\alpha/2$ として同じ式を適用するのが一般的である（ソフトウェアによって $\alpha$ の意味が異なるので確認すること）．
\begin{supbox}[補足：無効中止]
効果が見込めないときの早期中止（無効中止）には，$\beta$ 消費関数による下側の境界や，条件付き検出力を用いる．
無効中止の境界は第1種の過誤を増やさない（拘束的でなければ $\alpha$ の計算には入れない）．
\end{supbox}

\section{第II相の二段階デザイン}\label{sec:14-simon}
\Hinshutsu 単群で有効率 $\pi$ を評価する．$\pi_0$（これ以下なら無効），$\pi_1$（期待する有効率）とし，$H_0:\pi\le\pi_0$ を検定する．
\begin{itemize}
\item 第1段階：$n_1$ 人を治療し，有効数 $X_1\le r_1$ なら無効として中止．
\item 第2段階：$X_1>r_1$ なら $n_2$ 人を追加し，総有効数 $X_1+X_2>r$ なら有効と結論．
\end{itemize}
$b(x;\pi,n)$，$B(x;\pi,n)=\sum_{k\le x}b(k;\pi,n)$ として（$X_1,X_2$ は独立）
\begin{align}
\alpha&=\Prob(\text{有効と結論}\mid\pi_0)=\sum_{x=r_1+1}^{n_1}b(x;\pi_0,n_1)\{1-B(r-x;\pi_0,n_2)\},\label{eq:14-simonalpha}\\
\beta&=\Prob(\text{無効と結論}\mid\pi_1)=B(r_1;\pi_1,n_1)+\sum_{x=r_1+1}^{\min(n_1,r)}b(x;\pi_1,n_1)B(r-x;\pi_1,n_2),\label{eq:14-simonbeta}\\
\mathrm{PET}&=B(r_1;\pi_0,n_1),\qquad \E[N\mid\pi_0]=n_1+(1-\mathrm{PET})\,n_2.\label{eq:14-simonEN}
\end{align}
（$r-x<0$ なら $B(r-x)=0$．）\kakomon{2017}{3} の設計 $(n_1,r_1,n_2,r)=(5,0,7,3)$ は，$\alpha\le0.05$，$\beta\le0.1$ を満たす設計のうち最大症例数 $n_1+n_2=12$ を最小にし，その中で $\E[N\mid\pi_0]$ を最小にするもの（Simon のミニマックスデザイン）にあたる．
公式解答は「$\E[N\mid\pi_0]$ を最小にする」デザインと述べているが，同じ条件で $\E[N\mid\pi_0]$ を最小にする最適デザインは $(4,0,9,3)$（$\alpha=0.029$，$\beta=0.093$，$\E[N\mid\pi_0]=7.10$）である（二項確率の全探索による）．
\begin{supbox}[補足（出典：Simon \cite{simon1989}）]
Simon の最適（optimal）デザインは $\E[N\mid\pi_0]$ を，ミニマックス（minimax）デザインは最大症例数 $n_1+n_2$ を最小化する．
無効な治療に多くの患者を曝さないことが第1段階の目的である．
\end{supbox}

\section{多重性}\label{sec:14-mult}
\Hissu
\subsection{FWER}
\begin{teigi}[FWER]\label{def:14-fwer}
検定する帰無仮説の族 $\{H_1,\dots,H_m\}$ について，少なくとも1つの真の帰無仮説を誤って棄却する確率を FWER（familywise error rate）という．
どの帰無仮説が真であっても FWER $\le\alpha$ となる制御を\textbf{強い制御}，すべての帰無仮説が真のときだけ保証されるものを\textbf{弱い制御}という．
\end{teigi}
各検定を水準 $\alpha$ で行うと，独立な $m$ 個の検定では $\mathrm{FWER}=1-(1-\alpha)^m$．
$m=2$，$\alpha=0.05$ で $0.0975$（\kakomon{2025}{4}），$m=3$，$\alpha=0.025$ で $1-0.975^3=0.073$（\kakomon{2024}{4}）．

\subsection{単一段階の方法}
\begin{itemize}
\item \textbf{Bonferroni}：各検定を $\alpha/m$ で行う．Boole の不等式 $\Prob(\cup A_i)\le\sum\Prob(A_i)$ により，依存関係によらず強い制御．
\item \textbf{Šidák}：独立なら $1-(1-\alpha_1)(1-\alpha_2)=\alpha$，等しく配分して $\alpha_k=1-(1-\alpha)^{1/m}$（$m=2$，$\alpha=0.05$ で $0.0253$）．Bonferroni よりわずかに大きい．
\end{itemize}
\begin{meidai}[正の相関と FWER（\kakomon{2025}{4}〔2〕）]\label[meidai]{prop:14-rho}
$(Z_1,Z_2)$ が相関 $\rho\ge0$ の標準2変量正規分布に従い，各検定を $Z_k>z_\alpha$ で棄却するとき，
$\mathrm{FWER}(\rho)=2\alpha-\phi(\rho)$，$\phi(\rho)=\Prob(Z_1>z_\alpha,Z_2>z_\alpha)$ であり，$\mathrm{FWER}(\rho)\le\mathrm{FWER}(0)=2\alpha-\alpha^2$．
\end{meidai}
\begin{proof}
和事象の確率から $\Prob(A_1\cup A_2)=\alpha+\alpha-\phi(\rho)$．$\rho=0$ なら独立で $\phi(0)=\alpha^2$．$\rho\ge0$ では $\phi(\rho)\ge\alpha^2$（Slepian の不等式．試験では証明なしに使用可）なので $\mathrm{FWER}(\rho)\le\mathrm{FWER}(0)$．
\end{proof}
正の相関があると，独立を仮定した調整（Šidák）や Bonferroni は保守的になる．相関を利用すれば限界値を緩められる（Dunnett 法）．

\subsection{段階的方法}
\begin{itemize}
\item \textbf{Holm}（ステップダウン）：$p_{(1)}\le\dots\le p_{(m)}$ について，$p_{(k)}\le\alpha/(m-k+1)$ が $k=1,2,\dots$ と成り立つ限り棄却し，初めて成り立たなかったところで止める．強い制御で Bonferroni より常に強力．
\item \textbf{Simes 検定}（大域帰無仮説 $\cap H_k$ の検定）：ある $k$ で $p_{(k)}\le k\alpha/m$ なら棄却．大域帰無仮説の下で，独立なら棄却確率は $\alpha$ ちょうど，正の依存（PRDS）の下では $\alpha$ 以下．
      $m=2$ では「$p_{(1)}\le\alpha/2$ または $p_{(2)}\le\alpha$」で，Bonferroni（$p_{(1)}\le\alpha/2$）との差は正方形 $\alpha/2<p_1,p_2\le\alpha$ の部分（\kakomon{2025}{4}〔1-4〕）．
\item \textbf{Hochberg}（ステップアップ）：$k=m,m-1,\dots$ の順に $p_{(k)}\le\alpha/(m-k+1)$ となる最大の $k$ を探し，$H_{(1)},\dots,H_{(k)}$ を棄却．Simes 検定による閉手順（Hommel 法）の簡便版で，Hommel 法より保守的だが Holm 法より強力．
\end{itemize}

\begin{figure}[htbp]
\centering
\begin{tikzpicture}[scale=5]
\draw[->] (0,0)--(0.42,0) node[right,font=\scriptsize]{$p_1$};
\draw[->] (0,0)--(0,0.42) node[above,font=\scriptsize]{$p_2$};
\fill[bkblue!25] (0,0) rectangle (0.1,0.4);
\fill[bkblue!25] (0,0) rectangle (0.4,0.1);
\fill[bkorange!50] (0.1,0.1) rectangle (0.2,0.2);
\draw (0.1,0)--(0.1,-0.01) node[below,font=\scriptsize]{$\alpha/2$};
\draw (0.2,0)--(0.2,-0.01) node[below,font=\scriptsize]{$\alpha$};
\draw (0,0.1)--(-0.01,0.1) node[left,font=\scriptsize]{$\alpha/2$};
\draw (0,0.2)--(-0.01,0.2) node[left,font=\scriptsize]{$\alpha$};
\node[font=\scriptsize,align=left,anchor=west] at (0.24,0.3) {青：Bonferroni の棄却域\\橙：Simes で追加される棄却域};
\end{tikzpicture}
\caption{2つの $p$ 値に対する大域帰無仮説の棄却域（\kakomon{2025}{4}）}\label{fig:14-simes}
\end{figure}

\subsection{閉検定手順とゲートキーピング}
\begin{teiri}[閉検定手順]\label[teiri]{thm:14-closed}
族のすべての共通部分仮説 $H_I=\cap_{i\in I}H_i$ に水準 $\alpha$ の検定を用意し，
$H_i$ を「$i$ を含むすべての $H_I$ が棄却されたとき」に棄却すると，FWER は強い意味で $\alpha$ 以下に制御される．
\end{teiri}
\begin{proof}
真の帰無仮説の添字集合を $I_0$ とする．真の $H_i$ $(i\in I_0)$ が1つでも棄却されるには $H_{I_0}$（真）が棄却されなければならず，その確率は $\alpha$ 以下．
\end{proof}
Holm 法は各 $H_I$ を Bonferroni で検定する閉手順である．Simes 検定による閉手順は Hommel 法で，Hochberg 法はその保守的な簡便版である．
\paragraph{固定順序（階層的）検定}
あらかじめ順序を決め，前の仮説が棄却されたときだけ次を水準 $\alpha$ で検定する．閉手順の特別な場合で，主要評価項目→副次評価項目のゲートキーピングに使われる．
非劣性→優越性の順の検定もこれにあたる．
\paragraph{分散分析をゲートにした対比較（\kakomon{2024}{4}〔3〕）}
一元配置分散分析で全体の帰無仮説が棄却されたときだけ各対比較を水準 $\alpha$ で行う方法は，完全帰無仮説の下では FWER を制御する（弱い制御）．
しかし「$\mu_0=\mu_1=\mu_2<\mu_3$」のように真の帰無仮説が2つ以上ある部分的帰無仮説では，全体の検定は $\mu_3$ のためにほぼ確実に棄却され，
その後2つの真の帰無仮説をそれぞれ水準 $\alpha$ で検定することになり FWER は $\alpha$ を超えうる．
FWER が保証されるのは，すべての帰無仮説が真か，真の帰無仮説が高々1つの場合である．

\subsection{多対1比較と Dunnett 法}
\begin{meidai}[多対1比較の相関（\kakomon{2024}{4}〔4〕）]\label[meidai]{prop:14-dunnett}
$Z_j=(\bar Y_j-\bar Y_0)/\sqrt{\sigma^2(1/n_j+1/n_0)}$ について，すべての帰無仮説が真のとき
\[
\Corr(Z_j,Z_k)=\sqrt{\frac{n_jn_k}{(n_j+n_0)(n_k+n_0)}}\quad(j\ne k),\qquad n_j=n_k=n_0\ \text{なら}\ \frac12.
\]
\end{meidai}
\begin{proof}
$\Cov(\bar Y_j-\bar Y_0,\bar Y_k-\bar Y_0)=\V[\bar Y_0]=\sigma^2/n_0$（共通の対照群のみが共有される）．
各分散で割って $\frac{1/n_0}{\sqrt{(1/n_j+1/n_0)(1/n_k+1/n_0)}}=\sqrt{\frac{n_jn_k}{(n_j+n_0)(n_k+n_0)}}$．
\end{proof}
Dunnett 法はこの相関をもつ多変量正規（$t$）分布の最大値の分布から共通の限界値を求め，Bonferroni より緩い限界値で強い制御を行う．

\subsection{複数の主要評価項目}
\begin{supbox}[補足（出典：FDA Multiple Endpoints ガイダンス \cite{fda:multiple}）]
「複数のうちどれか1つで有効」を主張するなら多重性の調整が必要である．
「すべての評価項目で有効（co-primary）」を要件とするなら，IUT（\cref{thm:13-iut}）により各検定を水準 $\alpha$ で行ってよいが，
全体の検出力は各検出力の積程度に下がる（正の相関があれば少し改善）．
\end{supbox}

\section{用量反応と対比}\label{sec:14-contrast}
各群の平均 $\bar Y_j\sim\Normal(\mu_j,\sigma^2/n_j)$（$j=0,\dots,J$，独立）について，$\sum c_j=0$ を満たす対比係数 $\bm c$ による
\begin{equation}
Z_c=\frac{\sum_jc_j\bar Y_j}{\sqrt{\sigma^2\sum_jc_j^2/n_j}}\label{eq:14-contrast}
\end{equation}
は $H_0:\sum c_j\mu_j=0$ の検定統計量である．$\sum c_j=0$ なので全群の平均が等しければ $H_0$ が成り立つ．
\begin{itemize}
\item \textbf{直線的な用量反応}（等間隔の用量，各群同数）：$\bm c\propto(-3,-1,1,3)$（中心化した用量）．
      正規化すると $(-3,-1,1,3)/\sqrt{20}$（\kakomon{2024}{4}〔5〕）．真の $\mu_j$ が直線なら，この対比が最も検出力が高い．
\item \textbf{プラセボ対全用量の平均}：$(-3,1,1,1)$．
\item 2 値データでの直線対比の検定が Cochran--Armitage 傾向性検定（\cref{sec:04-trend}）である．
\end{itemize}
\begin{meidai}[対比検定の相関（\kakomon{2024}{4}〔6〕）]\label[meidai]{prop:14-contrastcorr}
$H_0$ の下で $\Corr(Z_c,Z_d)=\dfrac{\sum c_jd_j/n_j}{\sqrt{\sum c_j^2/n_j}\sqrt{\sum d_j^2/n_j}}$．各群同数なら $\bm c$ と $\bm d$ のなす角の余弦．
\end{meidai}
$(-3,-1,1,3)$ と $(-3,1,1,1)$ では $\frac{9-1+1+3}{\sqrt{20}\sqrt{12}}=\frac{12}{\sqrt{240}}=\frac{3}{\sqrt{15}}\approx0.775$．
複数の対比（用量反応の形の候補）を検定するときは，この相関を用いた多重性調整を行う（MCP-Mod などの方法）．

\section{仮定}\label{sec:14-assump}
\begin{itemize}
\item 群逐次：独立増分構造（ログランク・スコア統計量で近似的に成立），解析時期が結果に依存しない．
\item 二段階デザイン：患者の独立性，二項モデル．
\item 多重性の方法：Bonferroni・Holm は依存関係によらず有効．Simes・Hochberg は独立または正の依存を要する．Dunnett は正規性・等分散．
\end{itemize}

\section{結果の解釈}\label{sec:14-interp}
\begin{itemize}
\item 中間解析で早期に有効中止した試験は，効果を過大推定する傾向がある（早く境界を越えたのは偶然大きく出たため）．
\item 多重性の調整で有意にならなかった副次項目の名目上の $p<0.05$ は，探索的な結果として扱う．
\item 用量反応の対比が有意でも，どの用量が有効かは多対1比較や用量反応モデルで別に評価する．
\end{itemize}

\section{手法選択}\label{sec:14-choice}
\begin{itemize}
\item 中間解析：$\alpha$ 消費関数（通常は O'Brien--Fleming 型）．
\item 第II相単群：Simon の二段階デザイン．
\item 多用量対プラセボ：Dunnett 法，閉手順（固定順序・Holm），用量反応の対比．
\item 複数の主要評価項目：どれか1つで主張 → Bonferroni/Holm/Hochberg，すべてで主張 → 調整不要（IUT）．
\end{itemize}

\section{よくある誤り}\label{sec:14-err}
\begin{warnbox}
\begin{itemize}
\item 中間解析ごとに名目の $\alpha$ で検定する．
\item $\alpha$ 消費関数で，第 $k$ 回に消費する $\alpha$ を $\alpha(t_k)$ とする（正しくは増分 $\alpha(t_k)-\alpha(t_{k-1})$）．
\item $\Corr[Z(t_1),Z(t_2)]$ を $t_1/t_2$ とする（正しくは $\sqrt{t_1/t_2}$）．
\item 二段階デザインの $\beta$ で第1段階の中止を数え忘れる．
\item 分散分析をゲートにすれば常に FWER が制御されると考える．
\end{itemize}
\end{warnbox}

\section{数理統計との接続}\label{sec:14-link}
\begin{linkbox}
\begin{itemize}
\item $B(t)=\sqrt tZ(t)$ の独立増分構造は，独立な観測の部分和 $S_k=\sum_{i\le k}X_i$ の共分散 $\Cov(S_k,S_l)=\min(k,l)\sigma^2$ と同じである．
\item $p$ 値が帰無仮説の下で $\Unif(0,1)$ に従うこと（『現代数理統計学の基礎』定理7.22）が，Bonferroni・Šidák・Simes の FWER の計算の基礎である．
\item 2変量正規分布の上側同時確率が $\rho$ について増加すること（Slepian の不等式）が\cref{prop:14-rho}の背景にある．
\item \kakomonSM{2025}{4} は二項分布の確率化検定を扱い，離散分布では有意水準をちょうど $\alpha$ にできない（二段階デザインの $\alpha$ が $0.05$ より小さい）ことの背景となる．
\end{itemize}
\end{linkbox}

\section{例題}\label{sec:14-ex}
\begin{reidai}[$\alpha$ 消費関数]\label{reidai:14-1}
片側 $\alpha=0.025$，$\alpha(t)=0.025t^2$ で，情報時間 $t=0.5$ で中間解析，$t=1$ で最終解析を行う．
\begin{enumerate}[label=(\arabic*)]
\item 中間解析と最終解析で消費する $\alpha$ を求め，中間解析の棄却限界値を求めよ（$z_{0.00625}=2.50$）．
\item $\Corr[Z(0.5),Z(1)]$ を求め，最終解析の限界値 $c_2$ を定める式を書け（数値計算では $c_2\approx2.02$）．
\item 両解析を名目の $1.96$ で行った場合の全体の第1種の過誤確率は約 $0.042$ である．このことの意味を述べよ．
\end{enumerate}
\end{reidai}

\begin{reidai}[二段階デザイン]\label{reidai:14-2}
$\pi_0=0.2$，$\pi_1=0.5$ で，$n_1=6$，$r_1=1$，$n_2=6$，$r=4$ のデザインを考える（総有効数 $\ge5$ で有効）．
$\Bin(6,\pi)$ の確率関数 $b(x)$ と上側確率 $\Prob(X\ge x)$ は次のとおり．
\begin{center}\small
\begin{tabular}{@{}ccccc@{}}
\toprule
$x$ & $b(x;0.2,6)$ & $\Prob(X\ge x;0.2)$ & $b(x;0.5,6)$ & $\Prob(X\ge x;0.5)$\\
\midrule
0 & 0.2621 & 1 & 0.0156 & 1\\
1 & 0.3932 & 0.7379 & 0.0938 & 0.9844\\
2 & 0.2458 & 0.3446 & 0.2344 & 0.8906\\
3 & 0.0819 & 0.0989 & 0.3125 & 0.6562\\
4 & 0.0154 & 0.0170 & 0.2344 & 0.3438\\
5 & 0.0015 & 0.0016 & 0.0938 & 0.1094\\
6 & 0.0001 & 0.0001 & 0.0156 & 0.0156\\
\bottomrule
\end{tabular}
\end{center}
\begin{enumerate}[label=(\arabic*)]
\item 早期中止確率 PET と $\E[N\mid\pi_0]$ を求めよ．
\item $\alpha$ と検出力 $1-\beta$ を求めよ．
\item $\alpha\le0.05$ を満たすか．満たさない場合の改善策を述べよ．
\end{enumerate}
\end{reidai}

\begin{reidai}[多重性]\label{reidai:14-3}
3つの独立な仮説の $p$ 値が $0.010,\ 0.020,\ 0.030$ であった．$\alpha=0.05$ とする．
\begin{enumerate}[label=(\arabic*)]
\item 各検定を $0.05$ で行ったときの完全帰無仮説の下の FWER を求めよ．
\item Bonferroni，Holm，Hochberg の各方法で棄却される仮説を答えよ．
\end{enumerate}
\end{reidai}

\begin{reidai}[対比]\label{reidai:14-4}
プラセボと3用量（等間隔）の各群20人で，平均が $10,12,13,15$，$\sigma^2=16$（既知）であった．
\begin{enumerate}[label=(\arabic*)]
\item 直線対比 $(-3,-1,1,3)$ の検定統計量を求めよ．
\item 各用量対プラセボの $Z_j$ を求めよ．$Z_1$ と $Z_2$ の相関はいくらか．
\end{enumerate}
\end{reidai}

\section{解答}\label{sec:14-sol}
\begin{kaitou}{reidai:14-1}
(1) $\alpha(0.5)=0.025\times0.25=0.00625$，最終で $0.025-0.00625=0.01875$．$c_1=2.50$．\\
(2) $\sqrt{0.5/1}=0.707$．$\Prob(Z(0.5)\le2.50,\ Z(1)>c_2)=0.01875$，すなわち
$\int_{-\infty}^{2.50}\int_{c_2}^{\infty}\phi_2(x_1,x_2;0.707)\dd x_2\dd x_1=0.01875$．\\
(3) 名目の限界値で2回検定すると，第1種の過誤確率が0.025から0.042へ約1.7倍に増える．中間解析では限界値を調整しなければならない．
\end{kaitou}

\begin{kaitou}{reidai:14-2}
(1) $\mathrm{PET}=b(0)+b(1)=0.2621+0.3932=0.6553$．$\E[N]=6+(1-0.6553)\times6=8.07$．\\
(2) 有効の結論には $X_1=x\ge2$ かつ $X_2\ge5-x$．
$\alpha=0.2458\times0.0989+0.0819\times0.3446+0.0154\times0.7379+0.0015+0.0001=0.0243+0.0282+0.0114+0.0016=0.065$．
$1-\beta=0.2344\times0.6562+0.3125\times0.8906+0.2344\times0.9844+0.0938+0.0156=0.1538+0.2783+0.2307+0.1094=0.772$．\\
(3) $\alpha=0.065>0.05$ で満たさない．$r$ を5に上げる（総有効数 $\ge6$）と $\alpha$ は下がるが検出力も下がるので，
$n_1,n_2$ を増やしたデザインを探索する必要がある（Simon のデザインはこの探索を系統的に行う）．
\end{kaitou}

\begin{kaitou}{reidai:14-3}
(1) $1-0.95^3=0.143$．\\
(2) Bonferroni（$0.0167$）：$0.010$ のみ棄却．
Holm：$0.010\le0.0167$ 棄却，$0.020\le0.025$ 棄却，$0.030\le0.05$ 棄却で3つとも棄却．
Hochberg：最大の $p=0.030\le0.05$ なので3つとも棄却．
\end{kaitou}

\begin{kaitou}{reidai:14-4}
(1) $\sum c_j\bar Y_j=-30-12+13+45=16$，$\sqrt{16\times20/20}=4$，$Z_c=4.0$．\\
(2) $\SE=\sqrt{16\times2/20}=1.265$．$Z_1=2/1.265=1.58$，$Z_2=3/1.265=2.37$，$Z_3=5/1.265=3.95$．
相関は $\sqrt{20\cdot20/(40\cdot40)}=0.5$．
\end{kaitou}

\begin{summarybox}
\begin{itemize}
\item 情報時間 $t$，$B(t)=\sqrt tZ(t)$，$\Cov[B(t_1),B(t_2)]=\min(t_1,t_2)$，$\Corr[Z(t_1),Z(t_2)]=\sqrt{t_1/t_2}$．
\item $\alpha$ 消費：第1回 $c_1=z_{\alpha(t_1)}$，第 $k$ 回は増分 $\alpha(t_k)-\alpha(t_{k-1})$ を多変量正規で．
\item 二段階：$\alpha=\sum_{x>r_1}b(x;\pi_0,n_1)\{1-B(r-x;\pi_0,n_2)\}$，PET $=B(r_1;\pi_0,n_1)$，$\E[N]=n_1+(1-\mathrm{PET})n_2$．
\item FWER：独立なら $1-(1-\alpha)^m$．Bonferroni $\alpha/m$，Šidák $1-(1-\alpha)^{1/m}$，Holm（ステップダウン），Simes/Hochberg．正の相関で FWER は下がる．
\item 閉検定手順で強い制御．分散分析ゲートは弱い制御．
\item 多対1の相関 $\sqrt{n_jn_k/\{(n_j+n_0)(n_k+n_0)\}}$（同数なら1/2）．対比の相関は係数ベクトルの余弦．直線対比 $(-3,-1,1,3)$．
\end{itemize}
\end{summarybox}

% ===== ch15.tex =====
\chapter{メタアナリシス}\label{ch:15}

\begin{goalbox}
\begin{itemize}
\item 固定効果モデルと変量効果モデルを式で書き，推定対象の違いを説明できる．
\item 逆分散重み付き推定量を最尤法として導出し，その期待値・分散を求められる．
\item Cochran の $Q$ 統計量の期待値を導出し，DerSimonian--Laird のモーメント推定量 $\hat\tau^2$ と $I^2$ を求められる．
\item 変量効果モデルでの $\mu$ の推定量と信頼区間を計算し，フォレストプロットを読める．
\end{itemize}
\end{goalbox}

\begin{exambox}
\begin{itemize}
\item \kakomon{2021}{4}：脂質異常症治療薬の4試験の対数リスク比．変量効果モデル $Y_k=\mu+\delta_k+\varepsilon_k$ の周辺分布 $\Normal(\mu,\sigma_k^2+\tau^2)$，
      対数尤度，$\hat\mu=\sum\frac{y_k}{\sigma_k^2+\hat\tau^2}\big/\sum\frac{1}{\sigma_k^2+\hat\tau^2}$，$\hat\tau^2=0.013$ での数値（$\hat\mu\approx0.008$），$\tau^2=0.014$ を既知とした $\V[\hat\mu]$ と95\%CI．
\item \kakomon{2023}{3}：高脂血症治療薬の筋肉関連有害事象の4試験の対数オッズ比．固定効果の対数尤度と MLE，その期待値・分散，
      $Q$ の分解と $\E[Q]$，$\tau^2$ のモーメント推定量，$\sigma_k^2=\sigma^2$ の下での $\tau^2/(\tau^2+\sigma^2)$ の推定量 $(Q-K+1)/Q$，$Q=7.15$ で $0.58$．
\end{itemize}
\end{exambox}

\flowline{試験ごとの\\推定値と分散}{共通効果 $\theta$\\平均効果 $\mu$}{逆分散重み\\DL 法}{正規近似\\$\sigma_k^2$ 既知}{$Q$，$\hat\tau^2$，$I^2$}{異質性と\\一般化可能性}

\section{問題設定}\label{sec:15-setting}
$K$ 個の研究 $k=1,\dots,K$ から，効果の推定値 $Y_k$（対数 RR，対数 OR，対数 HR，平均差など）とその分散 $\sigma_k^2$ が得られる．
$\sigma_k^2$ は各研究の標本サイズから十分精度よく推定されているとして既知とみなし，重みを $w_k=1/\sigma_k^2$ とする．
比の指標は対数スケールで統合し，結果を指数変換して示す．

\section{固定効果モデル}\label{sec:15-fixed}
\Hissu
\begin{teigi}[固定効果（共通効果）モデル]\label{def:15-fe}
$Y_k\sim\Normal(\theta,\sigma_k^2)$，$k=1,\dots,K$ 独立．すべての研究が共通の真の効果 $\theta$ をもつ．
\end{teigi}
\begin{meidai}[\kakomon{2023}{3}〔1〕〔2〕]\label[meidai]{prop:15-fe}
対数尤度 $l(\theta)=-\frac12\sum_k\{\log(2\pi\sigma_k^2)+w_k(y_k-\theta)^2\}$ を最大化して
\begin{equation}
\hat\theta=\frac{\sum_kw_kY_k}{\sum_kw_k},\qquad \E[\hat\theta]=\theta,\qquad \V[\hat\theta]=\frac{1}{\sum_kw_k}.\label{eq:15-fe}
\end{equation}
\end{meidai}
\begin{proof}
$\partial l/\partial\theta=\sum w_k(y_k-\theta)=0$ より $\hat\theta$．
$\V[\hat\theta]=\sum w_k^2\sigma_k^2/(\sum w_k)^2=\sum w_k/(\sum w_k)^2$．
\end{proof}
逆分散重み付き平均は，$\theta$ の線形不偏推定量の中で分散最小である．精度の高い（大きな）研究ほど重みが大きい．

\section{異質性と \texorpdfstring{$Q$}{Q} 統計量}\label{sec:15-q}
\Hissu
\begin{teigi}[Cochran の $Q$]\label{def:15-q}
$Q=\sum_kw_k(Y_k-\hat\theta)^2$，$\hat\theta$ は固定効果推定量．
\end{teigi}
固定効果モデルの下で $Q\sim\chi^2_{K-1}$ であり，均一性（$\theta_1=\dots=\theta_K$）の検定に用いる．
\begin{teiri}[$Q$ の期待値（\kakomon{2023}{3}〔3〕）]\label[teiri]{thm:15-EQ}
任意の $\theta$ について $Q=\sum w_k(Y_k-\theta)^2-\sum w_k(\hat\theta-\theta)^2$ であり，$Y_k$ が独立で $\E[Y_k]=\theta$（変量効果モデルでは $\theta$ を $\mu$ と読む）のとき
\begin{equation}
\E[Q]=\sum_kw_k\V[Y_k]-\frac{\sum_kw_k^2\V[Y_k]}{\sum_kw_k}.\label{eq:15-EQ}
\end{equation}
\end{teiri}
\begin{proof}
$\sum w_k(Y_k-\theta)^2=\sum w_k\{(Y_k-\hat\theta)+(\hat\theta-\theta)\}^2=Q+2(\hat\theta-\theta)\sum w_k(Y_k-\hat\theta)+(\hat\theta-\theta)^2\sum w_k$ で，
中央の項は $\sum w_k(Y_k-\hat\theta)=0$ より消える．期待値をとると
$\E[\sum w_k(Y_k-\theta)^2]=\sum w_k\V[Y_k]$，$\E[(\hat\theta-\theta)^2]=\V[\hat\theta]=\sum w_k^2\V[Y_k]/(\sum w_k)^2$．
\end{proof}
固定効果モデルなら $\V[Y_k]=1/w_k$ で $\E[Q]=K-1$ である．

\section{変量効果モデル}\label{sec:15-random}
\Hinshutsu
\begin{teigi}[変量効果モデル]\label{def:15-re}
$Y_k=\theta_k+\varepsilon_k$，$\theta_k=\mu+\delta_k$，$\varepsilon_k\sim\Normal(0,\sigma_k^2)$，$\delta_k\sim\Normal(0,\tau^2)$，すべて独立．
$\mu$ は研究間の真の効果の分布の平均，$\tau^2$ は研究間分散．
\end{teigi}
\begin{meidai}[\kakomon{2021}{4}]\label[meidai]{prop:15-re}
(1) $Y_k\sim\Normal(\mu,\sigma_k^2+\tau^2)$．\\
(2) 対数尤度 $\log L(\mu,\tau^2)=-\frac12\sum_k\bigl\{\log2\pi+\log(\sigma_k^2+\tau^2)+\frac{(y_k-\mu)^2}{\sigma_k^2+\tau^2}\bigr\}$．\\
(3) $\tau^2$ を与えたときの MLE と分散は，$w_k^*=1/(\sigma_k^2+\tau^2)$ として
\begin{equation}
\hat\mu=\frac{\sum_kw_k^*Y_k}{\sum_kw_k^*},\qquad \V[\hat\mu]=\frac{1}{\sum_kw_k^*}.\label{eq:15-re}
\end{equation}
\end{meidai}
\begin{proof}
(1) 正規分布の和の再生性，$\E[Y_k]=\mu$，$\V[Y_k]=\tau^2+\sigma_k^2$．(3) $\partial\log L/\partial\mu=\sum w_k^*(y_k-\mu)=0$．分散は\cref{prop:15-fe}と同じ計算．
\end{proof}
実際には $\tau^2$ を推定値 $\hat\tau^2$ で置き換える．$\tau^2$ の推定の不確かさは $\V[\hat\mu]$ に含まれないので，研究数が少ないと区間は狭すぎる傾向がある．

\subsection{DerSimonian--Laird 推定量}\label{subsec:15-dl}
\begin{teiri}[モーメント推定量（\kakomon{2023}{3}〔4〕，\cite{dl1986}）]\label[teiri]{thm:15-dl}
変量効果モデルで $\V[Y_k]=\tau^2+1/w_k$ を\cref{eq:15-EQ}に代入すると
\[
\E[Q]=K-1+\tau^2\Bigl(\sum_kw_k-\frac{\sum_kw_k^2}{\sum_kw_k}\Bigr).
\]
$\E[Q]=Q$ とおいて
\begin{equation}
\hat\tau^2_{\mathrm{DL}}=\max\left\{0,\ \frac{Q-(K-1)}{\sum_kw_k-\sum_kw_k^2/\sum_kw_k}\right\}.\label{eq:15-dl}
\end{equation}
\end{teiri}
\begin{proof}
$\sum w_k(\tau^2+1/w_k)=\tau^2\sum w_k+K$，$\sum w_k^2(\tau^2+1/w_k)/\sum w_k=\tau^2\sum w_k^2/\sum w_k+1$ を\cref{eq:15-EQ}に代入する．
負になる場合は0に切り詰める．
\end{proof}
DL 推定量は反復不要で広く使われるが，ML や REML の推定値とは一般に異なる
（\kakomon{2021}{4} のデータでは $\hat\tau^2_{\mathrm{ML}}=0.013$ に対し $\hat\tau^2_{\mathrm{DL}}\approx0.032$）．

\subsection{\texorpdfstring{$I^2$}{I²}}\label{subsec:15-i2}
全研究で $\sigma_k^2=\sigma^2$（$w_k=w$）なら $\sum w-\sum w^2/\sum w=w(K-1)$ で $\hat\tau^2=\sigma^2\{Q-(K-1)\}/(K-1)$，これより
\begin{equation}
\widehat{\Bigl(\frac{\tau^2}{\tau^2+\sigma^2}\Bigr)}=\frac{Q-(K-1)}{Q}=:I^2\label{eq:15-i2}
\end{equation}
（\kakomon{2023}{3}〔5〕）．$I^2$ は「観測された推定値のばらつきのうち，偶然（研究内の誤差）ではなく研究間の真の差による部分の割合」と解釈される．ただしこれは一般には直感的な解釈であり，上のように研究内分散が等しい場合に限って $\tau^2/(\tau^2+\sigma^2)$ という分散比として導出できる（研究内分散が異なる場合は，$\sigma^2$ を「典型的な」研究内分散に置き換えた近似的な意味をもつ）．$I^2$ は
負なら0とする．目安として25\%，50\%，75\% がそれぞれ低・中・高の異質性とされる \cite{higgins2003}（$I^2$ の定義は \cite{higgins2002}）．
$Q=7.15$，$K=4$ なら $I^2=(7.15-3)/7.15=0.58$．

\section{推定対象の違い}\label{sec:15-target}
\begin{itemize}
\item 固定効果：「全研究に共通の効果 $\theta$」．異質性があるときは「各研究の効果の，精度による重み付き平均」を推定していると解釈できる．
\item 変量効果：「研究間で変動する真の効果の分布の平均 $\mu$」．研究を母集団からの標本とみなす．
\item 変量効果の重み $1/(\sigma_k^2+\tau^2)$ は，$\tau^2$ が大きいほど研究間で均等に近づき，小さな研究の相対的な重みが大きくなる．
\item 変量効果の CI は固定効果より広い（研究間変動を反映）．$\tau^2=0$ なら両者は一致する．
\end{itemize}
\begin{supbox}[補足：予測区間（出典：\cite{higgins2009}）]
新しい研究（あるいは新しい臨床現場）での真の効果の範囲を示すには，予測区間 $\hat\mu\pm t_{K-2}(0.025)\sqrt{\hat\tau^2+\V[\hat\mu]}$ を用いる．
$\mu$ の信頼区間が0を含まなくても，予測区間が0を含むことはよくある．
\end{supbox}

\section{フォレストプロット}\label{sec:15-forest}
各研究の点推定値と95\%CI を横棒で，重みを四角形の大きさで示し，最下段に統合推定値をひし形で示す図をフォレストプロットという．
横軸は通常，比の指標の対数目盛で，縦線が「効果なし」（RR $=1$）を表す．
\begin{figure}[htbp]
\centering
\begin{tikzpicture}[x=40mm,y=6mm]
\draw[gray] (0,0.3)--(0,-6.6);
\foreach \x/\l in {-1/-1.0,-0.5/-0.5,0/0,0.5/0.5}{\draw (\x,-6.6)--(\x,-6.75) node[below,font=\scriptsize]{\l};}
\node[font=\scriptsize] at (-0.25,-7.6) {対数オッズ比};
% studies: Y, se
\foreach \y/\m/\s/\w/\name in {-1/-0.40/0.2000/1.2/研究1,-2/-0.05/0.1414/1.7/研究2,-3/-0.60/0.3000/0.8/研究3,-4/0.15/0.1732/1.4/研究4}{
  \draw[thick] ({\m-1.96*\s},\y)--({\m+1.96*\s},\y);
  \fill[bkblue] ({\m-0.02*\w},{\y-0.133*\w})rectangle({\m+0.02*\w},{\y+0.133*\w});
  \node[anchor=east,font=\scriptsize] at (-1.25,\y) {\name};}
\fill[bkorange] (-0.298,-5)--(-0.1186,-4.7)--(0.0607,-5)--(-0.1186,-5.3)--cycle;
\node[anchor=east,font=\scriptsize] at (-1.25,-5) {固定効果};
\fill[bkred!70] (-0.4614,-6)--(-0.1684,-5.7)--(0.1247,-6)--(-0.1684,-6.3)--cycle;
\node[anchor=east,font=\scriptsize] at (-1.25,-6) {変量効果（DL）};
\end{tikzpicture}
\caption{\cref{reidai:15-1}のフォレストプロット}\label{fig:15-forest}
\end{figure}

\section{仮定}\label{sec:15-assump}
\begin{itemize}
\item 各研究の推定値の正規近似と，$\sigma_k^2$ を既知とみなせること（小さな研究やイベントの稀な研究では不正確）．
\item 研究の独立性．
\item 変量効果モデル：真の効果が正規分布に従い，研究の大きさと無関係であること．
\item 選択された研究が偏っていないこと（出版バイアス）．
\end{itemize}
\begin{supbox}[補足：出版バイアス]
有意な結果の研究ほど出版されやすいと，統合推定値は効果を過大評価する．
横軸に効果，縦軸に精度（$1/\SE$）をとったファネルプロットの非対称性で検討する．
\end{supbox}

\section{結果の解釈}\label{sec:15-interp}
\begin{itemize}
\item \kakomon{2021}{4}：$\tau^2=0.014$（既知）とすると $\hat\mu\approx0.007$，$\SE\approx0.105$，95\%CI $[-0.198,0.212]$（対数RR；$\hat\tau^2=0.013$ では $\hat\mu\approx0.008$）．高用量と標準用量で心血管以外の死亡リスクに差があるとは言えない．
\item $Q$ 検定は研究数が少ないと検出力が低い．有意でなくても異質性がないとは言えないので，$I^2$ と $\hat\tau^2$ で程度を示す．
\item 異質性が大きいときは統合値だけでなく，その原因（対象集団，用量，追跡期間）を部分集団解析やメタ回帰で探索する．
\end{itemize}

\section{手法選択}\label{sec:15-choice}
\Youchui $Q$ 検定の結果や $I^2$ の大きさだけで固定効果か変量効果かを機械的に選ばない．まず推定対象を決める：「全研究に共通の効果（または研究の精度で重み付けた平均）」を知りたいなら固定効果，「研究間で変動する真の効果の分布の平均」を知りたいなら変量効果である．そのうえで，以下を目安とする．
\begin{itemize}
\item 研究が同一のプロトコルで行われ異質性が小さい → 固定効果（MH 法，Peto 法を含む）．
\item 異なる集団・条件の研究を統合し平均的効果を知りたい → 変量効果（DL，REML）．
\item 2 値データでイベントが稀 → Peto 法や MH 法（逆分散法はゼロセルに弱い）．
\end{itemize}

\section{よくある誤り}\label{sec:15-err}
\begin{warnbox}
\begin{itemize}
\item 比の指標を元のスケールで平均する．
\item 変量効果の重みに $1/\sigma_k^2$ を使う（正しくは $1/(\sigma_k^2+\tau^2)$）．
\item $Q$ の自由度を $K$ とする（正しくは $K-1$）．$\hat\tau^2<0$ をそのまま使う．
\item $Q$ 検定が非有意なら固定効果でよいと機械的に判断する．
\end{itemize}
\end{warnbox}

\section{数理統計との接続}\label{sec:15-link}
\begin{linkbox}
\begin{itemize}
\item 逆分散重み付き平均が線形不偏推定量の中で最良であることは『現代数理統計学の基礎』第6章演習（問9）で，固定効果推定量の原型である．
\item 変量効果モデルは正規--正規の階層モデル（同書 4.2.4項，6.4.1項の経験ベイズ）であり，各研究の $\theta_k$ の事後平均は $Y_k$ を $\hat\mu$ の方へ縮小した値になる．
\item $\E[Q]$ の計算は平方和の分解 $\sum(X_i-\mu)^2=\sum(X_i-\bar X)^2+n(\bar X-\mu)^2$ の重み付き版で，不偏分散の除数 $n-1$ の導出と同じ構造である．
\end{itemize}
\end{linkbox}

\section{例題}\label{sec:15-ex}
\begin{reidai}[固定効果と変量効果]\label{reidai:15-1}
4研究の対数オッズ比 $Y_k$ と分散 $\sigma_k^2$ が $(-0.40,0.04)$，$(-0.05,0.02)$，$(-0.60,0.09)$，$(0.15,0.03)$ である．
\begin{enumerate}[label=(\arabic*)]
\item 固定効果推定値とその95\%CI を求めよ．
\item $Q$ を求め，均一性を検定せよ（$\chi^2_3(0.05)=7.815$）．$I^2$ を求めよ．
\item DL 法で $\hat\tau^2$ を求めよ．
\item 変量効果推定値とその95\%CI を求め，(1)と比べよ．
\end{enumerate}
\end{reidai}

\begin{reidai}[{$\E[Q]$ の応用}]\label{reidai:15-2}
$K=5$ 研究がすべて同じ分散 $\sigma^2=0.05$ をもつ．$Q=12.0$ のとき，$\hat\tau^2_{\mathrm{DL}}$ と $I^2$ を求めよ．
\end{reidai}

\section{解答}\label{sec:15-sol}
\enlargethispage{2\baselineskip}
\begin{kaitou}{reidai:15-1}
(1) $w=(25,50,11.11,33.33)$，$\sum w=119.44$，$\sum wY=-10-2.5-6.667+5=-14.17$．
$\hat\theta=-14.17/119.44=-0.119$，$\SE=1/\sqrt{119.44}=0.0915$，95\%CI $(-0.298,\,0.061)$（OR で $(0.74,1.06)$）．\\
(2) $Q=\sum w_k(Y_k+0.119)^2=1.980+0.235+2.575+2.405=7.19<7.815$．均一性は棄却されない．$I^2=(7.19-3)/7.19=0.58$．\\
(3) $\sum w^2=625+2500+123.5+1111.1=4359.6$，$\sum w-\sum w^2/\sum w=119.44-36.50=82.95$．$\hat\tau^2=(7.19-3)/82.95=0.0506$．\\
(4) $w^*=1/(\sigma_k^2+0.0506)=(11.04,14.17,7.11,12.41)$，$\sum w^*=44.74$，$\sum w^*Y=-7.53$．
$\hat\mu=-0.168$，$\SE=1/\sqrt{44.74}=0.150$，95\%CI $(-0.461,\,0.125)$．
変量効果では小さい研究（研究3）の相対的な重みが増えて点推定が負の側に移り，区間は研究間変動を反映して広くなった．
$Q$ 検定は有意でないが $I^2=58\%$ で中程度の異質性があり，研究数が少ないため検定の検出力が低いことに注意する．
\end{kaitou}

\begin{kaitou}{reidai:15-2}
$\hat\tau^2=\sigma^2\{Q-(K-1)\}/(K-1)=0.05\times8/4=0.10$．$I^2=(12-4)/12=0.667$（$=0.10/(0.10+0.05)$）．
\end{kaitou}

\begin{summarybox}
\begin{itemize}
\item 固定効果：$\hat\theta=\sum w_kY_k/\sum w_k$，$\V=1/\sum w_k$，$w_k=1/\sigma_k^2$．
\item $Q=\sum w_k(Y_k-\hat\theta)^2\sim\chi^2_{K-1}$（均一性の下）．$\E[Q]=\sum w\V[Y]-\sum w^2\V[Y]/\sum w$．
\item 変量効果：$Y_k\sim\Normal(\mu,\sigma_k^2+\tau^2)$，$\hat\mu=\sum w_k^*Y_k/\sum w_k^*$，$w_k^*=1/(\sigma_k^2+\tau^2)$，$\V=1/\sum w_k^*$．
\item DL：$\hat\tau^2=\max\{0,(Q-K+1)/(\sum w-\sum w^2/\sum w)\}$．$I^2=(Q-K+1)/Q$．
\item 固定効果は共通効果，変量効果は効果の分布の平均．変量効果の CI は広く，小研究の重みが相対的に大きい．
\end{itemize}
\end{summarybox}

% ===== ch16.tex =====
\chapter{欠測データと多重補完}\label{ch:16}

\begin{goalbox}
\begin{itemize}
\item 欠測メカニズム MCAR・MAR・MNAR を定義し，具体例を分類できる．
\item 完全ケース解析・単一補完・尤度に基づく解析・多重補完が妥当となる条件を説明できる．
\item 欠測の有無と群・共変量の関連を検定し，その結果から言えること・言えないことを述べられる．
\item 多重補完の結果を Rubin のルール $T=W+(1+1/M)B$ で統合し，補完間分散の意味を説明できる．
\item 補完モデルが与えられたとき，補完推定量の条件付き分散と補完間分散の期待値を導出できる．
\end{itemize}
\end{goalbox}

\begin{exambox}
\begin{itemize}
\item \kakomon{2024}{3}：改善割合を比較する試験で，欠測（$T$群30人，$C$群10人）と群の独立性の $\chi^2$ 検定（12.5，有意），
      3つの補完データセットでの改善割合の差の分散，Rubin のルールによる $\hat\V[\bar\theta]$，
      観測値を与えた下での $\hat\theta_k$ の条件付き分散と補完間分散 $B$ の条件付き期待値．
\item 統計数理では\kakomonSM{2016}{5}（MCAR の検定：有意でなくても MCAR とは言えない），\kakomonSM{2022}{5}（対応のあるデータで片方が欠測したときの検定）．
\item \kakomon{2019}{5}（共通）：ABO 血液型の遺伝子頻度の推定（EM アルゴリズム）．
\end{itemize}
\end{exambox}

\flowline{欠測を含む\\データ}{完全データでの\\推定対象}{CC・尤度\\MI・IPW}{MCAR/MAR\\補完モデル}{Rubin のルール}{感度分析\\（MNAR）}

\section{問題設定}\label{sec:16-setting}
完全データを $\bm Y=(\bm Y_{\mathrm{obs}},\bm Y_{\mathrm{mis}})$，欠測の指示変数を $\bm R$（観測なら1）とする．
解析の目的は，欠測がなかった場合に推定したかった量（estimand）を，観測されたデータから偏りなく推定することである．
臨床試験では，欠測の多くは中止・脱落によって起こり，脱落は治療の効果や副作用と関連しうる．

\section{欠測メカニズム}\label{sec:16-mech}
\Hissu
\begin{teigi}[欠測メカニズム（Rubin の分類）]\label{def:16-mech}
\begin{itemize}
\item \textbf{MCAR}（missing completely at random）：$\Prob(\bm R\mid\bm Y_{\mathrm{obs}},\bm Y_{\mathrm{mis}})=\Prob(\bm R)$．欠測はデータと無関係．
\item \textbf{MAR}（missing at random）：$\Prob(\bm R\mid\bm Y_{\mathrm{obs}},\bm Y_{\mathrm{mis}})=\Prob(\bm R\mid\bm Y_{\mathrm{obs}})$．欠測は観測された値（共変量・過去の測定値）だけに依存．
\item \textbf{MNAR}（missing not at random）：欠測が観測されていない値そのものに依存．
\end{itemize}
\end{teigi}
例：検体の破損による欠測は MCAR．前回の測定値が悪い患者が来院しなくなるのは（前回値が観測されていれば）MAR．
今回の症状が悪化したこと自体が来院しない理由なら MNAR．

\begin{meidai}[MAR の下での尤度の無視可能性]\label[meidai]{prop:16-ignor}
MAR で，データのパラメータ $\theta$ と欠測メカニズムのパラメータ $\psi$ が別個（distinct）なら，
観測データの尤度は $L(\theta,\psi)=f(\bm R\mid\bm Y_{\mathrm{obs}};\psi)\,f(\bm Y_{\mathrm{obs}};\theta)$ と分解され，
$\theta$ の推測は欠測メカニズムをモデル化せずに $f(\bm Y_{\mathrm{obs}};\theta)$ だけで行える．
\end{meidai}
\begin{proof}
$f(\bm Y_{\mathrm{obs}},\bm R)=\int f(\bm Y_{\mathrm{obs}},\bm Y_{\mathrm{mis}};\theta)f(\bm R\mid\bm Y_{\mathrm{obs}},\bm Y_{\mathrm{mis}};\psi)\,\dd\bm Y_{\mathrm{mis}}$．
MAR なら第2因子は $\bm Y_{\mathrm{mis}}$ によらず積分の外に出せて，残りは $\int f(\bm Y;\theta)\dd\bm Y_{\mathrm{mis}}=f(\bm Y_{\mathrm{obs}};\theta)$．
\end{proof}
反復測定の混合モデル（MMRM，\cref{ch:08}）が MAR の下で妥当なのはこのためである．

\paragraph{データから言えること}
欠測の有無と観測された変数（群，ベースライン値）の関連を調べれば，MCAR を\textbf{否定}することはできる．
\kakomon{2024}{3}〔1〕では欠測が $T$ 群30\%，$C$ 群10\%で，$\chi^2=\frac{200(30\times90-70\times10)^2}{100\cdot100\cdot40\cdot160}=12.5>3.84$ から，欠測は群と独立でない．
しかし関連がないことは MCAR の証明にならず（\kakomonSM{2016}{5}），MAR と MNAR は観測データからは区別できない．

\section{主な解析法}\label{sec:16-methods}
\begin{table}[htbp]
\centering\small
\caption{欠測への対処法と妥当となる条件}\label{tab:16-methods}
\begin{tabular}{@{}lL{48mm}L{52mm}@{}}
\toprule
方法 & 内容 & 妥当性\\
\midrule
完全ケース解析 & 欠測のある個体を除く & MCAR で不偏（精度は低下）．MAR では一般に偏る（欠測が解析モデルの共変量のみに依存する場合の回帰係数は例外）\\
単一補完（平均値，LOCF） & 欠測を1つの値で埋める & 推定対象により偏り（LOCF など），常に分散を過小評価（補完値を確定値として扱う）\\
尤度に基づく方法（MMRM など） & 観測データの尤度を最大化 & MAR で妥当（モデルが正しければ）\\
多重補完（MI） & 予測分布から複数回補完して統合 & MAR で妥当（補完モデルが正しければ）\\
逆確率重み付け（IPW） & 観測確率の逆数で重み付け & MAR で妥当（観測確率のモデルが正しければ）\\
\bottomrule
\end{tabular}
\end{table}
\Youchui LOCF（最終観測値の引き継ぎ）は「脱落後の値が変わらない」という強い仮定であり，保守的とも限らない．

\section{多重補完}\label{sec:16-mi}
\Hissu\Hinshutsu
\subsection{手順}
\begin{enumerate}
\item \textbf{補完}：欠測値を，観測データを条件とした予測分布（パラメータの不確かさを含む事後予測分布）から $M$ 回独立に生成し，$M$ 個の完全データセットを作る．
\item \textbf{解析}：各データセットで通常の解析を行い，推定値 $\hat\theta_k$ とその分散推定値 $\hat\V[\hat\theta_k]$ を得る．
\item \textbf{統合}：Rubin のルールで1つの推定値と分散にまとめる．
\end{enumerate}
\begin{teiri}[Rubin のルール \cite{rubin1987}]\label[teiri]{thm:16-rubin}
\begin{equation}
\bar\theta=\frac1M\sum_{k=1}^M\hat\theta_k,\quad W=\frac1M\sum_{k=1}^M\hat\V[\hat\theta_k],\quad B=\frac{1}{M-1}\sum_{k=1}^M(\hat\theta_k-\bar\theta)^2,\quad
T=W+\Bigl(1+\frac1M\Bigr)B.\label{eq:16-rubin}
\end{equation}
$W$ を\textbf{補完内分散}，$B$ を\textbf{補完間分散}という．
\end{teiri}
\paragraph{各項の意味}
$W$ は「データが完全だったとしても生じる標本変動」，$B$ は「欠測値がわからないことによる追加の不確かさ」である．
$B/M$ の項は，$\bar\theta$ を有限の $M$ 回の補完の平均で計算したことによるシミュレーション誤差である．
推測には $(\theta-\bar\theta)/\sqrt T$ を自由度 $\nu=(M-1)\bigl\{1+W/((1+1/M)B)\bigr\}^2$ の $t$ 分布で近似する．
$\hat\lambda=(1+1/M)B/T$ は欠測による情報の損失割合の目安である．

\kakomon{2024}{3} では，$\hat\theta_k=0.06,0.07,0.08$，各分散 $\hat\V[\hat\theta_k]=\frac{p_T(1-p_T)}{100}+\frac{p_C(1-p_C)}{100}$ が
$0.003732$，$0.003871$，$0.003718$ で，
$W=0.003774$，$B=\frac{0.01^2+0+0.01^2}{2}=0.0001$，$T=0.003774+\frac43\times0.0001=0.003907$．

\subsection{補完間分散の性質}\label{subsec:16-B}
以下の $M_T,M_C$ は欠測人数であり，補完回数 $M$ とは別である（過去問では補完回数を $K$ と書く）．
\begin{meidai}[\kakomon{2024}{3}〔4〕〔5〕]\label[meidai]{prop:16-B}
観測データを与えた下で $\hat\theta_1,\dots,\hat\theta_M$ が独立同分布なら $\E[B\mid\text{観測}]=\V[\hat\theta_k\mid\text{観測}]$．
2群の改善割合の差の例で，$T$ 群の観測 $N_T$ 人中改善 $y_T^o$，欠測 $M_T$ 人の改善数を $Y_{Tk}^m\sim\Bin(M_T,y_T^o/N_T)$（$C$ 群も同様，独立）で補完すると
\[
\V[\hat\theta_k\mid y_T^o,y_C^o]=\frac{M_T\hat p_T(1-\hat p_T)}{(N_T+M_T)^2}+\frac{M_C\hat p_C(1-\hat p_C)}{(N_C+M_C)^2},\quad \hat p_g=\frac{y_g^o}{N_g}.
\]
\end{meidai}
\begin{proof}
$\hat\theta_k=\frac{y_T^o+Y_{Tk}^m}{N_T+M_T}-\frac{y_C^o+Y_{Ck}^m}{N_C+M_C}$ で，条件付きで確率的なのは $Y^m$ だけ．二項分布の分散を代入する．
$B$ は条件付きで独立同分布な $\hat\theta_k$ の不偏分散なので，その期待値は $\V[\hat\theta_k\mid\text{観測}]$．
\end{proof}
\begin{supbox}[発展：適切な補完（proper imputation）]
上の補完は，改善確率を観測割合 $\hat p_g$ に固定しているので，$B$ は「$\hat p_g$ を与えたときの二項のばらつき」しか反映せず，
$\hat p_g$ 自体の推定誤差 $\hat p_g(1-\hat p_g)/N_g$ に由来する不確かさを取りこぼす．その結果 $T$ は真の分散を過小評価する．
Rubin のルールが正しい分散を与えるためには，各補完で最初にパラメータを事後分布（例：$p_g\sim\Be(y_g^o+1,N_g-y_g^o+1)$）から生成し，
次にそのパラメータで欠測値を生成する「適切な」補完が必要である．
\end{supbox}

\section{感度分析}\label{sec:16-sens}
MAR は検証できない仮定なので，MNAR の可能性に対する感度分析を行う．
\begin{itemize}
\item \textbf{デルタ調整}：MAR で補完した値に，脱落者は $\delta$ だけ悪いというずれを加え，$\delta$ を動かして結論の変化を見る．
\item \textbf{ティッピングポイント解析}：結論が覆る $\delta$ の大きさを求め，臨床的にもっともらしいかを議論する．
\item 最悪ケース補完（試験治療群の欠測を無効，対照群を有効など）は極端な下限を与える．
\end{itemize}
\begin{supbox}[補足（出典：ICH E9(R1) \cite{ich:e9r1}，EMA 欠測データガイドライン \cite{ema:missing}）]
欠測の扱いは estimand の定義（中間事象の戦略）と一体で，主解析の仮定と整合する感度分析をあらかじめ計画することが求められる．
欠測を減らすための試験の計画・実施上の工夫（中止後も追跡を続けるなど）が最も重要である．
\end{supbox}

\section{EM アルゴリズム}\label{sec:16-em}
\Hatten 観測データの尤度の最大化が難しいとき，完全データの対数尤度の条件付き期待値（Eステップ）と，その最大化（Mステップ）を交互に繰り返す．
\paragraph{例：ABO 血液型（\kakomon{2019}{5}）}
遺伝子頻度 $r,p,q$（O，A，B）で，表現型の度数 $n_O,n_A,n_B,n_{AB}$ だけが観測され，A型の遺伝子型 AA と AO の内訳 $f_{AA},f_{AO}$ は欠測である．
\begin{itemize}
\item Eステップ：$\E[f_{AA}\mid n_A]=n_A\dfrac{p^2}{p^2+2pr}$，$\E[f_{AO}\mid n_A]=n_A\dfrac{2pr}{p^2+2pr}$（B型も同様）．
\item Mステップ（遺伝子の数え上げ）：$\hat p=\dfrac{2f_{AA}+f_{AO}+n_{AB}}{2N}$，$\hat q=\dfrac{2f_{BB}+f_{BO}+n_{AB}}{2N}$，$\hat r=\dfrac{2n_O+f_{AO}+f_{BO}}{2N}$．
\end{itemize}
Eステップは「観測された A型の人数を，現在のパラメータでの AA と AO の条件付き確率で按分する」操作である．
各反復で観測データの尤度は減少しない．
\begin{supbox}[試験対策上の注意：公式解答の記法と標準的な EM]
\kakomon{2019}{5}〔4〕(ii) の公式解答は，観測されない度数の期待値を $\E[f_{AA}]=Np^2$，$\E[f_{AO}]=n_A-Np^2$（B型も同様）と書き，これを (i) と交互に用いる反復を示している．
一方，理論上の標準的な EM の Eステップは，上に示した観測値 $n_A$ を条件とする期待値 $n_A\,p^2/(p^2+2pr)$ である．
両者は一般には同じ反復にならないので，答案では，設問が求める式（公式解答の形）に沿って答えつつ，標準的な EM としてはどちらを用いるかを区別して理解しておくとよい．
\end{supbox}

\section{仮定}\label{sec:16-assump}
\begin{itemize}
\item MI・尤度法・IPW：MAR（検証不能），モデルの正しさ．補完モデルには解析モデルの変数（群，アウトカム，重要な共変量）をすべて含める．
\item Rubin のルール：補完が適切（パラメータの不確かさを反映）で，完全データの推定量が漸近正規．
\end{itemize}

\section{結果の解釈}\label{sec:16-interp}
\begin{itemize}
\item 群によって欠測割合が異なる（\kakomon{2024}{3}）なら，完全ケース解析は偏る可能性が高い．欠測の理由（有害事象，効果不十分）を群ごとに確認する．
\item $B$ が $W$ に比べて大きいほど欠測の影響が大きい．
\item MAR の下での結果と感度分析の結果が一貫していれば，結論は頑健である．
\end{itemize}

\section{手法選択}\label{sec:16-choice}
\begin{itemize}
\item 連続アウトカムの反復測定で単調な脱落 → MMRM（MAR）．
\item 2 値アウトカム，最終時点のみ → MI（ロジスティック補完），または脱落を無効とする複合戦略（estimand に依存）．
\item MNAR が疑われる → デルタ調整・ティッピングポイントの感度分析．
\end{itemize}

\section{よくある誤り}\label{sec:16-err}
\begin{warnbox}
\begin{itemize}
\item 欠測と群・共変量の関連が有意でないことを MCAR の証明とする．
\item 単一補完したデータを完全データとして分析し，標準誤差を過小評価する．
\item Rubin のルールで $(1+1/M)$ を落とす，$B$ の除数を $M$ とする．
\item 補完モデルからアウトカムや群を除く（関連が0に偏る）．
\end{itemize}
\end{warnbox}

\section{数理統計との接続}\label{sec:16-link}
\begin{linkbox}
\begin{itemize}
\item Rubin のルールは全分散の公式（『現代数理統計学の基礎』命題4.11）
      $\V[\theta\mid\bm Y_{\mathrm{obs}}]=\E[\V[\theta\mid\bm Y]\mid\bm Y_{\mathrm{obs}}]+\V[\E[\theta\mid\bm Y]\mid\bm Y_{\mathrm{obs}}]$ に対応する（第1項が $W$，第2項が $B$）．
\item $\E[B]=\V[\hat\theta_k]$ は，独立同分布の標本の不偏分散の不偏性そのものである．
\item \kakomonSM{2018}{3}（0切断二項分布の反復解法）は，観測されないゼロの個数を補う EM アルゴリズムの更新と一致する．
\end{itemize}
\end{linkbox}

\section{例題}\label{sec:16-ex}
\begin{reidai}[欠測メカニズムの分類]\label{reidai:16-1}
次の欠測を MCAR・MAR・MNAR に分類せよ（必要な仮定があれば述べよ）．
\begin{enumerate}[label=(\arabic*)]
\item 検査機器の故障で，ある日の全患者の測定値が欠測した．
\item 第4週の血圧が高かった患者ほど第8週に来院しなかった（第4週の値は観測済み）．
\item 抑うつ尺度の評価で，症状が最も悪い時期の患者が回答を拒否した．
\end{enumerate}
\end{reidai}

\begin{reidai}[Rubin のルール]\label{reidai:16-2}
$M=5$ 回の補完で，推定値 $1.20,1.35,1.10,1.28,1.22$，分散推定値 $0.040,0.042,0.039,0.041,0.043$ を得た．
統合推定値，$W$，$B$，$T$，標準誤差，欠測による情報損失割合 $\hat\lambda$，自由度 $\nu$ を求めよ（$\sqrt{0.05144}=0.2268$）．
\end{reidai}

\begin{reidai}[完全ケース解析の偏り]\label{reidai:16-3}
2 値の共変量 $X$（各50\%）で $\E[Y\mid X=0]=10$，$\E[Y\mid X=1]=20$．$Y$ の観測確率は $X=0$ で0.9，$X=1$ で0.5（$X$ は全員観測）．
\begin{enumerate}[label=(\arabic*)]
\item 欠測メカニズムは何か．
\item 完全ケースでの $Y$ の平均の期待値を求め，真の平均15と比べよ．
\item 観測確率の逆数で重み付けした平均の期待値を求めよ．
\end{enumerate}
\end{reidai}

\section{解答}\label{sec:16-sol}
\begin{kaitou}{reidai:16-1}
(1) MCAR（故障が患者の状態と無関係なら）．
(2) MAR（欠測が観測済みの第4週の値に依存し，それを与えれば第8週の値によらないなら）．第8週の値自体に依存するなら MNAR．
(3) MNAR（欠測が欠測値そのもの＝その時点の症状に依存）．
\end{kaitou}

\begin{kaitou}{reidai:16-2}
$\bar\theta=1.23$，$W=0.041$，$B=\frac{0.03^2+0.12^2+0.13^2+0.05^2+0.01^2}{4}=\frac{0.0348}{4}=0.0087$，
$T=0.041+1.2\times0.0087=0.05144$，$\SE=0.227$．$\hat\lambda=0.01044/0.05144=0.20$．
$\nu=4\{1+0.041/0.01044\}^2=4\times4.927^2=97$．
\end{kaitou}

\begin{kaitou}{reidai:16-3}
(1) 欠測は観測された $X$ だけに依存するので MAR．\\
(2) 観測者の構成は $X=0$：$0.45$，$X=1$：$0.25$．平均 $\frac{0.45\times10+0.25\times20}{0.70}=13.57$．真の15より小さく偏る．\\
(3) 重み $1/0.9$，$1/0.5$ で $\frac{0.45\times10/0.9+0.25\times20/0.5}{0.45/0.9+0.25/0.5}=\frac{5+10}{1}=15$．偏りが除かれる．
（$X$ ごとの平均を $X$ の分布で平均する，すなわち $X$ を含む補完・回帰でも同じ結果になる．）
\end{kaitou}

\begin{summarybox}
\begin{itemize}
\item MCAR：データと無関係．MAR：観測値のみに依存．MNAR：欠測値自体に依存．MAR と MNAR はデータで区別できない．
\item MAR＋パラメータの分離で尤度は欠測メカニズムを無視できる（MMRM が妥当）．
\item 完全ケースは MCAR で不偏．単一補完は分散を過小評価．
\item Rubin：$\bar\theta$，$W=\overline{\hat\V}$，$B=\frac{1}{M-1}\sum(\hat\theta_k-\bar\theta)^2$，$T=W+(1+1/M)B$．
\item $\E[B\mid\text{観測}]=\V[\hat\theta_k\mid\text{観測}]$．パラメータを固定した補完は不確かさを過小評価（適切な補完が必要）．
\item MNAR への感度分析（デルタ調整・ティッピングポイント）．
\end{itemize}
\end{summarybox}

% ===== ch17.tex =====
\chapter{ポアソン率と医療データのBayes解析}\label{ch:17}

\begin{goalbox}
\begin{itemize}
\item 人時間（観察時間・人口）を考慮したポアソンモデルで，率の最尤推定・信頼区間・2群比較（差・比・条件付き検定）ができる．
\item ガンマ分布のパラメータ化（形状・率／形状・尺度）を明示して扱える．
\item Beta--Binomial，Gamma--Poisson の共役事前分布による事後分布・事後平均・事後モード・事後確率を導出できる．
\item 事前予測分布（周辺分布）と事後予測分布を導出し，それぞれ離散一様分布・負の二項分布になる例を示せる．
\item 事後平均が「データの推定値と事前平均の重み付き平均」（縮小）になることを説明し，頻度論の推定との違いを述べられる．
\end{itemize}
\end{goalbox}

\begin{exambox}
\begin{itemize}
\item \kakomon{2022}{4}：有効例数 $Y\sim\Bin(n,\pi)$，一様事前分布．尤度と MLE，事後分布 $\Be(y+1,n-y+1)$，事後モード $=y/n$，
      事後平均 $(y+1)/(n+2)$ と MLE・事前平均の関係，$Y$ の周辺分布（離散一様），$\Prob(\pi\ge1/2\mid y)\ge0.8$ となる最小の $y$（$n=5$ で $y=4$）．
\item \kakomon{2025}{2}：2市の死亡数（人口2万・8万）．共通の死亡率の MLE（1.8），率の差の正規近似検定，
      ガンマ事前分布 $\Ga(1,1)$（形状・率）からの事後分布 $\Ga(19,11)$，人口5万の市の死亡数の事後予測分布（負の二項分布），
      その期待値 $95/11$ と MLE に基づく予測値9の比較．
\item 統計数理では\kakomonSM{2022}{3}（ガンマ--ポアソン混合＝負の二項，過分散），\kakomonSM{2021}{3}（ポアソン平均の十分性・MLE・スコア型区間），\kakomonSM{2024}{3}（Beta 事前分布の事後平均と同じ形の推定量）．
\end{itemize}
\end{exambox}

\flowline{件数＋人時間\\二項データ}{率・割合\\予測分布}{ポアソン・\\共役 Bayes}{ポアソン性\\事前分布}{事後 $=$ 事前 $\times$ 尤度}{縮小と\\事前の影響}

\section{ポアソン率}\label{sec:17-poisson}
\subsection{モデルと推定}
観察人時間（または人口）$T$ の集団で起きたイベント数を $Y$ とし，率 $\lambda$（単位人時間あたり）で
\[
Y\sim\Po(\lambda T),\qquad \Prob(Y=y)=\frac{e^{-\lambda T}(\lambda T)^y}{y!}
\]
とする．\kakomon{2025}{2} のように「1万人あたり $\Po(\theta)$」が独立なら，再生性により人口 $T$ 万人の死亡数は $\Po(T\theta)$ である．
\begin{meidai}\label[meidai]{prop:17-mle}
$\hat\lambda=Y/T$，$\V[\hat\lambda]=\lambda/T$（推定値 $Y/T^2$），$\V[\log\hat\lambda]\approx1/(\lambda T)\approx1/Y$．
複数の集団で共通の率なら $\hat\lambda=\sum Y_g/\sum T_g$．
\end{meidai}
\begin{proof}
$\log L=y\log\lambda-\lambda T+\text{const}$，$\partial/\partial\lambda=y/\lambda-T=0$．共通の率では $\log L=(\sum y_g)\log\lambda-\lambda\sum T_g$．分散はポアソン分布の性質とデルタ法．
\end{proof}
\kakomon{2025}{2}：$\hat\theta=(4+14)/(2+8)=1.8$（1万人あたり）．
この尤度 $\lambda^ye^{-\lambda T}$ は，打ち切りのある指数モデルの尤度 $\lambda^de^{-\lambda W}$（\cref{prop:10-expmle}）と同じ形である．

\subsection{2群の比較}\label{subsec:17-compare}
\begin{itemize}
\item \textbf{率の差}：$Z=\dfrac{\hat\lambda_A-\hat\lambda_B}{\sqrt{Y_A/T_A^2+Y_B/T_B^2}}$（\kakomon{2025}{2}〔2〕：$Z=0.25/\sqrt{1.219}=0.23$，有意でない）．
\item \textbf{率比}：$\log(\hat\lambda_A/\hat\lambda_B)\pm z_{\alpha/2}\sqrt{1/Y_A+1/Y_B}$ を指数変換．
\item \textbf{条件付き検定}：$Y_A+Y_B=n$ を与えると $Y_A\sim\Bin\bigl(n,\ \frac{\lambda_AT_A}{\lambda_AT_A+\lambda_BT_B}\bigr)$．
      $H_0:\lambda_A=\lambda_B$ の下で成功確率は $T_A/(T_A+T_B)$ となり，二項検定（exact）ができる．
\end{itemize}
\begin{supbox}[補足：過分散とポアソン回帰]
個体・施設ごとに率がばらつくと，件数の分散が平均を超える（過分散）．率がガンマ分布に従う混合モデルでは件数は負の二項分布に従い，
$\V=\mu+\mu^2/a$（\kakomonSM{2022}{3}）．共変量の調整には $\log\E[Y_i]=\log T_i+\bm x_i\transp\bm\beta$ のポアソン回帰（$\log T_i$ はオフセット）を用い，$e^{\beta_j}$ は率比である．
\end{supbox}

\section{ガンマ分布のパラメータ化}\label{sec:17-gamma}
\Youchui 本書ではガンマ分布を\textbf{形状 $a$・率 $b$} で表す．
\[
\Ga(a,b):\quad g(\theta)=\frac{b^a}{\Gamma(a)}\theta^{a-1}e^{-b\theta},\qquad \E=\frac ab,\quad \V=\frac{a}{b^2},\quad \text{モード}=\frac{a-1}{b}\ (a\ge1).
\]
\kakomon{2025}{2} も同じ流儀である．『現代数理統計学の基礎』は $\Ga(\alpha,\beta)$ を\textbf{形状・尺度}（密度 $\propto x^{\alpha-1}e^{-x/\beta}$，平均 $\alpha\beta$）で定義しており，
尺度 $\beta$ は率 $b$ の逆数 $1/b$ に当たる．答案では必ず密度を書いて流儀を明示すること．
指数分布 $\Ex(\lambda)$ は $\Ga(1,\lambda)$（形状1・率 $\lambda$）である．

\section{Bayes 推測の基本}\label{sec:17-bayes}
パラメータ $\theta$ に事前分布 $p(\theta)$ を置き，データ $y$ を観測した後の事後分布を
\[
p(\theta\mid y)=\frac{p(y\mid\theta)p(\theta)}{p(y)}\propto p(y\mid\theta)\,p(\theta),\qquad p(y)=\int p(y\mid\theta)p(\theta)\,\dd\theta
\]
とする．$p(y)$ は\emphx{事前予測分布}（周辺分布）であり，将来のデータ $\tilde y$ の\emphx{事後予測分布}は
\[
p(\tilde y\mid y)=\int p(\tilde y\mid\theta)\,p(\theta\mid y)\,\dd\theta
\]
（$\tilde y$ と $y$ が $\theta$ を与えて独立な場合）．
点推定には事後平均（2乗損失の Bayes 推定量），事後モード，事後中央値を，区間推定には事後分布の $2.5\%$ 点と $97.5\%$ 点による信用区間を用いる．
事後分布が事前分布と同じ分布族になる事前分布を\emphx{共役事前分布}という．

\section{Beta--Binomial モデル}\label{sec:17-betabin}
\Hissu
\begin{teiri}[Beta--Binomial]\label[teiri]{thm:17-betabin}
$Y\mid\pi\sim\Bin(n,\pi)$，$\pi\sim\Be(a,b)$ なら
\begin{equation}
\pi\mid y\sim\Be(a+y,\ b+n-y),\qquad \E[\pi\mid y]=\frac{a+y}{a+b+n}=\frac{n}{a+b+n}\cdot\frac yn+\frac{a+b}{a+b+n}\cdot\frac{a}{a+b}.\label{eq:17-betabin}
\end{equation}
事後モードは $(a+y-1)/(a+b+n-2)$．$Y$ の周辺分布は
$\Prob(Y=y)=\binom ny\dfrac{B(a+y,b+n-y)}{B(a,b)}$（ベータ二項分布）．
\end{teiri}
\begin{proof}
$p(\pi\mid y)\propto\pi^y(1-\pi)^{n-y}\cdot\pi^{a-1}(1-\pi)^{b-1}=\pi^{a+y-1}(1-\pi)^{b+n-y-1}$ はベータ分布の核．
平均は $\Be(\alpha,\beta)$ の平均 $\alpha/(\alpha+\beta)$，モードは $(\alpha-1)/(\alpha+\beta-2)$．
周辺分布は $\int\binom ny\pi^y(1-\pi)^{n-y}\frac{\pi^{a-1}(1-\pi)^{b-1}}{B(a,b)}\dd\pi$．
\end{proof}
事後平均は MLE $y/n$ と事前平均 $a/(a+b)$ の重み付き平均で，重みは「データの大きさ $n$」と「事前分布の大きさ $a+b$（事前の仮想的な標本サイズ）」の比である．
$n\to\infty$ で事前の影響は消える．

\paragraph{一様事前分布（\kakomon{2022}{4}）}
$a=b=1$ なら事後分布 $\Be(y+1,n-y+1)$，事後モード $y/n$（MLE と一致），事後平均 $\frac{y+1}{n+2}=\frac{n}{n+2}\frac yn+\frac{2}{n+2}\cdot\frac12$．
$Y$ の周辺分布は
\[
\Prob(Y=y)=\binom ny B(y+1,n-y+1)=\frac{n!}{y!(n-y)!}\cdot\frac{y!(n-y)!}{(n+1)!}=\frac{1}{n+1}\quad(y=0,\dots,n),
\]
すなわち離散一様分布である．

\paragraph{事後確率による判定}
「$\Prob(\pi\ge\pi_0\mid y)\ge\gamma$ なら有効」とする判定では，一様事前分布の下で次の関係が便利である．
\begin{meidai}\label[meidai]{prop:17-betatail}
$\pi\mid y\sim\Be(y+1,n-y+1)$ のとき $\Prob(\pi>p\mid y)=\Prob(W\le y)$，$W\sim\Bin(n+1,p)$．
\end{meidai}
\begin{proof}
$\Be(k,m-k+1)$ の分布関数と二項分布の関係 $\Prob(\Be(k,m-k+1)\le p)=\Prob(\Bin(m,p)\ge k)$（順序統計量 $U_{(k)}$ の分布：$m$ 個の一様乱数のうち $k$ 個以上が $p$ 以下）を
$k=y+1$，$m=n+1$ で用い，余事象をとる．
\end{proof}
\kakomon{2022}{4} では $n=5$，$p=1/2$ で $\Prob(\pi\ge\frac12\mid y=3)=\Prob(\Bin(6,\tfrac12)\le3)=42/64=0.656$，
$y=4$ で $57/64=0.891$．よって $0.8$ 以上となる最小の $y$ は4．

\paragraph{事後予測}
次の $\tilde n$ 人での有効数 $\tilde Y$ の事後予測分布は $\Be(a+y,b+n-y)$ を事前とするベータ二項分布で，
次の1人が有効である確率は事後平均 $\E[\pi\mid y]$ に等しい．

\section{Gamma--Poisson モデル}\label{sec:17-gampois}
\Hissu
\begin{teiri}[Gamma--Poisson]\label[teiri]{thm:17-gampois}
$Y\mid\theta\sim\Po(T\theta)$，$\theta\sim\Ga(a,b)$（形状・率）なら
\begin{equation}
\theta\mid y\sim\Ga(a+y,\ b+T),\qquad \E[\theta\mid y]=\frac{a+y}{b+T}=\frac{T}{b+T}\cdot\frac yT+\frac{b}{b+T}\cdot\frac ab.\label{eq:17-gampois}
\end{equation}
人時間 $\tilde T$ での将来の件数 $\tilde Y\mid\theta\sim\Po(\tilde T\theta)$ の事後予測分布は，$r=a+y$，$q=\dfrac{b+T}{b+T+\tilde T}$ として負の二項分布
\begin{equation}
\Prob(\tilde Y=k\mid y)=\frac{\Gamma(k+r)}{k!\,\Gamma(r)}q^r(1-q)^k\quad(k=0,1,\dots),\qquad
\E[\tilde Y\mid y]=\tilde T\,\frac{a+y}{b+T}.\label{eq:17-nb}
\end{equation}
\end{teiri}
\begin{proof}
$p(\theta\mid y)\propto e^{-T\theta}\theta^y\cdot\theta^{a-1}e^{-b\theta}=\theta^{a+y-1}e^{-(b+T)\theta}$．
事後予測は
\[
\int_0^\infty\frac{e^{-\tilde T\theta}(\tilde T\theta)^k}{k!}\frac{(b+T)^r}{\Gamma(r)}\theta^{r-1}e^{-(b+T)\theta}\dd\theta
=\frac{\tilde T^k(b+T)^r}{k!\,\Gamma(r)}\cdot\frac{\Gamma(k+r)}{(b+T+\tilde T)^{k+r}},
\]
これを整理すると\cref{eq:17-nb}．期待値は $\E[\tilde Y\mid y]=\E[\E[\tilde Y\mid\theta]\mid y]=\tilde T\E[\theta\mid y]$．
\end{proof}
事後予測分布の分散 $\tilde T\E[\theta\mid y]+\tilde T^2\V[\theta\mid y]$ は平均より大きく，
「$\theta$ の不確かさ」の分だけポアソン分布より広い．

\paragraph{\kakomon{2025}{2} の数値}
$a=b=1$（平均1の指数分布の事前），$y=18$，$T=10$ で $\theta\mid y\sim\Ga(19,11)$，事後平均 $19/11=1.73$．
人口5万の市（$\tilde T=5$）の死亡数の事後予測は $r=19$，$q=11/16$ の負の二項分布で，期待値 $5\times19/11=95/11\approx8.6$．
MLE による予測 $5\times1.8=9$ より小さいのは，事後平均が事前平均1の側へ縮小されているためである（$1.73=\frac{10}{11}\times1.8+\frac{1}{11}\times1$）．

\section{正規--正規モデル}\label{sec:17-normal}
$\bar X\mid\mu\sim\Normal(\mu,\sigma^2/n)$（$\sigma^2$ 既知），$\mu\sim\Normal(\xi,\tau^2)$ なら
\[
\mu\mid\bar x\sim\Normal\Bigl(\frac{(n/\sigma^2)\bar x+(1/\tau^2)\xi}{n/\sigma^2+1/\tau^2},\ \frac{1}{n/\sigma^2+1/\tau^2}\Bigr).
\]
事後平均は精度（分散の逆数）で重み付けた平均である．
平均への回帰の真値モデル（\cref{sec:08-rtm}）やメタアナリシスの変量効果モデル（\cref{ch:15}）で各研究の効果を推定する式と同じ形である．

\section{頻度論との違い}\label{sec:17-vs}
\begin{itemize}
\item 頻度論ではパラメータは定数で，確率は手順の繰り返しに関するもの．Bayes ではパラメータの不確かさを確率分布で表し，「$\Prob(\pi\ge0.5\mid\text{データ})$」を直接述べられる．
\item 信用区間は「データを与えたとき $\theta$ が区間に入る確率が95\%」と解釈できる．信頼区間はそうは解釈できない．
\item 事前分布が結論に影響する．小標本では影響が大きく，事前分布の選択の根拠と感度分析が必要である．
\item Bayes の判定規則（事後確率の閾値）も，頻度論的な性質（第1種の過誤確率・検出力）をシミュレーションで評価するのが臨床試験の実務である．
\end{itemize}
\begin{supbox}[補足：無情報事前分布]
二項比率では一様分布 $\Be(1,1)$ のほか，Jeffreys 事前分布 $\Be(1/2,1/2)$（Fisher 情報量の平方根に比例）がよく用いられる．
ポアソン率の Jeffreys 事前分布は $p(\lambda)\propto\lambda^{-1/2}$ で，全区間で積分が発散する非正則（improper）事前分布である（形式的に「$\Ga(1/2,0)$」と書かれることがあるが，率0のガンマ分布は正則な確率分布ではない）．データ $y$，人時間 $T$ を得た後の事後分布は $\Ga(y+1/2,\,T)$ で，$y\ge0$ なら正則になる．
\end{supbox}

\section{仮定}\label{sec:17-assump}
\begin{itemize}
\item ポアソンモデル：イベントが独立で，率が観察期間・集団内で一定．過分散があれば負の二項など．
\item Bayes：事前分布の妥当性．共役事前分布は計算の便宜であり，それ自体が正しい保証はない．
\item 交換可能性：事前分布を共有する単位（集団・患者）が事前に区別できないこと．
\end{itemize}

\section{結果の解釈}\label{sec:17-interp}
\begin{itemize}
\item 「事後平均 0.30，95\%信用区間 0.15--0.47」：データと事前情報を合わせると，有効率は95\%の確率でこの範囲にある．
\item 事後予測分布の期待値が MLE の予測値と違うのは，事前平均への縮小による．データが多いほど両者は近づく．
\item 率の差の検定が非有意でも，共通の率を仮定してよいという積極的な根拠にはならない．
\end{itemize}

\section{手法選択}\label{sec:17-choice}
\begin{itemize}
\item 観察期間が異なる件数データ → ポアソン（率）．群比較は率比の区間，条件付き二項検定．
\item 小標本の第I/II相で既存情報を使いたい → Beta--Binomial の事後確率による判定．
\item 将来の件数の予測 → 事後予測分布（Gamma--Poisson なら負の二項）．
\end{itemize}

\section{よくある誤り}\label{sec:17-err}
\begin{warnbox}
\begin{itemize}
\item ガンマ分布の形状・尺度と形状・率を混同する（平均 $a/b$ と $ab$ の取り違え）．
\item 事後予測分布を，事後平均を代入したポアソン分布（$\Po(\tilde T\hat\theta)$）とする（$\theta$ の不確かさを無視）．
\item 人時間を無視して件数どうしを比べる．
\item 事前予測分布（周辺分布）と事後予測分布を混同する．
\end{itemize}
\end{warnbox}

\section{数理統計との接続}\label{sec:17-link}
\begin{linkbox}
\begin{itemize}
\item ベータ--二項（『現代数理統計学の基礎』例4.14，例6.12），ガンマ--ポアソン＝負の二項（例4.15），ポアソン--ガンマの事後分布（第6章演習問8，形状・尺度で $\Ga(y+\alpha,\beta/(1+n\beta))$）．
\item 二項の累積確率とベータ分布の関係（同書第3章演習問19）が\cref{prop:17-betatail}，ポアソンの累積確率とガンマ分布の関係（問20）はそのポアソン過程版（二項→ポアソンの極限に対応する関係）である．
\item \kakomonSM{2024}{3} の MSE 一定の推定量 $(S_n+\sqrt n/2)/(n+\sqrt n)$ は $\Be(\sqrt n/2,\sqrt n/2)$ 事前分布の事後平均と同じ形である．
\item \kakomonSM{2021}{3} のポアソン平均の区間 $\frac1n(t+2\mp2\sqrt{t+1})$ は，正規近似 $(t-n\lambda)^2\le4n\lambda$ を解いたスコア型区間である．
\end{itemize}
\end{linkbox}

\section{例題}\label{sec:17-ex}
\begin{reidai}[ポアソン率の比較]\label{reidai:17-1}
病院Aでは3000患者日で院内感染12件，病院Bでは5000患者日で30件であった．
\begin{enumerate}[label=(\arabic*)]
\item 各病院の感染率（1000患者日あたり）を求めよ．
\item 率の差の $Z$ 統計量を求めよ．
\item 率比の95\%信頼区間を求めよ（$\log(2/3)=-0.405$，$e^{-1.075}=0.341$，$e^{0.264}=1.302$）．
\item 条件付き二項検定の正規近似による $Z$ を求めよ．
\end{enumerate}
\end{reidai}

\begin{reidai}[Beta--Binomial]\label{reidai:17-2}
先行研究から有効率の事前分布を $\Be(2,8)$ とし，新しい試験で20人中7人が有効であった．
\begin{enumerate}[label=(\arabic*)]
\item 事前平均，事後分布，事後平均，事後モードを求めよ．
\item 事後平均を MLE と事前平均の重み付き平均として表せ．
\item 次の2人の有効数の事後予測分布を求めよ．
\end{enumerate}
\end{reidai}

\begin{reidai}[事後確率による判定]\label{reidai:17-3}
一様事前分布の下，$n=9$ 人の試験で「$\Prob(\pi>0.5\mid y)\ge0.9$ なら有望」と判定する．
$\Bin(10,0.5)$ の累積確率 $\Prob(W\le y)$ は $y=5,6,7,8$ でそれぞれ $0.623,0.828,0.945,0.989$ である．
\begin{enumerate}[label=(\arabic*)]
\item 有望と判定される最小の有効数を求めよ．
\item 真の有効率が0.5のとき，有望と判定される確率を求めよ（$\Bin(9,0.5)$ で $\Prob(Y\ge7)=0.0898$）．
\end{enumerate}
\end{reidai}

\begin{reidai}[Gamma--Poisson]\label{reidai:17-4}
有害事象の発生率 $\theta$（100人年あたり）の事前分布を $\Ga(3,2)$（形状3・率2）とする．4（百人年）の観察で8件が起きた．
\begin{enumerate}[label=(\arabic*)]
\item 事後分布と事後平均を求め，MLE と比べよ．
\item 次の2（百人年）での件数の事後予測分布の平均と分散を求め，ポアソン分布と比べよ．
\end{enumerate}
\end{reidai}

\section{解答}\label{sec:17-sol}
\begin{kaitou}{reidai:17-1}
(1) A：$12/3=4.0$，B：$30/5=6.0$（1000患者日あたり）．\\
(2) $\hat\lambda_A-\hat\lambda_B=-0.002$/患者日，$\SE=\sqrt{12/3000^2+30/5000^2}=\sqrt{1.333\times10^{-6}+1.2\times10^{-6}}=0.00159$，$Z=-1.26$．\\
(3) $\SE(\log)=\sqrt{1/12+1/30}=0.342$．$-0.405\pm0.670=(-1.075,0.264)$ → $(0.34,1.30)$．\\
(4) $n=42$，$H_0$ で $Y_A\sim\Bin(42,0.375)$，$\E=15.75$，$\V=9.84$，$Z=(12-15.75)/3.14=-1.20$．いずれも有意でない．
\end{kaitou}

\begin{kaitou}{reidai:17-2}
(1) 事前平均 $0.2$．事後 $\Be(9,21)$，平均 $9/30=0.30$，モード $8/28=0.286$．\\
(2) $0.30=\frac{20}{30}\times0.35+\frac{10}{30}\times0.2$（事前の仮想標本10人，データ20人）．\\
(3) $\Prob(\tilde Y=0)=\frac{21}{30}\cdot\frac{22}{31}=\frac{462}{930}=0.497$，
$\Prob(\tilde Y=2)=\frac{9}{30}\cdot\frac{10}{31}=\frac{90}{930}=0.097$，
$\Prob(\tilde Y=1)=2\cdot\frac{9\cdot21}{30\cdot31}=\frac{378}{930}=0.406$．
\end{kaitou}

\begin{kaitou}{reidai:17-3}
(1) \cref{prop:17-betatail}より $\Prob(\pi>0.5\mid y)=\Prob(W\le y)$，$W\sim\Bin(10,0.5)$．$\ge0.9$ となる最小の $y$ は7（$0.945$）．\\
(2) $\Prob(Y\ge7\mid\pi=0.5)=0.0898$．この Bayes 判定の「第1種の過誤確率」は約9\%である．
\end{kaitou}

\begin{kaitou}{reidai:17-4}
(1) $\Ga(3+8,\,2+4)=\Ga(11,6)$，事後平均 $11/6=1.83$．MLE $8/4=2.0$ より事前平均 $1.5$ の側に縮小している
（$1.83=\frac46\times2.0+\frac26\times1.5$）．\\
(2) $r=11$，$q=6/8=0.75$．平均 $r(1-q)/q=11\times0.25/0.75=3.67$（$=2\times11/6$），
分散 $r(1-q)/q^2=4.89$．同じ平均のポアソン分布（分散3.67）より大きい．
\end{kaitou}

\begin{summarybox}
\begin{itemize}
\item $Y\sim\Po(\lambda T)$：$\hat\lambda=Y/T$，$\V\approx Y/T^2$，$\V[\log\hat\lambda]\approx1/Y$．$H_0:\lambda_A=\lambda_B$ の下で条件付きで $Y_A\mid n\sim\Bin(n,T_A/(T_A+T_B))$．
\item ガンマは形状・率 $\Ga(a,b)$：平均 $a/b$（教科書は形状・尺度：平均 $\alpha\beta$）．
\item Beta--Binomial：$\Be(a+y,b+n-y)$，事後平均は MLE と事前平均の重み付き平均．一様事前で周辺は離散一様，$\Prob(\pi>p\mid y)=\Prob(\Bin(n+1,p)\le y)$．
\item Gamma--Poisson：$\Ga(a+y,b+T)$，事後予測は負の二項（$r=a+y$，$q=(b+T)/(b+T+\tilde T)$），平均 $\tilde T(a+y)/(b+T)$．
\item 事後予測は $\theta$ の不確かさを含むので，プラグインのポアソンより分散が大きい．
\end{itemize}
\end{summarybox}

% ===== 90_appendix.tex =====
\appendix

% ============================================================
\chapter{手法選択フローチャート}\label[appendix]{app:A}
% ============================================================
アウトカムの型と観測の構造から手法を選ぶための一覧である．
まず\cref{fig:A-flow}で候補を絞り，\cref{tab:A-select}で推定対象と対応章を確認する．

\begin{figure}[!htbp]
\centering
\begin{tikzpicture}[
  typ/.style={draw=bkblue,fill=bkblue!8,rounded corners=2pt,font=\small\bfseries,text width=17mm,align=center,minimum height=10mm},
  lst/.style={draw=bkgray!60,fill=white,rounded corners=2pt,font=\small,text width=132mm,align=left,inner sep=4pt},
  hd/.style={font=\small\bfseries,text=bkblue},
  arr/.style={-{Stealth[length=2.2mm]},bkblue,thick}]
\node[lst,anchor=north west] (cl) at (0,0) {%
\textbf{独立2群}：正規・等分散 → 2標本 $t$，不等分散 → Welch，非正規・外れ値 → 順位和（Mann--Whitney）\\
\textbf{対応あり}：正規 → 対応のある $t$，対称 → 符号付き順位，形の仮定なし → 符号検定\\
\textbf{3群以上}：一元配置分散分析 / Kruskal--Wallis，用量の順序 → 対比（線形対比）\\
\textbf{ベースラインあり（RCT）}→ ANCOVA，\textbf{対数正規（PK）}→ 対数変換，\textbf{同等性} → TOST・90\%CI};
\node[lst,anchor=north west] (bl) at ([yshift=-5mm]cl.south west) {%
\textbf{独立2群}：$\chi^2$（大標本）/ Fisher（小標本），RD・RR・OR の区間\\
\textbf{対応あり}（前後・ペア・同一患者に2検査）→ McNemar，条件付き OR $n_{12}/n_{21}$\\
\textbf{層別・交絡} → MH 検定・MH 推定量・標準化，\textbf{多数の共変量} → ロジスティック回帰・傾向スコア\\
\textbf{順序のある $K$ 群} → Cochran--Armitage，\textbf{診断} → 感度・特異度・PPV・ROC/AUC};
\node[lst,anchor=north west] (pl) at ([yshift=-5mm]bl.south west) {%
\textbf{人時間つき件数} → ポアソン率 $Y/T$，率比の区間，条件付き二項検定\\
\textbf{過分散} → 負の二項，\textbf{共変量} → ポアソン回帰（オフセット $\log T$），\textbf{予測} → Gamma--Poisson};
\node[lst,anchor=north west] (sl) at ([yshift=-5mm]pl.south west) {%
\textbf{記述} → Kaplan--Meier（Greenwood），\textbf{累積ハザード} → Nelson--Aalen\\
\textbf{2群比較} → ログランク（比例ハザード型の対立で効率的），早期差重視 → 一般化 Wilcoxon\\
\textbf{調整・効果の大きさ} → Cox（HR），ハザード一定 → 指数モデル，\textbf{非比例} → RMST・層別 Cox\\
\textbf{競合イベント} → CIF（Aalen--Johansen），原因別ハザード Cox / Fine--Gray};
\node[lst,anchor=north west] (ml) at ([yshift=-5mm]sl.south west) {%
\textbf{連続} → 線形混合モデル・MMRM（共分散構造を AIC で選択），欠測は MAR なら尤度法\\
\textbf{2 値・カウント} → GEE / 一般化線形混合モデル};
\foreach \l/\t/\n in {cl/連続/c,bl/2 値/b,pl/{カウント\\・率}/p,sl/生存時間/s,ml/反復測定/m}{
  \node[typ,left=6mm of \l] (\n) {\t};
  \draw[arr] (\n.east)--(\l.west);}
\node[hd,anchor=south] at ([yshift=2mm]c.north) {アウトカム};
\node[hd,anchor=south west] at ([yshift=2mm]cl.north west) {状況 → 手法};
\end{tikzpicture}
\caption{アウトカムからの手法選択}\label{fig:A-flow}
\end{figure}
\clearpage

\begin{center}\small
\begin{longtable}{@{}L{36mm}L{30mm}L{52mm}c@{}}
\caption{データ構造・推定対象・手法の対応}\label{tab:A-select}\\
\toprule
データ構造 & 推定対象 & 手法 & 章\\
\midrule
\endfirsthead
\toprule
データ構造 & 推定対象 & 手法 & 章\\
\midrule
\endhead
\bottomrule
\endfoot
連続・対応あり & 差の平均・中央値 & 対応のある $t$，符号付き順位，符号 & \ref{ch:03}\\
連続・独立2群 & 平均差，$\Prob(X>Y)$ & 2標本 $t$，Welch，順位和 & \ref{ch:03}\\
連続・前後2時点（RCT） & 最終時点の群差 & ANCOVA（$\V\propto1-\rho^2$） & \ref{ch:08}\\
連続・多時点 & 時点別平均差 & 混合モデル（MMRM） & \ref{ch:08}\\
2 値・独立2群 & RD，RR，OR & $\chi^2$，Fisher，Woolf & \ref{ch:04}\\
2 値・$K$ 群 & 一様性・傾向 & $\chi^2_{K-1}$，Cochran--Armitage & \ref{ch:04}\\
2 値・対応あり & 周辺割合の差 & McNemar，LRT & \ref{ch:04}\\
2 値・層別 & 共通 OR/RR & MH，Peto，Breslow--Day & \ref{ch:07}\\
2 値・共変量 & 調整 OR，予測確率 & ロジスティック回帰，AIC & \ref{ch:06}\\
観察研究の交絡 & 因果効果 & 標準化，傾向スコア，IPW & \ref{ch:07}\\
診断検査 & 感度・特異度・PPV・AUC & ROC，Youden，McNemar，デルタ法 & \ref{ch:05}\\
件数＋人時間 & 率，率比 & ポアソン，条件付き二項 & \ref{ch:17}\\
生存時間 & $S(t)$，中央値 & KM，NA，Greenwood & \ref{ch:09}\\
生存時間・2群 & 分布の差 & ログランク & \ref{ch:09}\\
生存時間・共変量 & HR & Cox，部分尤度 & \ref{ch:10}\\
生存時間・非比例 & 平均生存時間 & RMST & \ref{ch:10}\\
競合リスク & CIF & Aalen--Johansen，Fine--Gray & \ref{ch:11}\\
同等性・非劣性 & 差がマージン内 & TOST，$(1-2\alpha)$CI & \ref{ch:13}\\
中間解析 & 早期中止 & $\alpha$ 消費関数 & \ref{ch:14}\\
多重比較・多重評価項目 & FWER 制御 & Bonferroni，Holm，Dunnett，閉手順 & \ref{ch:14}\\
複数研究 & 統合効果 & 固定効果，変量効果（DL） & \ref{ch:15}\\
欠測あり & 完全データでの効果 & MMRM，MI（Rubin），感度分析 & \ref{ch:16}\\
小標本＋事前情報 & 事後確率・予測 & Beta--Binomial，Gamma--Poisson & \ref{ch:17}\\
\end{longtable}
\end{center}

% ============================================================
\chapter{必須公式集}\label[appendix]{app:B}
% ============================================================
各公式の右欄は「いつ使うか」である．記号は本文の定義に従う．

\begin{center}\small
\renewcommand{\arraystretch}{1.5}
\begin{longtable}{@{}L{25mm}L{80mm}L{45mm}@{}}
\toprule
項目 & 公式 & いつ使うか\\
\midrule
\endfirsthead
\toprule
項目 & 公式 & いつ使うか\\
\midrule
\endhead
\bottomrule
\endfoot
\multicolumn{3}{@{}l}{\textbf{効果指標・分割表}}\\
RD & $\hat\pi_1-\hat\pi_0$，$\V=\frac{\hat\pi_1(1-\hat\pi_1)}{n_1}+\frac{\hat\pi_0(1-\hat\pi_0)}{n_0}$ & 絶対的な効果．NNT $=1/\mathrm{ARR}$（利益），NNH $=1/\mathrm{ARI}$（害）\\
RR & $\frac{a/n_1}{c/n_0}$，$\V[\log\widehat{\RR}]=\frac1a-\frac1{n_1}+\frac1c-\frac1{n_0}$ & コホート・RCTの相対効果\\
OR & $\frac{ad}{bc}$，$\V[\log\widehat{\OR}]=\frac1a+\frac1b+\frac1c+\frac1d$ & 症例対照研究，ロジスティック\\
OR と RR & $\OR=\RR\frac{1-\pi_0}{1-\pi_1}$ & 稀な疾患なら OR $\approx$ RR\\
Pearson $\chi^2$ & $\frac{N(ad-bc)^2}{n_1n_0m_1m_0}\sim\chi^2_1$ & 独立2群の比率の比較\\
一様性 & $\frac{\sum n_i(p_i-p)^2}{p(1-p)}\sim\chi^2_{K-1}$ & $K$ 群の比率\\
傾向性 & $\frac{\{\sum y_i(x_i-\bar x)\}^2}{p(1-p)\sum n_i(x_i-\bar x)^2}\sim\chi^2_1$ & 用量に順序がある比率\\
McNemar & $\frac{(n_{12}-n_{21})^2}{n_{12}+n_{21}}\sim\chi^2_1$，$\V[\hat\delta]=\frac{\pi_{12}+\pi_{21}-(\pi_{12}-\pi_{21})^2}{n}$ & 対応のある2 値\\
超幾何 & $\E[a]=\frac{n_1m_1}{N}$，$\V[a]=\frac{n_1n_0m_1m_0}{N^2(N-1)}$ & MH 検定，ログランク\\
MH & $\chi^2_{\mathrm{MH}}=\frac{\{\sum(a_k-E_k)\}^2}{\sum V_k}$，$\widehat{\OR}_{\mathrm{MH}}=\frac{\sum a_kd_k/N_k}{\sum b_kc_k/N_k}$ & 層別表・交絡の調整\\
Peto & $\widehat{\log\psi}=\frac{\sum(O-E)}{\sum V}$，$\V=\frac1{\sum V}$ & OR が1に近い層別・メタ\\
標準化 & $R_{xw}=\sum_kw_kp_{xk}$ & 交絡の調整（直接法）\\
傾向スコア & $e(\bm z)=\Prob(X=1\mid\bm z)$，$X\indep\bm Z\mid e(\bm Z)$ & 多数の交絡因子の要約\\
\multicolumn{3}{@{}l}{\textbf{診断}}\\
感度・特異度 & $\Prob(+\mid D)$，$\Prob(-\mid\bar D)$ & 検査性能（有病率に依存しない）\\
PPV・NPV & $\frac{Se\,p}{Se\,p+(1-Sp)(1-p)}$，$\frac{Sp(1-p)}{Sp(1-p)+(1-Se)p}$ & 有病率 $p$ の集団での的中率\\
AUC & $\Prob(X_1>X_0)$，経験 $\frac{1}{n_1n_0}\sum\sum I(x_{1i}>x_{0j})$，2正規 $\Phi\bigl(\frac{\mu_1-\mu_0}{\sqrt{\sigma_1^2+\sigma_0^2}}\bigr)$ & ROC 曲線の要約\\
Youden & $J=Se+Sp-1$ の最大化 & カットオフの選択\\
\multicolumn{3}{@{}l}{\textbf{回帰}}\\
ロジスティック & $\ell=\sum\{y_i\bm x_i\transp\bm\beta-\log(1+e^{\bm x_i\transp\bm\beta})\}$，$\bm U=\bm X\transp(\bm y-\bm\pi)$，$\bm I=\bm X\transp\bm W\bm X$ & 2 値アウトカムの調整 OR\\
LRT & $2(\ell_1-\ell_0)\sim\chi^2_{\text{パラメータ数の差}}$ & 入れ子モデルの比較\\
AIC & $-2\ell(\hat{\bm\theta})+2k$ & モデル・共分散構造の選択\\
ベースライン & $\V_d=\frac{2\sigma^2}{n}$，$\V_c=\frac{4\sigma^2(1-\rho)}{n}$，$\V_a=\frac{2\sigma^2(1-\rho^2)}{n}$ & 最終値・変化量・ANCOVA\\
混合モデル & $\bm V=\bm Z\bm G\bm Z\transp+\bm R$，$\hat{\bm\beta}=(\bm X\transp\bm V^{-1}\bm X)^{-1}\bm X\transp\bm V^{-1}\bm Y$ & 反復測定\\
\multicolumn{3}{@{}l}{\textbf{順位検定}}\\
符号 & $T\sim\Bin(n,\frac12)$，$\E=\frac n2$，$\V=\frac n4$ & 対応あり・分布の仮定なし\\
符号付き順位 & $\E[W^+]=\frac{n(n+1)}4$，$\V=\frac{n(n+1)(2n+1)}{24}$ & 対応あり・対称分布\\
順位和 & $\E[W]=\frac{m(N+1)}2$，$\V=\frac{mn(N+1)}{12}$，$W=U+\frac{m(m+1)}2$ & 独立2群・非正規\\
\multicolumn{3}{@{}l}{\textbf{生存時間}}\\
基本関係 & $h=\frac fS$，$H=-\log S$，$S=e^{-H}$，$f=hS$ & 分布の変換\\
KM & $\hat S(t)=\prod_{t_{(j)}\le t}\bigl(1-\frac{d_j}{n_j}\bigr)$ & 打ち切りのある生存曲線\\
Greenwood & $\hat S^2\sum\frac{d_j}{n_j(n_j-d_j)}$ & KM の分散\\
Nelson--Aalen & $\hat H=\sum\frac{d_j}{n_j}$，$e^{-\hat H}\ge\hat S$ & 累積ハザード\\
ログランク & $E_{1j}=\frac{d_jn_{1j}}{n_j}$，$V_j=\frac{d_jn_{1j}n_{0j}(n_j-d_j)}{n_j^2(n_j-1)}$，$\frac{(O-E)^2}{V}\sim\chi^2_1$ & 生存曲線の2群比較\\
Cox & $h_0(t)e^{\bm\beta\transp\bm x}$，$L=\prod_j\frac{e^{\bm\beta\transp\bm x_{(j)}}}{\sum_{R_j}e^{\bm\beta\transp\bm x_l}}$，$S=S_0^{\exp(\bm\beta\transp\bm x)}$ & HR の推定・調整\\
指数モデル & $L=\lambda^de^{-\lambda W}$，$\hat\lambda=\frac dW$，$I=\frac d{\lambda^2}$，$\V[\log\hat\lambda]\approx\frac1d$ & ハザード一定・率\\
RMST & $\int_0^\tau S$，$\E[X^2]=2\int_0^\tau tS$；指数 $\frac{1-e^{-\lambda\tau}}{\lambda}$ & 非比例ハザードの要約\\
CIF & $\lambda=\sum\lambda_j$，$F_j=\int_0^t\lambda_jS$，$\hat F_j=\sum\frac{d_{ji}}{n_i}\hat S(t_i^-)$ & 競合リスク\\
\multicolumn{3}{@{}l}{\textbf{設計・臨床試験}}\\
検出力 & $\mu=z_\alpha+z_\beta$，検出力 $\Phi(\mu-z_\alpha)$ & すべての症例数計算の基本\\
平均の差 & $n=\frac{2\sigma^2(z_{\alpha/2}+z_\beta)^2}{\Delta^2}$（各群） & 連続アウトカム\\
比率の差 & $n=\frac{\{z_{\alpha/2}\sqrt{2\bar\pi\bar q}+z_\beta\sqrt{\pi_1q_1+\pi_0q_0}\}^2}{(\pi_1-\pi_0)^2}$ & 2 値アウトカム\\
イベント数 & $D=\frac{4(z_{\alpha/2}+z_\beta)^2}{(\log\HR)^2}$（総数，1:1） & 生存時間試験\\
同等性 & $(1-2\alpha)$CI $\subset(-\Delta,\Delta)$，BE：GMR の90\%CI $\subset(0.80,1.25)$ & 同等性・生物学的同等性\\
対数正規 & 中央値 $e^\mu$，$\E=e^{\mu+\sigma^2/2}$，$\mathrm{CV}=\sqrt{e^{\sigma^2}-1}$ & 薬物動態パラメータ\\
$\alpha$ 消費 & $\Corr[Z(t_1),Z(t_2)]=\sqrt{t_1/t_2}$，増分 $\alpha(t_k)-\alpha(t_{k-1})$ & 中間解析\\
二段階 & $\alpha=\sum_{x>r_1}b(x;\pi_0,n_1)\{1-B(r-x;\pi_0,n_2)\}$，$\E[N]=n_1+(1-\mathrm{PET})n_2$ & 第II相単群\\
多重性 & $1-(1-\alpha)^m$，Bonferroni $\alpha/m$，多対1の相関 $\sqrt{\frac{n_jn_k}{(n_j+n_0)(n_k+n_0)}}$ & 複数の検定\\
\multicolumn{3}{@{}l}{\textbf{メタアナリシス・欠測・ベイズ}}\\
固定効果 & $\hat\theta=\frac{\sum w_kY_k}{\sum w_k}$，$\V=\frac1{\sum w_k}$ & 共通効果の統合\\
異質性 & $Q=\sum w_k(Y_k-\hat\theta)^2$，$\hat\tau^2=\max\bigl\{0,\frac{Q-(K-1)}{\sum w-\sum w^2/\sum w}\bigr\}$，$I^2=\max\bigl\{0,\frac{Q-(K-1)}{Q}\bigr\}$ & 研究間のばらつき\\
変量効果 & $\hat\mu=\frac{\sum w_k^*Y_k}{\sum w_k^*}$，$w_k^*=\frac1{\sigma_k^2+\tau^2}$ & 効果の分布の平均\\
Rubin & $T=W+\bigl(1+\frac1M\bigr)B$ & 多重補完の統合\\
ポアソン & $\hat\lambda=\frac YT$，$\V[\log\hat\lambda]\approx\frac1Y$，$Y_A\mid n\sim\Bin\bigl(n,\frac{T_A}{T_A+T_B}\bigr)$（$H_0$ 下） & 人時間つき件数\\
Beta--Binomial & $\Be(a+y,b+n-y)$，事後平均 $\frac{a+y}{a+b+n}$ & 割合の Bayes 推測\\
Gamma--Poisson & $\Ga(a+y,b+T)$（形状・率），予測は $\NB\bigl(a+y,\frac{b+T}{b+T+\tilde T}\bigr)$ & 率の Bayes 推測・予測\\
\end{longtable}
\end{center}

% ============================================================
\chapter{医薬生物学用語集}\label[appendix]{app:C}
% ============================================================
\begin{center}\small\renewcommand{\arraystretch}{1.05}
\begin{longtable}{@{}L{32mm}L{40mm}L{78mm}@{}}
\toprule
用語 & 英語 & 説明\\
\midrule
\endfirsthead
\toprule
用語 & 英語 & 説明\\
\midrule
\endhead
\bottomrule
\endfoot
評価項目 & endpoint & 治療効果を評価する変数．主要評価項目は事前に規定する．複数（co-primary など）とする場合は，成功の条件と多重性の処理も事前に規定する\\
代替評価項目 & surrogate endpoint & 真の評価項目（死亡など）の代わりに用いる検査値など\\
イベント & event & 生存時間解析で観察する事象（死亡・再発）\\
打ち切り & censoring & イベント前に観察が終わること．右側打ち切りが基本\\
独立な打ち切り & independent censoring & 打ち切りがその後の予後と無関係であること\\
リスク集合 & risk set & ある時点の直前にイベントも打ち切りもない人の集合\\
ハザード & hazard & 生存者の瞬間的なイベント発生率 $f/S$\\
累積ハザード & cumulative hazard & $H(t)=\int_0^th=-\log S(t)$\\
ハザード比 & hazard ratio (HR) & 2群のハザードの比．比例ハザードなら時間によらず一定\\
比例ハザード & proportional hazards & HR が時間によらず一定であること\\
中央生存時間 & median survival time & $S(t)=0.5$ となる時点\\
境界内平均生存時間 & restricted mean survival time (RMST) & $\int_0^\tau S(t)\dd t$\\
競合リスク & competing risk & あるイベントの発生により他のイベントが起こりえなくなる状況\\
原因別ハザード & cause-specific hazard & 無事な人での原因 $j$ の瞬間的発生率\\
累積発生関数 & cumulative incidence function (CIF) & $\Prob(T\le t,J=j)$\\
部分分布ハザード & subdistribution hazard & CIF に対応するハザード（Fine--Gray）\\
無作為化 & randomization & 治療を確率的に割り付けること．交絡を防ぐ\\
層別無作為化 & stratified randomization & 予後因子の層ごとに割り付けのバランスをとる\\
盲検化 & blinding & 割付を被験者・医師・評価者に知らせないこと\\
二重盲検 & double-blind & 被験者と医師の双方が割付を知らない\\
プラセボ & placebo & 有効成分を含まない偽薬\\
並行群間比較 & parallel group design & 各被験者が1つの治療のみを受けるデザイン\\
クロスオーバー & crossover design & 各被験者が時期を変えて複数の治療を受けるデザイン\\
休薬期間 & washout period & クロスオーバーで持ち越し効果を除くための期間\\
持ち越し効果 & carryover effect & 前の時期の治療効果が次の時期に残ること\\
優越性 & superiority & 対照より優れること\\
同等性 & equivalence & 差がマージン $(-\Delta,\Delta)$ 内にあること\\
非劣性 & non-inferiority & 対照より $\Delta$ を超えて劣らないこと\\
非劣性マージン & non-inferiority margin & 許容できる最大の劣り幅\\
生物学的同等性 & bioequivalence (BE) & 後発品と先発品の吸収の速度・量が同等であること\\
幾何平均比 & geometric mean ratio (GMR) & 対数スケールの平均差を逆変換した比\\
分析感度 & assay sensitivity & 試験が治療間の差を検出できる能力\\
ITT & intention-to-treat & 割り付けられた群のまま全例を解析する原則\\
推定対象 & estimand & 推定したい治療効果の定義（ICH E9(R1)）\\
中間事象 & intercurrent event & 治療開始後に起こり，評価項目の解釈または存在に影響する事象\\
感度 & sensitivity & 罹患者が陽性になる確率\\
特異度 & specificity & 非罹患者が陰性になる確率\\
陽性的中率 & positive predictive value (PPV) & 陽性者が罹患している確率\\
陰性的中率 & negative predictive value (NPV) & 陰性者が罹患していない確率\\
有病率 & prevalence & 集団で疾患を有する割合\\
尤度比 & likelihood ratio (LR) & $Se/(1-Sp)$ など．検査後オッズ＝LR×検査前オッズ\\
ROC 曲線 & receiver operating characteristic curve & (偽陽性率, 真陽性率) の軌跡\\
AUC & area under the curve & ROC 曲線下面積 $=\Prob(X_1>X_0)$．薬物動態では血中濃度曲線下面積\\
リスク & risk & 一定期間の発生確率（累積発生割合）\\
率 & rate & 人時間あたりの発生数\\
オッズ比 & odds ratio (OR) & オッズの比．症例対照研究で推定可能\\
リスク比 & risk ratio (RR) & リスクの比．相対リスク\\
リスク差 & risk difference (RD) & リスクの差\\
NNT & number needed to treat & $1/\mathrm{ARR}$（ARR：絶対リスク減少）．1人の発症を防ぐのに治療が必要な人数\\
NNH & number needed to harm & $1/\mathrm{ARI}$（ARI：絶対リスク増加）．治療により1人余分に害が生じる人数\\
コホート研究 & cohort study & 曝露で分けて追跡する観察研究\\
症例対照研究 & case-control study & 症例と対照を集めて過去の曝露を比べる研究\\
交絡 & confounding & 曝露と結果の両方に関連する因子による偏り\\
効果修飾 & effect modification & 効果の大きさが第3の変数で変わること\\
層別 & stratification & 交絡因子の値ごとに分けて解析すること\\
標準化 & standardization & 共通の重みで層別リスクを平均すること\\
マッチング & matching & 背景の似た個体を組にすること\\
傾向スコア & propensity score & 共変量を与えたときの治療を受ける確率\\
非可縮性 & non-collapsibility & 交絡がなくても層別と併合で値が変わる性質（OR，HR）\\
Simpson のパラドックス & Simpson's paradox & 併合した関連が層別の関連と逆転する現象\\
異質性 & heterogeneity & 研究間で真の効果が異なること\\
固定効果モデル & fixed-effect model & 全研究に共通の効果を仮定するモデル\\
変量効果モデル & random-effects model & 研究ごとの効果が分布に従うとするモデル\\
フォレストプロット & forest plot & 各研究と統合の推定値・区間を並べた図\\
出版バイアス & publication bias & 有意な結果ほど出版されやすいことによる偏り\\
欠測 & missing data & 予定された測定が得られないこと\\
MCAR・MAR・MNAR & missing completely at random / at random / not at random & 欠測がデータと無関係／観測値のみに依存／欠測値自体に依存\\
多重補完 & multiple imputation (MI) & 欠測値を複数回補完して統合する方法\\
完全ケース解析 & complete-case analysis & 欠測のない個体のみを用いる解析\\
感度分析 & sensitivity analysis & 仮定を変えたときの結論の頑健性の確認\\
中間解析 & interim analysis & 試験途中で行う解析\\
情報時間 & information time & 最終解析に対する情報量の割合\\
$\alpha$ 消費関数 & alpha spending function & 中間解析で使う第1種の過誤の配分を定める関数\\
早期中止 & early stopping & 有効・無効・安全性により試験を途中で終えること\\
多重性 & multiplicity & 複数の検定による第1種の過誤の増大\\
FWER & familywise error rate & 少なくとも1つの真の帰無仮説を棄却する確率\\
閉検定手順 & closed testing procedure & 共通部分仮説をすべて検定して強い FWER 制御を行う方法\\
ゲートキーピング & gatekeeping & 順序を決めて前の仮説の棄却を条件に次を検定する方法\\
用量反応 & dose-response & 用量とともに効果が変化する関係\\
対比 & contrast & 係数の和が0の平均の線形結合\\
共分散分析 & analysis of covariance (ANCOVA) & ベースラインなどの共変量で調整した群比較\\
平均への回帰 & regression to the mean & 極端な初回値の人の次回値が平均に近づく現象\\
混合効果モデル & mixed-effects model & 固定効果と変量効果を含むモデル\\
級内相関係数 & intraclass correlation (ICC) & 同一個体内の測定値の相関\\
事前分布・事後分布 & prior / posterior distribution & Bayes 推測でのパラメータの分布\\
共役事前分布 & conjugate prior & 事後分布が事前分布と同じ分布族になる事前分布\\
信用区間 & credible interval & 事後分布の分位点で作る区間\\
事前予測分布 & prior predictive distribution & データの周辺分布 $\int p(y\mid\theta)p(\theta)\dd\theta$\\
人時間 & person-time & 各人の観察期間の合計．率の分母\\
過分散 & overdispersion & 件数の分散が平均を超えること（ポアソンより大きいばらつき）\\
事後予測分布 & posterior predictive distribution & 将来のデータの予測分布\\
\end{longtable}
\end{center}

% ============================================================
\chapter{過去問対応表}\label[appendix]{app:D}
% ============================================================
必要能力の略号：定＝定義，計＝計算，導＝導出，検＝検定，モ＝モデル化，選＝手法選択，解＝解釈．
問5は4分野共通問題（「共」）で，医薬に関連する部分がある場合のみ対応章を示す．

\begin{center}\scriptsize
\setlength{\tabcolsep}{3pt}
\begin{longtable}{@{}ccL{36mm}L{58mm}L{16mm}c@{}}
\toprule
年 & 問 & 主テーマ & 副テーマ & 必要能力 & 章\\
\midrule
\endfirsthead
\toprule
年 & 問 & 主テーマ & 副テーマ & 必要能力 & 章\\
\midrule
\endhead
\bottomrule
\endfoot
2016 & 1 & 対応のある連続データ & 対応のある $t$，符号検定（exact），符号付き順位（正規近似），外れ値の影響 & 検 計 導 解 & \ref{ch:03}\\
2016 & 2 & 4群の比率 & 一様性 $\chi^2_3$，$P=0.077$ の解釈，Cochran--Armitage（重み付き最小二乗） & 検 導 解 選 & \ref{ch:04},\ref{ch:14}\\
2016 & 3 & 生存時間解析 & KM 曲線，ログランク（$\chi^2=5.62$），Cox 部分尤度 & 計 導 & \ref{ch:09},\ref{ch:10}\\
2016 & 4 & 対数正規と BE & 中央値，密度，平均・分散，CV から $\sigma^2$，クロスオーバーの90\%CI & 導 計 解 & \ref{ch:13}\\
2016 & 5共 & 2標本 $t$ 検定 & 群ごとの CI と検定，CI の重なり（約29\%） & 計 検 解 & \ref{ch:01}\\
2017 & 1 & ハザードの理論 & $\lambda=f/S$，$\Lambda=-\log S$，$\Lambda(T)\sim\Ex(1)$，Cox の2重対数表現，極値分布，乱数生成 & 導 モ & \ref{ch:09},\ref{ch:10}\\
2017 & 2 & 群逐次デザイン & $\alpha$ 消費関数，情報時間，$\Cov[Z(0.4),Z(1)]=\sqrt{0.4}$，限界値 & 計 導 & \ref{ch:14}\\
2017 & 3 & 第II相二段階 & $\alpha,\beta$ の二項和，PET，期待症例数（Simon） & 導 計 & \ref{ch:14}\\
2017 & 4 & 層別症例対照 & 曝露 OR，超幾何の平均分散，Peto 型推定，共通 OR の CI，均一性 & 計 導 解 & \ref{ch:04},\ref{ch:07}\\
2017 & 5共 & 割付法 & 単純無作為化と偏りを抑える割付，CI 幅から $n$ & 計 導 & \ref{ch:02},\ref{ch:12}\\
2018 & 1 & 指数生存モデル & 打ち切り尤度，MLE，Fisher 情報，デルタ法，必要イベント数 & 導 計 & \ref{ch:10},\ref{ch:12}\\
2018 & 2 & マッチドペア2 値 & 多項分布，$\V[\hat\delta]$，CI，制約付き MLE，LRT & 導 検 & \ref{ch:04}\\
2018 & 3 & ROC 解析 & 感度・偽陽性率，経験 ROC，Youden，AUC＝Mann--Whitney & 計 導 解 & \ref{ch:05}\\
2018 & 4 & ロジスティック回帰 & 対数尤度，調整 OR，AIC，Bernoulli の KL，リスクスコアの解釈 & 導 解 & \ref{ch:06},\ref{ch:05}\\
2018 & 5共 & 正規混合 & 平均・分散，二峰性の条件 & 導 & ---\\
2019 & 1 & KM・NA・RMST & KM 表，NA $\ge$ KM，RMST $=\int S$，指数での分散・推定 & 計 導 & \ref{ch:09},\ref{ch:10}\\
2019 & 2 & 交絡と傾向スコア & 粗 RD の CI，直接標準化，飽和 PS モデル，バランス特性の確認と証明 & 計 導 解 & \ref{ch:07}\\
2019 & 3 & 対応のある診断比較 & 感度・PPV，McNemar 型，8セル多項＋デルタ法 & 計 検 導 & \ref{ch:05},\ref{ch:04}\\
2019 & 4 & 小標本2群比較 & Student $t$ の仮定，順位和の exact 分布，離散性による問題 & 検 計 解 & \ref{ch:03}\\
2019 & 5共 & 適合度と EM & 適合度 $\chi^2$，自由度の根拠，ABO 血液型の EM & 検 導 & \ref{ch:16}\\
2021 & 1 & 競合リスク & $f=\lambda S$，$1-$KM，CIF の導出と推定，過大推定 & 導 計 解 & \ref{ch:09},\ref{ch:11}\\
2021 & 2 & 同等性と BE & 片側 $t$ 検定，IUT，90\%CI による BE 判定 & 導 検 解 & \ref{ch:13}\\
2021 & 3 & 2×2表と層別 & OR と Woolf CI，積二項尤度，スコア型 $X$，Cochran 検定，Simpson & 計 導 検 解 & \ref{ch:04},\ref{ch:07}\\
2021 & 4 & 変量効果メタ & 周辺分布，対数尤度，$\hat\mu$，$\V[\hat\mu]$，CI & 導 計 & \ref{ch:15}\\
2021 & 5共 & 診断の正規モデル & 感度・特異度，PPV・NPV，閾値，有病率の影響 & 計 解 & \ref{ch:05}\\
2022 & 1 & Cox 比例ハザード & 比例性，$S_0^{\exp(\beta x)}$，2重対数プロット，HR の CI と解釈 & 導 計 解 & \ref{ch:10}\\
2022 & 2 & マッチと MH & 粗 RR・OR，MH-RR・OR，非可縮性，間接比較 & 計 解 & \ref{ch:04},\ref{ch:07}\\
2022 & 3 & 線形混合モデル & $\bm V=\bm Z\bm G\bm Z\transp+\bm R$，OLS/GLS，複合対称，被験者内相関 & 導 計 モ & \ref{ch:08}\\
2022 & 4 & Beta--Binomial & 事後分布・モード・平均，周辺分布，事後確率による判定 & 導 計 解 & \ref{ch:17}\\
2022 & 5共 & 平均への回帰 & 2変量正規の条件付き分布，真値モデル & 導 計 解 & \ref{ch:08}\\
2023 & 1 & ロジスティック回帰 & 対数オッズ，集計データの尤度，MLE，LRT，PPV・NPV & 計 導 検 & \ref{ch:06},\ref{ch:05}\\
2023 & 2 & 単群指数試験の症例数 & MLE，Fisher 情報，$\mu=z_\alpha+z_\beta$，$n=32$ & 導 計 & \ref{ch:12},\ref{ch:10}\\
2023 & 3 & メタアナリシス & 固定効果 MLE，$Q$ の分解と期待値，DL，$I^2$ & 導 計 & \ref{ch:15}\\
2023 & 4 & 反復測定 & $\hat{\bm\mu}$ の分布，Toeplitz 対無構造の AIC，相関行列，変化量の検定 & 計 モ 検 & \ref{ch:08}\\
2023 & 5共 & 条件付き確率と二項検定 & Monty Hall，exact と正規近似 & 計 & ---\\
2024 & 1 & 診断・ROC・McNemar & 感度等の表，台形 AUC，対応のある2 値の分散と McNemar & 計 導 検 & \ref{ch:05},\ref{ch:04}\\
2024 & 2 & ベースライン調整 & 最終値・変化量・ANCOVA の分散，不偏条件，相対分散 & 導 計 解 & \ref{ch:08}\\
2024 & 3 & 欠測と多重補完 & 欠測と群の $\chi^2$，補完内分散，Rubin，$\E[B]$ & 計 導 & \ref{ch:16}\\
2024 & 4 & 多重性と用量反応 & FWER，Bonferroni，ゲート付き検定，多対1の相関，対比 & 計 導 選 & \ref{ch:14}\\
2024 & 5共 & 重回帰の抑制 & 正規方程式，決定係数 & 導 & ---\\
2025 & 1 & 順位和検定の検出力 & 位置ずれ，$\E[R_X]$，$\Prob(X>Y)$，検出力の式 & 導 計 & \ref{ch:03},\ref{ch:12}\\
2025 & 2 & ポアソン率と Bayes & 共通率 MLE，率の差の検定，Gamma 事前，事後予測（負の二項） & 導 計 解 & \ref{ch:17}\\
2025 & 3 & 競合リスクの理論 & $\lambda_j=f_j/S$，$1-F_j>S_j$，$\lambda=\sum\lambda_j$，Weibull 型，原因別 Cox の解釈 & 導 解 & \ref{ch:11}\\
2025 & 4 & 多重性 & 独立 $p$ 値の FWER，Bonferroni，Šidák，Simes，相関の影響 & 計 導 & \ref{ch:14}\\
2025 & 5共 & ランダム化回答法 & 誤分類の補正，MLE，$p$ の選択 & 導 解 & ---\\
\end{longtable}
\end{center}

% ============================================================
\chapter{総合チェックリスト}\label[appendix]{app:E}
% ============================================================
試験前に，各項目を\textbf{何も見ずに}説明・導出できるか確認する．できない項目は対応する章に戻る．
\newcommand{\chk}{\item[$\square$]}

\section*{第1〜2章　問題の読み方・効果指標}
\begin{itemize}
\chk 設問からデザイン・観測単位・アウトカムの型・推定対象を書き出せる．
\chk $P=0.077$，「CI が1を含む」，「有意差なし」を正しく解釈できる．
\chk 群ごとの CI の重なりと2標本検定の関係（約29\%）を導ける．
\chk $\OR=\RR(1-\pi_0)/(1-\pi_1)$ と，症例対照研究で OR だけが推定できる理由を示せる．
\chk OR の非可縮性を数値例で説明できる．estimand の5要素を挙げられる．
\end{itemize}
\section*{第3章　連続・順位検定}
\begin{itemize}
\chk 符号検定・符号付き順位・順位和の帰無仮説と仮定を区別できる．
\chk $\E[W^+],\V[W^+]$，$\E[W],\V[W]$ を導出できる（非復元抽出の修正を含む）．
\chk 小標本の exact 分布を数え上げ，最小の P 値を求められる．
\chk $W=U+m(m+1)/2$ と，$\Prob(X>Y)$ を用いた検出力を導ける．
\end{itemize}
\section*{第4章　分割表}
\begin{itemize}
\chk $X^2=N(ad-bc)^2/(n_1n_0m_1m_0)$ をスコア型統計量として導ける．
\chk デルタ法で $\V[\log\widehat{\RR}]$，$\V[\log\widehat{\OR}]$ を導ける．
\chk Cochran--Armitage の $\hat\beta$ と $\V[\hat\beta\mid H_0]$ を導き，一様性検定との使い分けを説明できる．
\chk McNemar の分散・検定・LRT（ラグランジュ乗数法）を導ける．
\chk 超幾何分布の平均・分散から MH 検定を書ける．
\end{itemize}
\section*{第5章　診断}
\begin{itemize}
\chk 感度・特異度・PPV・NPV を定義し，有病率から PPV を計算できる．
\chk 経験 ROC を描き，台形公式で AUC を計算し，Youden のカットオフを選べる．
\chk AUC $=\Prob(X_1>X_0)$ と経験 AUC＝Mann--Whitney を示せる．
\chk 同一被験者の2検査の感度を McNemar で，PPV を多項＋デルタ法で比較する方法を説明できる．
\end{itemize}
\section*{第6〜7章　ロジスティック・交絡}
\begin{itemize}
\chk ロジスティックモデルの対数尤度・スコア・情報行列を導ける．2群なら Woolf と一致することを示せる．
\chk LRT の自由度，AIC のパラメータ数を正しく数えられる．
\chk 交絡と効果修飾を区別し，Simpson のパラドックスを例示できる．
\chk MH-OR，Peto 推定，直接標準化を計算できる．
\chk 傾向スコアのバランス特性を繰り返し期待値で証明できる．
\end{itemize}
\section*{第8章　縦断・混合モデル}
\begin{itemize}
\chk 最終値・変化量・ANCOVA の分散 $2\sigma^2/n$，$4\sigma^2(1-\rho)/n$，$2\sigma^2(1-\rho^2)/n$ を導ける．
\chk 平均への回帰を2変量正規と真値モデルで説明できる．
\chk $\bm V=\bm Z\bm G\bm Z\transp+\bm R$，GLS 推定量，変量切片が複合対称を生むことを示せる．
\chk 共分散構造のパラメータ数を数え，AIC を計算できる．
\end{itemize}
\section*{第9〜11章　生存時間}
\begin{itemize}
\chk $h=f/S$，$H=-\log S$，$H(T)\sim\Ex(1)$ を証明できる．
\chk KM の表を作り，曲線を描き，中央値を読める．Greenwood を導ける．
\chk NA 推定量と $e^{-\hat H}\ge\hat S$ を示せる．
\chk ログランクの $E_{1j},V_j$ を超幾何分布から導き，手計算できる．
\chk Cox の部分尤度を書き下し，$S=S_0^{\exp(\beta x)}$ と2重対数プロットの平行性を示せる．
\chk 打ち切りのある指数分布の尤度・MLE・Fisher 情報を導き，デルタ法で HR の検定を作れる．
\chk RMST $=\int_0^\tau S$ を示し，指数モデルと KM で計算できる．
\chk $\lambda=\sum\lambda_j$，$F_j=\int\lambda_jS$ を導き，Aalen--Johansen 推定値を計算できる．$1-$KM の過大推定を示せる．
\chk 原因別 Cox と Fine--Gray の違いを説明できる．
\end{itemize}
\section*{第12〜14章　臨床試験デザイン}
\begin{itemize}
\chk $\mu=z_\alpha+z_\beta$ を導き，平均・比率・イベント数の症例数を計算できる．
\chk TOST の有意水準が $\alpha$ 以下であること，90\%CI との同値性を示せる．
\chk 対数正規の中央値・平均・分散・CV を導き，BE を判定できる．
\chk $\Corr[Z(t_1),Z(t_2)]=\sqrt{t_1/t_2}$ と $\alpha$ 消費関数による限界値の決め方を説明できる．
\chk 二段階デザインの $\alpha,\beta$，PET，$\E[N]$ を計算できる．
\chk FWER，Bonferroni，Holm，Simes，閉検定手順，多対1の相関，対比の相関を説明・計算できる．
\end{itemize}
\section*{第15〜17章　メタ・欠測・ベイズ}
\begin{itemize}
\chk 固定効果の MLE と分散，$\E[Q]$，DL の $\hat\tau^2$，$I^2$ を導ける．
\chk 変量効果の $\hat\mu$ と CI を計算し，固定効果との違いを説明できる．
\chk MCAR・MAR・MNAR を定義・分類でき，MAR での尤度の無視可能性を示せる．
\chk Rubin のルールを計算し，観測を与えた下で $\E[B]=\V[\hat\theta_k]$ を示せる．適切な補完の必要性を説明できる．
\chk ポアソン率の推定・比較（差・比・条件付き二項）ができる．
\chk Beta--Binomial，Gamma--Poisson の事後分布・事後予測分布を導ける．ガンマのパラメータ化を明示できる．
\chk $\Prob(\pi>p\mid y)=\Prob(\Bin(n+1,p)\le y)$ を用いて事後確率による判定ができる．
\end{itemize}

% ===== 99_bib.tex =====
\backmatter
\renewcommand{\bibname}{参考文献}
\begin{thebibliography}{99}

\item[] \textbf{■ ローカル教材（本書の執筆で参照したもの）}
\bibitem{kubokawa} 久保川達也．『現代数理統計学の基礎』．共立出版，2017．（第2〜7章を参照）
\bibitem{jss:past} 日本統計学会 編．『日本統計学会公式認定 統計検定1級 公式問題集』（2016〜2019年，2021〜2025年の各年版）．
      統計応用（医薬生物学）および統計数理の問題・解答例．
      （2025年の解答は公式略解に導出過程を補ったもの．本書では「2022年 医薬 問2」「2016年 数理 問2」の形で参照した．）

\item[] \textbf{■ ガイドライン・公的資料（Web で確認）}
\bibitem{ich:e9r1} ICH. E9(R1) Addendum on Estimands and Sensitivity Analysis in Clinical Trials to the Guideline on Statistical Principles for Clinical Trials. 2019.
      \url{https://www.fda.gov/media/108698/download}（FDA 版 2021）
\bibitem{ich:e9} ICH. E9 Statistical Principles for Clinical Trials. 1998.
\bibitem{ich:e10} ICH. E10 Choice of Control Group and Related Issues in Clinical Trials. 2000.
\bibitem{jp:be} 厚生労働省．後発医薬品の生物学的同等性試験ガイドライン（平成9年12月22日 医薬審第487号，最終改正 令和2年3月19日 薬生薬審発0319第1号）．
      \url{https://www.pmda.go.jp/files/000234565.pdf}
\bibitem{fda:ni} U.S. FDA. Non-Inferiority Clinical Trials to Establish Effectiveness: Guidance for Industry. 2016.
\bibitem{fda:multiple} U.S. FDA. Multiple Endpoints in Clinical Trials: Guidance for Industry. 2022.
\bibitem{ema:missing} EMA. Guideline on Missing Data in Confirmatory Clinical Trials. 2010.
\bibitem{web:lec} LEC東京リーガルマインド．統計検定1級合格講座（試験形式の説明）．\url{https://www.lec-jp.com/jssc/class1/}
\bibitem{web:academaid} Academaid．統計検定1級合格へのロードマップ．\url{https://academ-aid.com/qualif/jssc/roadmap-jssc1}
\bibitem{estimandprimer} Kahan, B.~C. et al. The estimands framework: a primer on the ICH E9(R1) addendum. \textit{BMJ}, 384, e076316, 2024.
      \url{https://pmc.ncbi.nlm.nih.gov/articles/PMC10802140/}

\item[] \textbf{■ 原論文・標準的教科書}
\bibitem{km1958} Kaplan, E.~L. and Meier, P. Nonparametric estimation from incomplete observations. \textit{JASA}, 53, 457--481, 1958.
\bibitem{greenwood1926} Greenwood, M. The natural duration of cancer. \textit{Reports on Public Health and Medical Subjects}, 33, 1--26, 1926.
\bibitem{cox1972} Cox, D.~R. Regression models and life-tables (with discussion). \textit{JRSS B}, 34, 187--220, 1972.
\bibitem{mantel1966} Mantel, N. Evaluation of survival data and two new rank order statistics arising in its consideration. \textit{Cancer Chemotherapy Reports}, 50, 163--170, 1966.
\bibitem{mh1959} Mantel, N. and Haenszel, W. Statistical aspects of the analysis of data from retrospective studies of disease. \textit{JNCI}, 22, 719--748, 1959.
\bibitem{rbg1986} Robins, J., Breslow, N. and Greenland, S. Estimators of the Mantel--Haenszel variance consistent in both sparse data and large-strata limiting models. \textit{Biometrics}, 42, 311--323, 1986.
\bibitem{armitage1955} Armitage, P. Tests for linear trends in proportions and frequencies. \textit{Biometrics}, 11, 375--386, 1955.
\bibitem{mcnemar1947} McNemar, Q. Note on the sampling error of the difference between correlated proportions or percentages. \textit{Psychometrika}, 12, 153--157, 1947.
\bibitem{rr1983} Rosenbaum, P.~R. and Rubin, D.~B. The central role of the propensity score in observational studies for causal effects. \textit{Biometrika}, 70, 41--55, 1983.
\bibitem{finegray1999} Fine, J.~P. and Gray, R.~J. A proportional hazards model for the subdistribution of a competing risk. \textit{JASA}, 94, 496--509, 1999.
\bibitem{schoenfeld1983} Schoenfeld, D.~A. Sample-size formula for the proportional-hazards regression model. \textit{Biometrics}, 39, 499--503, 1983.
\bibitem{royston2013} Royston, P. and Parmar, M.~K.~B. Restricted mean survival time: an alternative to the hazard ratio for the design and analysis of randomized trials with a time-to-event outcome. \textit{BMC Medical Research Methodology}, 13, 152, 2013.
\bibitem{landemets1983} Lan, K.~K.~G. and DeMets, D.~L. Discrete sequential boundaries for clinical trials. \textit{Biometrika}, 70, 659--663, 1983.
\bibitem{obf1979} O'Brien, P.~C. and Fleming, T.~R. A multiple testing procedure for clinical trials. \textit{Biometrics}, 35, 549--556, 1979.
\bibitem{simon1989} Simon, R. Optimal two-stage designs for phase II clinical trials. \textit{Controlled Clinical Trials}, 10, 1--10, 1989.
\bibitem{holm1979} Holm, S. A simple sequentially rejective multiple test procedure. \textit{Scandinavian Journal of Statistics}, 6, 65--70, 1979.
\bibitem{simes1986} Simes, R.~J. An improved Bonferroni procedure for multiple tests of significance. \textit{Biometrika}, 73, 751--754, 1986.
\bibitem{dunnett1955} Dunnett, C.~W. A multiple comparison procedure for comparing several treatments with a control. \textit{JASA}, 50, 1096--1121, 1955.
\bibitem{schuirmann1987} Schuirmann, D.~J. A comparison of the two one-sided tests procedure and the power approach for assessing the equivalence of average bioavailability. \textit{Journal of Pharmacokinetics and Biopharmaceutics}, 15, 657--680, 1987.
\bibitem{dl1986} DerSimonian, R. and Laird, N. Meta-analysis in clinical trials. \textit{Controlled Clinical Trials}, 7, 177--188, 1986.
\bibitem{higgins2002} Higgins, J.~P.~T. and Thompson, S.~G. Quantifying heterogeneity in a meta-analysis. \textit{Statistics in Medicine}, 21, 1539--1558, 2002.
\bibitem{higgins2003} Higgins, J.~P.~T., Thompson, S.~G., Deeks, J.~J. and Altman, D.~G. Measuring inconsistency in meta-analyses. \textit{BMJ}, 327, 557--560, 2003.
\bibitem{higgins2009} Higgins, J.~P.~T., Thompson, S.~G. and Spiegelhalter, D.~J. A re-evaluation of random-effects meta-analysis. \textit{JRSS A}, 172, 137--159, 2009.
\bibitem{rubin1987} Rubin, D.~B. \textit{Multiple Imputation for Nonresponse in Surveys}. Wiley, 1987.
\bibitem{littlerubin} Little, R.~J.~A. and Rubin, D.~B. \textit{Statistical Analysis with Missing Data}, 3rd ed. Wiley, 2019.
\bibitem{delong1988} DeLong, E.~R., DeLong, D.~M. and Clarke-Pearson, D.~L. Comparing the areas under two or more correlated receiver operating characteristic curves: a nonparametric approach. \textit{Biometrics}, 44, 837--845, 1988.
\bibitem{hanley1982} Hanley, J.~A. and McNeil, B.~J. The meaning and use of the area under a receiver operating characteristic (ROC) curve. \textit{Radiology}, 143, 29--36, 1982.
\bibitem{kp2002} Kalbfleisch, J.~D. and Prentice, R.~L. \textit{The Statistical Analysis of Failure Time Data}, 2nd ed. Wiley, 2002.
\bibitem{agresti2013} Agresti, A. \textit{Categorical Data Analysis}, 3rd ed. Wiley, 2013.
\bibitem{fitzmaurice2011} Fitzmaurice, G.~M., Laird, N.~M. and Ware, J.~H. \textit{Applied Longitudinal Analysis}, 2nd ed. Wiley, 2011.
\bibitem{rothman2008} Rothman, K.~J., Greenland, S. and Lash, T.~L. \textit{Modern Epidemiology}, 3rd ed. Lippincott Williams \& Wilkins, 2008.
\end{thebibliography}

\end{document}
